% Generated by build_chapter.py from the GY1977 files. Do not edit by hand; edit the files and rebuild.
\chapter{The academic year 1977--78}\label{ch:1977}
\noindent\textbf{Scope and editorial date.}
This chapter covers all 28 requested catalogue records and all 94 supplied
source pages, with status assessed on 10 November 2026. It indexes 305
mathematical occurrences, not 305 independent discoveries. The packet is
not a perfectly delimited academic year. Its record envelopes overlap;
record 422 is a presence/heading rather than a second Kirzhnits lecture;
the omitted catalogue record 437 survives on pages 753--754 and is treated
as an explicitly attributed precursor to 438. Records 445, 451 and 457
contain 42 occurrences from 1978--79. Their requested identifiers
are retained, with the actual notebook dates explained. The date field
preserves the supplied catalogue date; it must be read with the provenance
notice, especially for Manin's 21 November 1977 and Fuchs's 3 April 1978
and 5 March 1979 lectures. P000705 retains an unresolved April/September
date conflict. One truncated attachment, P000722, was recovered
from MCCME and checked against its intact neighbors and the uploaded prefix.

\begin{center}
\begin{tabular}{@{}lrrrrrrr@{}}
\toprule
Class & A & B & C & D & E & F & Total\\
\midrule
Occurrences & 55 & 11 & 4 & 0 & 210 & 25 & 305\\
\bottomrule
\end{tabular}
\end{center}
Here A means near-contemporary publication by the speaker or Gelfand;
B, later development by others; C, a precisely delimited partial result;
D, an undeveloped identifiable problem; E, prior established material;
and F, insufficiently identifiable content. These are historical fates,
not correctness ratings. A published thread can contain a false displayed
formula, and an F entry is not evidence of an orphan research idea.

\paragraph{The main lines.}
The strongest connected sequences concern local algebra models for
Harish--Chandra modules, self-dual gauge fields and their finite-dimensional
data, and spectral methods for nonlinear evolution equations. The local
representation-theory announcements match
\citet[Theorems~1--2]{bernsteingelfandgelfand1978local1} and
\citet[Theorem~1]{bernsteingelfandgelfand1978local2}; their small-domain
and regular-central-character qualifications are essential. Manin's talk
precedes the receipt of \citet{drinfeldmanin1978selfdual} by six days.
The instanton construction does not classify every finite-action
Yang--Mills critical point. The current comparison includes
\citet[Theorems~1.1--1.2]{dayapremawaldron2026parabolic}, with their stated
energy and moduli-space scopes kept separate.

Krichever's rational KP theorem, received in September 1977 and published
in 1978, is a close match to the November announcement
\citep[Theorem~3]{krichever1978rational}. Its pole dynamics require a
consistent choice of time and Hamiltonian normalization. Dickey's Poisson
bracket, Faddeev's classical and quantum scattering discussion, and the
Calogero--Degasperis lectures connect infinite-dimensional equations with
finite spectral or particle data. BGG's exterior-algebra account of bundles,
Aomoto's hypergeometric and Ising lectures, Buchstaber's substitution groups,
and Fuchs's vector-field cohomology supply additional substantial lines.
Margulis, Kirzhnits and the counterfeit-coin discussion contain much
established material whose hypotheses still need careful reconstruction.

\paragraph{Questions and development.}
The 13 problem entries include questions answered before the talk,
near-contemporary solutions and surviving restricted questions; they are
not a count of currently open problems. Nonarithmetic lattices in complex
hyperbolic dimension at least four remain open in the located literature
\citep[Introduction]{falbelpasquinelli2023hybrid}; an October 2026 preprint
still reports examples only in dimensions two and three
\citep[Introduction]{yu2026mirror}. The broad Yang--Mills
critical-point classification and the exact scope of soliton-resolution
statements must not be replaced by narrower solved assertions. Buchstaber's
integer-substitution question omits coefficients and cochain topology.
\citet[Theorems~1.2 and~1.5]{babenkobogatyi2008substitutions} resolve its
real finite-jet version and compute a closed-commutator quotient. For an
explicit integral finite-jet interpretation, this chapter proves degree
one is \(\mathbb Z^2\), all degree-two torsion is \((\mathbb Z/2)^2\),
and all products of degree-one classes vanish. These are elementary
consequences and explicit calculations, with no priority claim and no
assertion about the full discrete cohomology ring.

\paragraph{Continuity and critical findings.}
The supplied earlier index connects the spectral-curve material to the
1975 rank-one correspondence and particle systems, gauge theory to the
1975 connection lectures, and the cohomology material to earlier
formal-vector-field discussions. Higher-rank commuting operators and
Hamiltonian-vector-field cohomology are related but different problems;
their old problem IDs are not silently reused. The earlier index is
incomplete as a year-by-year inventory, so a missing match establishes
only that no exact recurrence was identifiable there.

There are 111 critical entries, separating transcription,
translation, ambiguity, provenance and faithful mathematical error.
Among the independently checked cases, thirteen coins is correct for
identity-only detection; the KP coefficients change with convention;
the Calogero equilibrium linearization needs an identity matrix; and two
weight-25 spectral labels coincide, so gap partitions cannot count
distinct Laplace eigenvalues. Exact computations and their outputs are
provided. No D item is inferred from an unsuccessful search. The detailed
report explains the search limits, verification levels and requirements
for presenting these distinctions on the companion website.
\nocite{mccme2026gelfand}

\GStalk{419}{1977-09-05}{V. G. Drinfeld}{P000705-P000707}{Self-dual Yang-Mills fields and twistor bundles}
The report presents the four-dimensional self-duality equation, its conformal invariance, and the twistor fibration \(\mathbf{CP}^{3}\to\mathbf{HP}^{1}\simeq S^{4}\). It states the instanton--holomorphic-bundle correspondence in abbreviated form and sketches how a connection produces a holomorphic structure. A proposed spectral-determinant path integral is much less specified. The opening reference to noncommutative algebraic geometry and BGG is a possible seminar theme, not a developed construction. The record ends before Krichever begins in the lower part of P000707. The catalogue date is retained: the handwritten month on P000705 appears to be IV, so the date is a provenance conflict, not an emendation of a clearly written IX.

\begin{GSitem}{1977-419.1}{E}{A unitary connection and the self-duality equation}
\GSfield{Kind; source}{definition; \GSpage{P000705-P000706}}
\GSfield{The notes}{The notes take a smooth complex bundle \(\xi\) on \(S^{4}\), a connection \(\nabla\), and its curvature \(K\), an \(\operatorname{End}(\xi)\)-valued two-form. They require a Hermitian metric and unitarity and write \(*K=K\).}
\GSfield{Statement}{On an oriented Riemannian four-manifold, a unitary connection \(A\) on a Hermitian bundle has \(F_A\in\Omega^2(\mathfrak u(E))\). Self-duality means \(*F_A=F_A\); reversing orientation interchanges this equation with anti-self-duality. This is a gauge-invariant condition.}
\GSfield{Reconstruction}{The distinction between self-dual and anti-self-dual conventions, and the gauge-equivalence qualification, are made explicit here.}
\GSfield{Fate}{This is the setup of the instanton correspondence reported in \citet[abstract and pp.~117--124]{atiyahward1977instantons}; that paper appeared in June 1977, before the catalogue date of this seminar.}
\GSfield{Present status}{As of November 2026, this is an established expository item, subject to the hypotheses stated above.}
\GSfield{Assessment}{literal statement: correct; publication match: partial; verification: publisher metadata checked; development: annotated; confidence: high}
\end{GSitem}
\begin{GSitem}{1977-419.2}{E}{Conformal invariance of four-dimensional self-duality}
\GSfield{Kind; source}{theorem; \GSpage{P000705-P000706}}
\GSfield{The notes}{The star on two-forms in dimension four is said not to change when the metric is multiplied by a number. The next page describes a conformal metric as defined up to multiplication by a positive function.}
\GSfield{Statement}{For \(\widehat g=e^{2\varphi}g\) on an oriented \(d\)-manifold, the Hodge star on \(p\)-forms satisfies \(*_{\widehat g}=e^{(d-2p)\varphi}*_g\). Thus in dimension four the self-duality equation for two-forms depends only on the oriented conformal class.}
\GSfield{Reconstruction}{The general scaling identity explains why a positive function, not merely a constant, is allowed; it follows directly by comparing the induced inner product on forms with the volume form.}
\GSfield{Fate}{The conformal formulation is part of the twistor description in \citet[abstract]{atiyahward1977instantons}; the scaling calculation here is elementary reconstruction rather than a new result.}
\GSfield{Present status}{As of November 2026, this is an established expository item, subject to the hypotheses stated above.}
\GSfield{Assessment}{literal statement: correct; publication match: partial; verification: publisher metadata checked; development: annotated; confidence: high}
\end{GSitem}
\begin{GSitem}{1977-419.3}{F}{Gluons and the proposed physical role of instantons}
\GSfield{Kind; source}{remark; \GSpage{P000706}}
\GSfield{The notes}{The notes say that the quantum analogues are gluons and that they hold quarks together.}
\GSfield{Statement}{The identifiable content is a motivation from non-Abelian gauge fields and the strong interaction. It supplies no mathematical confinement criterion or theorem relating the displayed self-dual classical fields to a mass gap or a quark potential.}
\GSfield{Reconstruction}{An identification with the precise confinement question of earlier item 1975-377.10 would add a claim not made on this page.}
\GSfield{Fate}{No exact mathematical assertion about confinement can be identified here; the neighboring instanton paper is a construction in classical gauge theory, not a proof of this physical slogan.}
\GSfield{Present status}{As of November 2026, this occurrence remains a physical motivation without a separately recoverable mathematical proposition.}
\GSfield{Assessment}{literal statement: not determinable; publication match: none found; verification: unresolved lead; development: none; confidence: high}
\end{GSitem}
\begin{GSitem}{1977-419.4}{F}{A spectral-determinant path-integral prescription}
\GSfield{Kind; source}{method; \GSpage{P000706}}
\GSfield{The notes}{To a pair \((\xi,\nabla)\) one is to assign a linear operator, form a regularized determinant from its spectrum, and integrate over solutions; the notes add that this is the denominator.}
\GSfield{Statement}{No operator, regularization, integration measure, gauge quotient, insertion or normalization is specified. Consequently there is no determined integral, determinant identity or renormalization claim to formalize.}
\GSfield{Reconstruction}{It is tempting to supply a one-loop gauge-fixed determinant, but the notebook does not choose that construction.}
\GSfield{Fate}{Searches around the cited 1977 instanton/twistor correspondence identify the geometric report but do not identify this incomplete quantization prescription as a particular published formula.}
\GSfield{Present status}{As of November 2026, the operator and measure remain unidentified from this source; this is not an open-problem entry.}
\GSfield{Assessment}{literal statement: not determinable; publication match: none found; verification: unresolved lead; development: none; confidence: low}
\end{GSitem}
\begin{GSitem}{1977-419.5}{E}{The twistor fibration over the four-sphere}
\GSfield{Kind; source}{construction; \GSpage{P000706}}
\GSfield{The notes}{Using \(\mathbf C\subset\mathbf H\), the notes pass from complex lines in \(\mathbf H^2\) to quaternionic lines and obtain \(\mathbf{CP}^3\to\mathbf{HP}^1=S^4\), with two-dimensional fibers.}
\GSfield{Statement}{Fix right scalar multiplication. The map sends a complex line \(\ell\subset\mathbf H^2\) to its quaternionic span \(\ell\mathbf H\). It is a smooth fiber bundle \(\pi:\mathbf{CP}^{3}\to\mathbf{HP}^{1}\simeq S^4\) with fiber \(\mathbf{CP}^1\simeq S^2\). Projectivization uses \(\mathbf H^2\setminus\{0\}\), so the zero vector must be excluded from the quotient.}
\GSfield{Reconstruction}{The notes use left quaternionic projectivization; the equivalent right-scalar convention is chosen explicitly here. Deletion of zero is supplied. The fiber dimension is real dimension two, not complex dimension two.}
\GSfield{Fate}{This is the twistor space underlying \citet[abstract]{atiyahward1977instantons}, already published before the catalogue date.}
\GSfield{Present status}{As of November 2026, this is an established expository item, subject to the hypotheses stated above.}
\GSfield{Assessment}{literal statement: correct with missing hypotheses; publication match: partial; verification: publisher metadata checked; development: annotated; confidence: high}
\end{GSitem}
\begin{GSitem}{1977-419.6}{E}{Quaternionic conjugation, invariant lines and the Klein quadric}
\GSfield{Kind; source}{construction; \GSpage{P000706-P000707}}
\GSfield{The notes}{An antilinear coordinate map pairs the four complex coordinates, and the notes describe invariant projective lines, the Grassmannian of lines, a quadric, and a real form with real points \(S^4\).}
\GSfield{Statement}{Let \(J(z_1,z_2,z_3,z_4)=(-\overline z_2,\overline z_1,-\overline z_4,\overline z_3)\). Then \(J^2=-1\), so its projectivization is a fixed-point-free antiholomorphic involution of \(\mathbf{CP}^3\). A complex two-plane invariant under \(J\) is a quaternionic line. Its locus in \(\operatorname{Gr}(2,4)\), identified by the Pl\"ucker embedding with the Klein quadric in \(\mathbf{CP}^5\), is \(\mathbf{HP}^1\simeq S^4\).}
\GSfield{Reconstruction}{The real points are points of the induced real structure on the Grassmannian, not fixed points in \(\mathbf{CP}^3\). This distinction removes a possible ambiguity in the compressed wording.}
\GSfield{Fate}{The construction belongs to the geometric background of \citet[pp.~117--124]{atiyahward1977instantons}; only the publisher abstract and bibliographic record were available for that reference in this audit.}
\GSfield{Present status}{As of November 2026, this is an established expository item, subject to the hypotheses stated above.}
\GSfield{Assessment}{literal statement: correct with missing hypotheses; publication match: partial; verification: publisher metadata checked; development: annotated; confidence: high}
\end{GSitem}
\begin{GSitem}{1977-419.7}{E}{The instanton--holomorphic-bundle correspondence}
\GSfield{Kind; source}{theorem; \GSpage{P000706-P000707}}
\GSfield{The notes}{Under the label Atiyah\textquoteright s theorem, the notes assert a one-to-one correspondence between bundles with self-dual connection on \(S^4\) and algebraic bundles on \(\mathbf{CP}^3\) trivial on every fiber of \(\pi\).}
\GSfield{Statement}{A safe rank-two version is a correspondence between gauge-equivalence classes of smooth \(SU(2)\) instantons on \(S^4\) and isomorphism classes of holomorphic rank-two bundles on \(\mathbf{CP}^3\), with trivial determinant, holomorphically trivial restriction to each real twistor line, and the compatible quaternionic reality structure required for a positive unitary real form. Choose the orientation so the notebook\textquoteright s sign \(*F=F\) corresponds to this convention.}
\GSfield{Reconstruction}{The notebook allows arbitrary complex rank and omits equivalence relations and reality conditions. The statement here deliberately specifies the \(SU(2)\) case identified by the publication; bare complex bundles trivial on the real lines do not by themselves encode a unitary connection.}
\GSfield{Fate}{\citet[abstract]{atiyahward1977instantons} identifies precisely the \(SU(2)\) construction relating instantons to algebraic bundles. The paper was received 31 May and published June 1977. Class E uses the supplied 5 September 1977 catalogue date; its pre-talk chronology is conditional because P000705 instead appears to read 5/IV--77. The handwritten date has not been resolved. The full theorem was not accessible in the official copy consulted, so exact matching of every reconstructed reality hypothesis is not claimed.}
\GSfield{Present status}{As of November 2026, this is the established Ward correspondence. The historical report is recognizable, but its literal all-rank formulation omits indispensable structure.}
\GSfield{Assessment}{literal statement: correct with missing hypotheses; publication match: partial; verification: publisher metadata checked; development: annotated; confidence: medium}
\end{GSitem}
\begin{GSitem}{1977-419.8}{E}{A holomorphic structure on the pulled-back bundle}
\GSfield{Kind; source}{construction; \GSpage{P000707}}
\GSfield{The notes}{The direction from a connection to an algebraic bundle is indicated by \(\xi\mapsto f^*\xi\), followed by the notice that details will be given next time. Triviality is specified to mean analytic, or complex, triviality.}
\GSfield{Statement}{The \((0,1)\)-part of the pulled-back connection defines a Dolbeault operator when its square vanishes. On twistor space, the appropriate half-flatness condition on the original four-dimensional curvature makes this \((0,2)\)-curvature vanish. Along each fiber the pullback is holomorphically trivial.}
\GSfield{Reconstruction}{The Dolbeault operator, its integrability criterion and the construction of the holomorphic structure are supplied here; they are not a proof recorded on this page. The orientation convention is inherited from the preceding item.}
\GSfield{Fate}{This is the forward construction of the correspondence identified by \citet[abstract]{atiyahward1977instantons}; the consulted publisher record alone does not supply a full proof. Its class-E chronology is conditional on the supplied 5 September 1977 catalogue date: the paper was received 31 May and published in June, while the unresolved manuscript heading appears to read 5/IV--77.}
\GSfield{Present status}{As of November 2026, this is an established expository item, subject to the hypotheses stated above.}
\GSfield{Assessment}{literal statement: correct with missing hypotheses; publication match: partial; verification: publisher metadata checked; development: annotated; confidence: medium}
\end{GSitem}
\begin{GSitem}{1977-419.9}{E}{The four-sphere has no complex structure}
\GSfield{Kind; source}{claim; \GSpage{P000707}}
\GSfield{The notes}{The notes observe that \(S^4\) is not a complex manifold, in contrast to its twistor space.}
\GSfield{Statement}{The smooth four-sphere admits no almost-complex structure, hence no integrable one. Indeed an almost-complex tangent bundle would have \(c_1=0\) because \(H^2(S^4;\mathbf Z)=0\), \(\int c_2=\chi(S^4)=2\), and \(p_1=c_1^2-2c_2\). But \(\int p_1=3\sigma(S^4)=0\), a contradiction.}
\GSfield{Reconstruction}{The stronger almost-complex obstruction and the characteristic-class argument are explanatory reconstruction; no proof is recorded on the page.}
\GSfield{Fate}{This is a classical topological obstruction, not a new consequence of the reported instanton construction. The argument above makes the precise scope of the remark independently checkable.}
\GSfield{Present status}{As of November 2026, this is an established expository item, subject to the hypotheses stated above.}
\GSfield{Assessment}{literal statement: correct; publication match: none found; verification: indirect; development: annotated; confidence: high}
\end{GSitem}

\GStalk{420}{1977-09-05}{I. M. Krichever}{P000707-P000708}{Algebraic spectral curves and rational solutions of Lax equations}
Krichever describes zero-curvature equations for ordinary differential operators, spectral curves and their degenerations, and rational KdV solutions governed by pole dynamics. The suggestion that rationality should replace an extra commuting operator for every Lax equation has no specified class of equations. These pages continue the finite-gap line represented by earlier items 1975-374.2--1975-374.5 and 1975-379.1. The actual transition to Margulis is at P000709, not at the ends of the overlapping catalogue page ranges.

\begin{GSitem}{1977-420.1}{E}{Zero curvature for differential-operator evolution equations}
\GSfield{Kind; source}{construction; \GSpage{P000707}}
\GSfield{The notes}{The notes write \([\partial_y-L_1,\partial_t-L_2]=0\), where \(L_1,L_2\) are scalar or matrix ordinary differential operators in \(x\) of orders \(n,m\).}
\GSfield{Statement}{For differential operators with sufficiently differentiable coefficients in \((x,y,t)\), the compatibility of \(\psi_y=L_1\psi\), \(\psi_t=L_2\psi\) is \(L_{1,t}-L_{2,y}+[L_1,L_2]=0\). Equality means that every coefficient of the differential operator vanishes.}
\GSfield{Reconstruction}{The expanded sign convention and the regularity of coefficients are supplied. Matrix coefficients require attention to noncommutative multiplication but not a change in the compatibility identity.}
\GSfield{Fate}{The construction predates the talk: \citet[pp.~291--294]{krichever1976zakharov} was received in October 1975 and published in 1976; \citet[\S1]{krichever1977methods} gives the spectral-curve construction. Earlier occurrence: 1975-379.1.}
\GSfield{Present status}{As of November 2026, this is an established expository item, subject to the hypotheses stated above.}
\GSfield{Assessment}{literal statement: correct; publication match: exact; verification: read the result; development: annotated; confidence: high}
\end{GSitem}
\begin{GSitem}{1977-420.2}{E}{Quasiperiodic families, rational degenerations and soliton separatrices}
\GSfield{Kind; source}{claim; \GSpage{P000707}}
\GSfield{The notes}{The constructed solutions are described as generally conditionally periodic, with rational specializations; multisoliton KdV solutions form a separatrix-type family for \(u_t=6uu_x+u_{xxx}\).}
\GSfield{Statement}{Algebro-geometric scalar solutions associated with a smooth compact spectral curve are expressed through a theta function and linear motion on its Jacobian. For suitable real, nonsingular spectral data this gives smooth real quasiperiodic solutions. Singular degenerations can produce soliton or rational solutions, but arbitrary complex spectral data need not yield bounded real quasiperiodic solutions.}
\GSfield{Reconstruction}{The reality and regularity conditions are essential additions. ``Separatrix is retained as a qualitative description of a degeneration, not as a proved global phase-space classification.}
\GSfield{Fate}{\citet[\S1 and Appendix~2]{krichever1977methods} gives the theta-function and rational-degeneration constructions in a survey based on the earlier seminar program; the printed survey is 1977. The generic real-regular qualification is not stated in the notebook.}
\GSfield{Present status}{As of November 2026, this is an established expository item, subject to the hypotheses stated above.}
\GSfield{Assessment}{literal statement: correct with missing hypotheses; publication match: partial; verification: read the result; development: annotated; confidence: high}
\end{GSitem}
\begin{GSitem}{1977-420.3}{E}{Algebraic type and an additional commuting operator}
\GSfield{Kind; source}{definition; \GSpage{P000707}}
\GSfield{The notes}{For an equation with no \(y\), the notes attribute to Manin a condition called algebraic type: there is \(L_3\) commuting with \(L_1\), with \(L_3\notin\mathbf C[L_1,L_2]\).}
\GSfield{Statement}{A standard finite-gap hypothesis for a scalar ordinary differential operator \(L\) is the existence of a positive-order differential operator \(M\) commuting with \(L\), not generated by \(L\) alone; suitable coprime-order or rank conditions determine the associated spectral data. In a Lax evolution \(L_t=[P,L]\), the time generator \(P\) need not commute with \(L\). Therefore \(\mathbf C[L_1,L_2]\) in the notebook cannot automatically be treated as a commutative spectral ring.}
\GSfield{Reconstruction}{The standard commuting-ring formulation is a reconstruction, not a silent replacement of the literal condition. The notes do not specify which operators the notation \(\mathbf C[L_1,L_2]\) is intended to generate.}
\GSfield{Fate}{\citet[Introduction and \S1]{krichever1978commutative} treats commuting rings and clarifies rank restrictions; it also warns of errors in the earlier arbitrary-rank description. Earlier related occurrences are 1975-374.2--1975-374.5, not a new precise Manin problem.}
\GSfield{Present status}{As of November 2026, this is an established expository item, subject to the hypotheses stated above.}
\GSfield{Assessment}{literal statement: not determinable; publication match: partial; verification: read the result; development: annotated; confidence: medium}
\end{GSitem}
\begin{GSitem}{1977-420.4}{E}{Spectral curves and genus under degeneration}
\GSfield{Kind; source}{construction; \GSpage{P000707}}
\GSfield{The notes}{The notes attach a curve through relations involving coefficients of fractional powers, say that its genus is generically constant and drops at special parameters, and box genus zero beside the multisoliton degeneration.}
\GSfield{Statement}{A commuting scalar differential-operator ring determines an affine algebraic spectral curve through polynomial relations among simultaneous eigenvalues. In a degeneration of such curves, the geometric genus of the normalization may fall; a flat family preserves arithmetic genus. Rational or multisoliton limits can have a rational normalization while the singular curve retains positive arithmetic genus.}
\GSfield{Reconstruction}{The distinction between geometric and arithmetic genus is absent in the notes. No specific fractional-power coefficient relation is recorded, so the individual equation of the curve cannot be recovered.}
\GSfield{Fate}{The spectral-curve and rational-limit mechanisms are discussed in \citet[\S1 and Appendix~2]{krichever1977methods}. This is a recognizable construction with compressed local details, not a theorem that all notions of genus decrease.}
\GSfield{Present status}{As of November 2026, this is an established expository item, subject to the hypotheses stated above.}
\GSfield{Assessment}{literal statement: correct with missing hypotheses; publication match: partial; verification: read the result; development: annotated; confidence: high}
\end{GSitem}
\begin{GSitem}{1977-420.5}{E}{Polynomial dependence of commuting scalar differential operators}
\GSfield{Kind; source}{theorem; \GSpage{P000707}}
\GSfield{The notes}{The notes assert algebraic dependence of \(L_1,L_2\), with a tentative Jacobson attribution and a date in the 1930s.}
\GSfield{Statement}{For commuting ordinary scalar differential operators of positive order over a differential field of characteristic zero with constant field \(\mathbf C\), the usual Burchnall--Chaundy setting gives a nonzero polynomial \(R\in\mathbf C[X,Y]\) with \(R(L_1,L_2)=0\). This assertion is not true merely because \(L_1,L_2\) occur in a zero-curvature equation: they must commute in the stationary spectral problem.}
\GSfield{Reconstruction}{The commutativity hypothesis and the distinction from the evolution generator are supplied. The notebook\textquoteright s historical attribution is not used as proof of priority.}
\GSfield{Fate}{\citet[Introduction]{krichever1978commutative} explicitly traces the commuting-operator relation to Burchnall and Chaundy in the 1920s. Earlier occurrence: 1975-374.2. The exact role of the vaguely dated Jacobson reference was not identified.}
\GSfield{Present status}{As of November 2026, this is an established expository item, subject to the hypotheses stated above.}
\GSfield{Assessment}{literal statement: correct with missing hypotheses; publication match: partial; verification: read the result; development: annotated; confidence: high}
\end{GSitem}
\begin{GSitem}{1977-420.6}{E}{Rational KdV solutions and many-particle pole dynamics}
\GSfield{Kind; source}{theorem; \GSpage{P000708}}
\GSfield{The notes}{Airault, McKean and Moser are credited with replacing rational KdV evolution by an explicitly solvable system for the poles.}
\GSfield{Statement}{For a rational KdV solution decaying at spatial infinity, away from times with pole collisions the potential has double poles with prescribed leading coefficients, and substitution into KdV gives differential equations and algebraic constraints for their positions. With the notebook\textquoteright s sign convention \(u_t=6uu_x+u_{xxx}\), the principal part is \(-2(x-x_j)^{-2}\). The number of poles, counted in the generic configuration, is triangular in the standard rational KdV family.}
\GSfield{Reconstruction}{The decay condition, generic-time qualification and sign of the principal part are supplied; the notebook gives no explicit pole equations. The sign follows by canceling the order-five pole in the displayed KdV equation.}
\GSfield{Fate}{\citet[Appendix~2, pp.~205--207]{krichever1977methods} explicitly relates the rational KdV pole system to Airault--McKean--Moser and describes its constrained configurations. This exact section was read; the original Airault--McKean--Moser publisher copy was inaccessible in this run.}
\GSfield{Present status}{As of November 2026, this is an established expository item, subject to the hypotheses stated above.}
\GSfield{Assessment}{literal statement: correct with missing hypotheses; publication match: exact; verification: read the result; development: annotated; confidence: high}
\end{GSitem}
\begin{GSitem}{1977-420.7}{E}{Rationality implies a stationary higher KdV relation}
\GSfield{Kind; source}{claim; \GSpage{P000708}}
\GSfield{The notes}{The notes say that, for KdV, requiring rationality can replace the requirement of a third commuting operator.}
\GSfield{Statement}{Within the standard rational KdV class decaying at infinity, the associated Schr\"odinger operator belongs to the degenerate finite-gap family and admits an additional commuting differential operator; equivalently the potential satisfies a stationary equation in the higher KdV hierarchy. Rationality without the decay, hierarchy and operator-class qualifications is not the statement verified here.}
\GSfield{Reconstruction}{The precise class is reconstructed from the printed rational KdV result, and is stated separately from the notebook\textquoteright s unqualified sentence.}
\GSfield{Fate}{\citet[Appendix~2, pp.~205--207]{krichever1977methods} states the connection of rational solutions with stationary higher equations. This resolves the identifiable KdV claim; the universal extension in the next item is distinct.}
\GSfield{Present status}{As of November 2026, this is an established expository item, subject to the hypotheses stated above.}
\GSfield{Assessment}{literal statement: correct with missing hypotheses; publication match: exact; verification: read the result; development: annotated; confidence: high}
\end{GSitem}
\begin{GSitem}{1977-420.8}{F}{The proposed extension from KdV to every Lax equation}
\GSfield{Kind; source}{conjecture; \GSpage{P000708}}
\GSfield{The notes}{The final suggestion is that replacing the additional commuting operator by rationality is apparently possible for every Lax equation.}
\GSfield{Statement}{The notes do not specify the class of Lax pairs, coefficient fields, rationality variable, boundary behavior, permitted commuting operators or meaning of equivalence. There is therefore no determined universal assertion to compare with later theorems.}
\GSfield{Reconstruction}{Restricting the sentence to rational KdV, rational KP, or a particular Calogero--Moser correspondence would create a different and much narrower problem.}
\GSfield{Fate}{The KdV result was located and read in \citet[Appendix~2]{krichever1977methods}. Searches for the broader rationality/commuting-operator formulation did not supply quantifiers missing from this occurrence; a later theorem on one hierarchy would not settle ``every Lax equation.}
\GSfield{Present status}{As of November 2026, the intended universal class is unidentified. No new problem identifier is assigned to this underdetermined extrapolation.}
\GSfield{Assessment}{literal statement: not determinable; publication match: thematic only; verification: unresolved lead; development: none; confidence: low}
\end{GSitem}

\GStalk{421}{1977-09-12}{G. A. Margulis}{P000709-P000711}{Arithmeticity and superrigidity of higher-rank lattices}
The talk defines lattices, arithmetic subgroups and irreducibility, states higher-real-rank arithmeticity and superrigidity, and sketches the passage from local-field representations to integral structure. It contrasts the theorem with the rank-one situation, then incompletely understood. The modern surviving question in complex hyperbolic dimension at least four is separated from the broad historical formulation. Local rigidity was already discussed in items 1964-8.1--1964-8.2 and 1966-93.1; those are related occurrences, not an earlier occurrence of the precise remaining rank-one problem.

\begin{GSitem}{1977-421.1}{E}{Lattices in semisimple Lie groups}
\GSfield{Kind; source}{definition; \GSpage{P000709}}
\GSfield{The notes}{A discrete subgroup \(\Gamma\subset G\), with \(G\) connected semisimple and for simplicity centerless, is called a lattice when \(\nu(G/\Gamma)<\infty\), where \(\nu\) is induced by Haar measure.}
\GSfield{Statement}{A lattice in a locally compact unimodular group is a discrete subgroup whose quotient carries finite invariant measure. In a connected semisimple Lie group this definition applies because the group is unimodular. Compactness of the quotient is a stronger condition, called cocompactness or uniformity.}
\GSfield{Fate}{The definition is the one used in \citet[\S1.1]{tits1977margulis}, a February 1976 exposition revised in April 1976 and printed in 1977.}
\GSfield{Present status}{As of November 2026, this is an established expository item, subject to the hypotheses stated above.}
\GSfield{Assessment}{literal statement: correct; publication match: exact; verification: read the result; development: annotated; confidence: high}
\end{GSitem}
\begin{GSitem}{1977-421.2}{E}{Algebraic groups over a field and integral points}
\GSfield{Kind; source}{definition; \GSpage{P000709}}
\GSfield{The notes}{The notes define an algebraic subgroup of \(GL(n,\mathbf C)\), a \(k\)-group by equations over \(k\), and \(G(k)\). They warn that integral matrices form a group when contained in \(SL(n,\mathbf Z)\).}
\GSfield{Statement}{For a linear algebraic group \(\mathbf H\) defined over \(\mathbf Q\) and a fixed faithful rational representation, set \(\mathbf H(\mathbf Z)=\mathbf H(\mathbf Q)\cap GL_N(\mathbf Z)\). Here \(GL_N(\mathbf Z)\) means matrices whose inverses also have integer entries, equivalently determinant \(\pm1\). Merely intersecting \(GL_N(\mathbf C)\) with integer-entry matrices is insufficient.}
\GSfield{Reconstruction}{Distinct symbols \(\mathbf H\) and \(G\) avoid the notebook\textquoteright s reuse of one letter for an algebraic group and a real Lie group.}
\GSfield{Fate}{\citet[Introduction and \S2]{borelharishchandra1962arithmetic} gives this rational/integral setup; the opened Russian translation is from 1964, of the 1962 original.}
\GSfield{Present status}{As of November 2026, this is an established expository item, subject to the hypotheses stated above.}
\GSfield{Assessment}{literal statement: correct; publication match: exact; verification: read the result; development: annotated; confidence: high}
\end{GSitem}
\begin{GSitem}{1977-421.3}{E}{Arithmetic subgroups and commensurability}
\GSfield{Kind; source}{definition; \GSpage{P000709}}
\GSfield{The notes}{A subgroup commensurable with \(G(\mathbf Z)\) is arithmetic. For a general real Lie group, the notes use an epimorphism from the real points of a \(\mathbf Q\)-group with compact kernel.}
\GSfield{Statement}{Two subgroups are commensurable if their intersection has finite index in each. For the semisimple Lie groups in this talk, arithmeticity means commensurability with the image of integral points of a suitable semisimple \(\mathbf Q\)-group under a continuous surjection on real points with compact kernel, allowing the standard finite-center identifications.}
\GSfield{Reconstruction}{Semisimplicity and the finite-center convention are kept explicit; the following finite-volume conclusion is not true for arbitrary algebraic groups with rational characters.}
\GSfield{Fate}{Compare \citet[\S2.1]{tits1977margulis} for the restriction-of-scalars and finite-center formulation. Arithmeticity is independent of the integral model up to commensurability.}
\GSfield{Present status}{As of November 2026, this is an established expository item, subject to the hypotheses stated above.}
\GSfield{Assessment}{literal statement: correct; publication match: exact; verification: read the result; development: annotated; confidence: high}
\end{GSitem}
\begin{GSitem}{1977-421.4}{E}{Finite covolume of arithmetic subgroups}
\GSfield{Kind; source}{theorem; \GSpage{P000710}}
\GSfield{The notes}{The notes cite Borel--Harish-Chandra for the statement that an arithmetic subgroup has finite covolume.}
\GSfield{Statement}{If \(\mathbf H\) is a semisimple algebraic group defined over \(\mathbf Q\), then \(\mathbf H(\mathbf Z)\) is a lattice in \(\mathbf H(\mathbf R)\). More generally, the connected rational-character obstruction must vanish; the semisimple case automatically satisfies that condition. Commensurable subgroups and compact-kernel images preserve the relevant finite-volume property.}
\GSfield{Reconstruction}{The general algebraic-group qualification prevents reading this as a theorem for \(\mathbf G_m(\mathbf Z)\subset\mathbf G_m(\mathbf R)\), which has infinite covolume.}
\GSfield{Fate}{\citet[Thms.~7.8 and 9.4]{borelharishchandra1962arithmetic} proves precisely these statements. Both theorem statements and the semisimple proof were read in the Russian translation.}
\GSfield{Present status}{As of November 2026, this is an established expository item, subject to the hypotheses stated above.}
\GSfield{Assessment}{literal statement: correct with missing hypotheses; publication match: exact; verification: read the result; development: annotated; confidence: high}
\end{GSitem}
\begin{GSitem}{1977-421.5}{E}{The lattice SL(n,Z) and the modular-area exercise}
\GSfield{Kind; source}{theorem; \GSpage{P000710}}
\GSfield{The notes}{\(SL(n,\mathbf Z)\) is given as an example. The audience is asked to check \(n=2\) by calculating a hyperbolic triangle\textquoteright s area.}
\GSfield{Statement}{\(SL_n(\mathbf Z)\) is a lattice in \(SL_n(\mathbf R)\). For \(PSL_2(\mathbf Z)\) acting on the upper half-plane with curvature \(-1\), the standard fundamental region \(\{|\operatorname{Re}z|\leq\tfrac12,\ |z|\geq1\}\) has area \(\pi/3\). Its finite area, together with the compact stabilizer \(SO(2)\), gives finite Haar covolume.}
\GSfield{Reconstruction}{The region and normalization are supplied. It is a union of two reflection triangles of area \(\pi/6\), so the triangle drawn in the notes need not itself be a fundamental domain for the orientation-preserving group.}
\GSfield{Fate}{This is a classical special case of \citet[Thm.~7.8]{borelharishchandra1962arithmetic}. The exercise is solved by \(\int_{-1/2}^{1/2}(1-x^2)^{-1/2}\,dx=\pi/3\); no modern research problem remains.}
\GSfield{Present status}{As of November 2026, this is an established expository item, subject to the hypotheses stated above.}
\GSfield{Assessment}{literal statement: correct; publication match: exact; verification: read the result; development: annotated; confidence: high}
\end{GSitem}
\begin{GSitem}{1977-421.6}{E}{Irreducibility of a lattice}
\GSfield{Kind; source}{definition; \GSpage{P000710}}
\GSfield{The notes}{A lattice is reducible if, for a nontrivial decomposition \(G=G_1\times G_2\), the product \((G_1\cap\Gamma)\times(G_2\cap\Gamma)\) has finite index in \(\Gamma\).}
\GSfield{Statement}{For a semisimple group with finite center and no compact factors, irreducibility means that no decomposition into two positive-dimensional normal factors splits the lattice into a product up to finite index. Passing to the adjoint group removes finite central overlaps. Equivalently, projections to proper factor quotients are dense in the usual irreducible-lattice setting.}
\GSfield{Reconstruction}{The finite-center convention distinguishes a direct product from an almost-direct product.}
\GSfield{Fate}{The finite-index definition appears in \citet[\S1.1]{tits1977margulis}; the notebook uses its centerless version.}
\GSfield{Present status}{As of November 2026, this is an established expository item, subject to the hypotheses stated above.}
\GSfield{Assessment}{literal statement: correct; publication match: exact; verification: read the result; development: annotated; confidence: high}
\end{GSitem}
\begin{GSitem}{1977-421.7}{E}{The overly broad arithmeticity conjecture outside PSL(2,R)}
\GSfield{Kind; source}{conjecture; \GSpage{P000710}}
\GSfield{The notes}{A historical conjecture attributed to Selberg says that every irreducible lattice is arithmetic except in \(PSL_2(\mathbf R)\). The notes immediately mention counterexamples of Makarov and Vinberg in real hyperbolic groups.}
\GSfield{Statement}{The literal all-groups assertion is false: rank-one real hyperbolic groups beyond \(PSL_2(\mathbf R)\) admit nonarithmetic lattices. The valid higher-rank theorem requires real rank at least two, irreducibility, finite center and no compact factors.}
\GSfield{Reconstruction}{The notebook itself retracts the broad formulation through its counterexamples; this is not classified as a conjecture still open in 1977.}
\GSfield{Fate}{\citet[\S2.2]{tits1977margulis} records the Makarov--Vinberg examples in dimensions three and four and states the correctly restricted Selberg--Piatetski-Shapiro conjecture. The notebook\textquoteright s additional dimension-five attribution was not checked in an original construction.}
\GSfield{Present status}{As of November 2026, the displayed broad conjecture remains false; it is historical context for higher-rank arithmeticity and the separately indexed rank-one remainder.}
\GSfield{Assessment}{literal statement: false as written; publication match: partial; verification: read the result; development: annotated; confidence: medium}
\end{GSitem}
\begin{GSitem}{1977-421.8}{E}{Margulis arithmeticity in higher real rank}
\GSfield{Kind; source}{theorem; \GSpage{P000710}}
\GSfield{The notes}{The answer is that every irreducible lattice is arithmetic when \(\operatorname{rank}G\geq2\). The notes point to an appendix to Raghunathan\textquoteright s book for a proof.}
\GSfield{Statement}{Let \(G\) be a connected semisimple real Lie group with finite center, without compact factors, and with \(\operatorname{rank}_{\mathbf R}G\geq2\). Then every irreducible lattice in \(G\) is arithmetic.}
\GSfield{Reconstruction}{Real rank and the structural hypotheses already present or inserted elsewhere on the pages are brought together. The unspecified edition of Raghunathan\textquoteright s book is not treated as an independently verified bibliographic citation.}
\GSfield{Fate}{\citet[Thm.~1]{tits1977margulis} already reported the full theorem in February 1976. The nonuniform real case had been published by \citet[main arithmeticity result]{margulis1974arithmetic} in 1974; that paper alone is not evidence for the full cocompact case.}
\GSfield{Present status}{As of November 2026, this is Margulis\textquoteright s established arithmeticity theorem. Class E records circulation before this talk, even though some complete expositions appeared later.}
\GSfield{Assessment}{literal statement: correct with missing hypotheses; publication match: exact; verification: read the result; development: annotated; confidence: high}
\end{GSitem}
\begin{GSitem}{1977-421.9}{E}{Rigidity and the special role of surface lattices}
\GSfield{Kind; source}{claim; \GSpage{P000710}}
\GSfield{The notes}{The notes say that parameters occur only for \(SL_2(\mathbf R)\), while in the other cases there is rigidity, attributed to Selberg and Weil.}
\GSfield{Statement}{For cocompact lattices in semisimple real Lie groups, the familiar local-rigidity theorem excludes an \(SL_2(\mathbf R)\) factor in the relevant irreducible case. Finite-volume strong rigidity and local rigidity of arbitrary representations are different statements: nonuniform hyperbolic three-manifold groups can deform in \(PSL_2(\mathbf C)\) through representations that cease to be lattices.}
\GSfield{Reconstruction}{The notebook does not distinguish deformations through lattices from deformations in the full representation variety. Its slogan cannot be used as the latter unrestricted assertion.}
\GSfield{Fate}{Local rigidity recurs from earlier items 1964-8.1--1964-8.2 and 1966-93.1. \citet[\S2.2]{tits1977margulis} uses the surface-moduli contrast for arithmeticity, but does not validate every possible reading of the slogan.}
\GSfield{Present status}{As of November 2026, this is an established expository item, subject to the hypotheses stated above.}
\GSfield{Assessment}{literal statement: correct with missing hypotheses; publication match: thematic only; verification: indirect; development: annotated; confidence: medium}
\end{GSitem}
\begin{GSitem}{1977-421.10}{E}{Algebraic coefficients after conjugation}
\GSfield{Kind; source}{claim; \GSpage{P000710}}
\GSfield{The notes}{As a consequence of rigidity, every lattice in \(SL(n,\mathbf R)\) is said to be conjugate to one with algebraic coefficients; no lower bound on \(n\) is written.}
\GSfield{Statement}{For \(n\geq3\), a lattice in \(SL_n(\mathbf R)\) can be conjugated to a group of matrices over a number field contained in \(\mathbf R\). The claim fails without qualification for \(n=2\): continuously varying surface lattices can have transcendental conjugacy-invariant traces.}
\GSfield{Reconstruction}{The missing hypothesis \(n\geq3\) is substantive, not a change to the scan. The number-field statement uses finite generation in addition to algebraicity of individual entries.}
\GSfield{Fate}{The algebraic-definition step is part of the deduction of arithmeticity discussed in \citet[\S5]{tits1977margulis}. Its unqualified formulation in the notebook must be restricted before use.}
\GSfield{Present status}{As of November 2026, this is an established expository item, subject to the hypotheses stated above.}
\GSfield{Assessment}{literal statement: false as written; publication match: partial; verification: read the result; development: annotated; confidence: high}
\end{GSitem}
\begin{GSitem}{1977-421.11}{E}{The archimedean superrigidity assertion}
\GSfield{Kind; source}{theorem; \GSpage{P000711}}
\GSfield{The notes}{The main lemma assumes real rank at least two, an irreducible lattice and no compact factors, all explicitly inserted in the manuscript. If \(T:\Gamma\to GL(n,k)\) has connected simple Zariski closure and non-relatively-compact image, then for \(k=\mathbf R,\mathbf C\) it extends to \(G\).}
\GSfield{Statement}{In the standard superrigidity formulation, an unbounded Zariski-dense homomorphism of an irreducible higher-real-rank lattice into a connected adjoint simple group over an archimedean local field extends continuously to the ambient real Lie group, with the standard finite-cover and finite-index qualifications when working with nonadjoint groups or finite central twists.}
\GSfield{Reconstruction}{The target-adjoint or finite-central qualifications are added; the stronger unqualified extension to the precise matrix representation need not follow without them.}
\GSfield{Fate}{\citet[Thm.~2, \S3]{tits1977margulis} reports superrigidity before the talk. Its statement, rather than a metadata-only attribution, was read.}
\GSfield{Present status}{As of November 2026, this is an established expository item, subject to the hypotheses stated above.}
\GSfield{Assessment}{literal statement: correct with missing hypotheses; publication match: exact; verification: read the result; development: annotated; confidence: high}
\end{GSitem}
\begin{GSitem}{1977-421.12}{E}{The nonarchimedean compactness assertion}
\GSfield{Kind; source}{theorem; \GSpage{P000711}}
\GSfield{The notes}{For \(k=\mathbf Q_p\), the same main lemma says that connected simple Zariski closure and noncompact image are incompatible. The opening definition of local field also includes finite extensions of \(\mathbf Q_p\).}
\GSfield{Statement}{For a characteristic-zero nonarchimedean local field \(k\), a Zariski-dense homomorphism of the same irreducible higher-real-rank lattice into a connected adjoint simple \(k\)-group has relatively compact image. Noncompactness would force a continuous extension from connected \(G\) to a totally disconnected target, which is impossible.}
\GSfield{Reconstruction}{The result applies to finite extensions of \(\mathbf Q_p\) as well, and the central/finite-index conventions are the same as in the preceding item.}
\GSfield{Fate}{This is the nonarchimedean consequence of \citet[Thm.~2 and \S5]{tits1977margulis}; it was already part of the 1976 report.}
\GSfield{Present status}{As of November 2026, this is an established expository item, subject to the hypotheses stated above.}
\GSfield{Assessment}{literal statement: correct with missing hypotheses; publication match: exact; verification: read the result; development: annotated; confidence: high}
\end{GSitem}
\begin{GSitem}{1977-421.13}{E}{From local compactness to arithmeticity}
\GSfield{Kind; source}{method; \GSpage{P000711}}
\GSfield{The notes}{Put \(\Gamma\) in matrices over a number field, list the finitely many primes in denominators of generators, and use relative \(p\)-adic compactness to bound all denominators. The notes conclude arithmeticity.}
\GSfield{Statement}{After obtaining a number field of definition, one must use its completions \(K_v\), not only \(\mathbf Q_p\), and control all archimedean embeddings. Boundedness at finite places gives preservation of suitable local lattices; restriction of scalars combines the embeddings into a rational algebraic group, whose unwanted archimedean factors are compact. This yields an ambient arithmetic lattice, and finite covolume then gives the required finite index.}
\GSfield{Reconstruction}{The displayed denominator argument is one step of the proof. Integral containment alone, without the other archimedean factors and finite-index step, is not arithmeticity.}
\GSfield{Fate}{\citet[\S5]{tits1977margulis} explains the deduction. The source sketch is recognizable but omits exactly the global-field steps listed here.}
\GSfield{Present status}{As of November 2026, this is an established expository item, subject to the hypotheses stated above.}
\GSfield{Assessment}{literal statement: correct with missing hypotheses; publication match: partial; verification: read the result; development: annotated; confidence: high}
\end{GSitem}
\begin{GSitem}{1977-421.14}{C}{The arithmeticity problem in the remaining rank-one groups}
\GSfield{Kind; source}{question; \GSpage{P000710}}
\GSfield{The notes}{After real hyperbolic counterexamples and the higher-rank theorem, the notes say that in the other rank-one groups the answer is unknown.}
\GSfield{Statement}{The historical question is which real-rank-one simple groups admit nonarithmetic lattices. A precise surviving subquestion is: for which integers \(n\geq4\), if any, does \(PU(n,1)\) contain a nonarithmetic lattice?}
\GSfield{Reconstruction}{The \(PU(n,1)\), \(n\geq4\), formulation is the modern remainder of the broader question, not a verbatim question confined to complex hyperbolic groups in the notebook.}
\GSfield{Fate}{Real hyperbolic groups in every dimension admit nonarithmetic lattices; quaternionic hyperbolic groups of quaternionic dimension at least two and the Cayley hyperbolic group have arithmetic lattices only. \citet[\S2.1]{parker2009lattices} surveys these resolutions. \citet[Introduction]{falbelpasquinelli2023hybrid} explicitly retains the complex-hyperbolic question for \(n\geq4\), while discussing constructions in dimensions two and three.}
\GSfield{Present status}{As of 10 November 2026, the search found no resolution of existence of nonarithmetic lattices in \(PU(n,1)\) for any \(n\geq4\). The introduction of \citet{yu2026mirror}, submitted 30 October 2026, still describes known examples as occurring only in dimensions two and three; its new example is in \(PU(2,1)\). This is a current research statement plus a dated search, not proof that an unlocated solution cannot exist.}
\GSfield{Assessment}{literal statement: correct; publication match: partial; verification: read the result; development: annotated; confidence: medium; problem \hyperref[prob:GP-1977-01]{GP-1977-01}}
\GSfield{Worth developing}{This is a precise and important surviving problem, but no bounded contribution is attempted here: producing a new lattice or an arithmeticity obstruction requires substantially more than the notebook supplies.}
\end{GSitem}

\GStalk{422}{1977-09-19}{D. A. Kirzhnits}{P000712}{Announcement of Kirzhnits’s presence}
The overlapping catalogue range contains the bracketed notice that the physicist Kirzhnits is present. It does not identify a second independent mathematical exposition. Drinfeld\textquoteright s coordinate formulas on P000712 are assigned to record 423; Kirzhnits\textquoteright s subsequent physics discussion to record 425, with Gelfand\textquoteright s explicit coin-weighing intervention assigned to record 424. Thus this heading has no duplicated occurrence items.


\GStalk{423}{1977-09-19}{V. G. Drinfeld}{P000712}{Coordinate curvature and the Yang-Mills equations}
A short continuation writes curvature, self-duality and the Yang--Mills equation in matrix coordinates. The formulas are the local version of record 419. The implication from self-duality to Yang--Mills is reconstructed using the Bianchi identity, rather than presented as a proof found in the notebook.

\begin{GSitem}{1977-423.1}{E}{Curvature in a local gauge}
\GSfield{Kind; source}{definition; \GSpage{P000712}}
\GSfield{The notes}{For matrix functions \(A_\mu\) of four variables the notes write \(F_{\mu\nu}=\partial_\mu A_\nu-\partial_\nu A_\mu+[A_\mu,A_\nu]\).}
\GSfield{Statement}{In a local trivialization of a smooth vector bundle, write \(D_\mu=\partial_\mu+A_\mu\). Then \([D_\mu,D_\nu]\) acts by multiplication by \(F_{\mu\nu}\), the displayed curvature matrix. For a unitary connection the matrices \(A_\mu\) are skew-Hermitian in a unitary frame.}
\GSfield{Fate}{This is the coordinate form of the connection in record 419 and the setup identified in \citet[abstract]{atiyahward1977instantons}.}
\GSfield{Present status}{As of November 2026, this is an established expository item, subject to the hypotheses stated above.}
\GSfield{Assessment}{literal statement: correct; publication match: partial; verification: publisher metadata checked; development: annotated; confidence: high}
\end{GSitem}
\begin{GSitem}{1977-423.2}{E}{The coordinate self-duality and Yang-Mills equations}
\GSfield{Kind; source}{definition; \GSpage{P000712}}
\GSfield{The notes}{The notes display \(F_{ij}=\tfrac12\varepsilon_{ijk\ell}F_{k\ell}\) and \([\partial_\mu+A_\mu,F_{\mu\nu}]=0\).}
\GSfield{Statement}{In oriented orthonormal Euclidean coordinates, self-duality is \(F_{ij}=\frac12\sum_{k,\ell}\varepsilon_{ijk\ell}F_{k\ell}\). The Yang--Mills equation is \(\sum_\mu D_\mu F_{\mu\nu}=0\), where \(D_\mu F=\partial_\mu F+[A_\mu,F]\). On a curved four-manifold these are \(*F=F\) and \(d_A^*F=0\), with the metric divergence understood.}
\GSfield{Reconstruction}{Summation over repeated indices and the use of a local orthonormal frame on a curved manifold are supplied. A plain partial-derivative divergence cannot be used in arbitrary coordinates on \(S^4\).}
\GSfield{Fate}{These are the local equations underlying the earlier instanton report; \citet[abstract]{atiyahward1977instantons} confirms the subject but is not cited as an exact source for this coordinate transcription.}
\GSfield{Present status}{As of November 2026, this is an established expository item, subject to the hypotheses stated above.}
\GSfield{Assessment}{literal statement: correct with missing hypotheses; publication match: partial; verification: publisher metadata checked; development: annotated; confidence: high}
\end{GSitem}
\begin{GSitem}{1977-423.3}{E}{Self-duality implies the Yang-Mills equation}
\GSfield{Kind; source}{theorem; \GSpage{P000712}}
\GSfield{The notes}{The self-duality and Yang--Mills equations are displayed in succession; their implication is not proved on the page.}
\GSfield{Statement}{For a smooth connection on an oriented Riemannian four-manifold, \(*F_A=\pm F_A\) implies \(d_A^*F_A=0\). Indeed the Bianchi identity gives \(d_AF_A=0\), and \(d_A^*F_A=-*d_A*F_A=\mp*d_AF_A=0\).}
\GSfield{Reconstruction}{The implication and its one-line proof are explanatory reconstruction of the juxtaposed equations, not an additional attributed proof by Drinfeld.}
\GSfield{Fate}{This standard curvature identity explains why the self-dual fields in \citet[abstract]{atiyahward1977instantons} solve Yang--Mills; no separate research claim is extracted.}
\GSfield{Present status}{As of November 2026, this is an established expository item, subject to the hypotheses stated above.}
\GSfield{Assessment}{literal statement: correct; publication match: partial; verification: indirect; development: annotated; confidence: high}
\end{GSitem}

\GStalk{424}{1977-09-19}{I. M. Gelfand}{P000713-P000714}{Counterfeit coins and entropy}
Gelfand interrupts the thermodynamic discussion with ternary weighing puzzles: 27 coins when the counterfeit is known to be lighter, and 13 when its sign is unknown. The number 13 is consistent with identifying the coin alone, leaving the sign of a never-weighed coin unresolved; requiring its sign as well gives the familiar 12-coin version. The ensuing entropy formulas include a missing minus sign in the binary case. Only the explicit intervention and the immediate combinatorial formulas are assigned here; the physical entropy discussion remains in record 425.

\begin{GSitem}{1977-424.1}{E}{A known lighter counterfeit among 27 coins}
\GSfield{Kind; source}{problem; \GSpage{P000713}}
\GSfield{The notes}{Gelfand asks for the counterfeit among 27 coins in three weighings when it is known to be lighter.}
\GSfield{Statement}{Assume exactly one coin is lighter and all others have the same mass. An ideal balance has three outcomes. Comparing equal groups of nine locates the group containing the counterfeit; repeat with groups of three and then single coins. In general \(3^r\) coins can be handled in \(r\) weighings, and the ternary decision-tree bound proves optimality.}
\GSfield{Reconstruction}{The assumptions on exactly one counterfeit and an ideal balance are stated explicitly; the recursive construction and counting lower bound are supplied.}
\GSfield{Fate}{This is a classical weighing puzzle, part of the subject of \citet[pp.~231--234]{dyson1946pennies}. Only Dyson\textquoteright s publication metadata was accessible; the elementary argument above independently checks this occurrence.}
\GSfield{Present status}{As of November 2026, this is an established expository item, subject to the hypotheses stated above.}
\GSfield{Assessment}{literal statement: correct; publication match: thematic only; verification: publisher metadata checked; development: annotated; confidence: high}
\end{GSitem}
\begin{GSitem}{1977-424.2}{E}{Thirteen coins when only the counterfeit identity is required}
\GSfield{Kind; source}{problem; \GSpage{P000713-P000714}}
\GSfield{The notes}{The unknown-sign version is answered by 13 coins in three weighings, then by \((3^N-1)/2\) coins in \(N\) weighings. The proposed first weighing is four against four, leaving five.}
\GSfield{Statement}{With exactly one counterfeit, no externally supplied genuine coin, equal coin counts on the pans, and only its identity required, \((3^N-1)/2\) coins are achievable and optimal for \(N\geq1\). If both identity and heavier/lighter sign are required, the capacity is \((3^N-3)/2\). For \(N=3\), the respective answers are 13 and 12.}
\GSfield{Reconstruction}{The Russian says to find the coin and does not require its sign. A balanced ternary code with 12 nonzero columns and a thirteenth zero column verifies the literal 13-coin reading. All-balance outcomes identify coin 13 but leave its sign undetermined. The general construction and bound are explained in the supplementary annotation, and the supplied code enumerates all signed hypotheses for \(N=1,\ldots,6\).}
\GSfield{Fate}{The classical literature includes \citet[pp.~231--234]{dyson1946pennies}; its full proof was not available in the publisher copy. This chapter\textquoteright s code and elementary reconstruction check the exact convention independently. No correction from 13 to 12 is warranted without changing the required output.}
\GSfield{Present status}{As of November 2026, this is an established expository item, subject to the hypotheses stated above.}
\GSfield{Assessment}{literal statement: correct with missing hypotheses; publication match: partial; verification: publisher metadata checked; development: annotated; confidence: high}
\end{GSitem}
\begin{GSitem}{1977-424.3}{E}{Choosing informative three-outcome weighings}
\GSfield{Kind; source}{method; \GSpage{P000714}}
\GSfield{The notes}{The algorithm is to divide the coins so that the entropy of the piles is equal, dividing the entropy into three equal parts and avoiding exhaustive search.}
\GSfield{Statement}{For a deterministic observation \(Y\) with three possible outcomes, information gained about the hidden hypothesis \(X\) is \(I(X;Y)=H(Y)\leq\log3\). Equality requires three equally likely outcomes. Thus one seeks a physically feasible weighing that balances posterior probability mass, not three literally equal entropy values or necessarily equal numbers of coins.}
\GSfield{Reconstruction}{For 13 equally likely counterfeit identities with equally likely signs, the initial four-versus-four weighing has signed branch counts \(8,8,10\), not three equal counts. Entropy maximization is a guiding criterion; it does not by itself prove a complete weighing algorithm exists.}
\GSfield{Fate}{\citet[\S6, properties 2, 3 and 5]{shannon1948communication} supplies the entropy bound and conditional-entropy identity. The application to the constrained weighing experiment is reconstruction.}
\GSfield{Present status}{As of November 2026, this is an established expository item, subject to the hypotheses stated above.}
\GSfield{Assessment}{literal statement: correct with missing hypotheses; publication match: partial; verification: read the result; development: annotated; confidence: high}
\end{GSitem}
\begin{GSitem}{1977-424.4}{E}{Shannon entropy and the binary sign error}
\GSfield{Kind; source}{definition; \GSpage{P000714}}
\GSfield{The notes}{The binary formula is written \(S=p\ln p+q\ln q\), while the following general formula is \(-\sum_i p_i\ln p_i\). For equiprobable outcomes the notes obtain \(S=\ln N\).}
\GSfield{Statement}{For a finite probability distribution, \(H(p)=-\sum_i p_i\log p_i\), with \(0\log0=0\). Thus \(H(p,1-p)=-p\log p-(1-p)\log(1-p)\), and \(H\) is maximized at the uniform distribution, where \(H=\log N\). Natural logarithms give nats.}
\GSfield{Reconstruction}{The minus sign in the binary expression is absent in the scan as well as the transcription. It is a faithful notebook error, not an English translation error.}
\GSfield{Fate}{\citet[\S6, Thm.~2 and property 2]{shannon1948communication} already gives the characterization and maximum. The formula and sign were checked in the scanned source and against the printed entropy properties.}
\GSfield{Present status}{As of November 2026, this is an established expository item, subject to the hypotheses stated above.}
\GSfield{Assessment}{literal statement: false as written; publication match: exact; verification: read the result; development: annotated; confidence: high}
\end{GSitem}
\begin{GSitem}{1977-424.5}{E}{Additivity and conditional information}
\GSfield{Kind; source}{claim; \GSpage{P000714}}
\GSfield{The notes}{The notes appeal to additivity and identify information with a change of entropy, writing \(\Delta I=-\ln(N_f/N_i)\).}
\GSfield{Statement}{Independent finite random variables satisfy \(H(X,Y)=H(X)+H(Y)\); in general \(H(X,Y)=H(X)+H(Y\mid X)\). The expected information supplied by an observation is \(I(X;Y)=H(X)-H(X\mid Y)\). The displayed count formula equals this entropy reduction on a branch only when the prior and posterior possibilities are uniform over sets of sizes \(N_i,N_f\).}
\GSfield{Reconstruction}{Independence, conditional probabilities and the uniformity qualification are supplied. Information gain is not an unconditional identity with an arbitrarily signed physical entropy change.}
\GSfield{Fate}{\citet[\S6, properties 3 and 5--6]{shannon1948communication} gives the identities. Earlier seminar entropy occurrences include 1964-10.1 and 1966-107.4, with different physical emphases.}
\GSfield{Present status}{As of November 2026, this is an established expository item, subject to the hypotheses stated above.}
\GSfield{Assessment}{literal statement: correct with missing hypotheses; publication match: exact; verification: read the result; development: annotated; confidence: high}
\end{GSitem}

\GStalk{425}{1977-09-19}{D. A. Kirzhnits}{P000713-P000716}{Entropy, black-hole thermodynamics and accelerated observers}
The physicist connects microscopic dynamics with entropy, black-hole horizons, the area law, Hawking radiation and the thermal response of accelerated observers. The talk also mentions temperature-dependent symmetry breaking and cosmology. Several statements are heuristic or numerically inaccurate: the solar Schwarzschild radius, the binary entropy sign in the adjoining intervention, and the normalization of the acceleration temperature require separate treatment. A causal loss of access to data is not silently upgraded to a precise quantum-information-loss conjecture. Vishik\textquoteright s closing Monte Carlo aphorism is assigned separately to record 427.

\begin{GSitem}{1977-425.1}{E}{Temperature can change a broken-symmetry phase}
\GSfield{Kind; source}{claim; \GSpage{P000713; P000716}}
\GSfield{The notes}{The opening connects the return of field theory with spontaneous symmetry breaking, macroscopic influences and cosmic history. The closing drawings show an evolving potential well as temperature changes.}
\GSfield{Statement}{A precise illustrative model is \(V_T(\phi)=\tfrac12 m^2(T)\phi^2+\tfrac\lambda4\phi^4\), \(\lambda>0\). If \(m^2(T)>0\), its unique minimum is \(0\); if \(m^2(T)<0\), the minima are \(\pm\sqrt{-m^2(T)/\lambda}\). Thus a temperature-dependent effective potential can change the symmetry of its minimizing state. No assertion that every theory restores symmetry at high temperature follows.}
\GSfield{Reconstruction}{The polynomial potential, parameter and explicit minima are an illustrative reconstruction; neither this formula nor a critical temperature is written in the notes.}
\GSfield{Fate}{The subject predates the talk in \citet{kirzhnitslinde1972macroscopic}. The author-institution publication record and publisher bibliographic search were checked, but the full article was inaccessible, so the drawings are matched thematically, not to a particular theorem.}
\GSfield{Present status}{As of November 2026, this is an established expository item, subject to the hypotheses stated above.}
\GSfield{Assessment}{literal statement: correct with missing hypotheses; publication match: thematic only; verification: publisher metadata checked; development: annotated; confidence: medium}
\end{GSitem}
\begin{GSitem}{1977-425.2}{E}{Dynamical instability and loss of coarse information}
\GSfield{Kind; source}{claim; \GSpage{P000713}}
\GSfield{The notes}{A region in phase space spreads because of instability, with Sinai billiards and positive entropy as examples; the notes describe this as losing information.}
\GSfield{Statement}{For a Hamiltonian flow preserving Liouville measure, the fine-grained Gibbs entropy of an absolutely continuous density is invariant under exact evolution. Instability can make a set highly stretched and fine-scaled, so information at finite observational resolution can be lost and a coarse-grained entropy can grow. Positive metric entropy is a property of a specified invariant dynamical system, not a consequence of every instability.}
\GSfield{Reconstruction}{The distinction between fine-grained conservation and coarse-grained loss is supplied. The page does not specify a billiard geometry or a partition and therefore does not support a new billiard theorem.}
\GSfield{Fate}{This is a classical dynamics-to-statistics motivation. Earlier items 1964-10.1 and 1964-10.4 already distinguish metric entropy and information-forgetting language. The claimed loss is valid only after an observational coarse graining is chosen.}
\GSfield{Present status}{As of November 2026, this is an established expository item, subject to the hypotheses stated above.}
\GSfield{Assessment}{literal statement: correct with missing hypotheses; publication match: none found; verification: indirect; development: annotated; confidence: high}
\end{GSitem}
\begin{GSitem}{1977-425.3}{E}{Boltzmann entropy of a macrostate}
\GSfield{Kind; source}{definition; \GSpage{P000713}}
\GSfield{The notes}{A macrostate is fixed by a small number of parameters but allows many microscopic configurations; the notes write \(S=\ln N_{\rm micro}\).}
\GSfield{Statement}{For a finite set of \(W\) equiprobable microstates compatible with specified macroscopic constraints, \(S=k_B\log W\). For a continuous phase space one needs a reference measure and coarse-graining convention; a raw cardinality is not finite. The notebook uses units \(k_B=1\).}
\GSfield{Reconstruction}{Equiprobability, the entropy unit and the continuous-state qualification are explicit additions.}
\GSfield{Fate}{This is the equiprobable case of the information-theoretic entropy discussed in \citet[\S6]{shannon1948communication}; its thermodynamic application long predates the talk. Earlier occurrence: 1964-9.3.}
\GSfield{Present status}{As of November 2026, this is an established expository item, subject to the hypotheses stated above.}
\GSfield{Assessment}{literal statement: correct with missing hypotheses; publication match: partial; verification: read the result; development: annotated; confidence: high}
\end{GSitem}
\begin{GSitem}{1977-425.4}{E}{Ideal-gas expansion and the entropy-information sign}
\GSfield{Kind; source}{claim; \GSpage{P000714}}
\GSfield{The notes}{Kirzhnits writes \(\Delta S=-\Delta I\), draws a gas expanding into a larger volume, and gives \(S=\ln(V/N)\) per molecule. The final sentence says that acquiring information decreases information.}
\GSfield{Statement}{For an ideal gas at fixed particle number and temperature, the volume-dependent entropy change per particle is \(\Delta S/N=k_B\log(V_f/V_i)\). In particular doubling accessible volume gives \(k_B\log2\). The full entropy also contains temperature, normalization and additive terms. If information is defined as a reduction of uncertainty in the same statistical model, its dimensionless change is opposite to the entropy change.}
\GSfield{Reconstruction}{The formula in the notes is only the volume contribution. The last sentence should say that acquiring information decreases entropy or uncertainty; the scan faithfully contains the wrong repeated word.}
\GSfield{Fate}{The information/entropy convention is clarified by \citet[\S6]{shannon1948communication}; the one-molecule gas comparison is developed explicitly in \citet[Fig.~12 and accompanying discussion]{bennett1982thermodynamics}. The gas entropy law itself is pre-1977 exposition.}
\GSfield{Present status}{As of November 2026, this is an established expository item, subject to the hypotheses stated above.}
\GSfield{Assessment}{literal statement: false as written; publication match: partial; verification: read the result; development: annotated; confidence: high}
\end{GSitem}
\begin{GSitem}{1977-425.5}{E}{Maxwell’s demon and the entropy of the whole apparatus}
\GSfield{Kind; source}{remark; \GSpage{P000714}}
\GSfield{The notes}{Maxwell\textquoteright s demon lowers entropy by following the particles.}
\GSfield{Statement}{An informed controller can reduce a subsystem\textquoteright s entropy by sorting particles or extracting work, but the complete cyclic process must include the controller\textquoteright s memory and its reset. A subsystem entropy decrease alone is not a violation of the second law for an isolated total system.}
\GSfield{Reconstruction}{The memory, cyclic-operation and total-system qualifications are absent in the notebook. No universal claim that every measurement itself dissipates \(k_BT\log2\) is inserted.}
\GSfield{Fate}{\citet[discussion of Maxwell\textquoteright s demon, especially Fig.~12]{bennett1982thermodynamics}, received in 1981 and published in 1982, distinguishes reversible measurement from resetting a record. It is a later clarification of a much older thought experiment, not the first appearance of the notebook\textquoteright s remark.}
\GSfield{Present status}{As of November 2026, this is an established expository item, subject to the hypotheses stated above.}
\GSfield{Assessment}{literal statement: correct with missing hypotheses; publication match: partial; verification: read the result; development: annotated; confidence: high}
\end{GSitem}
\begin{GSitem}{1977-425.6}{E}{A black-hole horizon as a causal barrier}
\GSfield{Kind; source}{definition; \GSpage{P000714-P000715}}
\GSfield{The notes}{The notes ask for forces that keep information in a region, name gravity because it is long-range and of one sign, and describe a black hole as a body that releases no information.}
\GSfield{Statement}{In an asymptotically flat spacetime with future null infinity \(\mathcal I^+\), the black-hole region is \(M\setminus J^-(\mathcal I^+)\), and its future boundary is the event horizon. Future-directed causal signals from this region cannot reach \(\mathcal I^+\). This is a statement about the causal metric, not a theorem deduced merely from the range and sign of a force.}
\GSfield{Reconstruction}{The global causal definition is supplied. Gravity\textquoteright s claimed uniqueness and always-attractive character are heuristic motivations, not hypotheses sufficient for the existence of an event horizon.}
\GSfield{Fate}{The classical causal setting and its difference from quantum emission are explicit in \citet[\S1]{hawking1975particle}. The report\textquoteright s black-hole description is recognizable once restricted to classical signals.}
\GSfield{Present status}{As of November 2026, this is an established expository item, subject to the hypotheses stated above.}
\GSfield{Assessment}{literal statement: correct with missing hypotheses; publication match: partial; verification: read the result; development: annotated; confidence: high}
\end{GSitem}
\begin{GSitem}{1977-425.7}{E}{The Schwarzschild radius and the incorrect solar value}
\GSfield{Kind; source}{claim; \GSpage{P000714-P000715}}
\GSfield{The notes}{Equating Newtonian escape speed to \(c\) gives \(R=2\gamma M/c^2=r_g\). The notes say this survives in exact theory and give seven kilometers for a solar mass.}
\GSfield{Statement}{For the Schwarzschild solution with mass \(M\), the event-horizon areal radius is \(r_s=2GM/c^2\). The Newtonian escape-speed argument reproduces the expression but does not derive relativistic causal structure. Using the nominal solar mass parameter \(GM_\odot=1.3271244\times10^{20}\,{\rm m^3s^{-2}}\) gives \(r_s\simeq2.95325\,{\rm km}\), rather than seven kilometers.}
\GSfield{Reconstruction}{The nonrotating, uncharged solution and the meaning of areal radius are specified. Seven is visible in the scan, so this is a faithful numerical notebook error. The quoted ordinary solar radius \(10^6\,{\rm km}\) is only an order-of-magnitude estimate.}
\GSfield{Fate}{The modern constant is verified in \citet[Table~1]{prsa2016nominal}; the radius is a direct substitution into the notebook\textquoteright s own formula. The classical Schwarzschild setting is also used in \citet[\S2]{hawking1975particle}.}
\GSfield{Present status}{As of November 2026, this is an established expository item, subject to the hypotheses stated above.}
\GSfield{Assessment}{literal statement: false as written; publication match: partial; verification: read the result; development: annotated; confidence: high}
\end{GSitem}
\begin{GSitem}{1977-425.8}{E}{The limited no-hair statement and the missing angular momentum}
\GSfield{Kind; source}{claim; \GSpage{P000715}}
\GSfield{The notes}{Only a few parameters of a black hole can be known: mass, charge and radius.}
\GSfield{Statement}{A precise established uniqueness theorem assumes a four-dimensional, stationary, asymptotically flat, analytic electrovacuum spacetime, the usual global regularity condition, and a connected nondegenerate horizon. Its domain of outer communications is Kerr--Newman. The family is parameterized by mass, angular momentum and charge; the horizon radius or area is derived, not an independent substitute for angular momentum.}
\GSfield{Reconstruction}{The field equations, asymptotics, stationarity and regularity are reconstructed. There are other matter theories and other asymptotics with black-hole hair, so the sentence is not a universal information-counting theorem.}
\GSfield{Fate}{\citet[Thm.~3.3]{chruscielcostaheusler2012stationary} gives the precise theorem read here, including the alternative rotating-horizon assumption. The no-hair program was already established before the seminar, but the notebook omits both angular momentum and the theorem\textquoteright s scope.}
\GSfield{Present status}{As of November 2026, this is an established expository item, subject to the hypotheses stated above.}
\GSfield{Assessment}{literal statement: false as written; publication match: partial; verification: read the result; development: annotated; confidence: high}
\end{GSitem}
\begin{GSitem}{1977-425.9}{E}{Classical increase of event-horizon area}
\GSfield{Kind; source}{theorem; \GSpage{P000715}}
\GSfield{The notes}{Hawking is credited with the assertion that the black-hole surface \(S\) increases in every process; merging two holes is presented as geometrically evident from the radius formula.}
\GSfield{Statement}{In classical general relativity, with the null energy condition and the global predictability/regularity assumptions needed to exclude future endpoints or exposed singularities of horizon generators, the event-horizon area cannot decrease toward the future. For a merger this concerns the final area compared with the sum of initial component areas. The radius formula alone does not prove the merger inequality, because energy can escape as radiation.}
\GSfield{Reconstruction}{The energy and causal hypotheses, and the distinction between classical area and generalized quantum entropy, are added.}
\GSfield{Fate}{\citet{hawking1971colliding} published the collision application in May 1971; only its abstract and metadata were read here. \citet[abstract and \S1]{hawking1975particle} explicitly explains that quantum emission can violate the classical area law. The quantum qualification therefore also predates this talk.}
\GSfield{Present status}{As of November 2026, the classical theorem is established under its stated hypotheses; the notebook\textquoteright s phrase every process is not valid for semiclassical evaporation.}
\GSfield{Assessment}{literal statement: correct with missing hypotheses; publication match: partial; verification: read the result; development: annotated; confidence: high}
\end{GSitem}
\begin{GSitem}{1977-425.10}{E}{Area entropy, Hawking temperature and the generalized second law}
\GSfield{Kind; source}{claim; \GSpage{P000715}}
\GSfield{The notes}{The area symbol \(S\) motivates a thermodynamic analogy; Bekenstein is said to remove the quotation marks from black-hole temperature by counting inaccessible information. The parenthesis incorrectly calls him Hawking\textquoteright s student.}
\GSfield{Statement}{In semiclassical Einstein gravity, \(S_{\rm BH}=k_Bc^3A/(4G\hbar)\) and \(T_H=\hbar\kappa/(2\pi c k_B)\) when \(\kappa\) is physical surface acceleration. The generalized second law concerns \(S_{\rm outside}+S_{\rm BH}\). Bekenstein proposed area entropy and the generalized law; Hawking\textquoteright s radiation calculation fixed the temperature and the standard entropy coefficient.}
\GSfield{Reconstruction}{The exact normalization and attribution are reconstructed from the publications. The causal information metaphor is not itself a microscopic derivation of the entropy in every theory of gravity.}
\GSfield{Fate}{\citet[abstract]{bekenstein1973entropy} was received in November 1972 and published in April 1973. \citet[abstract]{hawking1976thermodynamics} states the temperature/entropy normalization before the seminar. \citet[recollection of their Princeton student years]{wald2015bekenstein} confirms that Bekenstein\textquoteright s supervisor was John Wheeler.}
\GSfield{Present status}{As of November 2026, this is an established expository item, subject to the hypotheses stated above.}
\GSfield{Assessment}{literal statement: correct with missing hypotheses; publication match: partial; verification: read the result; development: annotated; confidence: high}
\end{GSitem}
\begin{GSitem}{1977-425.11}{E}{Comparing absorption of hot and cold bodies}
\GSfield{Kind; source}{claim; \GSpage{P000715}}
\GSfield{The notes}{A hot body thrown into a black hole is said to increase its area more than a cold body, in exact agreement with the increase in entropy.}
\GSfield{Statement}{For a quasistatic perturbation of a stationary hole, the first law is \(\delta(Mc^2)=\kappa c^2\delta A/(8\pi G)+\Omega_H\delta J+\Phi_H\delta Q\). At fixed angular-momentum and charge transfer, larger delivered energy gives larger area increase. Temperature or matter entropy alone does not determine the delivered energy, so the comparison needs specified bodies, trajectories and other conserved quantities.}
\GSfield{Reconstruction}{The first-law formulation states precisely which comparison is justified; the unconditional hot-versus-cold rule is not a mathematical theorem.}
\GSfield{Fate}{\citet[Eq.~(3.1)]{hawking1975particle} uses the classical first law in the radiation analysis. Bekenstein\textquoteright s generalized-law proposal \citep[abstract]{bekenstein1973entropy} concerns the combined entropy, not an equality between arbitrary matter entropy and horizon-area change.}
\GSfield{Present status}{As of November 2026, this is an established expository item, subject to the hypotheses stated above.}
\GSfield{Assessment}{literal statement: correct with missing hypotheses; publication match: partial; verification: read the result; development: annotated; confidence: high}
\end{GSitem}
\begin{GSitem}{1977-425.12}{E}{Hawking radiation and the thermal occupation factor}
\GSfield{Kind; source}{theorem; \GSpage{P000715}}
\GSfield{The notes}{The notes attribute radiation to Hawking and write the Planck factor \((e^{\hbar\omega/(kT)}-1)^{-1}\); the preceding phrase about pairs is only schematic.}
\GSfield{Statement}{For a free bosonic field on an asymptotically collapsing geometry settling to a Schwarzschild black hole, the late-time outgoing occupation number has the thermal factor \((e^{\hbar\omega/(k_BT_H)}-1)^{-1}\), multiplied by the mode\textquoteright s transmission or greybody factor. Fermions have a plus sign in the denominator. This is a semiclassical result with a prescribed background and specified incoming vacuum.}
\GSfield{Reconstruction}{The scattering factor, field statistics and semiclassical assumptions are supplied. A pair-creation cartoon is not the derivation, and the observed spectrum at infinity is not an exact unit-emissivity Planck spectrum in every mode.}
\GSfield{Fate}{\citet[Eq.~(2.29) and the following fermion discussion]{hawking1975particle} was received in April 1975 and published in 1975. The exact occupation formula and its absorption factor were read in the full paper.}
\GSfield{Present status}{As of November 2026, this is an established expository item, subject to the hypotheses stated above.}
\GSfield{Assessment}{literal statement: correct with missing hypotheses; publication match: exact; verification: read the result; development: annotated; confidence: high}
\end{GSitem}
\begin{GSitem}{1977-425.13}{F}{Causal inaccessibility versus quantum information loss}
\GSfield{Kind; source}{remark; \GSpage{P000715}}
\GSfield{The notes}{Information is said to be lost not through stochasticization but because it is mechanically cut off by the black hole.}
\GSfield{Statement}{The stated contrast concerns the exterior observer\textquoteright s causal access to the interior. It does not specify whether a complete quantum-gravitational evolution maps pure states to mixed states, what final state is compared, or which algebra of observables is used.}
\GSfield{Reconstruction}{Promoting this sentence to the precise black-hole unitarity problem would introduce quantifiers and a quantum model not given in the notebook.}
\GSfield{Fate}{\citet[\S1]{hawking1975particle} supplies the nearby semiclassical setting, but the notes do not isolate a complete quantum-evolution conjecture. No exact open problem is inferred merely from the word information.}
\GSfield{Present status}{As of November 2026, this occurrence is identifiable as a physical motivation, but not as a definite quantum-information-loss assertion whose resolution can be tracked independently.}
\GSfield{Assessment}{literal statement: not determinable; publication match: thematic only; verification: unresolved lead; development: none; confidence: medium}
\end{GSitem}
\begin{GSitem}{1977-425.14}{E}{Uniform acceleration and a Rindler horizon}
\GSfield{Kind; source}{construction; \GSpage{P000715}}
\GSfield{The notes}{A uniformly accelerated observer in vacuum receives light only from part of space, described as half of space; a spacetime diagram accompanies the statement.}
\GSfield{Statement}{In Minkowski space, the eternal worldline \(ct=(c^2/a)\sinh(a\tau/c)\), \(x=(c^2/a)\cosh(a\tau/c)\) has constant proper acceleration \(a>0\). Its adapted stationary region is the right Rindler wedge \(x>|ct|\), bounded by null horizons. The events from which it can ever receive a signal form the causal past of the worldline; this is not simply a fixed Euclidean half-space at all times.}
\GSfield{Reconstruction}{The eternal-acceleration assumption and distinction between the wedge and the entire causal past are supplied. This matters because a finite period of acceleration does not by itself produce the same eternal observer horizon.}
\GSfield{Fate}{\citet[\S~II.B, Eqs.~(2.34)--(2.37)]{crispinohiguchimatsas2008unruh} gives the geometry; the original 1976 detector paper is \citet{unruh1976notes}.}
\GSfield{Present status}{As of November 2026, this is an established expository item, subject to the hypotheses stated above.}
\GSfield{Assessment}{literal statement: correct with missing hypotheses; publication match: partial; verification: read the result; development: annotated; confidence: high}
\end{GSitem}
\begin{GSitem}{1977-425.15}{E}{The Unruh temperature and observer-dependent particles}
\GSfield{Kind; source}{theorem; \GSpage{P000715-P000716}}
\GSfield{The notes}{The inertial vacuum is described as a ground state of the electromagnetic Hamiltonian; the accelerated observer is said to find a Planck distribution. The diagram labels \(kT=ha/c\), without a \(2\pi\) factor, and the ground-state equation uses different ket letters on its two sides.}
\GSfield{Statement}{For the Minkowski vacuum restricted to a Rindler wedge, the boost-time state is thermal in the KMS sense. An observer of constant proper acceleration \(a\) has \(T_U=\hbar a/(2\pi c k_B)\). For free bosonic Rindler modes the mean occupation is \((e^{\hbar\omega/(k_BT_U)}-1)^{-1}\). A detector\textquoteright s detailed balance is thermal, but its entire response need not equal that of an inertial detector in a Minkowski thermal bath.}
\GSfield{Reconstruction}{The notation is standardized to reduced Planck\textquoteright s constant \(\hbar\); if the handwritten \(h\) means ordinary Planck\textquoteright s constant, the coefficient is instead \(h/(4\pi^2)\). The inertial ground-state equation must use the same vacuum vector on both sides.}
\GSfield{Fate}{\citet[abstract]{unruh1976notes} appeared in August 1976. \citet[Eq.~(2.125), \S~II.I and Introduction]{crispinohiguchimatsas2008unruh} explicitly verifies the occupation factor, KMS statement and detector qualification; its preprint circulated in October 2007, before 2008 publication.}
\GSfield{Present status}{As of November 2026, this is an established expository item, subject to the hypotheses stated above.}
\GSfield{Assessment}{literal statement: false as written; publication match: partial; verification: read the result; development: annotated; confidence: high}
\end{GSitem}
\begin{GSitem}{1977-425.16}{E}{The microwave background as evidence for a hot past}
\GSfield{Kind; source}{claim; \GSpage{P000716}}
\GSfield{The notes}{The notes ask why the universe is hot, quote roughly \(3\,\mathrm K\) from a Planckian cosmic radiation distribution, call it relic radiation, and say the universe is cooling.}
\GSfield{Statement}{The measured cosmic microwave background is a thermal radiation field. In a homogeneous expanding cosmology with collisionless propagation after thermalization, photon wavelengths scale with the scale factor \(a\), so a Planck spectrum remains Planckian with \(T\propto a^{-1}\). The notebook\textquoteright s \(3\,\mathrm K\) is a rough historical temperature, not an exact modern measurement.}
\GSfield{Reconstruction}{The cosmological assumptions and redshift argument are supplied. A single-frequency antenna temperature is not itself a full measurement of a Planck spectrum.}
\GSfield{Fate}{\citet[pp.~419--421]{penziaswilson1965antenna} reports the 1965 excess antenna temperature, with a May 1965 manuscript date, well before this seminar. That paper establishes the observational antecedent; it is not cited as a full-spectrum measurement or a derivation of cosmic initial conditions.}
\GSfield{Present status}{As of November 2026, this is an established expository item, subject to the hypotheses stated above.}
\GSfield{Assessment}{literal statement: correct with missing hypotheses; publication match: partial; verification: read the result; development: annotated; confidence: high}
\end{GSitem}
\begin{GSitem}{1977-425.17}{F}{Does a cosmological horizon explain the radiation temperature?}
\GSfield{Kind; source}{question; \GSpage{P000716}}
\GSfield{The notes}{The question why the universe is hot is followed by an inaccessible region, the phrase event horizon, the explanation that light has not yet arrived, and an arrow suggesting that this accounts for the temperature.}
\GSfield{Statement}{Light not yet received from a cosmological region concerns a particle horizon; an event horizon concerns signals that can never be received in the future. Neither causal boundary alone fixes a microwave-background temperature. A horizon-temperature calculation requires a specified spacetime, quantum state and time normalization, none of which are given here.}
\GSfield{Reconstruction}{The fragment cannot be uniquely identified with de Sitter temperature, a reheating mechanism, or a general explanation of cosmological initial entropy. These are distinct modern questions.}
\GSfield{Fate}{The nearby analogy to thermal horizons is discussed in \citet[Introduction]{crispinohiguchimatsas2008unruh}, but that does not identify this sketch with a particular cosmological theorem. The search found no exact formula whose hypotheses can be recovered from the page.}
\GSfield{Present status}{As of November 2026, the intended cosmological model and asserted temperature relation remain unidentified. No definite problem identifier is assigned.}
\GSfield{Assessment}{literal statement: not determinable; publication match: thematic only; verification: unresolved lead; development: none; confidence: low}
\end{GSitem}

\GStalk{426}{1977-10-03}{L. A. Dickey}{P000717-P000718}{Formal variational calculus and Hamiltonian differential operators}
Dickey defines differential polynomials, local functionals, variational derivatives and evolutionary derivations; he then gives an explicit Hamiltonian matrix and asks for a simpler proof of the Jacobi identity. The final comparison is with the third-order KdV Hamiltonian operator. P000718 belongs to this talk throughout, despite the supplied catalogue range ending at P000717; S. I. Gelfand begins on P000719. The principal theorem is linked explicitly in the notebook to Gelfand and Dickey\textquoteright s 1977 paper.

\begin{GSitem}{1977-426.1}{E}{The differential-polynomial algebra and its total derivative}
\GSfield{Kind; source}{definition; \GSpage{P000717}}
\GSfield{The notes}{The generators are \(u_k^{(i)}\), and \(A=K[u_k^{(i)}]\). The notes allow derivations given by infinite sums and define \(D=d/dx\) by \(D u_k^{(i)}=u_k^{(i+1)}\).}
\GSfield{Statement}{Over a field \(K\) of characteristic zero, let \(A\) be the polynomial algebra in the countably many jet variables, every polynomial involving only finitely many. Then \(D=\sum_{k,i}u_k^{(i+1)}\partial/\partial u_k^{(i)}\) is well defined. Likewise an arbitrary assignment of polynomial images to the generators determines a derivation; its displayed infinite sum acts finitely on each polynomial.}
\GSfield{Reconstruction}{Characteristic zero is imposed for the subsequent variational identities. Infinite summation here is locally finite on polynomials and does not invoke analytic convergence.}
\GSfield{Fate}{This formal setup is already explicit in \citet[\S1, pp.~13--14 of the Russian original]{gelfanddikii1977resolvent}, received in December 1976 and printed before the October 1977 talk.}
\GSfield{Present status}{As of November 2026, this is an established expository item, subject to the hypotheses stated above.}
\GSfield{Assessment}{literal statement: correct with missing hypotheses; publication match: exact; verification: read the result; development: annotated; confidence: high}
\end{GSitem}
\begin{GSitem}{1977-426.2}{E}{Local functionals modulo total derivatives}
\GSfield{Kind; source}{definition; \GSpage{P000717}}
\GSfield{The notes}{The normalization is \(u_n=1\), \(u_{n-1}=0\), and the space of functionals is \(\widetilde A=A/DA\), with the quotient map denoted \(f\mapsto\int f\,dx\).}
\GSfield{Statement}{After imposing the differential relations \(u_n=1\), \(u_{n-1}=0\), the free dependent variables are \(u_0,\ldots,u_{n-2}\). The local-functional space \(\mathcal F=A/D(A)\) is a vector-space quotient. Integration by parts means \(\int fDg\,dx=-\int (Df)g\,dx\). In general \(D(A)\) is not an ideal, so \(\mathcal F\) is not naturally a quotient algebra.}
\GSfield{Reconstruction}{The relations include derivatives of the normalized coefficients. The vector-space, rather than algebra, qualification is important for the later Lie bracket.}
\GSfield{Fate}{\citet[\S1]{gelfanddikii1977resolvent} gives this exact definition and its periodic/rapidly decreasing interpretations; \citet[\S1]{gelfanddikii1978jets} reuses it in the later elementary calculus.}
\GSfield{Present status}{As of November 2026, this is an established expository item, subject to the hypotheses stated above.}
\GSfield{Assessment}{literal statement: correct; publication match: exact; verification: read the result; development: annotated; confidence: high}
\end{GSitem}
\begin{GSitem}{1977-426.3}{E}{When does an identically zero local integral imply a total derivative?}
\GSfield{Kind; source}{problem; \GSpage{P000717}}
\GSfield{The notes}{The exercise asks whether \(\int_{\mathbf R}f[u(x)]\,dx=0\) for every Schwartz function is equivalent to \(f\in D\mathbf C[u]\).}
\GSfield{Statement}{Let \(A=\mathbf C[u_k^{(i)}:0\leq k\leq n-2,\ i\geq0]\), with no explicit \(x\), and assume the constant term \(f(0)\) is zero. Then \(\int_{\mathbf R}f[j^\infty u(x)]\,dx=0\) for every Schwartz tuple \(u\) if and only if \(f=Dg\) for some \(g\in A\). All evaluated nonconstant differential monomials are integrable under these hypotheses.}
\GSfield{Reconstruction}{The tuple interpretation and zero constant term make the analytic integral well defined. Algebraically, vanishing of all Euler derivatives gives \(f\in\mathbf C+D(A)\); the constant must not be dropped. Integration by parts proves one direction, while arbitrary compactly supported variations and the one-dimensional variational homotopy give the converse.}
\GSfield{Fate}{The periodic and Schwartz realizations and Euler-derivative lemma are discussed in \citet[\S1, pp.~13--14]{gelfanddikii1977resolvent}. Its printed shorthand for the kernel also requires separating constants. This is a classical null-Lagrangian exercise, not an undeveloped problem.}
\GSfield{Present status}{As of November 2026, this is an established expository item, subject to the hypotheses stated above.}
\GSfield{Assessment}{literal statement: correct with missing hypotheses; publication match: exact; verification: read the result; development: annotated; confidence: high}
\end{GSitem}
\begin{GSitem}{1977-426.4}{E}{The Euler variational derivative}
\GSfield{Kind; source}{construction; \GSpage{P000717}}
\GSfield{The notes}{The map \(\delta/\delta u:\widetilde A\to A^{n-1}\) is given by \((\delta\widetilde f/\delta u)_k=\sum_{\ell\geq0}(-D)^\ell\partial f/\partial u_k^{(\ell)}\).}
\GSfield{Statement}{For each \(k\), define \(E_k(f)=\sum_{\ell\geq0}(-D)^\ell(\partial f/\partial u_k^{(\ell)})\). The sum is finite on every differential polynomial, and \(E_k(Dg)=0\), so the vector \(E(f)\) is independent of the chosen representative of \(\int f\,dx\). It is the coefficient of an arbitrary variation after integration by parts.}
\GSfield{Reconstruction}{Finite support and well-definedness on the quotient are explained explicitly; the formula itself is faithfully transcribed.}
\GSfield{Fate}{\citet[\S1, variational derivative and following theorem]{gelfanddikii1977resolvent} gives the definition and the integration-by-parts characterization. This is the same formal calculus continued in \citet[\S1]{gelfanddikii1978jets}.}
\GSfield{Present status}{As of November 2026, this is an established expository item, subject to the hypotheses stated above.}
\GSfield{Assessment}{literal statement: correct; publication match: exact; verification: read the result; development: annotated; confidence: high}
\end{GSitem}
\begin{GSitem}{1977-426.5}{E}{The first Gelfand--Dickey Hamiltonian matrix}
\GSfield{Kind; source}{construction; \GSpage{P000718}}
\GSfield{The notes}{The notes define \((\ell f)_r=\sum_s\ell_{rs}f_s\) by a binomial sum involving \(u_{r+s+\gamma+1}\) and powers of \(\pm iD\).}
\GSfield{Statement}{With \(0\leq r,s\leq n-2\), \(u_n=1\), \(u_{n-1}=0\), and empty sums interpreted as zero, the operator is
\[\ell_{rs}=\sum_{\gamma=0}^{n-1-r-s}\left[\binom{\gamma+r}{r}u_{r+s+\gamma+1}(-iD)^\gamma-\binom{\gamma+s}{s}(iD)^\gamma\circ u_{r+s+\gamma+1}\right].\]
Here multiplication by a coefficient is an operator, so the right-hand derivative acts on its product with the argument. The matrix is formally skew-adjoint, \(\ell_{rs}^*=-\ell_{sr}\).}
\GSfield{Reconstruction}{Composition, empty-sum convention and the use of \(K\supset\mathbf Q(i)\) are supplied. The repeated minus at the manuscript line break denotes a single continued subtraction, not two negatives giving a plus. This is the first, coefficient-linear Gelfand--Dickey bracket; it must not be confused with the later quadratic bracket.}
\GSfield{Fate}{\citet[\S2, Eq.~(8)]{gelfanddikii1977resolvent} gives the matrix version. \citet[Eq.~(1) and \S5]{lebedevmanin1979volterra} identifies the same coefficient-linear operator, after the derivative convention is changed, with a coadjoint-orbit construction.}
\GSfield{Present status}{As of November 2026, this is an established expository item, subject to the hypotheses stated above.}
\GSfield{Assessment}{literal statement: correct; publication match: exact; verification: read the result; development: annotated; confidence: high}
\end{GSitem}
\begin{GSitem}{1977-426.6}{E}{The first nonconstant small-order example}
\GSfield{Kind; source}{remark; \GSpage{P000718}}
\GSfield{The notes}{The cases \(n=2,3\) are called easy, while \(n=4\) is the first substantial example.}
\GSfield{Statement}{Substitution into the displayed matrix gives \(\ell=-2iD\) for \(n=2\), and \(\ell=\left(\begin{smallmatrix}0&-3iD\\-3iD&0\end{smallmatrix}\right)\) for \(n=3\). Both have constant coefficients. For \(n=4\), \(\ell_{00}=2iD^3-i(2u_2D+u_2^{\prime})\), so coefficient dependence first appears.}
\GSfield{Reconstruction}{These explicit small-order substitutions are calculated here to explain the notebook\textquoteright s qualitative remark; they are not additional formulas copied from the page.}
\GSfield{Fate}{The calculation is a direct specialization of \citet[Eq.~(8)]{gelfanddikii1977resolvent} and agrees with the scan. The general Poisson theorem, rather than only these examples, was already proved in that paper.}
\GSfield{Present status}{As of November 2026, this is an established expository item, subject to the hypotheses stated above.}
\GSfield{Assessment}{literal statement: correct; publication match: exact; verification: read the result; development: annotated; confidence: high}
\end{GSitem}
\begin{GSitem}{1977-426.7}{E}{Evolutionary derivations commute with the total derivative}
\GSfield{Kind; source}{construction; \GSpage{P000718}}
\GSfield{The notes}{To \(f\in A^{n-1}\) the notes assign \(\partial_f=\sum_{k,\ell}f_k^{(\ell)}\partial/\partial u_k^{(\ell)}\), assert \([\partial_f,D]=0\), and hence let it act on \(\widetilde A\).}
\GSfield{Statement}{Define \(\partial_f(u_k^{(\ell)})=D^\ell f_k\). On each generator, \(\partial_fD=D\partial_f\); therefore the identity holds on all of \(A\). It follows that \(\partial_f(Dg)=D(\partial_fg)\), so \(\partial_f\) descends to \(A/D(A)\). Conversely, every derivation commuting with \(D\) has this form.}
\GSfield{Reconstruction}{The generator check and converse make the local-finiteness and descent arguments explicit.}
\GSfield{Fate}{This correspondence is already in \citet[\S1, Eq.~(5)]{gelfanddikii1977resolvent}.}
\GSfield{Present status}{As of November 2026, this is an established expository item, subject to the hypotheses stated above.}
\GSfield{Assessment}{literal statement: correct; publication match: exact; verification: read the result; development: annotated; confidence: high}
\end{GSitem}
\begin{GSitem}{1977-426.8}{E}{The Lie algebra of evolutionary characteristics}
\GSfield{Kind; source}{theorem; \GSpage{P000718}}
\GSfield{The notes}{The bracket is \([f,g]_r=\partial_fg_r-\partial_gf_r\), with \([\partial_f,\partial_g]=\partial_{[f,g]}\); this makes \(A^{n-1}\) a Lie algebra.}
\GSfield{Statement}{The map \(f\mapsto\partial_f\) is injective, since \(\partial_fu_k=f_k\), and identifies \(A^{n-1}\) with the Lie algebra of derivations of \(A\) commuting with \(D\). The displayed formula is therefore bilinear, alternating and satisfies Jacobi as the pullback of the commutator of derivations.}
\GSfield{Reconstruction}{Injectivity is the extra step justifying the inference from the commutator identity to Jacobi.}
\GSfield{Fate}{\citet[\S2, before Thm.~1]{gelfanddikii1977resolvent} uses exactly this Lie algebra.}
\GSfield{Present status}{As of November 2026, this is an established expository item, subject to the hypotheses stated above.}
\GSfield{Assessment}{literal statement: correct; publication match: exact; verification: read the result; development: annotated; confidence: high}
\end{GSitem}
\begin{GSitem}{1977-426.9}{E}{The Poisson bracket on local functionals}
\GSfield{Kind; source}{theorem; \GSpage{P000718}}
\GSfield{The notes}{The notes define \(\{\widetilde f,\widetilde g\}=\int\sum_{r,s}(\delta\widetilde f/\delta u_r)\ell_{rs}(\delta\widetilde g/\delta u_s)\,dx\) and state that it makes \(\widetilde A\) a Lie algebra.}
\GSfield{Statement}{On the differential-polynomial local-functional space over characteristic zero, with the normalization and operator of item 1977-426.5, the displayed bracket is well defined and skew-symmetric and satisfies the Jacobi identity. Since the quotient need not carry a product, the immediate precise assertion is a Lie bracket on local functionals; its Hamiltonian interpretation is formal variational calculus.}
\GSfield{Reconstruction}{The hypotheses are gathered together and the absence of a natural quotient-algebra product is kept visible.}
\GSfield{Fate}{\citet[\S2, Thm.~2 and its proof, pp.~16--18]{gelfanddikii1977resolvent} proves Jacobi. This was received 22 December 1976 and published in volume 11, issue 2 of 1977, before the October talk; the notebook itself cites that issue.}
\GSfield{Present status}{As of November 2026, this is an established expository item, subject to the hypotheses stated above.}
\GSfield{Assessment}{literal statement: correct; publication match: exact; verification: read the result; development: annotated; confidence: high}
\end{GSitem}
\begin{GSitem}{1977-426.10}{B}{A conceptual proof of the first Gelfand--Dickey Jacobi identity}
\GSfield{Kind; source}{problem; \GSpage{P000718}}
\GSfield{The notes}{Immediately after the theorem the notes ask: \ru{доказать её попроще} (prove it more simply).}
\GSfield{Statement}{Give a structural proof of the Jacobi identity for the coefficient-linear operator \(\ell\), explaining its origin instead of checking a large binomial identity term by term.}
\GSfield{Reconstruction}{``Structural spells out a reasonable interpretation of the request for simplicity. It does not claim that the original theorem was unproved, nor does it substitute the distinct second Gelfand--Dickey bracket.}
\GSfield{Fate}{\citet[Introduction, \S5 theorem and \S6]{lebedevmanin1979volterra} explicitly identifies the cumbersome direct Jacobi verification as the problem, realizes this first operator through the coadjoint representation of the formal Volterra group, and gives a new invariant proof. Their paper was received 10 November 1978, circulated as a 1978 ITEF preprint, and appeared in 1979; the introduction acknowledges independent work by Adler.}
\GSfield{Present status}{As of November 2026, the conceptual-proof request is fulfilled by the Volterra/coadjoint interpretation. The exact operator and the new proof were checked in the full primary article, so this is not inferred from a title alone.}
\GSfield{Assessment}{literal statement: correct; publication match: exact; verification: read the result; development: annotated; confidence: high; problem \hyperref[prob:GP-1977-02]{GP-1977-02}}
\end{GSitem}
\begin{GSitem}{1977-426.11}{E}{Novikov’s third-order Hamiltonian operator for KdV}
\GSfield{Kind; source}{claim; \GSpage{P000718}; S. P. Novikov}
\GSfield{The notes}{Novikov observes that \(A=D^3+uD+Du\) also gives a Lie algebra and can be used to write KdV.}
\GSfield{Statement}{The scalar operator \(P=D^3+uD+D\circ u=D^3+2uD+u_x\) is Hamiltonian on local functionals. With \(H=\tfrac12\int u^2\,dx\), the evolution \(u_t=P(\delta H/\delta u)\) is \(u_t=u_{xxx}+3uu_x\). Writing \(u=2v\) gives \(v_t=v_{xxx}+6vv_x\). This is the familiar second scalar KdV Poisson structure, distinct from the first-order \(n=2\) matrix above.}
\GSfield{Reconstruction}{Composition by \(u\), the Hamiltonian and the scaling are supplied. The notebook states no numerical KdV normalization at this point.}
\GSfield{Fate}{The exact operator is a specialization of the scalar Hamiltonian pair in \citet[\S7, p.~29 of the Russian original]{gelfanddorfman1979hamiltonian}. That section reports that the KdV Lenard scheme was known to the authors from a Moscow lecture by Lax in 1976, providing retrospective evidence of circulation before this talk. This does not establish the earliest priority for the second bracket. The displayed KdV flow is checked here by applying the operator to \(u\).}
\GSfield{Present status}{As of November 2026, this is an established expository item, subject to the hypotheses stated above.}
\GSfield{Assessment}{literal statement: correct; publication match: exact; verification: read the result; development: annotated; confidence: high}
\end{GSitem}

\GStalk{427}{1977-09-19}{S. M. Vishik}{P000716}{A Monte Carlo metaphor for a broad seminar}
The only identifiable contribution is the closing aphorism that integrating a large domain calls for the Monte Carlo method, offered as an explanation of the seminar\textquoteright s breadth. No measure, integrand, error estimate or mathematical problem is specified. It is preserved as a remark, not reconstructed as a separate numerical-analysis lecture.

\begin{GSitem}{1977-427.1}{F}{Monte Carlo as a metaphor for the seminar’s breadth}
\GSfield{Kind; source}{remark; \GSpage{P000716}}
\GSfield{The notes}{The notes attribute to Senya Vishik the explanation that, to integrate a large domain successfully, one should use the Monte Carlo method.}
\GSfield{Statement}{The sentence specifies no domain, integrand, probability measure, estimator or accuracy criterion. In its stated context it is a metaphor about broad intellectual sampling, not a recoverable numerical-analysis theorem.}
\GSfield{Reconstruction}{An unbiased sample-mean estimator and an \(N^{-1/2}\) error bound would be a possible illustration, but attributing either to this occurrence would add mathematical content absent from the source.}
\GSfield{Fate}{The attribution and metaphor are directly identifiable in the scan. No separate published technical proposal can be matched to the unquantified sentence, and no open numerical-integration problem is manufactured from it.}
\GSfield{Present status}{As of November 2026, the occurrence remains a contextual remark without a definite mathematical claim to resolve.}
\GSfield{Assessment}{literal statement: not determinable; publication match: none found; verification: unresolved lead; development: none; confidence: high}
\end{GSitem}

\GStalk{428}{1977-10-03}{S. I. Gelfand}{P000719-P000720}{Local algebra models for Harish–Chandra modules}
The opening lecture introduces admissible \((\mathfrak g,K)\)-modules, generalized infinitesimal characters and the problem of replacing suitable blocks by modules over explicit algebras. The finite-generation condition is supplied in the next lecture. The example is a matrix order for \(\operatorname{PSL}_2(\mathbb R)\). The heading begins on P000719; P000718 belongs to the preceding Dickey--Novikov material. The final nonmathematical report is not an occurrence-index item.

\begin{GSitem}{1977-428.1}{E}{The category of Harish–Chandra modules}
\GSfield{Kind; source}{definition; \GSpage{P000719}}
\GSfield{The notes}{The notes describe a compatible action of \(G\) and \(\mathfrak g\), with finite multiplicities on restriction to \(K\), and denote the category by \(\mathcal C\).}
\GSfield{Statement}{For a real linear connected semisimple group \(G\), maximal compact subgroup \(K\), and complexified Lie algebra \(\mathfrak g\), a Harish--Chandra module is a finitely generated \(U(\mathfrak g)\)-module with compatible locally finite algebraic \(K\)-action and finite multiplicity of every irreducible \(K\)-type.}
\GSfield{Reconstruction}{The notes here do not require a global action of \(G\) on the algebraic module; compatibility is between \(\mathfrak g\) and \(K\). Finite generation is imported from the next lecture, explicitly marked.}
\GSfield{Fate}{\citet[\S2]{bernsteingelfandgelfand1978local1} state the full definition. The manuscript supplies finite generation on P000721; it is needed in the theorem that follows.}
\GSfield{Present status}{As of 10 November 2026, this is an established result in the stated scope.}
\GSfield{Assessment}{literal statement: correct with missing hypotheses; publication match: exact; verification: read the result; development: none; confidence: high}
\end{GSitem}
\begin{GSitem}{1977-428.2}{A}{Algebraization of a bounded part of the Harish–Chandra category}
\GSfield{Kind; source}{problem; \GSpage{P000719}}
\GSfield{The notes}{Gelfand asks for a description of \(\mathcal C\) by representations of an associative algebra with finitely many generators and relations.}
\GSfield{Statement}{For a bounded connected open subset \(D\) of the spectrum \(\Theta\) of the center \(Z(U(\mathfrak g))\), seek an algebra \(Q\), finite as a module over the center, for which the category of Harish--Chandra modules with generalized infinitesimal characters in \(D\) is equivalent to the corresponding support-restricted finite-dimensional \(Q\)-module category.}
\GSfield{Reconstruction}{The bounded central-support restriction is supplied by the theorem actually proved; the opening formulation does not state it.}
\GSfield{Fate}{The precise affirmative theorem is \citet[\S3, Theorem~1]{bernsteingelfandgelfand1978local1}, presented to the Paris Academy on 9 January 1978 and printed in March 1978. It is local in central support, not an equivalence between the entire category and all modules over one finite-dimensional algebra.}
\GSfield{Present status}{As of 10 November 2026, this is an established result in the stated scope.}
\GSfield{Assessment}{literal statement: correct with missing hypotheses; publication match: exact; verification: read the result; development: none; confidence: high; problem \hyperref[prob:GP-1977-03]{GP-1977-03}}
\end{GSitem}
\begin{GSitem}{1977-428.3}{E}{The Lorentz-group matrix-pair precedent}
\GSfield{Kind; source}{remark; \GSpage{P000719}}
\GSfield{The notes}{The notes recall the Gelfand--Ponomarev reduction for the Lorentz group to two operators \(A,B\) satisfying \(AB=BA=0\).}
\GSfield{Statement}{The cited precedent concerns explicit finite-dimensional matrix problems describing suitable Harish--Chandra blocks for the Lorentz group. The relations \(AB=BA=0\) alone, without the block, support, nilpotence and equivalence conventions, do not specify the whole representation category.}
\GSfield{Fate}{The later account \citet[\S6, case \(\operatorname{SL}_2(\mathbb C)\)]{bernsteingelfandgelfand1978local2} confirms an explicit local matrix-order model. This verifies the direction of the precedent but not an unrestricted equivalence with the bare two-operator category.}
\GSfield{Present status}{As of 10 November 2026, this is an established result in the stated scope.}
\GSfield{Assessment}{literal statement: correct with missing hypotheses; publication match: partial; verification: read the result; development: none; confidence: medium}
\end{GSitem}
\begin{GSitem}{1977-428.4}{E}{Central characters and support}
\GSfield{Kind; source}{definition; \GSpage{P000719-P000720}}
\GSfield{The notes}{The center \(Z\) has spectrum \(\Theta\); the notes attach a subset \(\operatorname{supp}V\subset\Theta\) to a module.}
\GSfield{Statement}{For a Harish--Chandra module \(V\), its central support is the finite set of characters \(\theta:Z\to\mathbb C\) for which the generalized \(\theta\)-eigenspace is nonzero. Finite generation and admissibility imply that \(V\) is annihilated by a finite-codimensional ideal of \(Z\).}
\GSfield{Fate}{\citet[\S2]{bernsteingelfandgelfand1978local1} formulate this central-finiteness statement and use it to define the categories attached to subsets of \(\Theta\).}
\GSfield{Present status}{As of 10 November 2026, this is an established result in the stated scope.}
\GSfield{Assessment}{literal statement: correct with missing hypotheses; publication match: exact; verification: read the result; development: none; confidence: high}
\end{GSitem}
\begin{GSitem}{1977-428.5}{E}{Decomposition by generalized infinitesimal character}
\GSfield{Kind; source}{theorem; \GSpage{P000720}}
\GSfield{The notes}{The category attached to a set of characters is written as a sum of the single-character categories; distinct characters do not interact.}
\GSfield{Statement}{Every Harish--Chandra module is a finite direct sum of its generalized central-character submodules. A homomorphism between modules with disjoint generalized central supports is zero. The category is therefore the finite-support direct sum of its generalized-character blocks.}
\GSfield{Fate}{This is the central decomposition used in \citet[\S2]{bernsteingelfandgelfand1978local1}. It follows from the primary decomposition of a module over a finite-dimensional quotient of the commutative center.}
\GSfield{Present status}{As of 10 November 2026, this is an established result in the stated scope.}
\GSfield{Assessment}{literal statement: correct; publication match: exact; verification: read the result; development: none; confidence: high}
\end{GSitem}
\begin{GSitem}{1977-428.6}{A}{An algebra finite over the center}
\GSfield{Kind; source}{theorem; \GSpage{P000720}}
\GSfield{The notes}{For bounded \(D\), the notes state an equivalence \(\mathcal C_D\simeq\operatorname{mod}Q\), where \(Q\) is finite over \(Z\).}
\GSfield{Statement}{For bounded connected open \(D\subset\Theta\), there exists a \(Z\)-algebra \(Q\), finite over \(Z\), and a \(Z\)-linear equivalence \(\mathcal C_D\simeq\operatorname{Mod}_D Q\), where the right side consists of finite-dimensional modules whose central support lies in \(D\).}
\GSfield{Fate}{Exactly \citet[\S3, Theorem~1]{bernsteingelfandgelfand1978local1}. The later analytic localization \(Q\otimes_Z\mathcal O(D)\) permits a category of all finite-dimensional modules of that localized algebra; this is a different formulation.}
\GSfield{Present status}{As of 10 November 2026, this is an established result in the stated scope.}
\GSfield{Assessment}{literal statement: correct with missing hypotheses; publication match: exact; verification: read the result; development: none; confidence: high; problem \hyperref[prob:GP-1977-03]{GP-1977-03}}
\end{GSitem}
\begin{GSitem}{1977-428.7}{A}{A rank-three analytic matrix order for PSL2(R)}
\GSfield{Kind; source}{construction; \GSpage{P000720}}
\GSfield{The notes}{The example uses \(3\times3\) matrix-valued analytic functions with specified entries vanishing at a distinguished point.}
\GSfield{Statement}{Near the regular infinitesimal character of the trivial representation of \(\operatorname{PSL}_2(\mathbb R)\), the block category is represented by a matrix order \(R\subset M_3(\mathcal O(D))\). In the published left-module convention, with rows and columns labeled \(-1,0,1\), the entries \(a_{-1,0},a_{1,0},a_{-1,1},a_{1,-1}\) vanish at the centered parameter \(z=0\).}
\GSfield{Reconstruction}{The published parameter and left-module convention are stated explicitly; the notebook changes orientation in later diagrams.}
\GSfield{Fate}{\citet[\S5]{bernsteingelfandgelfand1978local2} give the exact three-type order in March 1978. The notebook\textquotesingle s row/column placement must be compared up to permutation and opposite-algebra convention, not silently identified entry by entry.}
\GSfield{Present status}{As of 10 November 2026, this is an established result in the stated scope.}
\GSfield{Assessment}{literal statement: correct with missing hypotheses; publication match: exact; verification: read the result; development: none; confidence: high}
\end{GSitem}

\GStalk{429}{1977-10-10}{S. I. Gelfand}{P000721-P000724}{Nice algebras and low-rank Harish–Chandra blocks}
The lecture makes the category and equivalence problem precise, explains Morita equivalence, and gives analytic matrix orders for rank-one groups and \(\operatorname{SL}_3(\mathbb R)\). The later Bernstein--Gelfand--Gelfand notes identify the precise restrictions: bounded central support for the first algebraization theorem, and regular central characters and sufficiently small neighborhoods for the nice-algebra theorem. Quiver and matrix conventions in the notebook are not fully consistent.

\begin{GSitem}{1977-429.1}{E}{Finite generation and admissibility}
\GSfield{Kind; source}{definition; \GSpage{P000721}}
\GSfield{The notes}{The definition now requires finite generation over \(U(\mathfrak g)\), compatible actions, and finite multiplicity for each irreducible \(K\)-type; Harish-Chandra\textquotesingle s theorem is invoked.}
\GSfield{Statement}{A Harish--Chandra module has the three conditions just stated. For irreducible admissible representations in the usual reductive representation-theoretic setting, the \(K\)-finite vectors form such a module. This is not a theorem about arbitrary representations without a topology or admissibility hypothesis.}
\GSfield{Fate}{\citet[\S2]{bernsteingelfandgelfand1978local1} use this precise algebraic category. The classical admissibility theorem is background to the construction, rather than the novel local-algebra statement.}
\GSfield{Present status}{As of 10 November 2026, this is an established result in the stated scope.}
\GSfield{Assessment}{literal statement: correct with missing hypotheses; publication match: exact; verification: read the result; development: none; confidence: high}
\end{GSitem}
\begin{GSitem}{1977-429.2}{E}{The support criterion is a character, not its value}
\GSfield{Kind; source}{definition; \GSpage{P000721}}
\GSfield{The notes}{The support condition is written with \(\theta(z)\), followed by a common central eigenvector condition.}
\GSfield{Statement}{A character \(\theta\in\Theta\) lies in the support of \(V\) if its generalized eigenspace is nonzero. Since the central action factors through a finite-dimensional commutative algebra, this is equivalent to the existence of a nonzero vector on which every \(z\in Z\) acts as \(\theta(z)\). The scalar \(\theta(z)\) itself is not a point of \(\Theta\).}
\GSfield{Fate}{The correctly typed support definition appears in \citet[\S2]{bernsteingelfandgelfand1978local1}.}
\GSfield{Present status}{As of 10 November 2026, this is an established result in the stated scope.}
\GSfield{Assessment}{literal statement: false as written; publication match: exact; verification: read the result; development: none; confidence: high}
\end{GSitem}
\begin{GSitem}{1977-429.3}{E}{Separation of disjoint central supports}
\GSfield{Kind; source}{theorem; \GSpage{P000721}}
\GSfield{The notes}{The notes repeat that modules with separated support have no morphisms between them and split into single-support pieces.}
\GSfield{Statement}{If \(\operatorname{supp}V\cap\operatorname{supp}W=\varnothing\), then \(\operatorname{Hom}_{(\mathfrak g,K)}(V,W)=0\); the primary decomposition of the central action is canonical and finite for each object.}
\GSfield{Fate}{The same construction is used in \citet[\S2]{bernsteingelfandgelfand1978local1} and recurs here as preparation for a local equivalence.}
\GSfield{Present status}{As of 10 November 2026, this is an established result in the stated scope.}
\GSfield{Assessment}{literal statement: correct; publication match: exact; verification: read the result; development: none; confidence: high}
\end{GSitem}
\begin{GSitem}{1977-429.4}{E}{Equivalence of categories}
\GSfield{Kind; source}{definition; \GSpage{P000721-P000722}}
\GSfield{The notes}{The notes explicitly say \ru{не совпад., а изоморфно!}: the composites are isomorphic to the identity, and the diagrams commute.}
\GSfield{Statement}{An equivalence consists of functors \(F:\mathcal C\to\mathcal D\), \(G:\mathcal D\to\mathcal C\), and natural isomorphisms \(GF\cong\operatorname{Id}_{\mathcal C}\), \(FG\cong\operatorname{Id}_{\mathcal D}\). Equality on the nose is unnecessary.}
\GSfield{Fate}{This is exactly the standard categorical meaning of equivalence used in \citet[\S3]{bernsteingelfandgelfand1978local1}; the recovered complete scan of P000722 confirms the distinction.}
\GSfield{Present status}{As of 10 November 2026, this is an established result in the stated scope.}
\GSfield{Assessment}{literal statement: correct; publication match: partial; verification: read the result; development: none; confidence: high}
\end{GSitem}
\begin{GSitem}{1977-429.5}{E}{Morita equivalence for a full matrix algebra}
\GSfield{Kind; source}{construction; \GSpage{P000722}}
\GSfield{The notes}{Finite-dimensional vector spaces are identified with modules over \(M_2(\mathbb C)\), by tensoring with the defining representation and recovering a multiplicity space.}
\GSfield{Statement}{The functors \(E\mapsto\mathbb C^2\otimes E\) and \(V\mapsto\operatorname{Hom}_{M_2(\mathbb C)}(\mathbb C^2,V)\) are mutually quasi-inverse equivalences between finite-dimensional vector spaces and finite-dimensional left \(M_2(\mathbb C)\)-modules.}
\GSfield{Fate}{The notebook gives an elementary Morita example; the general equivalence setting is that of \citet[\S3]{bernsteingelfandgelfand1978local1}. The recovered complete P000722 scan confirms the two categories.}
\GSfield{Present status}{As of 10 November 2026, this is an established result in the stated scope.}
\GSfield{Assessment}{literal statement: correct; publication match: thematic only; verification: read the result; development: none; confidence: high}
\end{GSitem}
\begin{GSitem}{1977-429.6}{A}{Nice analytic algebras}
\GSfield{Kind; source}{definition; \GSpage{P000722}}
\GSfield{The notes}{A \ru{приятная} algebra lies inside matrix-valued analytic functions and has the full matrix algebra as generic fiber.}
\GSfield{Statement}{Let \(E\) be a connected complex manifold and \(B\) a finite-dimensional semisimple algebra. A nice algebra is an \(\mathcal O(E)\)-subalgebra \(R\subset B\otimes\mathcal O(E)\), containing scalar functions and becoming \(B\otimes\mathcal M(E)\) after extension to the field of meromorphic functions; the published definition includes the required finiteness over \(\mathcal O(E)\).}
\GSfield{Fate}{\citet[\S2]{bernsteingelfandgelfand1978local2} give the definition and distinguish the stronger freeness condition called strongly nice.}
\GSfield{Present status}{As of 10 November 2026, this is an established result in the stated scope.}
\GSfield{Assessment}{literal statement: correct with missing hypotheses; publication match: exact; verification: read the result; development: none; confidence: high}
\end{GSitem}
\begin{GSitem}{1977-429.7}{A}{The local nice-algebra theorem}
\GSfield{Kind; source}{theorem; \GSpage{P000722}}
\GSfield{The notes}{The category is described by a nice algebra on a domain with dimension equal to the real split rank.}
\GSfield{Statement}{If \(\theta\) is a regular infinitesimal character, a sufficiently small neighborhood \(D\) of \(\theta\) has \(\mathcal C_D\) equivalent to the finite-dimensional modules of a nice analytic algebra. The parameter manifold has dimension equal to the real split rank; in the split case the formulation uses \(E=D\).}
\GSfield{Reconstruction}{Regularity and the sufficiently small neighborhood are essential hypotheses supplied from the published theorem, not visible in the notebook formulation.}
\GSfield{Fate}{\citet[\S2, Theorem~1]{bernsteingelfandgelfand1978local2} prove this with the regular-character and local-neighborhood hypotheses. The theorem is not stated there for arbitrary singular characters or arbitrary bounded domains.}
\GSfield{Present status}{As of 10 November 2026, this is an established result in the stated scope.}
\GSfield{Assessment}{literal statement: correct with missing hypotheses; publication match: exact; verification: read the result; development: none; confidence: high; problem \hyperref[prob:GP-1977-03]{GP-1977-03}}
\end{GSitem}
\begin{GSitem}{1977-429.8}{E}{The sl2 Casimir has a sign inconsistency}
\GSfield{Kind; source}{claim; \GSpage{P000722}}
\GSfield{The notes}{The quadratic central element is written with a negative \(H^2\) term.}
\GSfield{Statement}{For \([H,E_+]=2E_+\), \([H,E_-]=-2E_-\), \([E_+,E_-]=H\), the element \(C=4E_-E_++H^2+2H\) is central. Replacing \(+H^2\) by \(-H^2\) does not yield a central element in these conventions.}
\GSfield{Reconstruction}{The correction uses the positive standard matrix generators on the recovered complete P000722 scan and the differential realization on P000728. The recovered scan independently confirms the minus sign before \(H^2\).}
\GSfield{Fate}{The Casimir and rank-one principal-series facts here are classical background, not first results of the 1978 local-category papers. The local parameter in \citet[\S5]{bernsteingelfandgelfand1978local2} is a shifted normalized central character. The explicit sign diagnosis follows directly from the Lie brackets and is an editorial calculation, not a claim that the published paper printed the notebook formula.}
\GSfield{Present status}{As of 10 November 2026, this is an established result in the stated scope.}
\GSfield{Assessment}{literal statement: false as written; publication match: partial; verification: read the result; development: annotated; confidence: high}
\end{GSitem}
\begin{GSitem}{1977-429.9}{A}{Generators of the rank-one matrix order}
\GSfield{Kind; source}{construction; \GSpage{P000723}}
\GSfield{The notes}{Nine generators are given as matrix units, with selected off-diagonal units multiplied by the local coordinate \(z\).}
\GSfield{Statement}{If \(R\subset M_3(\mathcal O(D))\) is the indicated order, its \(\mathcal O(D)\)-module generators are diagonal matrix units, ordinary matrix units in unrestricted positions, and \(z e_{ij}\) in positions required to vanish at \(z=0\). Multiplication is inherited from \(e_{ij}e_{kl}=\delta_{jk}e_{il}\).}
\GSfield{Fate}{This is the explicit order of \citet[\S5]{bernsteingelfandgelfand1978local2}, up to the notebook\textquotesingle s changing row/column conventions. It determines all relations without relying on the inconsistent arrow labels.}
\GSfield{Present status}{As of 10 November 2026, this is an established result in the stated scope.}
\GSfield{Assessment}{literal statement: correct with missing hypotheses; publication match: exact; verification: read the result; development: none; confidence: high}
\end{GSitem}
\begin{GSitem}{1977-429.10}{A}{The three-vertex quiver and its relations}
\GSfield{Kind; source}{construction; \GSpage{P000723-P000724}}
\GSfield{The notes}{Three idempotent spaces \(L_1,L_2,L_3\), arrows \(e_{12},e_{21},e_{13},e_{31}\), product relations and a nilpotent central coordinate are drawn.}
\GSfield{Statement}{For a left module over the matrix order, \(e_{ij}\) maps the \(j\)-th idempotent space to the \(i\)-th. The arrows between the central vertex and the other two obey the matrix-unit multiplication relations, with the relevant two-step products equal to the central parameter and mixed products determined by the remaining generators. For the single-character block, that parameter acts nilpotently.}
\GSfield{Fate}{\citet[\S5]{bernsteingelfandgelfand1978local2} verify the order and hence its module presentation. The right-hand notebook diagram reverses the \(e_{13}\) and \(e_{31}\) directions relative to this convention, while retaining the other pair; it cannot be repaired by one global transpose.}
\GSfield{Present status}{As of 10 November 2026, this is an established result in the stated scope.}
\GSfield{Assessment}{literal statement: false as written; publication match: partial; verification: read the result; development: annotated; confidence: high}
\end{GSitem}
\begin{GSitem}{1977-429.11}{A}{The rank-two order for SL3(R)}
\GSfield{Kind; source}{construction; \GSpage{P000724}}
\GSfield{The notes}{The algebra \(M_4\oplus M_3\oplus M_3\oplus\mathbb C\) is restricted on the three walls \(z_1=0\), \(z_2=0\), \(z_1+z_2=0\) by vanishing and identification of matrix entries.}
\GSfield{Statement}{Near the regular character of the trivial representation of \(\operatorname{SL}_3(\mathbb R)\), the local category is represented by a nice order in \((M_4\oplus M_3\oplus M_3\oplus\mathbb C)\otimes\mathcal O(D)\). On \(z_1=0\), \(a_{13}=a_{14}=a_{23}=a_{24}=b_{12}=b_{13}=0\) and the lower right \(2\times2\) blocks of \(a,b\) agree; on \(z_2=0\), \(a_{12}=a_{14}=a_{32}=a_{34}=c_{12}=c_{13}=0\), with \(a_{22}=c_{22},a_{24}=c_{23},a_{42}=c_{32},a_{44}=c_{33}\); on \(z_1+z_2=0\), \(b_{13}=b_{23}=c_{13}=c_{23}=0\), \(b_{33}=c_{33}\).}
\GSfield{Fate}{\citet[\S6, case \(\operatorname{SL}_3(\mathbb R)\)]{bernsteingelfandgelfand1978local2} publish these explicit wall conditions. This is an exact match for the formulas, apart from the independent notebook diagram convention.}
\GSfield{Present status}{As of 10 November 2026, this is an established result in the stated scope.}
\GSfield{Assessment}{literal statement: correct with missing hypotheses; publication match: exact; verification: read the result; development: none; confidence: high}
\end{GSitem}

\GStalk{430}{1977-10-24}{S. I. Gelfand}{P000725-P000729}{Normalizer quotients and principal-series calculations}
The lecture constructs the algebra acting on finitely many \(K\)-types as a normalizer quotient and illustrates its polynomial realization through the principal series of \(\operatorname{PSL}_2(\mathbb R)\). The notes contain a Casimir-sign error and inconsistent parameter conventions; these are kept separate from the later published theorem. The proposed computations for additional groups are a research program, without a specified central block.

\begin{GSitem}{1977-430.1}{E}{Compact weights in the PSL2(R) example}
\GSfield{Kind; source}{definition; \GSpage{P000725-P000726}}
\GSfield{The notes}{For \(K\simeq\operatorname{SO}(2)\), the operator \(X_0=\left(\begin{smallmatrix}0&1\\-1&0\end{smallmatrix}\right)\) is said to have eigenvalues in \(2\mathbb Z\); admissibility, finite generation and bounded multiplicity are discussed.}
\GSfield{Statement}{For the displayed real infinitesimal generator, a complex algebraic \(K\)-module has eigenvalues in \(2i\mathbb Z\) when the central element \(-I\) acts trivially. Integer weights are eigenvalues of \(X_0/i\). Harish--Chandra modules have finite multiplicities and, in this rank-one setting, bounded multiplicities; the conventions must be fixed before writing the truncating polynomial.}
\GSfield{Fate}{\citet[\S5]{bernsteingelfandgelfand1978local2} use integral labels for the \(K\)-types without confusing the normalized weight operator with the real compact generator.}
\GSfield{Present status}{As of 10 November 2026, this is an established result in the stated scope.}
\GSfield{Assessment}{literal statement: false as written; publication match: partial; verification: read the result; development: none; confidence: high}
\end{GSitem}
\begin{GSitem}{1977-430.2}{E}{The Casimir and bounded central support}
\GSfield{Kind; source}{claim; \GSpage{P000726}}
\GSfield{The notes}{The notes again write \(4E_-E_+-H^2+2H\) as central and use its eigenvalues to localize the category.}
\GSfield{Statement}{For the displayed standard \(\mathfrak{sl}_2\) generators the central element is \(C=4E_-E_++H^2+2H\). Fixing a bounded region of its generalized spectrum is a rank-one instance of restricting central support.}
\GSfield{Fate}{The Casimir and rank-one principal-series facts here are classical background, not first results of the 1978 local-category papers. The central-support theorem is \citet[\S3, Theorem~1]{bernsteingelfandgelfand1978local1}. The sign error is visible on the complete scan of P000726 and is a faithful notebook error.}
\GSfield{Present status}{As of 10 November 2026, this is an established result in the stated scope.}
\GSfield{Assessment}{literal statement: false as written; publication match: partial; verification: read the result; development: annotated; confidence: high}
\end{GSitem}
\begin{GSitem}{1977-430.3}{E}{Morita functors in arbitrary matrix size}
\GSfield{Kind; source}{construction; \GSpage{P000726-P000727}}
\GSfield{The notes}{The equivalence with \(M_\ell(\mathbb C)\)-modules is written using tensor product and \(\operatorname{Hom}\).}
\GSfield{Statement}{For finite-dimensional left \(M_\ell(\mathbb C)\)-modules, evaluation gives a natural isomorphism \(\mathbb C^\ell\otimes\operatorname{Hom}_{M_\ell}(\mathbb C^\ell,V)\cong V\); these functors are inverse equivalences.}
\GSfield{Fate}{This is elementary Morita equivalence, used as the model for the functor in \citet[\S4]{bernsteingelfandgelfand1978local1}.}
\GSfield{Present status}{As of 10 November 2026, this is an established result in the stated scope.}
\GSfield{Assessment}{literal statement: correct; publication match: thematic only; verification: read the result; development: none; confidence: high}
\end{GSitem}
\begin{GSitem}{1977-430.4}{E}{A left ideal selecting finitely many compact weights}
\GSfield{Kind; source}{construction; \GSpage{P000727}}
\GSfield{The notes}{A left ideal \(J\) is generated by \(\prod_{n=-N}^{N}(X_0-2n)\).}
\GSfield{Statement}{With the normalized generator \(T=X_0/i\), let \(p_N(T)=\prod_{n=-N}^{N}(T-2n)\) and \(J=U(\mathfrak g)p_N(T)\). Then \(\operatorname{Hom}_{U(\mathfrak g)}(U(\mathfrak g)/J,V)\) identifies with \(\ker p_N(T)\), the sum of the selected \(K\)-weight spaces.}
\GSfield{Fate}{The module/ideal identity in this item is elementary algebraic background; the subsequent block-equivalence theorem is a separate new assertion. This is the rank-one version of the finite-\(K\)-type construction in \citet[\S4]{bernsteingelfandgelfand1978local1}. The factor \(i\) is required if one retains the notebook\textquotesingle s real generator.}
\GSfield{Present status}{As of 10 November 2026, this is an established result in the stated scope.}
\GSfield{Assessment}{literal statement: false as written; publication match: partial; verification: read the result; development: none; confidence: high}
\end{GSitem}
\begin{GSitem}{1977-430.5}{E}{The normalizer quotient algebra}
\GSfield{Kind; source}{construction; \GSpage{P000727}}
\GSfield{The notes}{The normalizer is called the \ru{наименьшая} subalgebra containing \(J\) in which \(J\) is two-sided, and \(Q=N/J\).}
\GSfield{Statement}{For a left ideal \(J\subset U\), the right normalizer \(N_U(J)=\{u\in U:Ju\subset J\}\) is the largest subalgebra of \(U\) containing \(J\) in which \(J\) is a two-sided ideal. The quotient \(N_U(J)/J\) is an algebra acting on vectors annihilated by \(J\), with the chosen left/right convention.}
\GSfield{Fate}{The module/ideal identity in this item is elementary algebraic background; the subsequent block-equivalence theorem is a separate new assertion. \citet[\S4]{bernsteingelfandgelfand1978local1} use this normalizer construction. The scan explicitly says smallest; this is a faithful mathematical error, not a translation issue.}
\GSfield{Present status}{As of 10 November 2026, this is an established result in the stated scope.}
\GSfield{Assessment}{literal statement: false as written; publication match: exact; verification: read the result; development: annotated; confidence: high}
\end{GSitem}
\begin{GSitem}{1977-430.6}{A}{The finite-K-type functor}
\GSfield{Kind; source}{construction; \GSpage{P000727}}
\GSfield{The notes}{The functor \(\tau\) retains the weight spaces with indices \(-N,\ldots,N\); \(Q\) acts on their sum.}
\GSfield{Statement}{For an admissible \((\mathfrak g,K)\)-module \(V\), the selected sum \(\tau_N(V)=\bigoplus_{|n|\leq N}V_{2n}\) is finite-dimensional. The normalizer quotient acts because \(Jv=0\) implies \(Juv=0\) for \(u\in N_U(J)\).}
\GSfield{Fate}{\citet[\S4]{bernsteingelfandgelfand1978local1} define the corresponding functor for a finite set of \(K\)-types.}
\GSfield{Present status}{As of 10 November 2026, this is an established result in the stated scope.}
\GSfield{Assessment}{literal statement: correct; publication match: exact; verification: read the result; development: none; confidence: high}
\end{GSitem}
\begin{GSitem}{1977-430.7}{A}{A sufficiently large finite set of K-types detects the block}
\GSfield{Kind; source}{theorem; \GSpage{P000727}}
\GSfield{The notes}{For bounded \(D\), taking sufficiently many compact weights is asserted to give the desired equivalence.}
\GSfield{Statement}{For a bounded central-support region \(D\), a sufficiently large finite set of \(K\)-types can be chosen such that the associated normalizer-quotient functor is an equivalence on \(\mathcal C_D\); the algebra is finite over the center.}
\GSfield{Fate}{\citet[\S4, Proposition~1]{bernsteingelfandgelfand1978local1} prove this and thereby establish their main theorem. A finite set independent of an arbitrary unbounded support region is not asserted.}
\GSfield{Present status}{As of 10 November 2026, this is an established result in the stated scope.}
\GSfield{Assessment}{literal statement: correct with missing hypotheses; publication match: exact; verification: read the result; development: none; confidence: high; problem \hyperref[prob:GP-1977-03]{GP-1977-03}}
\end{GSitem}
\begin{GSitem}{1977-430.8}{E}{Homogeneous-function principal series}
\GSfield{Kind; source}{construction; \GSpage{P000728}}
\GSfield{The notes}{Functions on \(\mathbb R^2\setminus\{0\}\) homogeneous of degree \(s\) realize \(T^s\); on the circle, even trigonometric polynomials give the \(K\)-finite part.}
\GSfield{Statement}{For \(f(tv)=|t|^s f(v)\), the even principal-series representation restricts to even Fourier modes on the circle. Differentiation gives \(E_+=x\partial_y\), \(E_-=y\partial_x\), \(H=x\partial_x-y\partial_y\), up to the overall action convention explicitly chosen.}
\GSfield{Fate}{The principal-series family underlies \citet[\S5, Proposition~2]{bernsteingelfandgelfand1978local1} and the rank-one calculation of \citet[\S5]{bernsteingelfandgelfand1978local2}.}
\GSfield{Present status}{As of 10 November 2026, this is an established result in the stated scope.}
\GSfield{Assessment}{literal statement: correct with missing hypotheses; publication match: exact; verification: read the result; development: none; confidence: high}
\end{GSitem}
\begin{GSitem}{1977-430.9}{A}{Polynomial dependence on the principal-series parameter}
\GSfield{Kind; source}{method; \GSpage{P000728}}
\GSfield{The notes}{The matrix elements of the Lie algebra action in selected compact types are polynomial in \(s\), producing an algebra of matrix polynomials.}
\GSfield{Statement}{After a fixed compact-picture realization and choice of finitely many \(K\)-types, the action of the normalizer quotient in principal-series families depends polynomially on the inducing parameter. The resulting representation embeds the quotient into a finite-dimensional endomorphism algebra tensored with the coordinate ring of the parameter space.}
\GSfield{Fate}{\citet[\S5, lemma and Theorem~2]{bernsteingelfandgelfand1978local1} prove precisely this polynomial realization.}
\GSfield{Present status}{As of 10 November 2026, this is an established result in the stated scope.}
\GSfield{Assessment}{literal statement: correct with missing hypotheses; publication match: exact; verification: read the result; development: none; confidence: high}
\end{GSitem}
\begin{GSitem}{1977-430.10}{E}{Subquotients of principal series}
\GSfield{Kind; source}{claim; \GSpage{P000728}}
\GSfield{The notes}{The notes suggest that every Harish--Chandra module appears as a subquotient of a principal-series representation.}
\GSfield{Statement}{Every irreducible Harish--Chandra module is a subquotient of a suitable principal-series representation. This assertion does not say that every finite-length module, including nonsplit extensions with nonsemisimple central action, is a subquotient of one ordinary principal series with a fixed scalar infinitesimal character.}
\GSfield{Fate}{\citet[\S5, Proposition~2(c)]{bernsteingelfandgelfand1978local1} state the irreducible result. The stronger literal formulation in the notes is unsupported and fails for modules with nonsemisimple central action.}
\GSfield{Present status}{As of 10 November 2026, this is an established result in the stated scope.}
\GSfield{Assessment}{literal statement: false as written; publication match: exact; verification: read the result; development: annotated; confidence: high}
\end{GSitem}
\begin{GSitem}{1977-430.11}{E}{Casimir eigenvalues in the homogeneous realization}
\GSfield{Kind; source}{claim; \GSpage{P000728}}
\GSfield{The notes}{The eigenvalue is written \(-s^2+2s\), and reducibility is discussed at \(s=0,2\).}
\GSfield{Statement}{For the displayed differential generators and homogeneity \((x\partial_x+y\partial_y)f=sf\), the correctly signed Casimir acts as \(s(s+2)\). Replacing the parameter by its negative gives \(s(s-2)\). The notebook\textquotesingle s scalar and reducibility parameter cannot both be used with its stated homogeneity convention.}
\GSfield{Fate}{The Casimir and rank-one principal-series facts here are classical background, not first results of the 1978 local-category papers. The published treatment \citet[\S5]{bernsteingelfandgelfand1978local2} instead uses a normalized parameter with central coordinate \(z=s^2-1\). The present identity is a direct editorial computation, not an attribution of the notebook formula to that paper.}
\GSfield{Present status}{As of 10 November 2026, this is an established result in the stated scope.}
\GSfield{Assessment}{literal statement: false as written; publication match: partial; verification: read the result; development: annotated; confidence: high}
\end{GSitem}
\begin{GSitem}{1977-430.12}{E}{The reducible principal series at the trivial character}
\GSfield{Kind; source}{construction; \GSpage{P000728-P000729}}
\GSfield{The notes}{At \(s=0\), the constant function gives the trivial constituent and two remaining pieces are indicated; a second, dual reducibility point is shown at \(s=2\).}
\GSfield{Statement}{At homogeneity degree zero, the even principal series contains the constants and has the two lowest/highest-weight discrete-series constituents as quotient. The contragredient extension occurs at degree \(-2\), not \(+2\), for this same homogeneity convention. The sequence is generally nonsplit.}
\GSfield{Fate}{The Casimir and rank-one principal-series facts here are classical background, not first results of the 1978 local-category papers. \citet[\S5]{bernsteingelfandgelfand1978local2} describe the three constituents and compute the local order from their selected \(K\)-types. Their normalized parameter differs from the notebook\textquotesingle s degree.}
\GSfield{Present status}{As of 10 November 2026, this is an established result in the stated scope.}
\GSfield{Assessment}{literal statement: false as written; publication match: partial; verification: read the result; development: none; confidence: high}
\end{GSitem}
\begin{GSitem}{1977-430.13}{E}{Generic Weyl reflection of the principal-series parameter}
\GSfield{Kind; source}{theorem; \GSpage{P000729}}
\GSfield{The notes}{The notes relate \(T^s\) and \(T^{2-s}\) away from exceptional parameters.}
\GSfield{Statement}{For the homogeneous degree convention above, the generic intertwining equivalence is between degrees \(s\) and \(-s-2\). With the opposite parameter convention it is \(s\leftrightarrow2-s\); reducibility points require separate treatment, and a generic isomorphism need not extend as an invertible holomorphic family there.}
\GSfield{Fate}{\citet[\S5, Proposition~2(b)]{bernsteingelfandgelfand1978local1} give the generic Weyl-orbit equivalence; \citet[\S5]{bernsteingelfandgelfand1978local2} exploit it in the explicit order.}
\GSfield{Present status}{As of 10 November 2026, this is an established result in the stated scope.}
\GSfield{Assessment}{literal statement: correct with missing hypotheses; publication match: exact; verification: read the result; development: none; confidence: high}
\end{GSitem}
\begin{GSitem}{1977-430.14}{A}{Three compact types near the trivial character}
\GSfield{Kind; source}{construction; \GSpage{P000729}}
\GSfield{The notes}{The three weights \(-1,0,1\) are enough for the explicit algebra near the trivial character.}
\GSfield{Statement}{With the normalized integral labels for the even \(K\)-types, retaining labels \(-1,0,1\) gives the local rank-three order for a sufficiently small neighborhood of the trivial infinitesimal character.}
\GSfield{Fate}{\citet[\S5]{bernsteingelfandgelfand1978local2} explicitly choose these three types and compute the order.}
\GSfield{Present status}{As of 10 November 2026, this is an established result in the stated scope.}
\GSfield{Assessment}{literal statement: correct; publication match: exact; verification: read the result; development: none; confidence: high}
\end{GSitem}
\begin{GSitem}{1977-430.15}{F}{Further groups proposed for explicit computation}
\GSfield{Kind; source}{remark; \GSpage{P000729}}
\GSfield{The notes}{After listing rank-one groups and \(\operatorname{SL}_3(\mathbb R)\), the notes suggest \(\operatorname{SL}_3(\mathbb C)\), \(\operatorname{SU}(2,1)\), \(\operatorname{SO}(4,1)\), \(\operatorname{Sp}_4(\mathbb C)\), and the harder \(\operatorname{SL}_4\).}
\GSfield{Statement}{This is a program to compute explicit local Harish--Chandra orders for additional groups. The central character, real form in the last example, neighborhood and equivalence convention are not all specified, so it is not one definite classification problem.}
\GSfield{Fate}{\citet[\S6]{bernsteingelfandgelfand1978local2} publish the listed low-rank examples, but do not identify or solve a single additional problem with all the notebook\textquotesingle s proposed groups and unspecified blocks.}
\GSfield{Present status}{As of 10 November 2026, the program is identifiable; the notes do not determine an exact additional problem whose full fate can be asserted.}
\GSfield{Assessment}{literal statement: not determinable; publication match: partial; verification: read the result; development: none; confidence: medium}
\end{GSitem}

\GStalk{431}{1977-10-24}{I. M. Krichever}{P000729-P000730}{Commuting differential operators and their earlier history}
Krichever relates the recently discussed spectral-curve method to the earlier papers of Burchnall, Chaundy and Baker. The short intervention distinguishes the existence of algebraic relations, the geometric classification of commuting rings, and explicit theta-function formulas. An unnamed historical reference cannot be securely identified.

\begin{GSitem}{1977-431.1}{E}{An algebraic relation between commuting ordinary differential operators}
\GSfield{Kind; source}{theorem; \GSpage{P000729-P000730}}
\GSfield{The notes}{Krichever reports that the problem \([L_1,L_2]=0\) was treated by Burchnall and Chaundy in papers listed as 1922, 1928 and 1931.}
\GSfield{Statement}{If scalar ordinary differential operators \(L_1,L_2\) with locally analytic coefficients commute, their orders are positive and their leading coefficients are normalized to one, then there is a nonzero polynomial \(Q\in\mathbb C[\lambda,\mu]\) such that \(Q(L_1,L_2)=0\). Coprimality of the orders is not required for this algebraic-dependence statement.}
\GSfield{Fate}{\citet[Theorem~1.1 and its proof]{krichever1978commutative} explicitly attribute and prove the Burchnall--Chaundy theorem. This confirms an old result, not a new theorem of the 1977 seminar. The earlier occurrence 1975-374.2 concerns the same spectral-curve line.}
\GSfield{Present status}{As of 10 November 2026, this is an established result in the stated scope.}
\GSfield{Assessment}{literal statement: correct with missing hypotheses; publication match: exact; verification: read the result; development: none; confidence: high}
\end{GSitem}
\begin{GSitem}{1977-431.2}{E}{The rank-one spectral-curve description}
\GSfield{Kind; source}{theorem; \GSpage{P000729-P000730}}
\GSfield{The notes}{The older work is said to have posed and solved the classification for scalar operators of relatively prime orders, using algebraic geometry.}
\GSfield{Statement}{A commuting rank-one ring of scalar ordinary differential operators is encoded by an algebraic spectral curve, a smooth marked point at infinity, a local parameter and rank-one sheaf data with a local trivialization. On a smooth curve in generic position this becomes a nonspecial divisor or line bundle; singular curves require torsion-free sheaves and cannot be omitted from an unrestricted classification.}
\GSfield{Reconstruction}{The marked-point and sheaf data spell out what a modern classification needs; they are not all enumerated in this short announcement.}
\GSfield{Fate}{\citet[introduction and \S1]{krichever1978commutative} confirm the local classification by Burchnall--Chaundy and distinguish it from the later effective reconstruction. This follows earlier seminar items 1975-374.5 and 1975-380.2; the latter\textquotesingle s smooth-curve-only scope remains a real qualification. The rank-one pointed-curve data also recur in 1975-383.9--383.12 and 1975-385.4--385.8, associated with GP-1975-14. The present historical classification remark does not itself restate the earlier inverse-construction question; GP-1975-26 concerns higher rank and is not an exact recurrence here.}
\GSfield{Present status}{As of 10 November 2026, this is an established result in the stated scope.}
\GSfield{Assessment}{literal statement: correct with missing hypotheses; publication match: partial; verification: read the result; development: none; confidence: high}
\end{GSitem}
\begin{GSitem}{1977-431.3}{E}{Theta formulas and linear motion on a Jacobian}
\GSfield{Kind; source}{method; \GSpage{P000729-P000730}}
\GSfield{The notes}{The historical construction is described as insufficiently effective: no theta functions, no straight-line Jacobian motion, and no recognition of conditional periodicity.}
\GSfield{Statement}{For nonsingular rank-one spectral data, the Baker--Akhiezer construction yields explicit theta-function formulas and linearizes the commuting flows on the Jacobian. Quasiperiodicity and smooth real solutions require additional reality and nonsingularity assumptions; they do not hold for arbitrary complex or singular data.}
\GSfield{Fate}{The underlying construction predates this talk: \citet[Theorem~1 and concluding discussion]{krichever1976zakharov}, received 8 October 1975 and published in 1976, construct the quasiperiodic solutions and announce theta formulas. \citet[\S1, pp.189--191]{krichever1977methods}, received 24 March 1977, give the explicit theta representation and Jacobian parameters. The 1978 restatement \citet[introduction and \S1]{krichever1978commutative} is the closest historical discussion, not the first appearance of the mathematical method.}
\GSfield{Present status}{As of 10 November 2026, this is an established result in the stated scope.}
\GSfield{Assessment}{literal statement: correct with missing hypotheses; publication match: partial; verification: read the result; development: none; confidence: high}
\end{GSitem}
\begin{GSitem}{1977-431.4}{F}{Unspecified historical reference and almost-commuting pairs}
\GSfield{Kind; source}{remark; \GSpage{P000730}}
\GSfield{The notes}{A 1939 paper is said to omit the older references, followed by the fragment that almost-commuting operators and Lax pairs exist.}
\GSfield{Statement}{The author, title and equation of the 1939 paper are absent. The phrase almost commuting specifies neither a commutator nor a deformation parameter, so no definite mathematical proposition can be reconstructed from it.}
\GSfield{Fate}{Comparison with \citet[introduction and bibliography]{krichever1978commutative} identifies the Burchnall--Chaundy line but does not resolve these two notebook fragments. No exact later result is assigned to them.}
\GSfield{Present status}{As of 10 November 2026, these fragments remain unidentified; this is a source-identification limit, not evidence of an undeveloped mathematical problem.}
\GSfield{Assessment}{literal statement: not determinable; publication match: none found; verification: unresolved lead; development: none; confidence: low}
\end{GSitem}

\GStalk{432}{1977-10-31}{S. P. Novikov}{P000730-P000734}{Instantons, their moduli and the twistor construction}
Novikov surveys finite-action Yang--Mills fields, topological charge, self-duality, dimension and connectedness questions for instanton moduli, and the twistor interpretation. Several numerical and convention-dependent assertions need qualification. The connectedness question is a definite problem, subsequently solved, whereas classification of all Yang--Mills critical points is broader than classification of instantons.

\begin{GSitem}{1977-432.1}{E}{Connections, curvature and local pure gauges}
\GSfield{Kind; source}{definition; \GSpage{P000730}}
\GSfield{The notes}{The connection form \(A=A_\mu dx^\mu\), curvature \(dA+[A,A]\), gauge transformation and the flat-connection form are introduced; chiral fields are mentioned.}
\GSfield{Statement}{For a compact matrix group, \(F_A=dA+A\wedge A\), or \(dA+\frac12[A\wedge A]\) when the bracket notation includes both antisymmetrized terms. A smooth flat connection is locally gauge-equivalent to the trivial connection, and globally so on a simply connected domain. A group-valued field produces a flat Maurer--Cartan form; a chiral field equation imposes additional conditions.}
\GSfield{Fate}{This is background to the gauge-theoretic constructions of \citet[\S1]{drinfeldmanin1978selfdual} and the definitions in \citet[\S2.1]{dayapremawaldron2026parabolic}. The earlier occurrence 1975-386.1 treats Yang--Mills connections and action.}
\GSfield{Present status}{As of 10 November 2026, this is an established result in the stated scope.}
\GSfield{Assessment}{literal statement: correct with missing hypotheses; publication match: exact; verification: read the result; development: none; confidence: high}
\end{GSitem}
\begin{GSitem}{1977-432.2}{E}{Topological charge requires a trace and normalization}
\GSfield{Kind; source}{definition; \GSpage{P000730}}
\GSfield{The notes}{The integral of \(\varepsilon_{abcd}F^{ab}F^{cd}\) is said to be an integer up to a factor, without a written trace.}
\GSfield{Statement}{For an \(\operatorname{SU}(2)\)-connection extending over \(S^4\), use anti-Hermitian matrices in the defining representation and define \(k=-(8\pi^2)^{-1}\int_{S^4}\operatorname{tr}(F\wedge F)\in\mathbb Z\), with the sign of the integer fixed by orientation. A matrix-valued integral without the invariant trace is not a characteristic number.}
\GSfield{Fate}{\citet[\S2.4]{dayapremawaldron2026parabolic} state a precise characteristic-number convention. The notebook\textquotesingle s coefficient cannot be compared without fixing the trace and component normalization.}
\GSfield{Present status}{As of 10 November 2026, this is an established result in the stated scope.}
\GSfield{Assessment}{literal statement: false as written; publication match: partial; verification: read the result; development: annotated; confidence: high}
\end{GSitem}
\begin{GSitem}{1977-432.3}{E}{Yang–Mills variation and the abelian case}
\GSfield{Kind; source}{definition; \GSpage{P000730}}
\GSfield{The notes}{The action is the integral of \(\operatorname{Sp}(F_{\mu\nu}F^{\mu\nu})\), and the equation is \(\delta S=0\). Electrodynamics is assigned \(SU(1)\) or \(\mathbb R\).}
\GSfield{Statement}{With a positive invariant inner product on the Lie algebra, the action is a positive multiple of \(\int|F_A|^2\), and its Euler--Lagrange equation is \(d_A^*F_A=0\). The compact abelian gauge group is \(U(1)\), whose Lie algebra is one-dimensional; \(SU(1)\) is trivial.}
\GSfield{Fate}{\citet[\S1]{dayapremawaldron2026parabolic} state the action and equation. The \(SU(1)\) reading is visible on P000730 and is a faithful notebook error.}
\GSfield{Present status}{As of 10 November 2026, this is an established result in the stated scope.}
\GSfield{Assessment}{literal statement: false as written; publication match: partial; verification: read the result; development: none; confidence: high}
\end{GSitem}
\begin{GSitem}{1977-432.4}{C}{Classifying all finite-action Yang–Mills critical points}
\GSfield{Kind; source}{problem; \GSpage{P000730}}
\GSfield{The notes}{The classical problem is to find all \(A\) satisfying the Yang--Mills equations and \(S(A)<\infty\).}
\GSfield{Statement}{Classify, modulo smooth gauge transformations and with a fixed compact gauge group, smooth Yang--Mills connections on \(\mathbb R^4\) of finite \(L^2\)-curvature. After removable-singularity compactification this is the classification of Yang--Mills critical connections on bundles over the round \(S^4\), not merely their self-dual minima.}
\GSfield{Fate}{The instanton part is explicitly classified by \citet[Theorem in \S4]{drinfeldmanin1978selfdual}. Nonminimal critical points exist, and \citet[\S1, Theorem~1.1]{dayapremawaldron2026parabolic} distinguish them from instantons and prove a low-energy deformation theorem. These results do not classify all finite-action critical points.}
\GSfield{Present status}{As of 10 November 2026, the instanton classification and substantial energy-gap results are established; the sources checked do not provide a full classification of all nonminimal Yang--Mills critical connections.}
\GSfield{Assessment}{literal statement: correct with missing hypotheses; publication match: partial; verification: read the result; development: none; confidence: medium; problem \hyperref[prob:GP-1977-04]{GP-1977-04}}
\GSfield{Worth developing}{A bounded energy or symmetry-restricted version is worth studying, because the full critical-point classification extends far beyond the instanton construction.}
\end{GSitem}
\begin{GSitem}{1977-432.5}{B}{Finite action and extension over the point at infinity}
\GSfield{Kind; source}{claim; \GSpage{P000730}}
\GSfield{The notes}{Requiring a field on \(S^4\) is called slightly stronger than finite energy on \(\mathbb R^4\).}
\GSfield{Statement}{For a smooth Yang--Mills connection with compact structure group in dimension four, finite \(\int|F_A|^2\) implies smooth extension, after a gauge change, across the point at infinity on a suitable bundle over \(S^4\). For an arbitrary connection not satisfying Yang--Mills, the same conclusion is not supplied by this theorem.}
\GSfield{Fate}{\citet[abstract and main removable-singularity result]{uhlenbeck1982removable} establish the equivalence for Yang--Mills fields. The result was announced in 1979 and fully published in 1982; the publisher page verifies both the result and the earlier announcement citation, although the full original proof was not accessible here.}
\GSfield{Present status}{As of 10 November 2026, this is an established result in the stated scope.}
\GSfield{Assessment}{literal statement: correct with missing hypotheses; publication match: exact; verification: publisher metadata checked; development: none; confidence: high}
\end{GSitem}
\begin{GSitem}{1977-432.6}{E}{Four-dimensional conformal invariance}
\GSfield{Kind; source}{theorem; \GSpage{P000730-P000731}}
\GSfield{The notes}{The Euclidean metric and round sphere are conformally equivalent; the conformal group is written \(O(5,1)\), with a light-cone construction.}
\GSfield{Statement}{In four dimensions the Hodge star on two-forms and the Yang--Mills action are conformally invariant. The conformal group of the round sphere is \(O(5,1)/\{\pm I\}\), with appropriate components for orientation preservation; \(S^4\) is the projectivized null cone in \(\mathbb R^{5,1}\). Pullback by smooth conformal diffeomorphisms preserves the equations. An arbitrary continuous map need not do so.}
\GSfield{Fate}{Conformal geometry is part of the construction in \citet{atiyahward1977instantons}; \citet[\S1]{drinfeldmanin1978selfdual} explicitly fix the standard conformal structure on \(S^4\).}
\GSfield{Present status}{As of 10 November 2026, this is an established result in the stated scope.}
\GSfield{Assessment}{literal statement: false as written; publication match: partial; verification: indirect; development: none; confidence: high}
\end{GSitem}
\begin{GSitem}{1977-432.7}{E}{The topological energy bound and self-dual minima}
\GSfield{Kind; source}{theorem; \GSpage{P000731}}
\GSfield{The notes}{The notes state \(S\geq n\), with a strict minimum at a self-dual field.}
\GSfield{Statement}{With action \(E=-\int\operatorname{tr}(F\wedge *F)\) and charge as above, \(E\geq8\pi^2|k|\); equality holds exactly for one of \(F=\pm *F\), with sign chosen according to \(k\). The minimum is not strict along a nontrivial instanton moduli space or along gauge orbits.}
\GSfield{Fate}{\citet[\S1 and \S2.4]{dayapremawaldron2026parabolic} state the minimizing role of instantons in a fixed bundle. The notebook suppresses the absolute value and normalization and overstates strictness.}
\GSfield{Present status}{As of 10 November 2026, this is an established result in the stated scope.}
\GSfield{Assessment}{literal statement: false as written; publication match: partial; verification: read the result; development: annotated; confidence: high}
\end{GSitem}
\begin{GSitem}{1977-432.8}{B}{Multi-instanton ansätze and the all-charge problem}
\GSfield{Kind; source}{construction; \GSpage{P000731}}
\GSfield{The notes}{The notes attribute a \(5n+4\)-parameter family to \textquotesingle t Hooft, state completeness for charge two, and contrast it with dimension \(8n-3\).}
\GSfield{Statement}{Explicit multi-instanton ansätze yield finite-dimensional subfamilies at each positive charge. Their raw parameter counts must be reduced by redundancies and gauge conventions: \(5n+4\) is not simultaneously a count of independent unframed moduli at \(n=2\), where \(8n-3=13\). Complete classification requires the more general linear algebra construction.}
\GSfield{Fate}{\citet[p.2662, paragraph preceding the theorem]{atiyahhitchinsinger1977deformations} explicitly attribute the \(5n+4\)-parameter family to Jackiw, Nohl and Rebbi, distinguishing it from the \(5n\)-parameter family of \textquotesingle t Hooft. Their theorem gives \(8n-3\) independent unframed moduli. \citet[\S4]{drinfeldmanin1978selfdual} subsequently give all instantons for the classical compact groups. The notebook conflates the ansatz attribution and raw versus independent parameter counts.}
\GSfield{Present status}{As of 10 November 2026, this is an established result in the stated scope.}
\GSfield{Assessment}{literal statement: false as written; publication match: partial; verification: read the result; development: none; confidence: high; problem \hyperref[prob:GP-1977-05]{GP-1977-05}}
\end{GSitem}
\begin{GSitem}{1977-432.9}{E}{Dimension of the SU2 instanton moduli space}
\GSfield{Kind; source}{theorem; \GSpage{P000731-P000732}}
\GSfield{The notes}{The dimension is \(8n-3\); an index calculation and removal of three gauge parameters are described.}
\GSfield{Statement}{For charge \(k\geq1\) and the round metric on \(S^4\), the moduli space of irreducible \(SU(2)\) anti-self-dual connections modulo the full gauge group is smooth of real dimension \(8k-3\). The framed moduli space has dimension \(8k\). At charge zero the trivial connection has a nontrivial stabilizer and the same formula does not apply.}
\GSfield{Fate}{\citet[theorem and proof, pp.2662--2663]{atiyahhitchinsinger1977deformations} prove smoothness and dimension \(8k-3\), using the self-dual deformation complex and vanishing of its obstruction space. Contributed 20 April 1977 and published in July 1977, this was available before the talk. Both pages of the primary scanned article were read.}
\GSfield{Present status}{As of 10 November 2026, this is an established result in the stated scope.}
\GSfield{Assessment}{literal statement: correct with missing hypotheses; publication match: exact; verification: read the result; development: none; confidence: high}
\end{GSitem}
\begin{GSitem}{1977-432.10}{E}{The elliptic deformation complex and obstruction vanishing}
\GSfield{Kind; source}{method; \GSpage{P000732}}
\GSfield{The notes}{The duality equation is linearized; the cokernel is said to vanish by Hodge-theoretic considerations.}
\GSfield{Statement}{For an anti-self-dual connection the deformation complex is \(\Omega^0(\operatorname{ad}E)\xrightarrow{d_A}\Omega^1(\operatorname{ad}E)\xrightarrow{d_A^+}\Omega^2_+(\operatorname{ad}E)\). On the round \(S^4\), the Weitzenböck formula gives \(H_A^2=0\); irreducibility gives \(H_A^0=0\), so the index is the local moduli dimension. Such cokernel vanishing is not automatic on arbitrary four-manifolds.}
\GSfield{Fate}{This is the deformation argument underlying \citet{atiyahhitchinsinger1977deformations}; \citet[proof in \S3.2]{dayapremawaldron2026parabolic} explicitly use the same vanishing from positive scalar curvature and zero self-dual Weyl tensor.}
\GSfield{Present status}{As of 10 November 2026, this is an established result in the stated scope.}
\GSfield{Assessment}{literal statement: correct with missing hypotheses; publication match: exact; verification: read the result; development: none; confidence: high}
\end{GSitem}
\begin{GSitem}{1977-432.11}{E}{Charge-one uniqueness and conformal symmetry}
\GSfield{Kind; source}{theorem; \GSpage{P000732}}
\GSfield{The notes}{For \(n=1\), all solutions are said to be spherically symmetric up to gauge; the five-dimensional moduli space and conformal stabilizer are discussed.}
\GSfield{Statement}{Every charge-one \(SU(2)\) instanton on the round four-sphere is gauge-equivalent to a conformal transform of the basic instanton. Its unframed moduli are center and scale, giving a five-dimensional space. This is symmetry after choosing the corresponding conformal center and scale, not invariance under all rotations about an arbitrarily fixed point.}
\GSfield{Fate}{\citet[p.2663, paragraph beginning with the conformal-group action]{atiyahhitchinsinger1977deformations} identify the charge-one moduli space with hyperbolic five-space, the orbit of the basic instanton. The same charge-one case is included in \citet[\S4]{drinfeldmanin1978selfdual}. The notebook\textquotesingle s group-dimension argument alone would not prove uniqueness, but the earlier primary theorem supplies the result.}
\GSfield{Present status}{As of 10 November 2026, this is an established result in the stated scope.}
\GSfield{Assessment}{literal statement: correct with missing hypotheses; publication match: exact; verification: read the result; development: none; confidence: high}
\end{GSitem}
\begin{GSitem}{1977-432.12}{B}{Connectedness of fixed-charge instanton moduli}
\GSfield{Kind; source}{question; \GSpage{P000731-P000732}}
\GSfield{The notes}{The notes ask whether the space of instanton solutions is connected, with charge two singled out as incompletely understood.}
\GSfield{Statement}{For each fixed positive charge \(k\), is the moduli space of \(SU(2)\) instantons on \(S^4\), modulo gauge, path-connected?}
\GSfield{Fate}{Taubes proved the \(SU(2)\) statement in 1984. \citet[Theorem~1.2 and \S3.3]{dayapremawaldron2026parabolic} give a new proof in 2026 and explicitly distinguish Taubes\textquotesingle s original rank restrictions from the general \(SU(r)\) extension. Thus the precise question is solved, not an open problem inherited from the notebook.}
\GSfield{Present status}{As of 10 November 2026, path-connectedness holds for every fixed-charge SU(2) instanton moduli space; the 2026 paper also states and proves the SU(r) version.}
\GSfield{Assessment}{literal statement: correct; publication match: exact; verification: read the result; development: none; confidence: high; problem \hyperref[prob:GP-1977-05]{GP-1977-05}}
\end{GSitem}
\begin{GSitem}{1977-432.13}{E}{A spectral-parameter formulation of self-duality}
\GSfield{Kind; source}{construction; \GSpage{P000732}}
\GSfield{The notes}{Two first-order covariant differential operators linear in \(\lambda\) are asserted to commute exactly when the duality equations hold; both displayed mixed terms have plus signs.}
\GSfield{Statement}{A consistent Euclidean pair, also recorded on P000741, is \(L_1=\lambda(\nabla_1-i\nabla_2)+(\nabla_3+i\nabla_4)\), \(L_2=\lambda(\nabla_3-i\nabla_4)-(\nabla_1+i\nabla_2)\). Vanishing of \([L_1,L_2]\) for every \(\lambda\) gives the three self-duality equations in the corresponding orientation. One cannot change the relative minus sign without changing the real structure or signature.}
\GSfield{Fate}{\citet[\S5 and references7--8]{drinfeldmanin1978selfdual} identify the Belavin--Zakharov construction used in the notes. The corrected explicit pair is checked by expanding the polynomial commutator, not read into the differently signed P000732 display.}
\GSfield{Present status}{As of 10 November 2026, this is an established result in the stated scope.}
\GSfield{Assessment}{literal statement: false as written; publication match: partial; verification: indirect; development: annotated; confidence: high}
\end{GSitem}
\begin{GSitem}{1977-432.14}{E}{Orthogonal complex structures and the twistor fibration}
\GSfield{Kind; source}{construction; \GSpage{P000732-P000733}}
\GSfield{The notes}{All linear complex structures on \(\mathbb R^4\) are said to be parameterized by one complex parameter, and \(S^7\to\mathbb CP^3\to S^4\) is drawn.}
\GSfield{Statement}{The orientation-compatible orthogonal complex structures on a Euclidean four-space form \(SO(4)/U(2)\cong S^2\cong\mathbb CP^1\). All linear complex structures form a larger space. The twistor projection \(\mathbb CP^3\to\mathbb HP^1\cong S^4\) has fiber \(\mathbb CP^1\); it factors the quaternionic Hopf map through the circle quotient \(S^7\to\mathbb CP^3\).}
\GSfield{Fate}{The twistor construction is the setting of \citet{atiyahward1977instantons} and \citet[\S1]{drinfeldmanin1978selfdual}. The notebook omits orthogonality and prints an inconsistent \(\mathbb CP^2\) fiber near the bottom of P000733.}
\GSfield{Present status}{As of 10 November 2026, this is an established result in the stated scope.}
\GSfield{Assessment}{literal statement: false as written; publication match: partial; verification: indirect; development: none; confidence: high}
\end{GSitem}
\begin{GSitem}{1977-432.15}{E}{Instantons as holomorphic bundles on twistor space}
\GSfield{Kind; source}{theorem; \GSpage{P000733}}
\GSfield{The notes}{The connection lifts to \(\mathbb CP^3\), defining an almost complex structure which is integrable precisely under the duality equations.}
\GSfield{Statement}{The pullback of a unitary anti-self-dual connection to twistor space defines an integrable \(\bar\partial\)-operator. For the inverse correspondence one requires a holomorphic bundle trivial on every real twistor line and a compatible real or quaternionic structure with positivity. An arbitrary holomorphic bundle on \(\mathbb CP^3\) does not determine a unitary instanton.}
\GSfield{Fate}{\citet[abstract]{atiyahward1977instantons} state the instanton/holomorphic-bundle identification; \citet[\S5]{drinfeldmanin1978selfdual} explicitly use this theorem in proving completeness. The full list of reality and line-triviality hypotheses is a modern reconstruction of the terse notes.}
\GSfield{Present status}{As of 10 November 2026, this is an established result in the stated scope.}
\GSfield{Assessment}{literal statement: correct with missing hypotheses; publication match: partial; verification: indirect; development: none; confidence: high}
\end{GSitem}
\begin{GSitem}{1977-432.16}{E}{Integrability is involutivity of the complex tangent distribution}
\GSfield{Kind; source}{claim; \GSpage{P000733-P000734}}
\GSfield{The notes}{An almost complex structure is described by local vector fields, which are required to commute.}
\GSfield{Statement}{For a smooth almost complex structure \(J\), integrability is equivalent to involutivity \([T^{0,1},T^{0,1}]\subset T^{0,1}\), or vanishing Nijenhuis tensor. In holomorphic coordinates there is a commuting local frame, but an arbitrary frame of an integrable distribution need not commute. The base \(S^4\) is not thereby given a global complex structure.}
\GSfield{Fate}{The notebook is invoking the classical complex integrability criterion in its twistor application; \citet[\S5]{drinfeldmanin1978selfdual} rely on the resulting holomorphic-bundle correspondence. This reference is thematic for the general criterion, not a claimed original proof of it.}
\GSfield{Present status}{As of 10 November 2026, this is an established result in the stated scope.}
\GSfield{Assessment}{literal statement: false as written; publication match: thematic only; verification: indirect; development: annotated; confidence: high}
\end{GSitem}

\GStalk{433}{1977-11-15}{Béla Sz.-Nagy}{P000735-P000736}{Common inner divisors and matrix reduction over Hardy functions}
The talk develops a generic linear-combination lemma for inner divisors in \(H^\infty\), starting from inner--outer factorization and the failure of an unrestricted Bézout identity. It is closely matched by Sz.-Nagy\textquotesingle s 1976 paper, which credits the two-function lemma to Sherman. The polynomial-combination formulation and matrix-reduction application must be distinguished from a general Smith normal form over \(H^\infty\).

\begin{GSitem}{1977-433.1}{E}{Bounded analytic and inner functions}
\GSfield{Kind; source}{definition; \GSpage{P000735}}
\GSfield{The notes}{\(H^\infty\) is the algebra of bounded analytic functions on the disk; inner functions have boundary modulus one almost everywhere.}
\GSfield{Statement}{\(H^\infty(\mathbb D)\) consists of bounded holomorphic functions on the unit disk. An inner function has radial boundary values of modulus one for Lebesgue-almost every boundary point. The notation means \(|u(e^{it})|=1\), not multiplication of \(u\) by \(e^{it}\).}
\GSfield{Fate}{\citet[\S1, opening definitions]{sznagy1976diagonalization} use precisely this setting, before the November 1977 talk.}
\GSfield{Present status}{As of 10 November 2026, this is an established result in the stated scope.}
\GSfield{Assessment}{literal statement: correct; publication match: exact; verification: read the result; development: none; confidence: high}
\end{GSitem}
\begin{GSitem}{1977-433.2}{E}{Canonical factorization and the singular inner factor}
\GSfield{Kind; source}{theorem; \GSpage{P000735-P000736}}
\GSfield{The notes}{A bounded analytic function is decomposed as an inner factor times an outer factor, and the inner factor as a Blaschke product times an exponential built from a singular measure.}
\GSfield{Statement}{Every nonzero \(u\in H^\infty\) has an inner--outer factorization. Its inner part is, up to a unimodular constant, a Blaschke product times \(S_\mu(z)=\exp\{-\int_{0}^{2\pi}(e^{it}+z)(e^{it}-z)^{-1}\,d\mu(t)\}\), for a finite positive singular measure \(\mu\). A positive exponent with positive measure gives the reciprocal and is generally unbounded.}
\GSfield{Fate}{\citet[\S1 and the proof of Lemma~2]{sznagy1976diagonalization} invoke canonical factorization. The minus sign omitted in the notebook is mathematically necessary for the stated class.}
\GSfield{Present status}{As of 10 November 2026, this is an established result in the stated scope.}
\GSfield{Assessment}{literal statement: false as written; publication match: partial; verification: read the result; development: annotated; confidence: high}
\end{GSitem}
\begin{GSitem}{1977-433.3}{E}{Greatest common inner divisor}
\GSfield{Kind; source}{definition; \GSpage{P000735-P000736}}
\GSfield{The notes}{The wedge \(\bigwedge_\alpha u_\alpha\) is described by a product over zeros and the lower bound of the singular measures.}
\GSfield{Statement}{The greatest common inner divisor of a family of nonzero bounded analytic functions is determined, up to a unimodular factor, by the minimum zero multiplicities at each point and the greatest common lower bound of the singular measures of their inner factors. Zeros belonging to only one member do not contribute. For an all-zero family a separate convention is needed.}
\GSfield{Fate}{\citet[\S1]{sznagy1976diagonalization} explicitly define the divisor for \(H^\infty\) families and use the all-zero value zero.}
\GSfield{Present status}{As of 10 November 2026, this is an established result in the stated scope.}
\GSfield{Assessment}{literal statement: correct with missing hypotheses; publication match: exact; verification: read the result; development: none; confidence: high}
\end{GSitem}
\begin{GSitem}{1977-433.4}{E}{Failure of an unrestricted Bézout identity in H-infinity}
\GSfield{Kind; source}{claim; \GSpage{P000736}}
\GSfield{The notes}{The notes deny that the greatest common inner divisor always has the form \(xu+yv\) with \(x,y\in H^\infty\).}
\GSfield{Statement}{Having no common nonconstant inner divisor does not imply that \(xu+yv=1\) has bounded analytic solutions. Such an identity forces \(\inf_{z\in\mathbb D}(|u(z)|+|v(z)|)>0\), a condition stronger than inner coprimality. Thus \(H^\infty\) is not a Bézout domain in this sense.}
\GSfield{Fate}{\citet[introduction and \S3]{sznagy1976diagonalization} explain why ordinary matrix equivalence must be replaced by quasi-equivalence. This supports the algebraic obstruction; the notebook gives no explicit counterexample.}
\GSfield{Present status}{As of 10 November 2026, this is an established result in the stated scope.}
\GSfield{Assessment}{literal statement: correct; publication match: partial; verification: read the result; development: none; confidence: high}
\end{GSitem}
\begin{GSitem}{1977-433.5}{F}{The stronger outer-factor denial is not a complete statement}
\GSfield{Kind; source}{remark; \GSpage{P000736}}
\GSfield{The notes}{A second denied identity has an outer factor \(h\) on one side and a heavily crossed-out mark on the other.}
\GSfield{Statement}{The visible formula suggests a proposed Bézout identity after multiplication by an outer function, but its quantifiers, the crossed-out coefficient, and whether the factor is required to be invertible are not fixed. An exact strengthened negative assertion cannot be supplied responsibly.}
\GSfield{Fate}{\citet[\S1, Lemma~2]{sznagy1976diagonalization} give the precise weaker generic-combination statement used next, but do not identify this damaged displayed denial.}
\GSfield{Present status}{As of 10 November 2026, only the following Sherman lemma is securely identified; the stronger denial remains ambiguous.}
\GSfield{Assessment}{literal statement: not determinable; publication match: none found; verification: unresolved lead; development: none; confidence: low}
\end{GSitem}
\begin{GSitem}{1977-433.6}{E}{Sherman’s generic linear-combination lemma}
\GSfield{Kind; source}{theorem; \GSpage{P000736}}
\GSfield{The notes}{For a fixed inner function \(\omega\), constants can be chosen so that \(u+tv=h_t(u\wedge v)\), with \(h_t\wedge\omega=1\); at most countably many \(t\) are exceptional.}
\GSfield{Statement}{Let \(u,v\in H^\infty\) be not both zero, let \(d=u\wedge v\) and let \(\omega\) be inner. After writing \(u=da\), \(v=db\), the functions \(a,b\) have no common nonconstant inner divisor. For all but at most countably many complex \(t\), \(\omega\wedge(a+tb)=1\). Consequently \(u+tv=d h_t\) with \(h_t\wedge\omega=1\). Applying the lemma to the normalized pair is essential; the unnormalized identity \(\omega\wedge(u+tv)=\omega\wedge u\wedge v\) alone is a weaker statement.}
\GSfield{Fate}{Exactly \citet[\S1, Lemma~2, pp.224--225]{sznagy1976diagonalization}, credited there to M. Sherman. The paper was received 1 November 1975 and published in 1976, so the talk is expository.}
\GSfield{Present status}{As of 10 November 2026, this is an established result in the stated scope.}
\GSfield{Assessment}{literal statement: correct; publication match: exact; verification: read the result; development: none; confidence: high; problem \hyperref[prob:GP-1977-06]{GP-1977-06}}
\end{GSitem}
\begin{GSitem}{1977-433.7}{E}{Finite polynomial combinations and simultaneous rows}
\GSfield{Kind; source}{theorem; \GSpage{P000736}}
\GSfield{The notes}{The notes extend the lemma to \(u_1+t u_2+\cdots+t^{r-1}u_{r+1}\), and then to finitely many rows.}
\GSfield{Statement}{For a finite list \(u_1,\ldots,u_r\), divide by its common inner divisor. For any fixed inner \(\omega\), the polynomial combination \(u_1+t u_2+\cdots+t^{r-1}u_r\) is then inner-coprime to \(\omega\) outside a countable set of \(t\). One may avoid the union of the exceptional sets for finitely or countably many such rows. The last subscript in the notebook must be \(r\), not \(r+1\).}
\GSfield{Fate}{\citet[\S1, Lemma~2 and \S2, Main Lemma]{sznagy1976diagonalization} establish the generic-combination mechanism and a stronger simultaneous \(\ell^1\)-coefficient theorem. The exact moment-curve wording in the notebook is a finite version; the paper\textquotesingle s Main Lemma does not insist on coefficients \(1,t,t^2,\ldots\).}
\GSfield{Present status}{As of 10 November 2026, this is an established result in the stated scope.}
\GSfield{Assessment}{literal statement: false as written; publication match: partial; verification: read the result; development: annotated; confidence: high; problem \hyperref[prob:GP-1977-06]{GP-1977-06}}
\end{GSitem}
\begin{GSitem}{1977-433.8}{E}{Matrix reduction by controlled inner divisors}
\GSfield{Kind; source}{method; \GSpage{P000736}}
\GSfield{The notes}{The final line mentions elementary transformations of a matrix.}
\GSfield{Statement}{The generic-combination lemma can select combinations of rows or columns with prescribed common inner divisors and supports a quasi-equivalence diagonalization theory for finite and semi-finite \(H^\infty\)-matrices. It does not provide a Smith normal form under arbitrary invertible \(H^\infty\) row and column operations, nor settle all doubly infinite matrices.}
\GSfield{Fate}{\citet[\S3]{sznagy1976diagonalization} carry out this program and explicitly separate the convergence difficulty for doubly infinite matrices. The notebook does not itself pose that larger question.}
\GSfield{Present status}{As of 10 November 2026, this is an established result in the stated scope.}
\GSfield{Assessment}{literal statement: correct with missing hypotheses; publication match: partial; verification: read the result; development: none; confidence: high}
\end{GSitem}

\GStalk{434}{1977-11-15}{I. M. Krichever}{P000736-P000737}{Rational KP solutions and particle dynamics}
The talk identifies rational decaying KP solutions with pole dynamics of a Calogero-type many-particle system and a commuting higher flow. It also records a Lax compatibility representation and a brief proposal about commutative subalgebras in a Weyl skew field. The published announcement was received in September 1977 and appeared in 1978; numerical coefficients in the notebook require checking against its normalization.

\begin{GSitem}{1977-434.1}{E}{The KP normalization in the notes}
\GSfield{Kind; source}{definition; \GSpage{P000736}}
\GSfield{The notes}{The equation is \(\frac34u_{yy}+\partial_x\{u_t+\frac14(6uu_x+u_{xxx})\}=0\).}
\GSfield{Statement}{The recorded equation is \(3u_{yy}+(4u_t+6uu_x+u_{xxx})_x=0\), over complex variables when rational poles are considered. Rescalings or imaginary changes of \(y\) change the sign and strength of the corresponding particle interaction.}
\GSfield{Fate}{The equation is already part of the earlier integrable-equation theory; see \citet[\S1, KP example]{krichever1977methods}. \citet[\S3]{krichever1978rational} print this normalization in the rational-pole context. The latter paper was received 6 September 1977 and appeared in the first issue of 1978; receipt is not by itself evidence of public circulation.}
\GSfield{Present status}{As of 10 November 2026, this is an established result in the stated scope.}
\GSfield{Assessment}{literal statement: correct; publication match: exact; verification: read the result; development: none; confidence: high}
\end{GSitem}
\begin{GSitem}{1977-434.2}{A}{The double-pole form of rational decaying KP solutions}
\GSfield{Kind; source}{theorem; \GSpage{P000737}}
\GSfield{The notes}{Every rational solution decreasing as \(x\to\infty\) is written as \(-2\sum_i(x-x_i(y,t))^{-2}\).}
\GSfield{Statement}{On a parameter domain where poles are distinct, a rational KP solution tending to zero as \(x\to\infty\) has only double poles, zero simple-pole coefficients and leading coefficient \(-2\): \(u=-2\sum_{i=1}^N(x-q_i(y,t))^{-2}\). The pole coordinates can collide at exceptional parameters; a global assertion must allow the resulting limiting rational function.}
\GSfield{Fate}{\citet[\S3, Theorem~3]{krichever1978rational} give the rational-solution characterization together with the associated Baker--Akhiezer function.}
\GSfield{Present status}{As of 10 November 2026, this is an established result in the stated scope.}
\GSfield{Assessment}{literal statement: correct with missing hypotheses; publication match: exact; verification: read the result; development: none; confidence: high; problem \hyperref[prob:GP-1977-07]{GP-1977-07}}
\end{GSitem}
\begin{GSitem}{1977-434.3}{A}{The particle Hamiltonian needs a normalization correction}
\GSfield{Kind; source}{claim; \GSpage{P000737}}
\GSfield{The notes}{The \(y\)-motion is assigned to \(H=\frac12\sum p_i^2+\sum_{i<j}(x_i-x_j)^{-2}\).}
\GSfield{Statement}{For precisely the displayed KP equation, the ansatz \(u=-2\sum(x-q_i)^{-2}\), and canonical momenta \(p_i=q_{i,y}\), Laurent coefficients give \(q_{i,yy}=-8\sum_{j\ne i}(q_i-q_j)^{-3}\). Thus one consistent canonical Hamiltonian is \(H=\frac12\sum p_i^2-4\sum_{i<j}(q_i-q_j)^{-2}\). Other standard Calogero conventions follow by rescaling or complexifying time and momenta; they must be stated.}
\GSfield{Reconstruction}{The canonical convention \(p=q_y\) and the corrected coefficient are an editorial Laurent calculation. A different convention might explain the speaker\textquotesingle s numerical coefficient, but none is recorded.}
\GSfield{Fate}{The pole correspondence is established by \citet[\S2, Theorem~2; \S3, Corollary~1]{krichever1978rational}. Their Hamiltonian notation must be compared together with their time and Poisson conventions; no claim of an error in that published theorem is made here. The displayed notebook equation, ansatz and canonical convention \(p=q_y\) give the coefficients in the statement by direct Laurent expansion.}
\GSfield{Present status}{As of 10 November 2026, this is an established result in the stated scope.}
\GSfield{Assessment}{literal statement: false as written; publication match: partial; verification: read the result; development: annotated; confidence: high; problem \hyperref[prob:GP-1977-07]{GP-1977-07}}
\end{GSitem}
\begin{GSitem}{1977-434.4}{A}{The KP time is a higher commuting particle flow}
\GSfield{Kind; source}{theorem; \GSpage{P000737}}
\GSfield{The notes}{The \(t\)-dependence is governed by a second commuting integral of the same particle system.}
\GSfield{Statement}{In the same canonical convention, pole cancellation gives \(q_{i,t}=\frac34p_i^2-3\sum_{j\ne i}(q_i-q_j)^{-2}\). It is the position component of the flow generated by \(K=\frac14\sum_i p_i^3-3\sum_{i<j}(p_i+p_j)(q_i-q_j)^{-2}\), which commutes with the Hamiltonian of the preceding item. These compatible flows reconstruct the rational KP solutions where poles remain distinct.}
\GSfield{Reconstruction}{The normalized Hamiltonian \(K\) is an editorial rewriting of the cubic flow; the notebook only names the additional commuting integral.}
\GSfield{Fate}{\citet[\S3, Theorem~3 and Corollary~1]{krichever1978rational} identify the higher flow by a cubic trace invariant. The explicit coefficients here retain the notebook\textquotesingle s KP signs and canonical momenta rather than silently using a differently normalized \(F_3\).}
\GSfield{Present status}{As of 10 November 2026, this is an established result in the stated scope.}
\GSfield{Assessment}{literal statement: correct with missing hypotheses; publication match: partial; verification: read the result; development: annotated; confidence: high; problem \hyperref[prob:GP-1977-07]{GP-1977-07}}
\end{GSitem}
\begin{GSitem}{1977-434.5}{E}{The differential-operator compatibility representation}
\GSfield{Kind; source}{construction; \GSpage{P000737}}
\GSfield{The notes}{The notes write \(L_1=\partial_x^2+u\), \(L_2=\partial_x^3+\frac32u\partial_x+u\), and a zero-curvature relation.}
\GSfield{Statement}{One consistent pair for the notebook\textquotesingle s KP normalization is \(L=\partial_x^2+u\), \(M=\partial_x^3+\frac32u\partial_x+w\), with \(w=\frac34u_x+\frac34\partial_x^{-1}u_y\), and \([\partial_y-L,\partial_t+M]=0\). The integration constant in \(\partial_x^{-1}u_y\) is fixed by boundary conditions or absorbed in a scalar gauge. The lower coefficient of the third-order operator is a new function \(w\), not simply \(u\).}
\GSfield{Reconstruction}{The coefficient \(w\) and time sign are supplied by direct compatibility; the source has insufficient normalization information to infer them silently.}
\GSfield{Fate}{The compatibility method and KP example already occur in \citet[Theorem~1]{krichever1976zakharov} and \citet[\S1]{krichever1977methods}. \citet[\S3]{krichever1978rational} use a separate scalar coefficient \(w\) in the rational-pole setting. The displayed pair is normalized directly to the notebook equation; the modern verification is not dated as a new 1978 theorem.}
\GSfield{Present status}{As of 10 November 2026, this is an established result in the stated scope.}
\GSfield{Assessment}{literal statement: false as written; publication match: partial; verification: read the result; development: annotated; confidence: high}
\end{GSitem}
\begin{GSitem}{1977-434.6}{F}{Commutative subalgebras of Weyl skew fields}
\GSfield{Kind; source}{remark; \GSpage{P000737}}
\GSfield{The notes}{A final proposal concerns commutative subalgebras of \(D_n\), the skew field associated with the Weyl algebra.}
\GSfield{Statement}{The fragment suggests studying commutative subalgebras or subfields inside a Weyl skew field by spectral geometry. It specifies neither a class of subalgebras, a rank, a maximality condition, nor an equivalence relation; it is not a stated classification theorem for all \(n\).}
\GSfield{Fate}{\citet{drinfeld1977commutative} is a contemporary primary lead, received 14 April 1976 and published in 1977, but its title and metadata alone do not verify a general Weyl-skew-field assertion. The notebook proposal remains too short to assign an exact fate.}
\GSfield{Present status}{As of 10 November 2026, no definite new problem with specified quantifiers can be extracted from this fragment.}
\GSfield{Assessment}{literal statement: not determinable; publication match: thematic only; verification: publisher metadata checked; development: none; confidence: low}
\end{GSitem}

\GStalk{435}{1977-09-21}{Yu. I. Manin}{P000738-P000745}{Linear algebraic construction of all instantons}
The catalogue date is retained in the date field, but the manuscript heading is unambiguously 21/XI/1977 (21 November), not 21 September. Manin gives the twistor and finite-dimensional linear-algebra description of unitary instantons, including orthogonal projection, Hermitian Gram matrices and a local connection formula. The matching Drinfeld--Manin announcement was received 27 November 1977 and published in 1978. The question about arbitrary connections already had an affirmative universal-connection theorem.

\begin{GSitem}{1977-435.1}{E}{Curvature and the self-duality equations}
\GSfield{Kind; source}{definition; \GSpage{P000738}}
\GSfield{The notes}{Manin starts with skew-Hermitian matrix potentials, their curvature, and the equations \(F_{\mu\nu}=\varepsilon^{\mu\nu\rho\sigma}F_{\rho\sigma}\).}
\GSfield{Statement}{For a unitary connection \(\nabla=d+A\), \(F=dA+A\wedge A\). In an oriented Euclidean four-manifold the self-duality equation is \(F=*F\), or \(F_{\mu\nu}=\frac12\sum_{\rho,\sigma}\varepsilon_{\mu\nu\rho\sigma}F_{\rho\sigma}\). The factor is absent only if one sums over a single complementary ordered pair.}
\GSfield{Reconstruction}{The summation convention is supplied explicitly; the notebook first discusses complementary pairs and then suppresses the convention in its epsilon formula.}
\GSfield{Fate}{This is the established differential-geometric setting recalled in \citet[\S1]{drinfeldmanin1978selfdual}. Its use in the later classification does not make the definition a new 1978 result.}
\GSfield{Present status}{As of 10 November 2026, this is an established result in the stated scope.}
\GSfield{Assessment}{literal statement: correct with missing hypotheses; publication match: exact; verification: read the result; development: none; confidence: high}
\end{GSitem}
\begin{GSitem}{1977-435.2}{E}{Gauge changes and removable frame singularities}
\GSfield{Kind; source}{construction; \GSpage{P000738}}
\GSfield{The notes}{Changing a local frame gives \(A\mapsto u^{-1}du+u^{-1}Au\); apparent finite-point singularities are required to disappear in a suitable gauge. The covariant-derivative sum incorrectly repeats \(i\) as its summation index.}
\GSfield{Statement}{In a local unitary frame, \(\nabla_\mu s_i=\sum_j A_\mu^{ij}s_j\), with a consistent row or column convention. Smoothness is a property of the connection on the bundle, not of one singular trivialization. Gauge equivalence preserves curvature up to conjugation and preserves self-duality.}
\GSfield{Fate}{\citet[\S1]{drinfeldmanin1978selfdual} formulate regularity and equivalence in terms of bundles with connections. The displayed gauge law is standard; the index repair is elementary and scan-confirmed.}
\GSfield{Present status}{As of 10 November 2026, this is an established result in the stated scope.}
\GSfield{Assessment}{literal statement: false as written; publication match: partial; verification: read the result; development: annotated; confidence: high}
\end{GSitem}
\begin{GSitem}{1977-435.3}{E}{Regularity at infinity and compactification}
\GSfield{Kind; source}{definition; \GSpage{P000738-P000739}}
\GSfield{The notes}{The connection and its bundle must extend over \(S^4=\mathbb R^4\cup\{\infty\}\); the extended bundle may be nontrivial.}
\GSfield{Statement}{A regular instanton here means a smooth unitary self-dual connection on a bundle over the round four-sphere, expressed on the affine chart \(\mathbb R^4\). Extension over infinity is an imposed regularity condition in this talk. It should not be confused with a proof that arbitrary finite-action fields extend.}
\GSfield{Fate}{This is precisely the sphere formulation in \citet[\S1]{drinfeldmanin1978selfdual}. The separate removable-singularity theorem is treated under record 432.}
\GSfield{Present status}{As of 10 November 2026, this is an established result in the stated scope.}
\GSfield{Assessment}{literal statement: correct; publication match: exact; verification: read the result; development: none; confidence: high}
\end{GSitem}
\begin{GSitem}{1977-435.4}{E}{Clutching classification and charge}
\GSfield{Kind; source}{theorem; \GSpage{P000739}}
\GSfield{The notes}{Clutching along the equator gives \(\pi_3(U(r))=\mathbb Z\), generated by the block embedding of \(SU(2)\), and the sign of charge is reversed with duality.}
\GSfield{Statement}{For \(r\geq2\), complex rank-\(r\) bundles over \(S^4\) are classified by an integer clutching class in \(\pi_3(U(r))\), equivalently the second Chern number after fixing orientation. For \(r=1\), \(\pi_3(U(1))=0\). This classifies underlying bundles, not their connections; a fixed charge supports a moduli space of instantons.}
\GSfield{Fate}{\citet[\S1 and \S6]{drinfeldmanin1978selfdual} use the charge \(n\) and rank \(r\) as discrete parameters. The omitted rank-one exception and the distinction from gauge equivalence are essential.}
\GSfield{Present status}{As of 10 November 2026, this is an established result in the stated scope.}
\GSfield{Assessment}{literal statement: correct with missing hypotheses; publication match: partial; verification: read the result; development: annotated; confidence: high}
\end{GSitem}
\begin{GSitem}{1977-435.5}{A}{Projection from a finite-dimensional trivial bundle}
\GSfield{Kind; source}{construction; \GSpage{P000739}}
\GSfield{The notes}{The announced construction places \(L\) inside \(H\times S^4\), with \(\dim H=2n+r\), and obtains the connection by differentiating and orthogonally projecting.}
\GSfield{Statement}{For the admissible instanton data described below, \(H\times S^4=L\oplus L^\perp\), with \(\operatorname{rank}L=r\), and \(\nabla s=P_L(ds)\). A Hermitian orthogonal projection gives a unitary connection. The sharp ambient dimension and the self-duality conditions belong to the instanton construction, not to an arbitrary connection embedding theorem.}
\GSfield{Fate}{\citet[\S3 and \S6]{drinfeldmanin1978selfdual} give this construction; the announcement was received 27 November 1977, six days after the source date, and published in 1978.}
\GSfield{Present status}{As of 10 November 2026, this is an established result in the stated scope.}
\GSfield{Assessment}{literal statement: correct; publication match: exact; verification: read the result; development: none; confidence: high}
\end{GSitem}
\begin{GSitem}{1977-435.6}{E}{Can every connection be induced from a trivial one?}
\GSfield{Kind; source}{question; \GSpage{P000739}}
\GSfield{The notes}{A bracketed question asks whether a bundle with connection can always be embedded into a trivial bundle with connection, comparing this with a Euclidean embedding of a Riemannian manifold.}
\GSfield{Statement}{For a smooth Hermitian vector bundle over a finite-dimensional smooth manifold and a compatible unitary connection, does there exist a finite-dimensional trivial Hermitian bundle and an isometric bundle embedding for which orthogonal projection of the trivial connection is the prescribed connection?}
\GSfield{Reconstruction}{The unitary compatibility and finite-dimensional smooth-base assumptions make the intended projection question precise.}
\GSfield{Fate}{Affirmative before the talk: \citet[p.286, universal-connection theorem and the unitary case]{narasimhan1962universal} announce the joint Narasimhan--Ramanan theorem at the September 1961 Prague symposium; the proceedings appeared in 1962. The universal Stiefel connection gives exactly the projection formulation. Earlier occurrence 1975-382.5 describes the tautological universal connection and its base; it has no earlier problem identifier. That construction and this explicit universality question are closely related, with the quantifiers stated here.}
\GSfield{Present status}{As of 10 November 2026, this is an established result in the stated scope.}
\GSfield{Assessment}{literal statement: correct; publication match: exact; verification: read the result; development: none; confidence: high; problem \hyperref[prob:GP-1977-08]{GP-1977-08}}
\end{GSitem}
\begin{GSitem}{1977-435.7}{E}{Quaternionic real structure and the twistor projection}
\GSfield{Kind; source}{construction; \GSpage{P000739-P000740}}
\GSfield{The notes}{The map \(\sigma(z_1,z_2,z_3,z_4)=(-\bar z_2,\bar z_1,-\bar z_4,\bar z_3)\) has no invariant complex lines; its invariant complex two-planes are identified with \(S^4\). The following page calls \(\sigma\) an involution.}
\GSfield{Statement}{On \(\mathbb C^4\), this antilinear map satisfies \(\sigma^2=-1\). It therefore induces a fixed-point-free involution on \(\mathbb{CP}^3\). Its invariant two-planes are quaternionic lines, parametrized by \(\mathbb{HP}^1\simeq S^4\). Sending \([z]\) to \(\operatorname{span}_{\mathbb C}(z,\sigma z)\) is the twistor fibration with fiber \(\mathbb{CP}^1\).}
\GSfield{Fate}{The twistor setting is already present in \citet{atiyahward1977instantons}; the precise real structure is written in \citet[\S0.2 and \S1.6]{drinfeldmanin1979instantons}. The latter is a detailed later exposition, not the first appearance of twistor space.}
\GSfield{Present status}{As of 10 November 2026, this is an established result in the stated scope.}
\GSfield{Assessment}{literal statement: false as written; publication match: partial; verification: read the result; development: annotated; confidence: high}
\end{GSitem}
\begin{GSitem}{1977-435.8}{E}{Affine coordinates of the twistor map}
\GSfield{Kind; source}{construction; \GSpage{P000740}}
\GSfield{The notes}{Two complex quotients give \(x_2+ix_1\) and \(x_4+ix_3\), and a remark says that three coordinates suffice.}
\GSfield{Statement}{On \((z_3,z_4)\ne(0,0)\), put \(x_2+ix_1=(z_1\bar z_3+\bar z_2z_4)/(|z_3|^2+|z_4|^2)\) and \(x_4+ix_3=(\bar z_2z_3-z_1\bar z_4)/(|z_3|^2+|z_4|^2)\). These are constant on the complex projective line associated with a quaternionic line. Three complex affine coordinates on \(\mathbb{CP}^3\) suffice; the statement is not a claim that three real coordinates parametrize \(S^4\).}
\GSfield{Fate}{\citet[\S1.6]{drinfeldmanin1979instantons} print the coordinate construction in the detailed proof. It is an elementary coordinate expression of the earlier twistor map.}
\GSfield{Present status}{As of 10 November 2026, this is an established result in the stated scope.}
\GSfield{Assessment}{literal statement: correct with missing hypotheses; publication match: partial; verification: read the result; development: none; confidence: high}
\end{GSitem}
\begin{GSitem}{1977-435.9}{E}{Self-duality as a type condition on pulled-back curvature}
\GSfield{Kind; source}{theorem; \GSpage{P000740}}
\GSfield{The notes}{The lemma says self-dual two-forms pull back to type \((1,1)\), while three complementary \((2,0)\) components are independent. Corollary (8) omits the component index on its bracket.}
\GSfield{Statement}{For the indicated orientation, \(F=*F\) if and only if \((D^*F)^{2,0}=0\); for a unitary connection the conjugate \((0,2)\) condition follows. The three independent components detect the opposite-duality part. The condition \(D^*F=0\) without projection would force \(F=0\), since \(D\) is a surjective submersion.}
\GSfield{Fate}{This is the local twistor criterion underlying \citet{atiyahward1977instantons} and made explicit in \citet[\S0.2 and \S1.6--1.9]{drinfeldmanin1979instantons}. The missing projection is a faithful notebook error.}
\GSfield{Present status}{As of 10 November 2026, this is an established result in the stated scope.}
\GSfield{Assessment}{literal statement: false as written; publication match: partial; verification: read the result; development: annotated; confidence: high}
\end{GSitem}
\begin{GSitem}{1977-435.10}{A}{Belavin’s covariant compatibility equations}
\GSfield{Kind; source}{construction; \GSpage{P000741}; A. A. Belavin}
\GSfield{The notes}{Belavin rewrites the criterion as commuting holomorphic covariant derivatives and then as \([L_1,L_2]=0\) for a spectral parameter \(\lambda\).}
\GSfield{Statement}{With \(\nabla_\mu=\partial_\mu+A_\mu\), the pair \(L_1=\lambda(\nabla_1-i\nabla_2)+(\nabla_3+i\nabla_4)\), \(L_2=\lambda(\nabla_3-i\nabla_4)-(\nabla_1+i\nabla_2)\) has zero commutator for all \(\lambda\) exactly when the three self-duality equations hold in this convention. Equivalently the \((1,0)\) part of \(D^*\nabla\) is flat in its holomorphic directions.}
\GSfield{Fate}{\citet[\S1.6--1.9]{drinfeldmanin1979instantons} use this compatibility and explicitly identify the Belavin--Zakharov route. Here the relative minus sign is retained; it differs from the problematic pair in the earlier Novikov notes.}
\GSfield{Present status}{As of 10 November 2026, this is an established result in the stated scope.}
\GSfield{Assessment}{literal statement: correct; publication match: partial; verification: read the result; development: none; confidence: high}
\end{GSitem}
\begin{GSitem}{1977-435.11}{A}{From a tensor subspace to the bundle fibers}
\GSfield{Kind; source}{construction; \GSpage{P000741-P000742}}
\GSfield{The notes}{Choose \(I\subset H\otimes\mathbb C^4\), \(\dim I=n\). If \(\ell_1,\ell_2\) cut out a \(\sigma\)-invariant two-plane, take the orthogonal complement of \(I_{\ell_1}+I_{\ell_2}\).}
\GSfield{Statement}{Write \(I_\ell=(1\otimes\ell)(I)\subset H\). For a quaternionic line \(D\subset\mathbb C^4\), let \(D^\circ\) be its annihilator and set \(L_D=(\sum_{\ell\in D^\circ}I_\ell)^\perp\). Choosing any basis \(\ell_1,\ell_2\) of \(D^\circ\) gives the same subspace, so the construction is independent of the two defining equations.}
\GSfield{Fate}{\citet[\S2--\S3 and the unitary version in \S6]{drinfeldmanin1978selfdual} supply the linear-algebra construction. The independence is immediate from linearity of contraction.}
\GSfield{Present status}{As of 10 November 2026, this is an established result in the stated scope.}
\GSfield{Assessment}{literal statement: correct; publication match: exact; verification: read the result; development: none; confidence: high}
\end{GSitem}
\begin{GSitem}{1977-435.12}{A}{Rank, positivity, and quaternionic orthogonality conditions}
\GSfield{Kind; source}{theorem; \GSpage{P000742-P000743}}
\GSfield{The notes}{Conditions 1 and 2 prescribe fiber dimension \(r\) and positivity. Condition 0 is inserted later as \(Q(I_\ell,I_{\ell\circ\sigma})=0\).}
\GSfield{Statement}{Require the contraction sum over every quaternionic line to have rank \(2n\), take a positive Hermitian form on \(H\), and impose the quaternionic orthogonality constraint. Since \(\ell\circ\sigma\) is antilinear, the second contraction must use the complex-linear form \(z\mapsto\overline{\ell(\sigma z)}\). With these admissibility conditions the orthogonal complements form a smooth rank-\(r\) bundle and the projected connection is self-dual.}
\GSfield{Reconstruction}{Positivity on the ambient Hermitian space is adopted as the convenient modern sufficient hypothesis. The source\textquotesingle s claim that positivity follows from its abbreviated conditions is not independently strengthened here.}
\GSfield{Fate}{The admissibility and projection theorem are in \citet[\S2--\S4 and \S6]{drinfeldmanin1978selfdual}. The detailed discussion of positivity appears in \citet[\S3--\S4]{drinfeldmanin1979instantons}. A bare indefinite form without the full conditions is insufficient.}
\GSfield{Present status}{As of 10 November 2026, this is an established result in the stated scope.}
\GSfield{Assessment}{literal statement: false as written; publication match: partial; verification: read the result; development: annotated; confidence: high}
\end{GSitem}
\begin{GSitem}{1977-435.13}{A}{Completeness and the equivalence relation on instanton data}
\GSfield{Kind; source}{theorem; \GSpage{P000742-P000743}}
\GSfield{The notes}{Theorem 1 says that the construction gives all self-dual Yang--Mills solutions and that nonisomorphic bundles correspond to inequivalent solutions. Isomorphisms of tensor diagrams are then described.}
\GSfield{Statement}{Every regular unitary instanton on \(S^4\) is obtained from admissible finite-dimensional data. Two data sets yield gauge-equivalent connections precisely under the corresponding isomorphisms of diagrams, including an isometry on \(H\) and a change of basis in \(I\). The classification concerns bundles with connection; nonisomorphism of the underlying topological bundle is far too coarse.}
\GSfield{Fate}{Exactly the completeness and uniqueness theorem of \citet[\S4]{drinfeldmanin1978selfdual}, received 27 November 1977 and published in 1978; \citet[\S3]{drinfeldmanin1979instantons} provide the fuller proof. The upper tensor factor in the notebook diagram needs \(\mathbb C^4\).}
\GSfield{Present status}{As of 10 November 2026, this is an established result in the stated scope.}
\GSfield{Assessment}{literal statement: correct with missing hypotheses; publication match: exact; verification: read the result; development: annotated; confidence: high; problem \hyperref[prob:GP-1977-05]{GP-1977-05}}
\end{GSitem}
\begin{GSitem}{1977-435.14}{F}{The quadruple-of-subspaces intervention}
\GSfield{Kind; source}{remark; \GSpage{P000743}; I. M. Gelfand and V. A. Ponomarev}
\GSfield{The notes}{Gelfand invokes unpublished work with Ponomarev on quadruples with an automorphism. Ponomarev takes the four images of \(H\) in \((H^{\oplus4})/I\), then asks for the quadruples compatible with an orthogonal-complement operation.}
\GSfield{Statement}{The quotient construction indeed associates four subspaces to a tensor inclusion. But the notes do not specify the additional automorphism, nondegeneracy hypotheses, duality functor or equivalence relation needed for the claimed Hermitian classification; no exact classification theorem can be reconstructed.}
\GSfield{Fate}{\citet{gelfandponomarev1971quadruples} is the earlier complex quadruple-classification lead, received 21 October 1970 and published in 1971. Its metadata do not verify the additional unpublished Hermitian assertion. The instanton diagram theorem of \citet[\S4]{drinfeldmanin1978selfdual} solves its own specified classification, not every possible quadruple-with-duality problem.}
\GSfield{Present status}{As of 10 November 2026, the extra unpublished assertion remains unidentified; the earlier quadruple work is thematic evidence only.}
\GSfield{Assessment}{literal statement: not determinable; publication match: thematic only; verification: publisher metadata checked; development: none; confidence: low}
\end{GSitem}
\begin{GSitem}{1977-435.15}{A}{Gram matrices and the absence of horizontal sections}
\GSfield{Kind; source}{construction; \GSpage{P000743}}
\GSfield{The notes}{A basis of \(I\) produces \(4n\) vectors \(e_{ij}\in H\) and their Gram matrix \(\widetilde Q\). The assumption \(\operatorname{rank}\widetilde Q=2n+r\) is equated with absence of horizontal sections.}
\GSfield{Statement}{For \(i=1,\ldots,n\), write the basis tensor as \(e_i=\sum_{j=1}^4e_{ij}\otimes f_j\). The Gram matrix determines these vectors up to unitary equivalence on their span. Full rank \(2n+r\) means they span \(H\). For the instanton data, failure of spanning gives a constant orthogonal summand and hence a flat factor; the converse uses the admissibility theorem.}
\GSfield{Fate}{\citet[\S6--\S7]{drinfeldmanin1978selfdual} state the full-span criterion and the horizontal-section interpretation. The notebook\textquotesingle s basis endpoint \(e_r\) must be \(e_n\).}
\GSfield{Present status}{As of 10 November 2026, this is an established result in the stated scope.}
\GSfield{Assessment}{literal statement: false as written; publication match: partial; verification: read the result; development: annotated; confidence: high}
\end{GSitem}
\begin{GSitem}{1977-435.16}{A}{The constrained four-by-four Gram blocks}
\GSfield{Kind; source}{construction; \GSpage{P000744}}
\GSfield{The notes}{Every block has the displayed six-parameter pattern; Hermitian symmetry and rank are required, while Conditions 1 and 2 are expressly omitted.}
\GSfield{Statement}{The block pattern is \(\left(\begin{smallmatrix}a&0&c&d\\0&a&e&f\\f&-d&b&0\\-e&c&0&b\end{smallmatrix}\right)\). It expresses the quaternionic orthogonality equations. Global Hermitian symmetry relates the block at \((k,\ell)\) to the adjoint of the block at \((\ell,k)\). Positivity and the pointwise rank conditions remain necessary; the linear block pattern alone does not classify regular instantons.}
\GSfield{Fate}{\citet[\S6]{drinfeldmanin1978selfdual} give the Gram-matrix formulation and explicitly acknowledge further regularity inequalities. Thus the notebook is an announcement of constraints, not an assertion that its listed linear equations suffice.}
\GSfield{Present status}{As of 10 November 2026, this is an established result in the stated scope.}
\GSfield{Assessment}{literal statement: correct with missing hypotheses; publication match: partial; verification: read the result; development: none; confidence: high}
\end{GSitem}
\begin{GSitem}{1977-435.17}{A}{Affine generators of the orthogonal complement}
\GSfield{Kind; source}{construction; \GSpage{P000744}}
\GSfield{The notes}{The vectors \(h_{k1}\) and \(h_{k2}\) are called \(2n\) global generators of \(L^\perp\) over \(S^4\), although their formulas use affine coordinates.}
\GSfield{Statement}{On the chart \(\mathbb R^4\), the vectors \(h_{k1}=(x_2+ix_1)e_{k1}+(x_4-ix_3)e_{k2}+e_{k3}\) and \(h_{k2}=-(x_4+ix_3)e_{k1}+(x_2-ix_1)e_{k2}+e_{k4}\) span \(L_x^\perp\). They are not a global frame on \(S^4\) for nonzero charge: a rank-\(2n\) bundle globally generated by exactly \(2n\) independent sections would be trivial, contradicting \(c(L)c(L^\perp)=1\).}
\GSfield{Fate}{\citet[\S6]{drinfeldmanin1978selfdual} use these vectors in affine Euclidean coordinates. The local construction extends geometrically across infinity after changing trivialization, not by retaining one global affine frame.}
\GSfield{Present status}{As of 10 November 2026, this is an established result in the stated scope.}
\GSfield{Assessment}{literal statement: false as written; publication match: partial; verification: read the result; development: annotated; confidence: high}
\end{GSitem}
\begin{GSitem}{1977-435.18}{A}{Computing the connection from a kernel and orthonormal frame}
\GSfield{Kind; source}{method; \GSpage{P000744-P000745}}
\GSfield{The notes}{Solve \(2n\) homogeneous equations in \(2n+r\) variables, choose a fundamental system, orthonormalize, and differentiate to obtain \(A_\mu^{ij}=Q(\partial_\mu y^{(i)},y^{(j)})\).}
\GSfield{Statement}{On a chart where a chosen minor is nonzero, solve \(Q(v,h_{ka}(x))=0\), \(k=1,\ldots,n\), \(a=1,2\), for a local basis of \(L_x\). Its coefficients are rational functions of the real coordinates. Gram--Schmidt introduces square roots of positive rational functions. If \(Y\) is the resulting orthonormal column frame, the connection matrix is \(Y^*dY\), up to the notebook\textquotesingle s transposed frame convention.}
\GSfield{Fate}{\citet[\S6]{drinfeldmanin1978selfdual} describe this explicit algorithm. Its formula is local and uses nonvanishing minors; no global rational orthonormal frame is implied.}
\GSfield{Present status}{As of 10 November 2026, this is an established result in the stated scope.}
\GSfield{Assessment}{literal statement: correct with missing hypotheses; publication match: exact; verification: read the result; development: none; confidence: high}
\end{GSitem}
\begin{GSitem}{1977-435.19}{A}{The cohomological monad classification}
\GSfield{Kind; source}{theorem; \GSpage{P000745}}
\GSfield{The notes}{The final geometric theorem says diagrams classify bundles on \(\mathbb{CP}^3\) satisfying \(H^i(L(-2))=0\), the same for \(L^*\), and \(H^2(L(-1))=H^0(L(-1))=0\); the range of \(i\) is not stated.}
\GSfield{Statement}{A precise published version defines an admissible locally free sheaf \(E\) by \(H^i(E(k))=0\) for \(i=0,1\) with \(i+k\leq-1\), and for \(i=2,3\) with \(i+k\geq0\). The middle cohomology of a special monad \(0\to I\otimes\mathcal O(-1)\to H\otimes\mathcal O\to I^{\prime}\otimes\mathcal O(1)\to0\), with fiberwise injective and surjective end maps, gives an equivalence with admissible sheaves. Instantons additionally require the real and positivity conditions.}
\GSfield{Reconstruction}{The full admissibility definition and monad exactness conditions are taken from the later theorem and are explicitly separated from the notebook\textquotesingle s abbreviated cohomology conditions.}
\GSfield{Fate}{\citet[\S0.3--0.4 and Theorem~2.1]{drinfeldmanin1979instantons} prove this equivalence; the paper was received 26 May 1978 and published in 1979. The short notebook vanishings are a precursor, not an independently complete list with all quantifiers.}
\GSfield{Present status}{As of 10 November 2026, this is an established result in the stated scope.}
\GSfield{Assessment}{literal statement: correct with missing hypotheses; publication match: partial; verification: read the result; development: none; confidence: high}
\end{GSitem}

\GStalk{436}{1977-11-28}{L. D. Faddeev}{P000746-P000752}{Classical and quantum solitons}
Faddeev compares particle and soliton descriptions, formulates asymptotic decomposition heuristics, and discusses sine-Gordon quantization, exact scattering, semiclassical determinants and mass corrections. The scope of the general solvability and resolution claims is not specified. Breathers are explicitly recognized in the notes, so they cannot be treated as a refutation of every formulation the speaker might have intended. P000753--P000754 begin an adjacent February 1978 Gelfand talk, treated with its continuation under record 438, rather than as Faddeev material.

\begin{GSitem}{1977-436.1}{E}{Hamiltonian formulation of a nonlinear scalar field}
\GSfield{Kind; source}{definition; \GSpage{P000746}}
\GSfield{The notes}{The nonlinear Klein--Gordon equation is written with its Lagrangian, momentum \(\pi=u_t\), canonical bracket and energy.}
\GSfield{Statement}{For a smooth real potential \(U(u)=m^2u^2/2+\varepsilon V(u)\), the formal Hamiltonian is \(H=\int[\pi^2/2+|\nabla u|^2/2+U(u)]\,dx\). The convention \(\{u(x),\pi(y)\}=\delta(x-y)\), \(u_t=\{u,H\}\), gives \(u_{tt}-\Delta u+U^{\prime}(u)=0\). Domain, boundary and regularity assumptions are needed for an actual infinite-dimensional Hamiltonian flow.}
\GSfield{Fate}{\citet[introduction and \S1]{takhtajanfaddeev1974nonlinear} give the sine-Gordon instance before the talk. The general variational identity is standard background.}
\GSfield{Present status}{As of 10 November 2026, this is an established result in the stated scope.}
\GSfield{Assessment}{literal statement: correct with missing hypotheses; publication match: exact; verification: read the result; development: none; confidence: high}
\end{GSitem}
\begin{GSitem}{1977-436.2}{E}{Free-field Fock energy and occupation densities}
\GSfield{Kind; source}{construction; \GSpage{P000746}}
\GSfield{The notes}{The free Hamiltonian is diagonalized by Fourier modes; the occupation density in (7) contains \(a^*(k^{\prime})a(k)\) and the energy is a sum of relativistic one-particle energies.}
\GSfield{Statement}{After normal ordering, the free bosonic Hamiltonian is \(d\Gamma(\omega)\), \(\omega(k)=\sqrt{|k|^2+m^2}\). On generalized momentum states its energy is \(\sum_j\omega(k_j)\). The formal occupation density is \(a^*(k)a(k)\), or rigorously its smeared number operator; sharp-momentum states and delta-function occupations are distributions, not normalizable eigenvectors.}
\GSfield{Fate}{\citet[\S1, pp.148--149 of the Russian text]{korepinfaddeev1975quantization} use exactly this particle interpretation. The prime on the first argument in the notebook is a scan-confirmed error.}
\GSfield{Present status}{As of 10 November 2026, this is an established result in the stated scope.}
\GSfield{Assessment}{literal statement: false as written; publication match: partial; verification: read the result; development: annotated; confidence: high}
\end{GSitem}
\begin{GSitem}{1977-436.3}{E}{Nonspreading finite-energy solutions and failure of free scattering}
\GSfield{Kind; source}{definition; \GSpage{P000747}}
\GSfield{The notes}{Solitons are first called everything failing to approach a free solution, then more narrowly finite-energy solutions that do not spread out.}
\GSfield{Statement}{A localized stationary or traveling solution can obstruct scattering to a free Klein--Gordon wave. Non-scattering alone is not a definition of a soliton: it also includes other recurrent dynamics and, outside a global-existence class, breakdown of the solution. The intended soliton notion here is a localized coherent structure whose energy remains concentrated in a moving bounded region.}
\GSfield{Fate}{\citet[introduction, pp.4--6]{faddeevkorepin1978quantum} make the same distinction between spreading waves and localized structures; \citet[\S4]{takhtajanfaddeev1974nonlinear} supply explicit sine-Gordon kinks and breathers.}
\GSfield{Present status}{As of 10 November 2026, this is an established result in the stated scope.}
\GSfield{Assessment}{literal statement: correct with missing hypotheses; publication match: partial; verification: read the result; development: none; confidence: high}
\end{GSitem}
\begin{GSitem}{1977-436.4}{E}{Lorentz boosts and the mass-shell relation}
\GSfield{Kind; source}{construction; \GSpage{P000747}}
\GSfield{The notes}{Boosting a static finite-energy solution gives \(u_s((x-x_0-vt)/\sqrt{1-v^2})\) and \(H_s(v)=M/\sqrt{1-v^2}=\sqrt{p^2+M^2}\).}
\GSfield{Statement}{For a localized static solution of a Lorentz-invariant scalar equation in one spatial dimension, a boost with \(|v|<1\) gives the displayed traveling profile. If \(M\) is its rest energy, the conserved energy-momentum vector gives \(E=M/\sqrt{1-v^2}\) and \(p=Mv/\sqrt{1-v^2}\). Boundary stress must vanish so the integrated charges transform correctly.}
\GSfield{Fate}{\citet[\S1, equation~(2)]{korepinfaddeev1975quantization} state this explicitly for the sine-Gordon kink. The notebook\textquotesingle s schematic higher-dimensional opening does not turn the one-coordinate boost formula into a general \(d\)-dimensional expression.}
\GSfield{Present status}{As of 10 November 2026, this is an established result in the stated scope.}
\GSfield{Assessment}{literal statement: correct with missing hypotheses; publication match: exact; verification: read the result; development: none; confidence: high}
\end{GSitem}
\begin{GSitem}{1977-436.5}{C}{Resolution into radiation and finitely many localized structures}
\GSfield{Kind; source}{conjecture; \GSpage{P000747-P000748}}
\GSfield{The notes}{Equation (12) proposes \(u\to u_0+\sum u_s\) at large positive and negative time, calling it a hypothesis in general and a theorem in models. The next page explicitly adds periodic localized solutions.}
\GSfield{Statement}{A precise model question is whether every global finite-energy sine-Gordon solution admits a large-time decomposition into finitely many localized coherent structures, including kinks, antikinks and breathers, plus radiation, in a specified topology. The notebook supplies neither a convergence norm nor genericity assumptions; it first writes a free-wave term and then explicitly adds periodic localized structures.}
\GSfield{Reconstruction}{The all-energy sine-Gordon formulation makes the model and quantifiers explicit. The inclusion of breathers is supported by P000748; the exact radiation topology remains to be specified rather than attributed to the notebook.}
\GSfield{Fate}{\citet[Theorem~1.1, Definition~2.6 and Theorem~10.1]{chenliulu2020sinegordon} prove resolution for generic data in \(H^{2,s}_{\sin}\times H^{1,s}\), \(s>1/2\); their radiation has modified scattering inside the light cone. The preprint first appeared 8 September 2020; the version read is 15 February 2022. This is a partial match to the unrestricted wording, not a theorem of free scattering for all energy data.}
\GSfield{Present status}{As of 10 November 2026, the cited generic weighted-data result is established. This search did not verify an all-finite-energy theorem matching the notebook wording. The verified theorem uses modified scattering for its radiation component.}
\GSfield{Assessment}{literal statement: not determinable; publication match: partial; verification: read the result; development: annotated; confidence: high; problem \hyperref[prob:GP-1977-09]{GP-1977-09}}
\GSfield{Worth developing}{Worth developing in a carefully bounded form, because identifying the admissible coherent structures and the convergence topology prevents a false soliton-resolution assertion.}
\end{GSitem}
\begin{GSitem}{1977-436.6}{E}{Additivity of asymptotic energies}
\GSfield{Kind; source}{claim; \GSpage{P000747}}
\GSfield{The notes}{The notes pass from separated soliton asymptotics to \(H(u)=H(u_0)+\sum H_s(v_i)\).}
\GSfield{Statement}{Energy additivity follows if the asymptotic decomposition converges strongly enough in the energy quantities and the interaction terms vanish as the localized profiles separate. Pointwise convergence alone is insufficient. In a long-range model the radiation energy may have a nonlinear scattering description even though its conserved energy is well-defined.}
\GSfield{Fate}{\citet[introduction, equation~(1.13)]{faddeevkorepin1978quantum} explicitly assume separation of the wave packet and localized solutions; \citet[\S4, equations~(32)--(33)]{takhtajanfaddeev1974nonlinear} obtain an exact spectral energy decomposition in the integrable sine-Gordon setting.}
\GSfield{Present status}{As of 10 November 2026, this is an established result in the stated scope.}
\GSfield{Assessment}{literal statement: correct with missing hypotheses; publication match: partial; verification: read the result; development: none; confidence: high}
\end{GSitem}
\begin{GSitem}{1977-436.7}{E}{Several particle species from one classical field}
\GSfield{Kind; source}{claim; \GSpage{P000747-P000748}}
\GSfield{The notes}{Quantization is expected to give sums of ordinary-particle and soliton energies, although only one scalar field occurs in the Lagrangian; soliton contributions persist in the classical limit.}
\GSfield{Statement}{The semiclassical particle interpretation assigns species to stable localized classical solutions and their quantized internal excitations. In sine-Gordon the kink mass scales as inverse coupling. The displayed additive spectrum describes asymptotic multiparticle states; it is not a proof that an arbitrary nonlinear field has a rigorously constructed Fock decomposition or that all species are independent.}
\GSfield{Fate}{\citet[\S1 and conclusion]{korepinfaddeev1975quantization} already formulate this semiclassical interpretation. In the exact sine-Gordon spectrum the lightest breather is identified with the perturbative neutral particle, so counting both as independent species would duplicate a state.}
\GSfield{Present status}{As of 10 November 2026, this is an established result in the stated scope.}
\GSfield{Assessment}{literal statement: correct with missing hypotheses; publication match: partial; verification: read the result; development: none; confidence: high}
\end{GSitem}
\begin{GSitem}{1977-436.8}{E}{Periodic solitons and their internal parameters}
\GSfield{Kind; source}{construction; \GSpage{P000748}}
\GSfield{The notes}{A periodic localized solution has center, velocity, period parameter \(\beta\) and initial phase \(\alpha\). Equation (18) then writes \(\sqrt{p_k^2+M(\alpha_k)}\).}
\GSfield{Statement}{A breather has a frequency or action parameter controlling its rest energy and a separate phase parameter. Relativistic energy is \(\sqrt{p^2+M(\beta)^2}\) if \(\beta\) denotes that internal parameter. It is independent of the initial time phase. Thus (18) needs both the mass square and a parameter-name correction relative to the sentence immediately before it.}
\GSfield{Fate}{\citet[\S1, equation~(3) and subsequent discussion]{korepinfaddeev1975quantization} distinguish phase from the internal action and give the breather mass. Both defects in the notebook formula are visible in the scan.}
\GSfield{Present status}{As of 10 November 2026, this is an established result in the stated scope.}
\GSfield{Assessment}{literal statement: false as written; publication match: partial; verification: read the result; development: annotated; confidence: high}
\end{GSitem}
\begin{GSitem}{1977-436.9}{E}{Classical sine-Gordon kink and breather masses}
\GSfield{Kind; source}{theorem; \GSpage{P000748}}
\GSfield{The notes}{For \(u_{tt}-u_{xx}+m^2\sin u=0\), introducing coupling \(\gamma\) gives kink mass \(8m/\gamma\) and breather mass \((16m/\gamma)\sin\alpha\).}
\GSfield{Statement}{Use \(H=\gamma^{-1}\int[\frac12u_t^2+\frac12u_x^2+m^2(1-\cos u)]dx\) and \(\{u(x),u_t(y)\}=\gamma\delta(x-y)\). The kink has energy \(8m/\gamma\). A breather with frequency \(m\cos\alpha\), \(0<\alpha<\pi/2\), has rest energy \((16m/\gamma)\sin\alpha\). The endpoint \(\alpha=\pi/2\) is a separatrix limit, not a finite-period breather.}
\GSfield{Fate}{\citet[\S4]{takhtajanfaddeev1974nonlinear} give these masses and explicit solutions; \citet[\S1]{korepinfaddeev1975quantization} retain them. The source reverses the bracket order in (21) without changing its sign, unlike (3); the canonical convention is fixed here.}
\GSfield{Present status}{As of 10 November 2026, this is an established result in the stated scope.}
\GSfield{Assessment}{literal statement: correct with missing hypotheses; publication match: partial; verification: read the result; development: annotated; confidence: high}
\end{GSitem}
\begin{GSitem}{1977-436.10}{E}{Bohr--Sommerfeld breather levels and the exact coupling}
\GSfield{Kind; source}{claim; \GSpage{P000748}}
\GSfield{The notes}{The internal angle is quantized as \(\alpha=\gamma n/16\), attributed approximately to Takhtajan--Faddeev, 1974.}
\GSfield{Statement}{This is a semiclassical level rule with an upper bound on \(n\). In the exact sine-Gordon scattering convention with Lagrangian coupling \(\beta^2\), write \(\Gamma=\beta^2/(1-\beta^2/(8\pi))\). Bound-state masses satisfy \(m_n=2M\sin(n\Gamma/16)\), \(1\leq n<8\pi/\Gamma\). The bare or classical coupling must not silently be identified with \(\Gamma\).}
\GSfield{Fate}{\citet[\S1 and \S5]{korepinfaddeev1975quantization} discuss the semiclassical rule in 1975. \citet[equation~(2)]{zamolodchikov1977jetp} state the exact mass-ratio formula in May 1977. The notebook\textquotesingle s approximate 1974 attribution is not used to date the full quantum result.}
\GSfield{Present status}{As of 10 November 2026, this is an established result in the stated scope.}
\GSfield{Assessment}{literal statement: correct with missing hypotheses; publication match: partial; verification: read the result; development: none; confidence: high}
\end{GSitem}
\begin{GSitem}{1977-436.11}{E}{Rigorous construction of exponential interactions at small coupling}
\GSfield{Kind; source}{claim; \GSpage{P000749}}
\GSfield{The notes}{Glimm, Jaffe and Fröhlich are invoked for rigorous small-coupling quantization, beside the expression \(m^2\gamma^{-1}(1-\cos\gamma u)\). Equation (26) writes \([\pi(x),u(y)]=i\delta(x-y)\), in the opposite order to the earlier Poisson bracket.}
\GSfield{Statement}{Rigorous construction requires a specified normal ordering, ultraviolet limit, volume limit and normalization. The notebook\textquotesingle s potential has quadratic term \(m^2\gamma u^2/2\); it is not the canonically normalized potential \(m^2\beta^{-2}(1-\cos\beta\varphi)\) with small-oscillation mass \(m\). A theorem for a massive reference field is not automatically an identical theorem for every massless-reference sine-Gordon formulation.}
\GSfield{Reconstruction}{With the convention used in item 1977-436.1, canonical quantization gives \([u(x),\pi(y)]=i\delta(x-y)\). The source changes the bracket order without specifying a corresponding sign convention.}
\GSfield{Fate}{\citet[abstract]{froehlich1975quantized} announced a Wightman theory for a sine-Gordon interaction with a nonvanishing mass term in March 1975. The accessible primary abstract supports the historical construction claim, but does not verify every omitted hypothesis of the notebook\textquotesingle s Hamiltonian statement.}
\GSfield{Present status}{As of 10 November 2026, this is an established result in the stated scope.}
\GSfield{Assessment}{literal statement: correct with missing hypotheses; publication match: partial; verification: publisher metadata checked; development: none; confidence: medium}
\end{GSitem}
\begin{GSitem}{1977-436.12}{E}{Short-range wave operators and a unitary scattering operator}
\GSfield{Kind; source}{theorem; \GSpage{P000749}}
\GSfield{The notes}{For \(H=-\Delta/(2m)+v\), with rapidly decaying \(v\), the double-time expression defines a unitary \(S\)-operator.}
\GSfield{Statement}{For a real bounded short-range potential, for example \(|v(x)|\leq C\langle x\rangle^{-1-\epsilon}\), the wave operators \(W_\pm=\operatorname{s-lim}_{t\to\pm\infty}e^{itH}e^{-itH_0}\) exist and have range \(\mathcal H_{\mathrm{ac}}(H)\). Their completeness implies \(S=W_+^*W_-\) is unitary on the free Hilbert space. Bound states are not additional scattering channels; decay without self-adjointness assumptions is insufficient.}
\GSfield{Fate}{The established short-range theorem is restated with proof in \citet[Theorem~2.4 and \S1]{yafaev2004lectures}, lectures delivered in 2001 and posted in 2004. This later exposition is used to check precise hypotheses, not to date the old scattering theorem.}
\GSfield{Present status}{As of 10 November 2026, this is an established result in the stated scope.}
\GSfield{Assessment}{literal statement: correct with missing hypotheses; publication match: exact; verification: read the result; development: none; confidence: high}
\end{GSitem}
\begin{GSitem}{1977-436.13}{A}{Extracting scattering amplitudes from propagator asymptotics}
\GSfield{Kind; source}{method; \GSpage{P000749-P000750}}
\GSfield{The notes}{Equations (30)--(31) evaluate the propagator along free classical rays and divide by free propagators to recover a momentum-space scattering kernel.}
\GSfield{Statement}{The useful assertion is a distributional large-time extraction after testing against momentum wave packets and retaining the stationary-phase normalization. A pointwise quotient of kernels with omitted normalization is not, by itself, a well-defined equality for a scattering distribution. The required free rays are \(x=pt/m\).}
\GSfield{Reconstruction}{Only the distributional, wave-packet interpretation is asserted here; the exact normalization cannot be recovered from the notebook\textquotesingle s formula alone.}
\GSfield{Fate}{\citet[\S2.1, equations~(1.3)--(1.11)]{faddeevkorepin1978quantum} present the corresponding prescription with its normalization and illustrate it by spectral expansion. The primary article was received 4 July 1977 and published in June 1978; the notebook\textquotesingle s ratio is an abbreviated precursor.}
\GSfield{Present status}{As of 10 November 2026, this is an established result in the stated scope.}
\GSfield{Assessment}{literal statement: correct with missing hypotheses; publication match: partial; verification: read the result; development: none; confidence: medium}
\end{GSitem}
\begin{GSitem}{1977-436.14}{E}{The stationary-phase exponent and amplitude}
\GSfield{Kind; source}{claim; \GSpage{P000750}}
\GSfield{The notes}{Equation (32) gives \(\widetilde\Psi(mx/t)e^{ix^2/(2mt)}\) for an integral with phase \(px-p^2t/(2m)\).}
\GSfield{Statement}{For the convention \(\Psi(x,t)=(2\pi)^{-1}\int\widehat\Psi(p)e^{i(px-p^2t/(2m))}dp\), \(m>0\), and Schwartz amplitude, the leading term as \(t\to+\infty\) on bounded velocity ranges is \(\sqrt{m/(2\pi t)}e^{-i\pi/4}e^{imx^2/(2t)}\widehat\Psi(mx/t)\). Both the mass position in the exponent and the power of \(t\) matter.}
\GSfield{Fate}{This follows directly from the critical point \(p=mx/t\) and Hessian \(-t/m\). Compare \citet[equation~(2.8)]{yafaev2004lectures} in the normalization \(H_0=-\Delta\). The notebook error is faithfully transcribed.}
\GSfield{Present status}{As of 10 November 2026, this is an established result in the stated scope.}
\GSfield{Assessment}{literal statement: false as written; publication match: partial; verification: read the result; development: annotated; confidence: high}
\end{GSitem}
\begin{GSitem}{1977-436.15}{E}{Path integrals with soliton boundary data}
\GSfield{Kind; source}{method; \GSpage{P000750-P000751}}
\GSfield{The notes}{Time-sliced mechanical and field path integrals are written, and sums of solitons are imposed as initial and final configurations to define transition quantities.}
\GSfield{Statement}{These oscillatory path integrals encode propagators and, after asymptotic normalization, scattering amplitudes. They require regularization, renormalization and appropriate boundary conditions; an unqualified infinite product of Lebesgue measures does not define a continuum measure. A complex amplitude is not itself a transition probability.}
\GSfield{Fate}{\citet[\S2--\S3]{korepinfaddeev1975quantization} already use this stationary-phase and renormalization scheme in 1975; \citet[\S2.2]{faddeevkorepin1978quantum} develop the general soliton boundary prescription in the 1978 review.}
\GSfield{Present status}{As of 10 November 2026, this is an established result in the stated scope.}
\GSfield{Assessment}{literal statement: correct with missing hypotheses; publication match: partial; verification: read the result; development: none; confidence: high}
\end{GSitem}
\begin{GSitem}{1977-436.16}{A}{A quantum inverse-scattering method}
\GSfield{Kind; source}{question; \GSpage{P000750}}
\GSfield{The notes}{The notes say that there is not yet a good quantum version of the inverse problem.}
\GSfield{Statement}{Construct quantum scattering-data operators and exchange relations that generate commuting quantum integrals and eigenstates for a specified integrable field model. This is a definite method problem when the model and operator algebra are fixed; it is not an existence theorem for every possible quantization.}
\GSfield{Fate}{\citet[\S2--\S4]{sklyaninfaddeev1978quantum} give an early quantum inverse-scattering construction, received 28 July 1978 and published in December 1978. They explicitly construct the nonlinear-Schrödinger operators, while stating that the corresponding comparison with classical sine-Gordon variables had not yet been completed. Thus the early result has a precise model-dependent scope.}
\GSfield{Present status}{As of 10 November 2026, this is an established result in the stated scope.}
\GSfield{Assessment}{literal statement: correct; publication match: partial; verification: read the result; development: none; confidence: high; problem \hyperref[prob:GP-1977-10]{GP-1977-10}}
\end{GSitem}
\begin{GSitem}{1977-436.17}{E}{The exact sine-Gordon scattering matrix has two charge channels}
\GSfield{Kind; source}{theorem; \GSpage{P000750-P000752}}
\GSfield{The notes}{Zamolodchikov\textquotesingle s guessed exact formula is cited. Equation (46) then writes a single phase times momentum-conserving delta functions for a soliton and an antisoliton.}
\GSfield{Statement}{The sine-Gordon two-particle matrix has transmission and reflection amplitudes in the soliton--antisoliton sector, with unitarity, crossing and prescribed bound-state poles. A single scalar phase describes an appropriate eigenchannel or reflectionless specialization; it is not the generic full charge-channel matrix.}
\GSfield{Fate}{\citet{zamolodchikov1977exact} was received 4 January 1977 and published in June 1977. The separately opened primary announcement \citet[equations~(4)--(10)]{zamolodchikov1977jetp}, submitted 14 April and published 20 May 1977, explicitly displays both amplitudes. This was available before the talk.}
\GSfield{Present status}{As of 10 November 2026, this is an established result in the stated scope.}
\GSfield{Assessment}{literal statement: correct with missing hypotheses; publication match: partial; verification: read the result; development: annotated; confidence: high}
\end{GSitem}
\begin{GSitem}{1977-436.18}{F}{Classical solvability versus quantum solvability}
\GSfield{Kind; source}{conjecture; \GSpage{P000750}}
\GSfield{The notes}{The hypothesis says that every classically solvable problem is solvable in its quantum version.}
\GSfield{Statement}{The source specifies neither solvability, quantization, preservation of integrals, domains, nor a class of models. Quantum anomalies and ordering choices are relevant to any precise formulation, so no universal yes-or-no proposition can be extracted without adding substantial hypotheses.}
\GSfield{Fate}{\citet[\S2]{sklyaninfaddeev1978quantum} provide a successful model construction, not a theorem that all classically integrable systems admit the required quantization. No exact fate is assigned to the unrestricted sentence.}
\GSfield{Present status}{As of 10 November 2026, the general sentence remains under-specified. The precise quantum inverse-scattering method is treated separately.}
\GSfield{Assessment}{literal statement: not determinable; publication match: thematic only; verification: read the result; development: none; confidence: low}
\end{GSitem}
\begin{GSitem}{1977-436.19}{E}{The fluctuation determinant and the sine-Gordon zero mode}
\GSfield{Kind; source}{method; \GSpage{P000751}}
\GSfield{The notes}{Expanding around \(U_{cl}\) produces a Hessian and determinant. For a stationary soliton the notes give \(q(x)=1/\cosh^2x\), call it a Bargmann potential, and refer to a zero eigenfunction.}
\GSfield{Statement}{For the sine-Gordon kink, the spatial Hessian is \(K=-\partial_x^2+m^2(1-2\operatorname{sech}^2(mx))\). Its translation mode \(u_s^{\prime}\) has eigenvalue zero. A determinant must remove this mode and account for the collective coordinate; the relevant object is a renormalized determinant ratio. The isolated positive \(\operatorname{sech}^2x\) is not this Hessian potential.}
\GSfield{Fate}{\citet[\S2, equations~(12)--(18)]{korepinfaddeev1975quantization} explicitly exclude the zero mode, use the reflectionless spatial operator and express the determinant through spectral/scattering information. This is a precise earlier match after the source\textquotesingle s missing sign, coefficient and constant are restored.}
\GSfield{Present status}{As of 10 November 2026, this is an established result in the stated scope.}
\GSfield{Assessment}{literal statement: false as written; publication match: partial; verification: read the result; development: annotated; confidence: high}
\end{GSitem}
\begin{GSitem}{1977-436.20}{F}{Explicit higher-loop functionals of the resolvent kernel}
\GSfield{Kind; source}{question; \GSpage{P000751}}
\GSfield{The notes}{The next-loop discussion calls for \(R=A^{-1}\), its kernel, and \(\int R^3(x,x)\,d^2x\), with the hope that such quantities can be calculated explicitly.}
\GSfield{Statement}{The expression is ambiguous between the diagonal of an operator cube and the cube of a coincident kernel; zero modes, boundary conditions, ultraviolet subtraction and counterterms are unspecified. It therefore does not define one finite invariant whose exact evaluation can be tracked.}
\GSfield{Fate}{\citet[\S3 and \S4]{korepinfaddeev1975quantization} discuss renormalized higher-loop diagrams and explicitly report their computational difficulty. This identifies the setting but not an exact value or theorem for the notebook\textquotesingle s ambiguous functional.}
\GSfield{Present status}{As of 10 November 2026, the displayed functional remains insufficiently specified; it is not entered as an allegedly untouched open problem.}
\GSfield{Assessment}{literal statement: not determinable; publication match: thematic only; verification: read the result; development: none; confidence: medium}
\end{GSitem}
\begin{GSitem}{1977-436.21}{E}{The one-loop sine-Gordon mass correction}
\GSfield{Kind; source}{theorem; \GSpage{P000751}}
\GSfield{The notes}{The stated expansion is \(M=8m/\gamma-m/\pi+O(\gamma)\).}
\GSfield{Statement}{In the semiclassical convention for weak sine-Gordon coupling, with renormalized mass parameter \(m_r\), the one-loop kink mass is \(M=8m_r/\gamma-m_r/\pi+O(\gamma)\). A finite renormalization of the mass parameter changes higher coefficients, so the symbol \(m\) must be specified when comparing schemes.}
\GSfield{Fate}{\citet[\S2, following equation~(18), and \S3]{korepinfaddeev1975quantization} already give this result in 1975, with \(m_r\) the renormalized mass. The 1978 review is not its first publication.}
\GSfield{Present status}{As of 10 November 2026, this is an established result in the stated scope.}
\GSfield{Assessment}{literal statement: correct with missing hypotheses; publication match: exact; verification: read the result; development: none; confidence: high}
\end{GSitem}
\begin{GSitem}{1977-436.22}{F}{The no-higher-corrections claim depends on the mass convention}
\GSfield{Kind; source}{conjecture; \GSpage{P000751}}
\GSfield{The notes}{The notebook conjectures the exact identity \(M=8m/\gamma-m/\pi\), attributed to Zamolodchikov\textquotesingle s justification.}
\GSfield{Statement}{Without defining the renormalized parameter \(m\), the all-orders statement is not determinate. For the physical lightest-breather mass \(m_1\) and the Lagrangian coupling \(g=\beta^2\), the exact ratio is \(M/m_1=[2\sin(\pi\xi/2)]^{-1}\), \(\xi=g/(8\pi-g)\). Its expansion is \(8/g-1/\pi+g/192+O(g^2)\), so the two-term expression is not exact in that physical-mass convention.}
\GSfield{Reconstruction}{The Taylor expansion in the physical mass convention is an editorial algebraic consequence of the published ratio; it is not a claim that the notebook used that convention.}
\GSfield{Fate}{The mass-ratio identity follows from \citet[equation~(2)]{zamolodchikov1977jetp}. \citet[\S2--\S3]{korepinfaddeev1975quantization} explicitly frame the earlier conjecture as the possibility that higher corrections reduce to renormalization. The source does not fix the convention needed to decide its literal all-orders wording.}
\GSfield{Present status}{As of 10 November 2026, the exact mass ratios are known in the specified sine-Gordon convention. The notebook\textquotesingle s bare claim cannot be labeled convention-independently true or false.}
\GSfield{Assessment}{literal statement: not determinable; publication match: partial; verification: read the result; development: annotated; confidence: high}
\end{GSitem}
\begin{GSitem}{1977-436.23}{E}{The coupling threshold is not simply soliton disappearance}
\GSfield{Kind; source}{claim; \GSpage{P000751}}
\GSfield{The notes}{The last sentence says everything is fine for \(\gamma\leq8\pi\), beyond which the soliton disappears.}
\GSfield{Statement}{The threshold \(\beta^2=8\pi\) concerns the quantum sine-Gordon interaction and its renormalization. It does not follow from extrapolating a one-loop mass to zero. Breather bound states already cease to exist when \(\xi\geq1\), equivalently \(\beta^2\geq4\pi\), while soliton and antisoliton scattering persist in the repulsive regime below \(8\pi\). Endpoint claims require separate hypotheses.}
\GSfield{Fate}{\citet[abstract]{coleman1975equivalence} discuss unbounded energy density beyond \(8\pi\) in the stated convention; \citet[equation~(2)]{zamolodchikov1977jetp} distinguish the bound-state threshold. Neither supports the notebook\textquotesingle s unrestricted disappearance sentence.}
\GSfield{Present status}{As of 10 November 2026, this is an established result in the stated scope.}
\GSfield{Assessment}{literal statement: false as written; publication match: partial; verification: read the result; development: annotated; confidence: high}
\end{GSitem}
\begin{GSitem}{1977-436.24}{E}{Semiclassical phase integrals and reflectionless couplings}
\GSfield{Kind; source}{method; \GSpage{P000752}}
\GSfield{The notes}{The abbreviated action gives an integral for the scattering phase. The notes replace it by a sum at \(\gamma=8\pi/N\) and describe interpolation to arbitrary coupling.}
\GSfield{Statement}{The semiclassical phase is obtained from the two-soliton action after subtracting asymptotic external contributions. A finite-product exact scattering formula occurs at reflectionless values \(\Gamma=8\pi/N\), where \(\Gamma\) is the renormalized coupling of the exact scattering formula. In terms of \(\beta^2\), these values are \(8\pi/(N+1)\). A logarithm branch and analytic continuation must be fixed in the phase integral.}
\GSfield{Fate}{\citet[\S4, equations~(27)--(30)]{korepinfaddeev1975quantization} derive the semiclassical phase; \citet[equations~(7)--(10)]{zamolodchikov1977jetp} explicitly discuss the reflectionless interpolation. The notebook uses one symbol for distinct coupling conventions.}
\GSfield{Present status}{As of 10 November 2026, this is an established result in the stated scope.}
\GSfield{Assessment}{literal statement: correct with missing hypotheses; publication match: partial; verification: read the result; development: annotated; confidence: high}
\end{GSitem}
\begin{GSitem}{1977-436.25}{E}{The hyperbolic-secant scattering analogy}
\GSfield{Kind; source}{remark; \GSpage{P000752}}
\GSfield{The notes}{Sine-Gordon scattering is said to resemble a Schrödinger problem in one space dimension \(-d^2/dx^2+a/\cosh^2(\alpha x)\); a possible eigenvalue formula is crossed out.}
\GSfield{Statement}{This is an analogy with the explicitly solvable hyperbolic-secant potential. At special negative strengths \(a=-\ell(\ell+1)\alpha^2\), \(\ell\) a nonnegative integer, the Schrödinger potential is reflectionless. The notebook gives no parameter map identifying its entire Schrödinger spectrum with the quantum field theory, and the deleted eigenvalue is not restored as a theorem.}
\GSfield{Fate}{The exact reflectionless fluctuation operator in \citet[\S2, equation~(16)]{korepinfaddeev1975quantization} and reflectionless quantum scattering in \citet[equation~(7)]{zamolodchikov1977jetp} explain the comparison. They do not identify a literal equivalence of the two Hamiltonians.}
\GSfield{Present status}{As of 10 November 2026, this is an established result in the stated scope.}
\GSfield{Assessment}{literal statement: correct with missing hypotheses; publication match: thematic only; verification: read the result; development: none; confidence: medium}
\end{GSitem}
\begin{GSitem}{1977-436.26}{E}{The constrained sphere sigma model and its classical Lax representation}
\GSfield{Kind; source}{theorem; \GSpage{P000752}}
\GSfield{The notes}{The closing example is \(n:\mathbb R^{1,1}\to S^n\), with Lagrangian \(\frac12\int\partial n\cdot\partial n\,dx\), and a Lax pair attributed to Pohlmeyer and to Zakharov--Mikhailov.}
\GSfield{Statement}{For a unit vector field \(n\in S^n\), the constrained Euler--Lagrange equation is \(\Box n+(\partial_\mu n\cdot\partial^\mu n)n=0\). Its nonlinearity comes from the constraint. The classical sphere sigma models admit a spectral-parameter zero-curvature formulation and conserved quantities; this does not by itself construct their quantum theory.}
\GSfield{Fate}{\citet{pohlmeyer1976integrable} appeared in 1976, received 28 August 1975. \citet[introduction, \S3 and \S6]{zakharovmikhailov1978relativistic}, submitted 4 January 1978 and published June 1978, explicitly identify Pohlmeyer\textquotesingle s sphere-model result and develop the reduction and integration scheme. Thus the classical claim was already known; the latter article is not treated as a theorem of integrability for all homogeneous spaces.}
\GSfield{Present status}{As of 10 November 2026, this is an established result in the stated scope.}
\GSfield{Assessment}{literal statement: correct with missing hypotheses; publication match: exact; verification: read the result; development: none; confidence: high}
\end{GSitem}
\begin{GSitem}{1977-436.27}{B}{Massive sigma-model particles scatter like solitons}
\GSfield{Kind; source}{claim; \GSpage{P000752}}
\GSfield{The notes}{The scan reads \ru{Массивные частицы рассеиваются как солитоны}; both transcriptions instead say that particles appear as solitons.}
\GSfield{Statement}{The intended comparison is elastic, factorized scattering of massive quantum particles: multiparticle amplitudes are assembled from two-particle amplitudes, with no particle production in the integrable description. It is not a claim that all these particles are classical topological kink solutions.}
\GSfield{Reconstruction}{The reading of the Russian verb was checked independently against the scan by both this record editor and the chapter assembler.}
\GSfield{Fate}{\citet[\S1--\S3]{zamolodchikovzamolodchikov1978on} construct an \(O(N)\)-invariant factorized matrix for \(N\geq3\) and support its sigma-model identification by perturbative comparisons, expressly distinguishing that identification from the algebraic construction. The article was received 22 August 1977, revised 14 November 1977, and published in February 1978. This closely matches the scattering remark but is not a proof of a fully axiomatized field theory.}
\GSfield{Present status}{As of 10 November 2026, this is an established result in the stated scope.}
\GSfield{Assessment}{literal statement: correct with missing hypotheses; publication match: partial; verification: read the result; development: annotated; confidence: high}
\end{GSitem}

\GStalk{438}{1978-02-27}{S. I. Gelfand}{P000753-P000754; P000756-P000762}{Exterior algebra modules, projective bundles, and derived categories}
The principal lecture starts on P000757, dated 27 February 1978. P000753--P000754 and P000756 preserve a preceding lecture of 13 February 1978, absent as a separate record from the supplied list; they are included here as explicitly adjacent evidence, not assigned to Faddeev. Both lectures explain the construction later called the Bernstein--Gelfand--Gelfand correspondence: a finite graded exterior-algebra module gives a linear complex of twists on projective space, and suitable modules describe vector bundles modulo free summands. The second lecture develops symmetric/exterior Koszul functors and the language of derived categories. Several displayed formulas are faithfully defective: the total differential lacks a sign, and the proposed unit map vanishes on a trivial module.

\begin{GSitem}{1977-438.1}{E}{Topoi and derived categories as organizing concepts}
\GSfield{Kind; source}{remark; \GSpage{P000753}; I. M. Gelfand}
\GSfield{The notes}{I. M. Gelfand singles out the topos and the derived category as two primary categorical concepts.}
\GSfield{Statement}{This is a methodological comparison, not a classification or universal-property theorem: a topos organizes sheaf-theoretic localization, while a derived category records complexes up to quasi-isomorphism.}
\GSfield{Reconstruction}{The explanatory distinction is editorial; the notebook only lists the two concepts.}
\GSfield{Fate}{The latter framework was already developed in the Harvard seminar of 1963--1964, published as \citet[Chapter~I]{hartshorne1966residues}.}
\GSfield{Present status}{As of 10 November 2026: established in the stated scope; the notebook is exposition of this result.}
\GSfield{Assessment}{literal statement: correct; publication match: partial; verification: publisher metadata checked; development: annotated; confidence: high}
\end{GSitem}
\begin{GSitem}{1977-438.2}{A}{Graded exterior modules and evaluation complexes}
\GSfield{Kind; source}{definition; \GSpage{P000753; P000757-P000758}}
\GSfield{The notes}{\(\Lambda=\Lambda(\Xi)\), \(\dim\Xi=n+1\), \(\deg\xi_i=-1\), and \(V=\bigoplus V_j\) is finite-dimensional. Multiplication by any \(0\ne\xi\in\Xi\) makes \(\cdots\to V_1\to V_0\to V_{-1}\to\cdots\) a complex because \(\xi^2=0\).}
\GSfield{Statement}{Over \(\mathbb C\), the exterior algebra of an \((n+1)\)-dimensional vector space has \(\dim\Lambda_{-k}=\binom{n+1}{k}\). For a finite graded left module, the maps \(v\mapsto\xi v\) have degree \(-1\) and square zero. A module is called proper here when its evaluation complex has homology only in degree zero for every nonzero \(\xi\).}
\GSfield{Fate}{The same definitions precede Theorem~1 of \citet{bernsteingelfandgelfand1978bundles}; submitted 10 January 1978, published in Russian later in 1978.}
\GSfield{Present status}{As of 10 November 2026: the stated construction is established; the qualifications above are essential.}
\GSfield{Assessment}{literal statement: correct; publication match: exact; verification: read the result; development: annotated; confidence: high}
\end{GSitem}
\begin{GSitem}{1977-438.3}{A}{Constant rank and the bundle of evaluation homology}
\GSfield{Kind; source}{claim; \GSpage{P000754; P000759}}
\GSfield{The notes}{For proper \(V\), \(\dim H_0(L_\xi(V))\) is constant by the Euler characteristic; these spaces form a bundle.}
\GSfield{Statement}{For a finite complex with all fiber homology except degree zero vanishing, the remaining dimension is the alternating sum of the dimensions of the terms, hence is independent of the projective point. The cohomology sheaf of the associated locally free complex is locally free and commutes with taking fibers.}
\GSfield{Reconstruction}{The local-freeness and base-change assertion is stated explicitly here; a constant Euler characteristic by itself would not suffice without the other fiber cohomology vanishing.}
\GSfield{Fate}{This is the well-definedness step for \(\Phi(V)\) in \citet[Section~3]{bernsteingelfandgelfand1978bundles}.}
\GSfield{Present status}{As of 10 November 2026: the stated construction is established; the qualifications above are essential.}
\GSfield{Assessment}{literal statement: correct with missing hypotheses; publication match: exact; verification: read the result; development: annotated; confidence: high}
\end{GSitem}
\begin{GSitem}{1977-438.4}{A}{Free exterior modules produce the zero bundle}
\GSfield{Kind; source}{claim; \GSpage{P000754}}
\GSfield{The notes}{The regular module and its shifts have zero evaluation homology. The subsequent sentence says that free modules give a \ru{тривиальное расслоение}.}
\GSfield{Statement}{For \(0\ne\xi\), choose \(\ell\in\Xi^*\) with \(\ell(\xi)=1\); contraction by \(\ell\) contracts multiplication by \(\xi\). Thus every finite graded free \(\Lambda\)-module gives an acyclic complex, hence the rank-zero bundle, not an arbitrary positive-rank trivial bundle.}
\GSfield{Reconstruction}{The contraction proof and explicit rank-zero wording clarify the notebook. No new theorem is claimed.}
\GSfield{Fate}{Acyclicity of free modules is used in \citet[Section~3]{bernsteingelfandgelfand1978bundles}.}
\GSfield{Present status}{As of 10 November 2026: the stated construction is established; the qualifications above are essential.}
\GSfield{Assessment}{literal statement: false as written; publication match: exact; verification: read the result; development: annotated; confidence: high}
\end{GSitem}
\begin{GSitem}{1977-438.5}{A}{Stable exterior-module description of algebraic vector bundles}
\GSfield{Kind; source}{theorem; \GSpage{P000754; P000759}}
\GSfield{The notes}{Every algebraic bundle is \(\Phi(V)\), and \(\Phi(V)\cong\Phi(V')\) iff \(V\oplus P\cong V'\oplus P'\) for free modules. The notes credit Bernstein and both Gelfands, and mention an independent related result of Beilinson.}
\GSfield{Statement}{For \(\mathbb P^n_{\mathbb C}\), isomorphism classes of algebraic vector bundles are in bijection with stable isomorphism classes of finite graded proper \(\Lambda(\mathbb C^{n+1})\)-modules. Stabilization is by finite sums of grading shifts of the regular module.}
\GSfield{Fate}{Exact match: \citet[Theorem~1]{bernsteingelfandgelfand1978bundles}. This is a reduction to module classification, not a list of all bundle moduli.}
\GSfield{Present status}{As of 10 November 2026: the stated construction is established; the qualifications above are essential.}
\GSfield{Assessment}{literal statement: correct; publication match: exact; verification: read the result; development: annotated; confidence: high; problem \hyperref[prob:GP-1977-11]{GP-1977-11}}
\end{GSitem}
\begin{GSitem}{1977-438.6}{E}{The twisting sheaf and its local rank-one property}
\GSfield{Kind; source}{definition; \GSpage{P000754; P000758}}
\GSfield{The notes}{Three definitions of \(O(j)\) are offered: transition functions, homogeneous degree-\(j\) rational functions, and powers of the anti-Hopf bundle. An exercise asks why the resulting bundle is one-dimensional.}
\GSfield{Statement}{For a Zariski-open \(U\subset\mathbb P^n\), sections of \(O(j)\) are homogeneous rational functions of degree \(j\) regular on the inverse image of \(U\) in \(\mathbb C^{n+1}\setminus\{0\}\). On \(x_i\ne0\), multiplication by \(x_i^{-j}\) identifies this module with \(O_U\), proving local rank one. Negative powers mean dual tensor powers.}
\GSfield{Reconstruction}{The topology and omission of the cone vertex are necessary. Regular algebraic sections, rather than arbitrary analytic sections, are intended.}
\GSfield{Fate}{The twisting sheaves and homogeneous-module construction are classical: \citet[Chapter~III]{serre1955coherent}.}
\GSfield{Present status}{As of 10 November 2026: established in the stated scope; the notebook is exposition of this result.}
\GSfield{Assessment}{literal statement: correct with missing hypotheses; publication match: exact; verification: read the result; development: annotated; confidence: high; problem \hyperref[prob:GP-1977-12]{GP-1977-12}}
\end{GSitem}
\begin{GSitem}{1977-438.7}{A}{The linear sheaf complex associated with an exterior module}
\GSfield{Kind; source}{construction; \GSpage{P000754; P000758}}
\GSfield{The notes}{\(\mathcal L^j=V_{-j}\otimes O(j)\) and \((df)(\xi)=\xi f(\xi)\).}
\GSfield{Statement}{The tautological section defines maps \(V_{-j}\otimes O(j)\to V_{-j-1}\otimes O(j+1)\). The exterior relations give \(d^2=0\); the construction is independent of a basis of \(\Xi\).}
\GSfield{Fate}{This is precisely \(L(V)\) in \citet[Section~3]{bernsteingelfandgelfand1978bundles}.}
\GSfield{Present status}{As of 10 November 2026: the stated construction is established; the qualifications above are essential.}
\GSfield{Assessment}{literal statement: correct; publication match: exact; verification: read the result; development: annotated; confidence: high}
\end{GSitem}
\begin{GSitem}{1977-438.8}{A}{A truncated Koszul kernel followed by syzygies}
\GSfield{Kind; source}{method; \GSpage{P000756}}
\GSfield{The notes}{The page starts mid-proof: \(C_N\otimes\Lambda\to C_{N+1}\otimes\Lambda\), \(c\otimes\lambda\mapsto\sum X_i(c)\otimes\xi_i\lambda\); take its kernel and then \(N\) free-kernel steps with grading shifts.}
\GSfield{Statement}{If \(C\) is a graded module over \(S=\mathbb C[x_0,\ldots,x_n]\), multiplication by the coordinate functions gives the displayed basis-independent linear map. For sufficiently high truncation and suitable grading conventions, a Koszul-dual module may be recovered using successive syzygies.}
\GSfield{Reconstruction}{The missing preceding P000755 is not attached. C, the admissible bound on N, and the exact shift cannot be recovered from P000756 alone; no unique explicit inverse is asserted.}
\GSfield{Fate}{The structural inverse belongs to the BGG equivalence, \citet[Theorem~2]{bernsteingelfandgelfand1978bundles}; this page's precise normalization is only partially matched.}
\GSfield{Present status}{As of 10 November 2026: the stated construction is established; the qualifications above are essential.}
\GSfield{Assessment}{literal statement: correct with missing hypotheses; publication match: partial; verification: read the result; development: annotated; confidence: medium}
\end{GSitem}
\begin{GSitem}{1977-438.9}{E}{Serre's quotient-category description of coherent sheaves}
\GSfield{Kind; source}{theorem; \GSpage{P000756; P000759}}
\GSfield{The notes}{Gelfand urges replacing bundles by sheaves, since bundles do not form an abelian category. Coherent sheaves on projective space correspond to finitely generated graded polynomial modules modulo finitely many components.}
\GSfield{Statement}{For \(S=\mathbb C[x_0,\ldots,x_n]\), sheafification induces \(\operatorname{grmod}(S)/\operatorname{tors}(S)\simeq\operatorname{Coh}(\mathbb P^n)\), where torsion means finite-dimensional graded modules. A representative is \(\bigoplus_{j\ge N}H^0(\mathcal E(j))\) for sufficiently large \(N\).}
\GSfield{Reconstruction}{Truncation avoids claiming the unbounded section module is finitely generated for every coherent sheaf; e.g. a skyscraper has sections in every negative twist.}
\GSfield{Fate}{\citet[Chapter~III]{serre1955coherent} is the relevant 1955 publication; the notebook's 1953 bibliographic date is inaccurate.}
\GSfield{Present status}{As of 10 November 2026: established in the stated scope; the notebook is exposition of this result.}
\GSfield{Assessment}{literal statement: correct with missing hypotheses; publication match: exact; verification: read the result; development: annotated; confidence: high}
\end{GSitem}
\begin{GSitem}{1977-438.10}{E}{Sheafification by degree-zero localization}
\GSfield{Kind; source}{construction; \GSpage{P000756}}
\GSfield{The notes}{On \(E_i=\{x_i\ne0\}\), sections are fractions \(m/x_i^s\) with \(\deg m=s\). The grading shift of \(S\) gives \(O(N)\).}
\GSfield{Statement}{For a graded \(S\)-module \(M\), define \(\widetilde M(D_+(x_i))=(M_{x_i})_0\) and glue by the canonical localization maps. With \(M(N)_j=M_{j+N}\), \(\widetilde{S(N)}=O(N)\).}
\GSfield{Fate}{This is Serre's sheafification construction; see \citet[Chapter~III]{serre1955coherent}.}
\GSfield{Present status}{As of 10 November 2026: established in the stated scope; the notebook is exposition of this result.}
\GSfield{Assessment}{literal statement: correct; publication match: exact; verification: read the result; development: annotated; confidence: high}
\end{GSitem}
\begin{GSitem}{1977-438.11}{A}{The symmetric-algebra linear complex}
\GSfield{Kind; source}{construction; \GSpage{P000759}}
\GSfield{The notes}{\(W^k=V_{-k}\otimes S\), \(d(v\otimes s)=\sum\xi_i v\otimes x_i s\), and \((W^k)_j=V_{-k}\otimes S_{j+k}\).}
\GSfield{Statement}{The differential is homogeneous of internal degree zero and cohomological degree one; its square vanishes because the \(x_i\) commute and the \(\xi_i\) anticommute. Sheafifying this complex recovers \(\mathcal L(V)\).}
\GSfield{Fate}{This is the polynomial form of the functor in \citet[Sections~3--4]{bernsteingelfandgelfand1978bundles}.}
\GSfield{Present status}{As of 10 November 2026: the stated construction is established; the qualifications above are essential.}
\GSfield{Assessment}{literal statement: correct; publication match: exact; verification: read the result; development: annotated; confidence: high}
\end{GSitem}
\begin{GSitem}{1977-438.12}{E}{Totalization and the missing Koszul sign}
\GSfield{Kind; source}{construction; \GSpage{P000760}}
\GSfield{The notes}{For a complex \((V^i,\partial)\), the written total differential is \(\partial\otimes1+\sum\xi_i\otimes x_i\), with a visible plus but no sign.}
\GSfield{Statement}{For commuting bicomplex differentials \(d_h,d_v\), totalize using \(d_{\rm tot}=d_h+(-1)^p d_v\) on horizontal degree \(p\). Equivalently first impose an anticommuting bicomplex convention. Without either convention the displayed sum need not square to zero in characteristic zero.}
\GSfield{Reconstruction}{The required sign is an editorial repair, not a newly read mark in the manuscript.}
\GSfield{Fate}{Totalization conventions are part of the pre-existing derived-category apparatus; \citet[Chapter~I]{hartshorne1966residues}.}
\GSfield{Present status}{As of 10 November 2026: established in the stated scope; the notebook is exposition of this result.}
\GSfield{Assessment}{literal statement: false as written; publication match: partial; verification: publisher metadata checked; development: annotated; confidence: high}
\end{GSitem}
\begin{GSitem}{1977-438.13}{A}{The Hom Koszul-dual functor}
\GSfield{Kind; source}{construction; \GSpage{P000760}}
\GSfield{The notes}{\(V=\operatorname{Hom}_{\mathbb C}(\Lambda,W)\), with a regrading, and \(\partial(v)(\lambda)=d(v(\lambda))+\sum x_i v(\xi_i\lambda)\).}
\GSfield{Statement}{The reverse Koszul functor uses the graded Hom complex from \(\Lambda\) into a complex of graded \(S\)-modules, with differential combining the original differential and the contraction induced by \(\sum x_i\otimes\xi_i\). The totalization must use compatible Koszul signs and boundedness conditions.}
\GSfield{Reconstruction}{The full grading/sign convention is reconstructed; the displayed unsigned expression is not endorsed as a chain differential.}
\GSfield{Fate}{The induced derived equivalence is \citet[Theorem~2]{bernsteingelfandgelfand1978bundles}; the notebook's explicit signs are not an exact formula-level match.}
\GSfield{Present status}{As of 10 November 2026: the stated construction is established; the qualifications above are essential.}
\GSfield{Assessment}{literal statement: correct with missing hypotheses; publication match: partial; verification: read the result; development: annotated; confidence: high}
\end{GSitem}
\begin{GSitem}{1977-438.14}{A}{The unit comparison and its erroneous value at one}
\GSfield{Kind; source}{claim; \GSpage{P000761}}
\GSfield{The notes}{A natural quasi-isomorphism \(i\to GF\) is asserted. Its displayed value is \(\varphi(v)(1)=\sum x_i\otimes\xi_i v\); for the one-dimensional trivial module the target is said to be the Koszul complex.}
\GSfield{Statement}{The Koszul adjunction has a natural unit represented, in compatible conventions, by the action map \(v\mapsto[\lambda\mapsto\lambda v]\), including the degree-zero component \(1\mapsto1\otimes v\). The notebook's value alone is zero on a nonzero trivial module and therefore cannot induce the claimed cohomology isomorphism.}
\GSfield{Reconstruction}{The repair specifies the essential nonzero constant component; it is not a full reconstruction of every grading convention.}
\GSfield{Fate}{The equivalence is established in \citet[Theorem~2]{bernsteingelfandgelfand1978bundles}; the scanned formula is a faithful notebook error.}
\GSfield{Present status}{As of 10 November 2026: the stated construction is established; the qualifications above are essential.}
\GSfield{Assessment}{literal statement: false as written; publication match: partial; verification: read the result; development: annotated; confidence: high}
\end{GSitem}
\begin{GSitem}{1977-438.15}{E}{Quasi-isomorphisms and the derived category}
\GSfield{Kind; source}{definition; \GSpage{P000761-P000762}}
\GSfield{The notes}{The lecture defines complexes, commuting chain maps and cohomology. A quasi-isomorphism induces isomorphisms on cohomology; \(D(A)\) is obtained by inverting these maps.}
\GSfield{Statement}{For an abelian category \(\mathcal A\), the derived category is the localization \(K(\mathcal A)[S^{-1}]\) of the homotopy category at quasi-isomorphisms, with a chosen boundedness convention. Objects are complexes and cohomology sends quasi-isomorphisms to isomorphisms.}
\GSfield{Reconstruction}{The intermediate homotopy category and size conventions are made explicit; strict maps alone do not support the naive Ore formulas.}
\GSfield{Fate}{The notebook itself cites \citet[Chapter~I]{hartshorne1966residues}, published 1966, and Verdier's earlier work.}
\GSfield{Present status}{As of 10 November 2026: established in the stated scope; the notebook is exposition of this result.}
\GSfield{Assessment}{literal statement: correct with missing hypotheses; publication match: partial; verification: publisher metadata checked; development: annotated; confidence: high}
\end{GSitem}
\begin{GSitem}{1977-438.16}{E}{Roofs, Ore squares, and cancellation}
\GSfield{Kind; source}{method; \GSpage{P000761-P000762}}
\GSfield{The notes}{Morphisms are roofs \(X\leftarrow Z\to Y\); common refinements identify roofs. Composition closure, Ore squares, and cancellation are listed.}
\GSfield{Statement}{In the homotopy category of complexes, quasi-isomorphisms admit a calculus of fractions: roofs are identified by common refinements commuting up to homotopy. Identity maps belong to the multiplicative system. An arbitrary chain of refinements is not itself a morphism or a substitute for the equivalence relation.}
\GSfield{Reconstruction}{The notebook suppresses homotopies and identity axioms; these are necessary for its purported precise definition.}
\GSfield{Fate}{The established localization framework is the one cited in \citet[Chapter~I]{hartshorne1966residues}.}
\GSfield{Present status}{As of 10 November 2026: established in the stated scope; the notebook is exposition of this result.}
\GSfield{Assessment}{literal statement: correct with missing hypotheses; publication match: partial; verification: publisher metadata checked; development: annotated; confidence: high}
\end{GSitem}
\begin{GSitem}{1977-438.17}{A}{Reduction of Koszul comparisons by finite filtrations}
\GSfield{Kind; source}{method; \GSpage{P000762}}
\GSfield{The notes}{Check the natural comparison on the one-dimensional trivial module, then introduce filtrations with standard factors.}
\GSfield{Statement}{A finite graded exterior module has a finite filtration with one-dimensional trivial graded quotients. If two exact functors and their natural transformation preserve the relevant triangles, the quasi-isomorphism statement extends from such quotients by the long exact cohomology sequence.}
\GSfield{Reconstruction}{Exactness of the functors and the finite filtration are the hypotheses implicit in the proof strategy.}
\GSfield{Fate}{This is a proof method for the correspondence of \citet[Theorem~2]{bernsteingelfandgelfand1978bundles}, not an additional open problem.}
\GSfield{Present status}{As of 10 November 2026: the stated construction is established; the qualifications above are essential.}
\GSfield{Assessment}{literal statement: correct with missing hypotheses; publication match: partial; verification: read the result; development: annotated; confidence: high}
\end{GSitem}

\GStalk{439}{1978-03-06}{K. Aomoto}{P000763-P000766}{Neutral massive fields, Clifford operators, and Painleve equations}
The lecture constructs neutral fermions and twist operators in two dimensions and relates their correlation functions to holonomic systems and Painlev\'e equations. The decisive later comparison is Sato--Miwa--Jimbo's neutral-field construction, not merely the general phrase ``isomonodromy.'' Gelfand first asks about the smooth ideal of a point and ends with Garnier's isomonodromic description of Painlev\'e VI. P000767 begins the next lecture and is not assigned here.

\begin{GSitem}{1977-439.1}{E}{The smooth ideal of a point as a line-bundle module}
\GSfield{Kind; source}{question; \GSpage{P000763}; I. M. Gelfand}
\GSfield{The notes}{Gelfand asks whether the sheaf of smooth functions on \((-1,1)\) vanishing at zero is a sheaf of sections of a vector bundle, and draws multiplication by \(x\) as an isomorphism.}
\GSfield{Statement}{For \(I=(-1,1)\), let \(\mathcal I(U)=\{f\in C^\infty(U):f(0)=0\text{ if }0\in U\}\). Multiplication by \(x\) is an isomorphism \(C_I^\infty\to\mathcal I\), so \(\mathcal I\) is locally free of rank one, despite its inclusion having zero evaluation at the origin.}
\GSfield{Reconstruction}{The formula proves surjectivity; continuity and vanishing off a single point prove injectivity. The associated bundle fiber is a module fiber, not the image of pointwise evaluation of the ideal.}
\GSfield{Fate}{This is the one-variable smooth division lemma: near zero, \(f(x)=x\int_0^1 f'(tx)\,dt\); away from zero divide by \(x\). It parallels the algebraic affine-line example indexed as 1965-70.10, but that earlier entry is not an identical smooth problem.}
\GSfield{Present status}{As of 10 November 2026: established in the stated scope; the notebook is exposition of this result.}
\GSfield{Assessment}{literal statement: correct; publication match: partial; verification: indirect; development: annotated; confidence: high}
\end{GSitem}
\begin{GSitem}{1977-439.2}{B}{Rapidity fermions and their canonical anticommutation relations}
\GSfield{Kind; source}{definition; \GSpage{P000763}}
\GSfield{The notes}{\(\psi(-u)=\psi(u)^+\), \(\{\psi(u),\psi(u')\}=2\pi|u|\delta(u+u')\), with invariant measure \(du/(2\pi|u|)\) on \(\mathbb R^*\).}
\GSfield{Statement}{These are operator-valued distributions satisfying the real neutral CAR after smearing against admissible test functions. Positive and negative rapidity modes are adjoints. The singular point \(u=0\), a common domain, and the vacuum representation must be specified before treating these expressions as operators.}
\GSfield{Fate}{The neutral free-fermion conventions and their analytic realization appear in \citet[Section~4.1]{satomiwajimbo1979holonomicIV}, received 20 October 1978 and published 1979.}
\GSfield{Present status}{As of 10 November 2026: the stated construction is established; the qualifications above are essential.}
\GSfield{Assessment}{literal statement: correct with missing hypotheses; publication match: partial; verification: read the result; development: annotated; confidence: high}
\end{GSitem}
\begin{GSitem}{1977-439.3}{B}{Massive Dirac solutions and the conserved pairing}
\GSfield{Kind; source}{construction; \GSpage{P000763-P000764}}
\GSfield{The notes}{\(\partial_{x^\pm}w_\pm\mp mw_\mp=0\), \(x^\pm=(x^0\pm x^1)/2\), and \(\langle w,w'\rangle=\int(w_+w'_++w_-w'_-)\,dx^1\). A rapidity integral with phase \(-im(x^-u+x^+/u)\) gives the fields \(\psi_\pm\).}
\GSfield{Statement}{For sufficiently decaying smooth solutions the bilinear pairing is time independent: the time derivative is a spatial divergence and its integral vanishes. The rapidity amplitudes must satisfy the two Dirac equations, with compatible square-root boundary values and reciprocal powers. Smearing the fields against a solution then gives a time-independent Clifford generator.}
\GSfield{Reconstruction}{The decay assumption, the distributional interpretation, and the requirement on reciprocal amplitudes are supplied. No unresolved radical mark is silently converted into a verified exact formula.}
\GSfield{Fate}{Compare the explicit neutral-Dirac realization in \citet[Section~4.1]{satomiwajimbo1979holonomicIV}. The notebook's superscripts on the square root do not by themselves fix every boundary-value convention.}
\GSfield{Present status}{As of 10 November 2026: the stated construction is established; the qualifications above are essential.}
\GSfield{Assessment}{literal statement: correct with missing hypotheses; publication match: partial; verification: read the result; development: annotated; confidence: medium}
\end{GSitem}
\begin{GSitem}{1977-439.4}{F}{Support decomposition versus isotropic Fock polarization}
\GSfield{Kind; source}{claim; \GSpage{P000764}}
\GSfield{The notes}{The solution space is decomposed orthogonally into solutions vanishing in opposite spacelike wedges. The same symbols \(W_x^\pm\) subsequently satisfy \(\{\psi^+_\mu,\psi^-_\nu\}=\delta_{\mu\nu}\) and same-sign anticommutators zero.}
\GSfield{Statement}{A real support decomposition is orthogonal for the conserved symmetric pairing. In contrast, a creation/annihilation decomposition of a complex Clifford space is into dual isotropic subspaces. These cannot be the same decomposition with the same nondegenerate pairing without a further change of polarization.}
\GSfield{Reconstruction}{Both assertions are retained separately. The omitted change of basis or complexification is not recoverable uniquely from these pages.}
\GSfield{Fate}{\citet[Sections~1 and~4.1--4.2]{satomiwajimbo1979holonomicIV} distinguishes Clifford polarizations and field localization. It identifies the surrounding construction but does not validate the notebook's conflation.}
\GSfield{Present status}{As of 10 November 2026: the two standard constructions are understood, but their identification in this notebook is not a valid statement as written.}
\GSfield{Assessment}{literal statement: false as written; publication match: thematic only; verification: read the result; development: annotated; confidence: medium}
\end{GSitem}
\begin{GSitem}{1977-439.5}{E}{Clifford normal ordering and its exterior-algebra symbol}
\GSfield{Kind; source}{construction; \GSpage{P000764}}
\GSfield{The notes}{The product map is inverted by Wick notation \({:}\varphi{:}\), giving \(A(W)\longleftrightarrow\Lambda W_x^+\wedge\Lambda W_x^-\).}
\GSfield{Statement}{For a finite-dimensional complex quadratic space with complementary dual isotropic subspaces \(W^+,W^-\), multiplication gives a vector-space isomorphism \(\Lambda W^+\otimes\Lambda W^-\simeq\operatorname{Cl}(W)\). It is not an isomorphism for the ordinary exterior product; the transported product contains contractions. Infinite-dimensional versions need an appropriate algebraic restriction or completion.}
\GSfield{Reconstruction}{The category of isomorphism is made explicit. The missing minus on the second factor of (10) is inferred from the following line, not read as a clear scanned mark.}
\GSfield{Fate}{The normal-symbol calculus is the established input to \citet[Section~1]{satomiwajimbo1979holonomicIV}; the paper's later date does not make the Clifford PBW identity new.}
\GSfield{Present status}{As of 10 November 2026: established in the stated scope; the notebook is exposition of this result.}
\GSfield{Assessment}{literal statement: correct with missing hypotheses; publication match: partial; verification: read the result; development: annotated; confidence: high}
\end{GSitem}
\begin{GSitem}{1977-439.6}{B}{The neutral twist operator and its quadratic kernel}
\GSfield{Kind; source}{construction; \GSpage{P000765}}
\GSfield{The notes}{Conjugation by \(\varphi_F(x)\) changes \(w^++w^-\) into \(w^+-w^-\). The notes acknowledge scalar ambiguity and write \(\varphi_F={:}\exp L_F{:}\), where the quadratic kernel is \(-i(u-u')/(u+u'-i0)\) times the two-particle rapidity phase.}
\GSfield{Statement}{A Clifford implementer of a fixed orthogonal transformation is only determined up to the kernel of the conjugation action. A vacuum normalization and a completed normal-ordered exponential specify a neutral twist field; the singular kernel is a distributional boundary value. The normalizer of \(W\) in all invertible Clifford elements is a Clifford group, not the normalized spin group without parity and norm restrictions.}
\GSfield{Reconstruction}{The notebook definition of Spin is repaired only at the group-theoretic level. Its infinite-dimensional implementation and exact kernel conventions are not claimed to follow from the finite-dimensional spin double cover.}
\GSfield{Fate}{The neutral twist-field construction is developed explicitly in \citet[Section~4.2]{satomiwajimbo1979holonomicIV}; publication followed the seminar in 1979.}
\GSfield{Present status}{As of 10 November 2026: the stated construction is established; the qualifications above are essential.}
\GSfield{Assessment}{literal statement: false as written; publication match: partial; verification: read the result; development: annotated; confidence: high}
\end{GSitem}
\begin{GSitem}{1977-439.7}{B}{The companion odd field and vacuum modules}
\GSfield{Kind; source}{construction; \GSpage{P000765}}
\GSfield{The notes}{\(\varphi^F(x)=\Psi_0(x)\exp L_F(x)\), with \(\Psi_0=\int\underline{du}\,\Psi(u)e^{-im(x^-u+x^+/u)}\). Left and right quotients of the Clifford algebra are called \(|\mathrm{vac}\rangle\) and \(\langle\mathrm{vac}|\).}
\GSfield{Statement}{An odd neutral field is obtained by multiplying the Gaussian twist construction by a linear fermionic insertion, with compatible normal ordering. The quotients \(A/AW^-\) and \(W^+A\backslash A\) are cyclic modules; their classes of \(1\) are vacuum vectors annihilated by the specified subspaces. They are not themselves one-dimensional vectors.}
\GSfield{Reconstruction}{Module, cyclic vector, and pairing are separated; the notebook compresses them into one symbol.}
\GSfield{Fate}{Compare \citet[Sections~1 and~4.2]{satomiwajimbo1979holonomicIV}. The missing Wick colons in (15) leave its literal operator-product normalization uncertain.}
\GSfield{Present status}{As of 10 November 2026: the stated construction is established; the qualifications above are essential.}
\GSfield{Assessment}{literal statement: correct with missing hypotheses; publication match: partial; verification: read the result; development: annotated; confidence: medium}
\end{GSitem}
\begin{GSitem}{1977-439.8}{B}{Mixed correlation functions and holonomic deformation}
\GSfield{Kind; source}{construction; \GSpage{P000765-P000766}}
\GSfield{The notes}{Vacuum expectations of fermions and twist fields are formed. For one moving fermion and one odd insertion, \(W^i_{F,n}(x)\) satisfies a system involving \(M_F=z\partial_z-\bar z\partial_{\bar z}+\tfrac12\operatorname{diag}(1,-1)\); the system is called maximally overdetermined and a Gauss--Manin connection.}
\GSfield{Statement}{After fixing ordered, distinct insertion points and a branch and a normalization, these correlations are governed by a finite holonomic system and isomonodromic deformation equations. The matrices in the notebook must be determined from the insertion data; calling them constant in \(x\) does not mean independent of all deformation parameters.}
\GSfield{Reconstruction}{The notebook does not provide all matrix entries or integrability equations. The modern identification is structural rather than a verified equality with every displayed coefficient.}
\GSfield{Fate}{\citet[Section~4.4]{satomiwajimbo1979holonomicIV} develops the deformation equations and \citet{satomiwajimbo1978holonomic} supplies the general holonomic-field framework, submitted September 1977.}
\GSfield{Present status}{As of 10 November 2026: the stated construction is established; the qualifications above are essential.}
\GSfield{Assessment}{literal statement: correct with missing hypotheses; publication match: partial; verification: read the result; development: annotated; confidence: high}
\end{GSitem}
\begin{GSitem}{1977-439.9}{F}{The defective second equation in the holonomic system}
\GSfield{Kind; source}{claim; \GSpage{P000766}}
\GSfield{The notes}{Equation (23) literally reads \((m-\Gamma)w=0\), \(\Gamma=\begin{pmatrix}0&1\\1&0\end{pmatrix}\), with no differential operator before \(\Gamma\).}
\GSfield{Statement}{Taken literally with scalar mass \(m\), this forces \(w=0\) whenever \(m\ne\pm1\); it is not the massive Dirac equation for arbitrary mass. A derivative-containing operator has been omitted or a symbol was misrecorded in the original notebook.}
\GSfield{Reconstruction}{No conjectural replacement is promoted to a source reading.}
\GSfield{Fate}{The actual neutral Dirac equations in \citet[Section~4.1]{satomiwajimbo1979holonomicIV} confirm the context but do not uniquely repair (23). The scan agrees with both transcriptions.}
\GSfield{Present status}{As of 10 November 2026: the intended precise second operator is not identifiable from the supplied page.}
\GSfield{Assessment}{literal statement: false as written; publication match: thematic only; verification: read the result; development: annotated; confidence: high}
\end{GSitem}
\begin{GSitem}{1977-439.10}{E}{The two-point Ising scaling function and Painleve III}
\GSfield{Kind; source}{theorem; \GSpage{P000766}}
\GSfield{The notes}{The twist correlation \(\tau=\langle\varphi_F(a_1)\cdots\varphi_F(a_n)\rangle\) is called a tau function; for \(n=2\), the notes refer to Painlev\'e III and cite Wu, McCoy, and Tracy, Physical Review B 13 (1976).}
\GSfield{Statement}{For the specified scaling limit, phase, and physical boundary conditions of the zero-field planar Ising two-point correlation, the radial correlation is expressed through a distinguished Painlev\'e III solution. The tau function is generally an exponential integral built from that solution, not the Painlev\'e dependent variable itself.}
\GSfield{Reconstruction}{The scaling and boundary conditions, and the distinction between a Painleve solution and its tau function, are necessary qualifications.}
\GSfield{Fate}{The intended reference is \citet{wumccoytracybarouch1976ising}, published 1 January 1976; E. Barouch is a fourth author. The neutral-field extension is in \citet[Sections~4.4--4.6]{satomiwajimbo1979holonomicIV}.}
\GSfield{Present status}{As of 10 November 2026: established in the stated scope; the notebook is exposition of this result.}
\GSfield{Assessment}{literal statement: correct with missing hypotheses; publication match: partial; verification: publisher metadata checked; development: annotated; confidence: high}
\end{GSitem}
\begin{GSitem}{1977-439.11}{E}{Fuchsian isomonodromy and Painleve VI}
\GSfield{Kind; source}{theorem; \GSpage{P000766}; I. M. Gelfand}
\GSfield{The notes}{Gelfand writes a scalar second-order Fuchsian equation with singularities \(0,1,t,\infty\) and an apparent point \(\lambda\). Constancy of monodromy is related to the displayed second-order equation for \(\lambda(t)\), with a corrected, unclear sign.}
\GSfield{Statement}{For a nonresonant rank-two four-pole Fuchsian system on \(\mathbb P^1\), after reduction to scalar form with one apparent singularity, the Schlesinger isomonodromy equations yield Painlev\'e VI for its location. Accessory parameters must satisfy the apparentness and fixed local-exponent constraints; arbitrary coefficients in the notebook's scalar equation do not suffice.}
\GSfield{Reconstruction}{The rank-two formulation and nonresonance are a modern reconstruction. The obscured sign and the parameter dictionary in (B) are not claimed checked term by term against Garnier.}
\GSfield{Fate}{Garnier's memoir treats the relation between fixed critical points and isomonodromic equations: \citet[Introduction, pp.~1--6, especially p.~5]{garnier1912equations}.}
\GSfield{Present status}{As of 10 November 2026: established in the stated scope; the notebook is exposition of this result.}
\GSfield{Assessment}{literal statement: correct with missing hypotheses; publication match: partial; verification: read the result; development: annotated; confidence: medium}
\end{GSitem}
\begin{GSitem}{1977-439.12}{E}{What fixed singularities mean in the Painleve classification}
\GSfield{Kind; source}{claim; \GSpage{P000766}; I. M. Gelfand}
\GSfield{The notes}{The notes say that six second-order types have fixed singular points and that only these have that property.}
\GSfield{Statement}{The Painlev\'e property excludes movable critical branch points; it does not exclude movable poles. The six irreducible Painlev\'e families arise within a specified classification of second-order algebraic differential equations after allowed transformations, alongside equations reducible to classical functions.}
\GSfield{Reconstruction}{Both the distinction between a critical point and a pole and the restricted class of equations are supplied.}
\GSfield{Fate}{The historical fixed-critical-point problem is the one studied in \citet[Introduction, pp.~1--3]{garnier1912equations}. The unrestricted sentence in the notes is broader than that classification.}
\GSfield{Present status}{As of 10 November 2026: established in the stated scope; the notebook is exposition of this result.}
\GSfield{Assessment}{literal statement: false as written; publication match: partial; verification: read the result; development: annotated; confidence: high}
\end{GSitem}

\GStalk{440}{1978-03-13}{K. Aomoto}{P000767-P000771}{Hypergeometric integrals and twisted cohomology of arrangements}
Aomoto passes from Gauss and Appell series to integrals of products of powers of affine linear forms. The lecture identifies the parameter discriminant, twisted de Rham differential, logarithmic complex, and dual twisted cycles. The positive-integer nonresonance condition as explicitly defined in the notes is insufficient for the stated vanishing theorem; zero weights provide a counterexample. An omitted Euler-characteristic sign is independent of that defect. Gelfand adds a fractional-calculus viewpoint.

\begin{GSitem}{1977-440.1}{E}{Gauss and Appell hypergeometric series}
\GSfield{Kind; source}{definition; \GSpage{P000767-P000768}}
\GSfield{The notes}{The notes give \({}_2F_1(\alpha,\beta;\gamma;x)\), Appell \(F_1\), rising factorials \((\alpha,n)\), and a putative \(F_4\) series missing its variables.}
\GSfield{Statement}{Writing \((a)_n=a(a+1)\cdots(a+n-1)\), \[F_1(a;b,b';c;x,y)=\sum_{m,n\ge0}\frac{(a)_{m+n}(b)_m(b')_n}{(c)_{m+n}m!n!}x^my^n,\] and \[F_4(a,b;c,c';x,y)=\sum_{m,n\ge0}\frac{(a)_{m+n}(b)_{m+n}}{(c)_m(c')_n m!n!}x^my^n.\] Denominator parameters avoid nonpositive integers unless a limiting or terminating interpretation is specified. The series define analytic germs in their convergence regions.}
\GSfield{Reconstruction}{The absent factor in F4 is an emendation; F1 initially converges for both moduli less than 1, and F4 for square-root moduli summing to less than 1.}
\GSfield{Fate}{These are classical functions used as motivating examples for the integral framework of \citet[Introduction]{aomoto1977structure}. No new series definition is claimed in the seminar.}
\GSfield{Present status}{As of 10 November 2026: established in the stated scope; the notebook is exposition of this result.}
\GSfield{Assessment}{literal statement: false as written; publication match: thematic only; verification: indirect; development: annotated; confidence: high}
\end{GSitem}
\begin{GSitem}{1977-440.2}{E}{Fractional differentiation as a hypergeometric construction}
\GSfield{Kind; source}{method; \GSpage{P000768}; I. M. Gelfand}
\GSfield{The notes}{Gelfand starts with \(\frac{d^n}{dx^n}(x^2-1)^n\), generalizes the exponents, and writes \(I_\alpha f=\Gamma(\alpha)^{-1}\int_0^x(x-t)^{\alpha-1}f(t)\,dt\), \(D^\alpha=I_{-\alpha}\) by analytic continuation.}
\GSfield{Statement}{For \(\Re\alpha>0\) and suitable functions on a fixed ray, the Riemann--Liouville integral is defined by that formula. Extending the distribution \(x_+^{\alpha-1}/\Gamma(\alpha)\) gives fractional derivatives with fixed endpoint conventions. Euler's beta integral then gives, for \(\Re a,\Re c>0\), \[\int_0^x t^{a-1}(1-t)^{b-1}(x-t)^{c-1}dt=x^{a+c-1}B(a,c){}_2F_1(1-b,a;a+c;x).\]}
\GSfield{Reconstruction}{Branches, endpoint assumptions, and the normalization factor are made explicit.}
\GSfield{Fate}{The displayed identity follows by substituting \(t=xs\) and integrating the binomial series term by term in its initial domain; analytic continuation extends it. Rodrigues' expression is \(2^nn!P_n(x)\), so its normalization is suppressed in the notes.}
\GSfield{Present status}{As of 10 November 2026: established in the stated scope; the notebook is exposition of this result.}
\GSfield{Assessment}{literal statement: correct with missing hypotheses; publication match: partial; verification: indirect; development: annotated; confidence: high}
\end{GSitem}
\begin{GSitem}{1977-440.3}{E}{Appell integrals from line and conic configurations}
\GSfield{Kind; source}{construction; \GSpage{P000768-P000769}}
\GSfield{The notes}{Four affine lines in general position represent \(F_1\); degenerate line arrangements give \(F_2,F_3\), and replacing two lines by a quadratic form gives \(F_4\). Three lines are normalized by an affine change of variables.}
\GSfield{Statement}{After choosing powers, a branch, and a suitable twisted integration cycle, periods of products of affine linear forms give multivariable hypergeometric functions. Changes in the configuration and a quadratic factor produce other Appell-type integral representations. Integration coordinates must be distinguished from the coefficients serving as parameters.}
\GSfield{Reconstruction}{The repeated letters a,b for both integration coordinates and fourth-line parameters are separated editorially; no exact formula for F2,F3 or the conic integral is reconstructed from the drawing alone.}
\GSfield{Fate}{\citet[Introduction and Section~4]{aomoto1977structure} studies this period construction. The supplied pictures alone do not determine a unique parameter dictionary for all four Appell functions.}
\GSfield{Present status}{As of 10 November 2026: established in the stated scope; the notebook is exposition of this result.}
\GSfield{Assessment}{literal statement: correct with missing hypotheses; publication match: partial; verification: indirect; development: annotated; confidence: medium}
\end{GSitem}
\begin{GSitem}{1977-440.4}{E}{Rational ratios of neighboring series coefficients}
\GSfield{Kind; source}{remark; \GSpage{P000769}}
\GSfield{The notes}{A general remark says that ratios of adjacent coefficients are rational functions of \(m,n\).}
\GSfield{Statement}{For the listed double hypergeometric series, \(c_{m+1,n}/c_{m,n}\) and \(c_{m,n+1}/c_{m,n}\) are rational in the indices wherever defined. Multiplying their denominator polynomials converts these recurrences into differential equations in Euler operators \(x\partial_x,y\partial_y\). This criterion by itself does not ensure convergence or a prescribed holonomic rank for an arbitrary formal series.}
\GSfield{Reconstruction}{Only the consequence justified by the coefficient recurrences is asserted.}
\GSfield{Fate}{For the displayed F1 and F4 coefficients the assertion follows directly from \((a)_{r+1}=(a+r)(a)_r\). It is the classical Horn-series viewpoint, not a new claim of rank determination.}
\GSfield{Present status}{As of 10 November 2026: established in the stated scope; the notebook is exposition of this result.}
\GSfield{Assessment}{literal statement: correct; publication match: partial; verification: indirect; development: annotated; confidence: high}
\end{GSitem}
\begin{GSitem}{1977-440.5}{E}{The universal arrangement and its discriminant}
\GSfield{Kind; source}{construction; \GSpage{P000769-P000770}}
\GSfield{The notes}{The matrix of \(m\) affine forms is augmented by \(f_{m+1}=1\). A complement family \(\mathcal M\) projects to coefficient space, and degeneracy is described by minors \([i_1,\ldots,i_{n+1}]\).}
\GSfield{Statement}{Homogenize the affine forms and the hyperplane at infinity in \(\mathbb P^n\). Over the open set where every relevant \((n+1)\)-row determinant is nonzero, the hyperplanes are in projective general position and the fibers are complements \(M_A=\mathbb C^n\setminus\bigcup_i\{f_i=0\}\). The discriminant is the union of these determinant-zero loci, not their common intersection.}
\GSfield{Reconstruction}{The final determinant row requires index \(i_{n+1}\), and the rank condition concerns selected submatrices. The notes conflate a parameter discriminant with its inverse image in the total space.}
\GSfield{Fate}{This is the configuration space used in \citet[Introduction and Section~4]{aomoto1977structure}.}
\GSfield{Present status}{As of 10 November 2026: established in the stated scope; the notebook is exposition of this result.}
\GSfield{Assessment}{literal statement: false as written; publication match: partial; verification: indirect; development: annotated; confidence: high}
\end{GSitem}
\begin{GSitem}{1977-440.6}{E}{The rank-one twisted de Rham complex}
\GSfield{Kind; source}{definition; \GSpage{P000770}}
\GSfield{The notes}{\(U=\prod f_i^{\lambda_i}\), \(\omega=\sum\lambda_i\,df_i/f_i\), and \(\nabla_\omega=d+\omega\wedge\). The identity \(d(U\varphi)=U\nabla_\omega\varphi\) yields \(\nabla_\omega^2=0\).}
\GSfield{Statement}{On the arrangement complement, \(\omega\) is a closed rational one-form and \(\omega\wedge\omega=0\); hence \((\Omega^\bullet(M_A),\nabla_\omega)\) is a complex. Its horizontal local system is generated by \(U^{-1}\), while the multiplier in the integral is \(U\). The residue at infinity is \(\lambda_{m+1}=-\sum_{i=1}^m\lambda_i\).}
\GSfield{Reconstruction}{The inverse local-system convention is explicit, preventing confusion in the later integration pairing.}
\GSfield{Fate}{This twisted de Rham setup predates the talk; see \citet[Section~1]{aomoto1975vanishing} and \citet[Introduction]{aomoto1977structure}.}
\GSfield{Present status}{As of 10 November 2026: established in the stated scope; the notebook is exposition of this result.}
\GSfield{Assessment}{literal statement: correct; publication match: exact; verification: read the result; development: annotated; confidence: high}
\end{GSitem}
\begin{GSitem}{1977-440.7}{E}{Top-degree concentration and the missing zero-weight exclusion}
\GSfield{Kind; source}{theorem; \GSpage{P000770}}
\GSfield{The notes}{For every subset of at most \(n\) residues, the sum is required not to lie in \(\mathbb Z^+=\{k\in\mathbb Z:k>0\}\). The notes conclude \(H^p=0\) for \(p\ne n\).}
\GSfield{Statement}{For a general-position affine arrangement and generic nonresonant rank-one monodromy, twisted cohomology is concentrated in degree \(n\), of dimension \(\binom{m-1}{n}\). A sufficient strengthened generic assumption is that all sums in the notebook, including the residues at infinity, are nonintegers. The written condition alone permits \(\lambda_i=0\) for every \(i\); then \(H^0(M_A;\mathbb C)=\mathbb C\), contradicting the conclusion for \(n>0\).}
\GSfield{Reconstruction}{The strengthened condition is sufficient and deliberately not claimed optimal. A logarithmic-complex comparison alone must not be confused with concentration in top degree.}
\GSfield{Fate}{\citet[Theorems~1--2]{aomoto1975vanishing} proves concentration under geometric and local-system hypotheses; the later logarithmic comparison theorem of \citet{esnaultschechtmanviehweg1992local} separates nonresonance from the further vanishing question. The scan's strictly positive set makes the counterexample unambiguous.}
\GSfield{Present status}{As of 10 November 2026: established in the stated scope; the notebook is exposition of this result.}
\GSfield{Assessment}{literal statement: false as written; publication match: partial; verification: read the result; development: annotated; confidence: high}
\end{GSitem}
\begin{GSitem}{1977-440.8}{E}{The finite logarithmic Aomoto complex}
\GSfield{Kind; source}{construction; \GSpage{P000770-P000771}}
\GSfield{The notes}{The algebra is generated by \(df_i/f_i\) modulo their relations, and its differential is multiplication by \(\omega\). The top cohomology is written as logarithmic \(n\)-forms modulo \(\omega\wedge\Omega^{n-1}_{\log}\).}
\GSfield{Statement}{Let \(A^\bullet\) be the finite-dimensional algebra of global logarithmic forms generated by the \(df_i/f_i\). Since these forms are closed, the restriction of \(\nabla_\omega\) is \(\omega\wedge\). Under the appropriate nonresonance conditions, its inclusion into the full twisted de Rham complex is a quasi-isomorphism; in top degree its cohomology is \(A^n/(\omega\wedge A^{n-1})\).}
\GSfield{Reconstruction}{The braces mean linear span, not a set quotient. The original arrangement relations and nonresonance conditions are part of the statement.}
\GSfield{Fate}{Aomoto's earlier framework leads to this finite model; \citet{esnaultschechtmanviehweg1992local} proves a general comparison theorem. The separately opened \citet{esnaultschechtmanviehweg1993erratum} corrects the numerical characterization of bad subspaces in Section~3 and explicitly leaves that theorem and its proof unaffected.}
\GSfield{Present status}{As of 10 November 2026: established in the stated scope; the notebook is exposition of this result.}
\GSfield{Assessment}{literal statement: correct with missing hypotheses; publication match: partial; verification: read the result; development: annotated; confidence: high}
\end{GSitem}
\begin{GSitem}{1977-440.9}{E}{Euler characteristic, bounded chambers, and solution rank}
\GSfield{Kind; source}{claim; \GSpage{P000769-P000771}}
\GSfield{The notes}{The notes identify the number of independent solutions with bounded real chambers and write \(\chi(M_A)=\binom{m-1}{n}\), also equating \(H^n\) itself with \(\chi\).}
\GSfield{Statement}{For \(m\) real affine hyperplanes in general position in \(\mathbb R^n\), the number of bounded chambers is \(\binom{m-1}{n}\). The complex complement has \(\chi(M_A)=(-1)^n\binom{m-1}{n}\). Under top-degree concentration the twisted dimension, and the associated generic period-system rank, is \((-1)^n\chi(M_A)\).}
\GSfield{Reconstruction}{Reality, general position, generic local system, the dimension sign, and the difference between a vector space and its dimension are supplied.}
\GSfield{Fate}{The chamber-cycle interpretation is part of \citet[Introduction and Section~4]{aomoto1977structure}; the sign also follows directly from inclusion-exclusion, \(\sum_{j=0}^n(-1)^j\binom mj=(-1)^n\binom{m-1}n\).}
\GSfield{Present status}{As of 10 November 2026: established in the stated scope; the notebook is exposition of this result.}
\GSfield{Assessment}{literal statement: false as written; publication match: partial; verification: indirect; development: annotated; confidence: high}
\end{GSitem}
\begin{GSitem}{1977-440.10}{E}{Twisted cycles, duality, and chamber regularization}
\GSfield{Kind; source}{construction; \GSpage{P000771}}
\GSfield{The notes}{Analytic continuation of \(U\) gives \(\pi_1(M_A)\to\mathbb C^*\). The notes pair \(H^n(M_A,\zeta)\) with \(H_n(M_A,\zeta^*)\), and associate a class to each bounded chamber.}
\GSfield{Statement}{Twisted de Rham integration pairs cohomology with homology carrying the dual local system. An open bounded chamber has boundary on the deleted arrangement; a chosen branch and nonresonant regularization, or a locally finite homology interpretation, are needed to turn it into an admissible twisted cycle. Under genericity these chamber cycles give a basis.}
\GSfield{Reconstruction}{The final star is indistinct in the scan. Its position depends on whether \(\zeta\) denotes horizontal sections \(U^{-1}\) or branches of \(U\); the pairing must always be dual, so no unsupported paleographic repair is made.}
\GSfield{Fate}{\citet[Section~4 and Appendix]{aomoto1977structure} develops visible cycles and their connection formulas; \citet[final section]{aomoto1975vanishing} discusses the corresponding period pairing.}
\GSfield{Present status}{As of 10 November 2026: established in the stated scope; the notebook is exposition of this result.}
\GSfield{Assessment}{literal statement: correct with missing hypotheses; publication match: partial; verification: read the result; development: annotated; confidence: medium}
\end{GSitem}

\GStalk{441}{1978-03-20}{V. M. Buchstaber, with A. V. Shokurov}{P000772-P000773}{Landweber-Novikov operations and formal substitutions}
The lecture identifies the Landweber--Novikov algebra with invariant differential operators on the formal substitution group, translates stably complex bordism into genus polynomials, and states a Hattori--Stong integrality criterion using the Todd translate of integral substitutions. It ends with the genuine question of computing cohomology of the integral substitution group. That question requires separating finite-jet, continuous, and abstract group cohomology; later real finite-jet results and the chapter's bounded integral calculation do not determine all of these theories.

\begin{GSitem}{1977-441.1}{A}{Formal substitutions, jets, and the positive Witt algebra}
\GSfield{Kind; source}{definition; \GSpage{P000772}}
\GSfield{The notes}{\(\operatorname{Diff}_1\) consists of \(\varphi(t)=t+a_1t^2+a_2t^3+\cdots\), with \(\varphi(0)=0\), \(\varphi'(0)=1\), and is the inverse limit of its jet groups. Its Lie algebra is denoted \(L_1\), and polynomial functions in the coefficients form \(P\).}
\GSfield{Statement}{Over a characteristic-zero field, composition gives a pro-unipotent group of tangent-to-identity formal substitutions. Its complete Lie algebra is \(t^2\mathbb R[[t]]\partial_t\); the graded polynomial subalgebra has basis \(e_i=t^{i+1}\partial_t\), \(i\ge1\), and \([e_i,e_j]=(j-i)e_{i+j}\). Polynomial coordinate functions depend on finitely many jets.}
\GSfield{Reconstruction}{The distinction between the complete formal Lie algebra and its graded polynomial part is explicit; cochains later require the corresponding continuous or graded convention.}
\GSfield{Fate}{The group and coefficient Hopf algebra are the framework of \citet[Section~1]{buchstabershokurov1978landweber}, submitted 5 May 1977 and published 1978.}
\GSfield{Present status}{As of 10 November 2026: the stated construction is established; the qualifications above are essential.}
\GSfield{Assessment}{literal statement: correct; publication match: exact; verification: read the result; development: annotated; confidence: high}
\end{GSitem}
\begin{GSitem}{1977-441.2}{A}{A stably complex manifold as a genus polynomial}
\GSfield{Kind; source}{construction; \GSpage{P000772-P000773}}
\GSfield{The notes}{Put \(\psi(t)=t/\varphi(t)\); take the degree-\(n\) coefficient of \(\prod_j\psi(\beta_jt)\), express it in elementary symmetric functions, substitute Chern classes, and pair with \(M\).}
\GSfield{Statement}{For a closed stably complex manifold of real dimension \(2n\), evaluate the multiplicative characteristic sequence with power series \(t/\varphi(t)\) on its stable tangent Chern roots. The degree-\(2n\) component paired with the fundamental class is a polynomial in \(a_1,\ldots,a_n\). Stabilization adds zero roots and leaves the value unchanged.}
\GSfield{Reconstruction}{Closedness, orientation, tangent-bundle convention, and cohomological degree are supplied.}
\GSfield{Fate}{The genus-polynomial realization underlies \citet[Sections~1--2]{buchstabershokurov1978landweber}; its manuscript predates this lecture.}
\GSfield{Present status}{As of 10 November 2026: the stated construction is established; the qualifications above are essential.}
\GSfield{Assessment}{literal statement: correct with missing hypotheses; publication match: exact; verification: read the result; development: annotated; confidence: high}
\end{GSitem}
\begin{GSitem}{1977-441.3}{A}{Geometric characteristic operations and their scope}
\GSfield{Kind; source}{construction; \GSpage{P000773}}
\GSfield{The notes}{Products of manifolds and cycles dual to Chern polynomials illustrate operations. A first Chern class is represented by zeros of a section; Landweber--Novikov operations are then called all natural cohomological operations.}
\GSfield{Statement}{A transverse section of a complex line bundle gives a codimension-two submanifold dual to its first Chern class. General characteristic operations in complex bordism require the bordism construction and stable normal data; an arbitrary Chern polynomial does not canonically produce a single embedded manifold. The Landweber--Novikov algebra is the operation subalgebra \(S\); the full stable-operation algebra also involves coefficient operations.}
\GSfield{Reconstruction}{The distinction between geometric examples, S, and the entire operation algebra is essential. The notebook word canonical does not determine whether its bundle is the canonical line bundle or a tautological characteristic bundle.}
\GSfield{Fate}{\citet[Theorem~1 and surrounding definitions]{buchstabershokurov1978landweber} identifies the precise algebra \(S\). It does not justify an unrestricted assertion covering every natural operation in every cohomology theory.}
\GSfield{Present status}{As of 10 November 2026: the stated construction is established; the qualifications above are essential.}
\GSfield{Assessment}{literal statement: false as written; publication match: partial; verification: read the result; development: annotated; confidence: high}
\end{GSitem}
\begin{GSitem}{1977-441.4}{A}{Operations as left-invariant differential operators}
\GSfield{Kind; source}{theorem; \GSpage{P000773}}
\GSfield{The notes}{Theorem~1 associates \(c_\omega:M\mapsto M_\omega\) with a left-invariant differential operator \(D_\omega\) on genus polynomials, asserting a bijective correspondence.}
\GSfield{Statement}{The integral Landweber--Novikov algebra is identified with the appropriately integral left-invariant differential-operator algebra on the formal substitution group. After tensoring with \(\mathbb Q\), it is the universal enveloping algebra of the positive graded vector-field Lie algebra. The integral structure is not the naive enveloping algebra over \(\mathbb Z\); divided-power phenomena matter.}
\GSfield{Reconstruction}{The rational versus integral qualification prevents replacing the operation lattice by an incorrect integral enveloping algebra.}
\GSfield{Fate}{Exact match to \citet[Theorem~1 and its corollary]{buchstabershokurov1978landweber}, received May 1977, published 1978.}
\GSfield{Present status}{As of 10 November 2026: the stated construction is established; the qualifications above are essential.}
\GSfield{Assessment}{literal statement: correct with missing hypotheses; publication match: exact; verification: read the result; development: annotated; confidence: high}
\end{GSitem}
\begin{GSitem}{1977-441.5}{A}{The Todd translate and bordism integrality}
\GSfield{Kind; source}{theorem; \GSpage{P000773}}
\GSfield{The notes}{A polynomial comes from a manifold precisely when it takes integral values on the right coset of \(\varphi_T(t)=1-e^{-t}\) by \(\operatorname{Diff}_1(\mathbb Z)\); the theorem is compared with Stong--Hattori.}
\GSfield{Statement}{Within the graded rational genus-polynomial realization of complex bordism used in the paper, the bordism lattice is characterized by integrality on the Todd translate \(\varphi_T\operatorname{Diff}_1(\mathbb Z)\), with composition in the stated order. This is a statement about the identified characteristic-number space, not about all arbitrary functions on the substitution group.}
\GSfield{Reconstruction}{The graded polynomial domain and the right-coset composition convention are part of the theorem.}
\GSfield{Fate}{\citet[Theorem~2]{buchstabershokurov1978landweber} gives precisely the arithmetic characterization and relates it to the Hattori--Stong theorem. Submission preceded the lecture; publication was in 1978.}
\GSfield{Present status}{As of 10 November 2026: the stated construction is established; the qualifications above are essential.}
\GSfield{Assessment}{literal statement: correct with missing hypotheses; publication match: exact; verification: read the result; development: annotated; confidence: high}
\end{GSitem}
\begin{GSitem}{1977-441.6}{A}{Rational Ext and vector-field cohomology}
\GSfield{Kind; source}{theorem; \GSpage{P000773}}
\GSfield{The notes}{\(\operatorname{Ext}_S^*(\mathbb Q,\mathbb Q)=H^*(L_1)\). The notes also say that all homotopy invariants of continuous maps are expressed through this cohomology.}
\GSfield{Statement}{Under the rational enveloping-algebra identification, \(\operatorname{Ext}_{S\otimes\mathbb Q}^*(\mathbb Q,\mathbb Q)\cong H^*_{\rm CE}(L_1;\mathbb Q)\), with the internal grading retained. Complex-bordism operations enter Adams--Novikov spectral sequences, but their rational Ext with trivial coefficients is not by itself a classification of all continuous maps or all stable homotopy groups.}
\GSfield{Reconstruction}{Rationalization, augmentation coefficients, and the restricted role of the spectral sequence are made explicit.}
\GSfield{Fate}{The Ext comparison follows from \citet[Theorem~1 and cohomological consequences]{buchstabershokurov1978landweber}. The broad sentence about all invariants is motivational shorthand and has no exact theorem identified in this form.}
\GSfield{Present status}{As of 10 November 2026: the stated construction is established; the qualifications above are essential.}
\GSfield{Assessment}{literal statement: correct with missing hypotheses; publication match: partial; verification: read the result; development: annotated; confidence: high}
\end{GSitem}
\begin{GSitem}{1977-441.7}{C}{Cohomology of formal substitutions with integer coefficients}
\GSfield{Kind; source}{problem; \GSpage{P000773}}
\GSfield{The notes}{\ru{Задача: вычислить когомологии группы} \(\operatorname{Diff}_1(\mathbb Z)\). The immediately preceding definition specifies integral coefficients of a tangent-to-identity formal series. No coefficient module or cochain topology is stated.}
\GSfield{Statement}{Let \(J=\{t+a_1t^2+a_2t^3+\cdots:a_i\in\mathbb Z\}\) under composition. A precise finite-jet version is \(H^*_{\rm jet}(J;A)=\varinjlim_NH^*(J/(t^{N+1});A)\) with a trivial discrete abelian coefficient group and inflation maps. This is distinct from unrestricted cohomology of the abstract group \(J\). The full integral finite-jet cohomology ring is a natural precise refinement; the original wording does not choose between these theories.}
\GSfield{Reconstruction}{The integral finite-jet formulation is an explicit refinement, not a claim about the unstated original coefficients. The chapter proves bounded consequences of known abelianization, without asserting a full solution or novelty.}
\GSfield{Fate}{\citet[Theorem~1.5]{babenkobogatyi2008substitutions} computes real cohomology with uniformly locally constant supports as the cohomology of \(L_1\), resolving that natural real version. Their Theorem~1.2 computes the continuous abelianization \(\mathbb Z^2\oplus(\mathbb Z/2)^2\). \citet{pavezcornejo2025abelianization}, first circulated in 2024, also uses the closure of the commutator subgroup; neither statement identifies the abstract discrete abelianization.}
\GSfield{Present status}{As of 10 November 2026: the real finite-jet version is known. The chapter proves \(\operatorname{tors}H^2_{\rm jet}(J;\mathbb Z)\cong(\mathbb Z/2)^2\) and that the cup product of the two integral degree-one generators vanishes, by an explicit four-jet coboundary. The full integral ring and unrestricted abstract group cohomology are not determined by these results; no unsupported assertion that either is wholly unstudied is made.}
\GSfield{Assessment}{literal statement: correct with missing hypotheses; publication match: partial; verification: read the result; development: developed in this chapter; confidence: high; problem \hyperref[prob:GP-1977-13]{GP-1977-13}}
\GSfield{Worth developing}{Worth developing in a specified cochain category: the real theory is known, while integral torsion and multiplication require additional information that real coefficients discard.}
\end{GSitem}

\GStalk{442}{1978-03-20}{F. Calogero and A. Degasperis}{P000774-P000776}{Spectral evolution, conserved quantities, and an adjacent talk on caustics}
The main lecture studies KdV recursion operators and their nonisospectral extensions, linear evolution of scattering data, and trace-type conserved quantities. Several normalization claims fail elementary low-order checks. The bottom of P000776 explicitly starts a different lecture: V. P. Palamodov at the Mathematical Society, 21 March 1978, on lacunas, sharp fronts, and singularity deformations. That adjacent material is retained with its actual speaker. P000777 starts the next Calogero lecture.

\begin{GSitem}{1977-442.1}{A}{The recursion operator and the nonisospectral hierarchy}
\GSfield{Kind; source}{construction; \GSpage{P000774}}
\GSfield{The notes}{\(Lf=f_{xx}-4uf+2u_x\int_x^\infty f\), and \(u_t=\alpha(L,t)u_x+\beta(L,t)(xu_x+2u)\). For \(\alpha=-L,\beta=0\) this is KdV; powers of \(L\) give the higher hierarchy.}
\GSfield{Statement}{For smooth rapidly decaying \(u\), with \(\partial^{-1}f=-\int_x^\infty f\), \(L=\partial^2-4u-2u_x\partial^{-1}\). Polynomial \(\alpha,\beta\) define the indicated hierarchy on a domain stable under these operations; more general functions require functional calculus. Since \(Lu_x=u_{xxx}-6uu_x\), the stated sign gives \(u_t=6uu_x-u_{xxx}\).}
\GSfield{Reconstruction}{Decay, the inverse-derivative convention, and polynomial rather than arbitrary undefined functional calculus are supplied.}
\GSfield{Fate}{The nonisospectral extension was submitted 9 January 1978 and published in May: \citet{calogerodegasperis1978extension}. Its publisher record points to a June 1977 Rome lecture precursor. The exact operator-level identity here is also checked directly from the supplied formula.}
\GSfield{Present status}{As of 10 November 2026: the stated construction is established; the qualifications above are essential.}
\GSfield{Assessment}{literal statement: correct with missing hypotheses; publication match: partial; verification: publisher metadata checked; development: annotated; confidence: high}
\end{GSitem}
\begin{GSitem}{1977-442.2}{E}{The Fourier analogy and the nonlinear spectral transform}
\GSfield{Kind; source}{method; \GSpage{P000774}}
\GSfield{The notes}{A constant-coefficient linear equation becomes an ordinary differential equation after Fourier transform; an extra term linear in \(x\) becomes differentiation in the spectral variable. The same pattern motivates the scattering transform.}
\GSfield{Statement}{Fourier transformation turns \(\partial_x\) into multiplication by \(ik\) and \(x\) into \(i\partial_k\), on a suitable function space. For the one-dimensional linear Schr\"odinger spectral problem, the scattering map depends nonlinearly on the potential; it is not the ordinary Fourier transform of \(u\); its admissible inverse data and reconstruction equation must be retained.}
\GSfield{Fate}{The analogy and Wronskian method were already used by \citet[abstract]{calogerodegasperis1976inverse}. The page's inserted ``linear'' is essential to the scope of the Fourier sentence.}
\GSfield{Present status}{As of 10 November 2026: established in the stated scope; the notebook is exposition of this result.}
\GSfield{Assessment}{literal statement: correct with missing hypotheses; publication match: partial; verification: publisher metadata checked; development: annotated; confidence: high}
\end{GSitem}
\begin{GSitem}{1977-442.3}{A}{Scattering data and their linear transport}
\GSfield{Kind; source}{construction; \GSpage{P000774-P000775}}
\GSfield{The notes}{\(u\leftrightarrow\{R(k,t);p_n,c_n^2\}\), with inverse reconstruction by Gelfand--Levitan. \(R_t=2ik\alpha(-4k^2,t)R-k\beta(-4k^2,t)R_k\); \(p_t=p\beta(-4p^2,t)\); the norming-constant equation is omitted.}
\GSfield{Statement}{For a real sufficiently decaying potential, the complete one-dimensional scattering data consist of a reflection coefficient, negative eigenvalues, and norming constants. The displayed first-order transport equation is the proposed continuous-data evolution. If a bound state is parametrized by \(k=i\kappa\), its argument \(-4k^2\) is \(+4\kappa^2\); one must know whether the notebook's \(p\) is a complex spectral point or a positive decay constant before applying its last formula.}
\GSfield{Reconstruction}{The literal p/rho ambiguity is preserved. The statement about changing soliton amplitudes and velocities is appropriate for nonisospectral motion, subject to the chosen real or complex flow.}
\GSfield{Fate}{\citet{calogerodegasperis1978extension} is the closely matching nonisospectral publication; publisher metadata and editorial footnotes, not the full normalization proof, were obtained. The omitted norming equation prevents a complete algorithm from this page alone.}
\GSfield{Present status}{As of 10 November 2026: the stated construction is established; the qualifications above are essential.}
\GSfield{Assessment}{literal statement: correct with missing hypotheses; publication match: partial; verification: publisher metadata checked; development: annotated; confidence: medium}
\end{GSitem}
\begin{GSitem}{1977-442.4}{A}{Conserved moments from the scaling seed}
\GSfield{Kind; source}{theorem; \GSpage{P000774-P000775}}
\GSfield{The notes}{\(c^{(n)}=\int_{\mathbb R}L^n(xu_x+2u)\,dx\), and \(\beta=0\) implies \(\dot c^{(n)}=0\).}
\GSfield{Statement}{For an isospectral KdV-hierarchy flow with sufficiently smooth rapid decay, these recursion expressions represent trace invariants, up to the chosen normalization. The first two are \(c^{(0)}=\int u\,dx\) and \(c^{(1)}=-3\int u^2\,dx\), obtained by integration by parts. Boundary terms must vanish.}
\GSfield{Reconstruction}{The two low-order identities are direct checks of the notebook normalization, not a claim to have read its complete hierarchy in the later paper.}
\GSfield{Fate}{\citet{calogerodegasperis1978conservation}, published June 1978, explicitly studies conserved quantities for these old and new spectral-transform classes. The abstract supports the identification; a complete formula-by-formula publication match is not claimed.}
\GSfield{Present status}{As of 10 November 2026: the stated construction is established; the qualifications above are essential.}
\GSfield{Assessment}{literal statement: correct with missing hypotheses; publication match: partial; verification: publisher metadata checked; development: annotated; confidence: high}
\end{GSitem}
\begin{GSitem}{1977-442.5}{E}{Hamiltonian operators and the order of recursion factors}
\GSfield{Kind; source}{construction; \GSpage{P000775}}
\GSfield{The notes}{\(I_0=\int u^2dx\), \(u_t=A_n\delta I_0/\delta u\), \(A_n=L^n\partial\). Novikov then writes \(L=\partial^{-1}L_3\), with inverse differentiation on the left.}
\GSfield{Statement}{For the convention on P000774, \(J_1=\partial\), \(J_2=\partial^3-2(u\partial+\partial u)\), and the recursion on evolutionary vector fields is \(L=J_2J_1^{-1}\). The operator \(J_1^{-1}J_2\) instead acts on variational gradients. These factor orders cannot be identified on the same space. Also \(\delta I_0/\delta u=2u\), introducing a factor two in the time normalization.}
\GSfield{Reconstruction}{The vector-field/gradient distinction explains the two legitimate orders. Higher nonlocal Poisson brackets require a specified domain and inverse derivative; the notes suppress these.}
\GSfield{Fate}{The Hamiltonian spectral framework is pre-existing: \citet{zakharovfaddeev1971kdv}. The earlier occurrence index includes the recursive-hierarchy problem GP-1973-53; this lecture uses the recursion and is not a new occurrence of an unresolved classification problem.}
\GSfield{Present status}{As of 10 November 2026: established in the stated scope; the notebook is exposition of this result.}
\GSfield{Assessment}{literal statement: false as written; publication match: partial; verification: publisher metadata checked; development: annotated; confidence: high}
\end{GSitem}
\begin{GSitem}{1977-442.6}{A}{The transmission phase and its inconsistent leading coefficient}
\GSfield{Kind; source}{claim; \GSpage{P000775}}
\GSfield{The notes}{The physical scattering convention is \(\psi\sim T e^{-ikx}\) at negative infinity and \(\psi\sim e^{-ikx}+R e^{ikx}\) at positive infinity. With \(T=|T|e^{2i\varphi}\), the notes claim \(\varphi(k)\sim\sum k^{-2n-1}\varphi^{(n)}\) and \(\varphi^{(n)}=(-1)^{n+1}c^{(n)}/[2(2n+1)]\).}
\GSfield{Statement}{For \(-\psi''+u\psi=k^2\psi\) with real \(u\in C(\mathbb R)\cap L^1(\mathbb R)\) and the written physical incoming-wave convention, the left-Jost Volterra estimate gives \(\log T(k)=-\frac{i}{2k}\int u+O(k^{-2})\). Consequently \(\varphi^{(0)}=-\tfrac14\int u=-c^{(0)}/4\), whereas the notebook gives \(-c^{(0)}/2\). Higher trace identities exist, but this displayed normalization is already false at \(n=0\).}
\GSfield{Reconstruction}{Only the decisive leading coefficient is repaired; a guessed all-order numerical factor is not substituted for the source.}
\GSfield{Fate}{The trace-invariant subject matches \citet{calogerodegasperis1978conservation}; metadata does not verify its agreement with equation (8). The leading-coefficient test is an independent mathematical emendation of the notebook.}
\GSfield{Present status}{As of 10 November 2026: the stated construction is established; the qualifications above are essential.}
\GSfield{Assessment}{literal statement: false as written; publication match: partial; verification: publisher metadata checked; development: annotated; confidence: high}
\end{GSitem}
\begin{GSitem}{1977-442.7}{A}{Moment evolution, triangularity, and periodicity}
\GSfield{Kind; source}{claim; \GSpage{P000776}}
\GSfield{The notes}{For \(\beta(z,t)=\sum\beta_m(t)z^m\), equation (11) says \(\dot C^{(n)}=n\sum_m\beta_m c^{(n+m)}\). Polynomial \(\beta\) makes this triangular; \(\beta_0=i\omega\) allegedly makes all quantities periodic.}
\GSfield{Statement}{For the pure constant scaling flow \(u_t=\beta_0(xu_x+2u)\), \(u(x,t)=e^{2\beta_0t}u_0(e^{\beta_0t}x)\), and \(\dot c^{(n)}=(2n+1)\beta_0c^{(n)}\). Thus the recorded factor \(n\) is inconsistent even for \(n=0\). Polynomial \(\beta\) gives finite coupling in each row of an infinite upper system, not a closed finite system. Complex periodic scaling does not automatically imply a globally regular periodic real-space solution.}
\GSfield{Reconstruction}{The capital/lowercase C/c ambiguity is retained, but no alternative moment normalization is defined that would justify the written equation.}
\GSfield{Fate}{The nonisospectral conservation-law problem belongs to \citet{calogerodegasperis1978conservation}. The all-order nonconstant-beta equation was not verified from full text; the pure-scaling counterexample is conclusive and does not need such verification.}
\GSfield{Present status}{As of 10 November 2026: the stated construction is established; the qualifications above are essential.}
\GSfield{Assessment}{literal statement: false as written; publication match: partial; verification: publisher metadata checked; development: annotated; confidence: high}
\end{GSitem}
\begin{GSitem}{1977-442.8}{E}{Adjacent Palamodov lecture: caustics and wave fronts}
\GSfield{Kind; source}{definition; \GSpage{P000776}; V. P. Palamodov (Mathematical Society, 21 March 1978)}
\GSfield{The notes}{Under a fresh heading dated 21 March 1978, Palamodov illustrates parabolic focusing and the envelope of rays reflected from a spherical mirror, calling the latter a caustic. A surface has horizontal sections equal to successive wave fronts.}
\GSfield{Statement}{A caustic is the locus of critical values of the projection of a ray congruence; in the planar drawing it is the envelope to which reflected rays are tangent. A time-dependent front may be encoded as a surface in space-time, with its time slices giving front positions. These definitions do not by themselves establish a universal singularity type.}
\GSfield{Fate}{This is expository geometric material in an adjacent Mathematical Society lecture. The supplied fragment contains no lacuna criterion or classification theorem; none is attributed to Calogero.}
\GSfield{Present status}{As of 10 November 2026: established in the stated scope; the notebook is exposition of this result.}
\GSfield{Assessment}{literal statement: correct with missing hypotheses; publication match: partial; verification: indirect; development: annotated; confidence: high}
\end{GSitem}
\begin{GSitem}{1977-442.9}{F}{The illustrated four-thirds front singularity}
\GSfield{Kind; source}{claim; \GSpage{P000776}; V. P. Palamodov}
\GSfield{The notes}{At the instant of caustic formation the notes underline \(y=x^{4/3}\).}
\GSfield{Statement}{The plane branch \((x,y)=(s^3,s^4)\) has that real graph and can occur in a special front degeneration. It differs from an ordinary semicubical cusp \((s^2,s^3)\). Which front family, coordinate equivalence, and genericity assumptions Palamodov intended are not recorded.}
\GSfield{Fate}{No exact theorem identifying this drawing and exponent was located. It is preserved as a specific adjacent assertion rather than generalized into a false claim that every caustic has this exponent.}
\GSfield{Present status}{As of 10 November 2026: the algebraic curve is identifiable, but the unrecorded geometric hypotheses prevent an exact historical theorem match.}
\GSfield{Assessment}{literal statement: not determinable; publication match: none found; verification: unresolved lead; development: annotated; confidence: low}
\end{GSitem}
\begin{GSitem}{1977-442.10}{E}{The forward Cauchy fundamental solution}
\GSfield{Kind; source}{definition; \GSpage{P000776}; V. P. Palamodov}
\GSfield{The notes}{For a hyperbolic operator \(\mathcal L\) of order \(m\), \(\mathcal LE=0\) for \(t>0\), and \(\partial_t^jE|_{t=0}=\delta(x')\delta_{j,m-1}\), \(0\le j<m\). The wave example uses \(\partial_t^2-c^{-2}\Delta\).}
\GSfield{Statement}{For an operator monic in \(\partial_t^m\), with a noncharacteristic initial hypersurface and hypotheses ensuring the distributional Cauchy problem is well posed, these data define the causal propagation kernel for the last initial datum. Extending it by zero to \(t<0\) gives the corresponding fundamental solution with a point source, subject to the operator normalization.}
\GSfield{Reconstruction}{Monicity and Cauchy well-posedness are necessary. The reciprocal c is not corrected: the text does not define c as velocity rather than slowness.}
\GSfield{Fate}{The page supplies the definition, not a new existence proof. The lacuna theory named in the heading is not actually developed in the surviving material.}
\GSfield{Present status}{As of 10 November 2026: established in the stated scope; the notebook is exposition of this result.}
\GSfield{Assessment}{literal statement: correct with missing hypotheses; publication match: partial; verification: indirect; development: annotated; confidence: high}
\end{GSitem}

\GStalk{443}{1978-03-27}{F. Calogero, jointly with A. Degasperis}{P000777-P000781}{Backlund transformations and polynomial-root dynamics}
Wronskian identities connect nonlinear transformations of potentials with multiplicative changes in scattering data. The lecture obtains the KdV recursion by an infinitesimal limit, an elementary soliton transformation, a Riccati equation and nonlinear superposition. A second topic converts a linear PDE for a polynomial into a solvable interacting-root system, with Hermite-zero equilibria. The final linearized equation is missing an identity matrix; its stated matrix spectrum itself is correct.

\begin{GSitem}{1977-443.1}{E}{Complete scattering data cannot generally be replaced by reflection alone}
\GSfield{Kind; source}{claim; \GSpage{P000777}}
\GSfield{The notes}{After initially including discrete spectrum data, the notes say that eigenvalues are poles of \(R(k)\), with residues \(ic_n^2\), and that \(q\) is determined by \(R\).}
\GSfield{Statement}{A real decaying potential is determined by admissible complete scattering data: reflection on the real axis together with bound-state eigenvalues and norming constants. Reflection alone is insufficient in this class: the zero potential and nonzero reflectionless solitons all have \(R=0\). Special analyticity or support restrictions can encode more data, but no such class is specified here.}
\GSfield{Reconstruction}{The undeciphered Schrodinger sign is fixed to \(q-k^2\) by the plane-wave asymptotics, an inference rather than a confident reading of the overwritten mark.}
\GSfield{Fate}{The standard spectral-transform setup is used in \citet{calogerodegasperis1976inverse}. The notebook's later assertion that a soliton is added while multiplying R itself provides an internal warning against unconditional uniqueness from R.}
\GSfield{Present status}{As of 10 November 2026: established in the stated scope; the notebook is exposition of this result.}
\GSfield{Assessment}{literal statement: false as written; publication match: partial; verification: publisher metadata checked; development: annotated; confidence: high}
\end{GSitem}
\begin{GSitem}{1977-443.2}{E}{Two-potential Wronskian identities}
\GSfield{Kind; source}{method; \GSpage{P000777-P000778}}
\GSfield{The notes}{For two potentials \(q,\widetilde q\), primes denote the second potential, not differentiation. Equations (5)--(6) express \(2ik f(-4k^2)(\widetilde R-R)\) and \((2ik)^2g(-4k^2)(\widetilde R+R)\) as integrals against \(\widetilde\Psi\Psi\).}
\GSfield{Statement}{Define \(Ih(x)=\int_x^\infty h(s)ds\), \(\Gamma h=(\widetilde q_x+q_x)h+(\widetilde q-q)I((\widetilde q-q)h)\), and \(\Lambda h=h_{xx}-2(\widetilde q+q)h+\Gamma Ih\). With sufficient decay and polynomial f,g, the generalized Wronskian identities pair these operators with products of scattering solutions. Boundary asymptotics supply the reflection differences and sums.}
\GSfield{Reconstruction}{The second Psi(x,t) in (6) must be Psi(x,k) in this stationary identity; the operator notation is made unambiguous.}
\GSfield{Fate}{The generalized Wronskian/B\"acklund method is already the subject of \citet[abstract]{calogerodegasperis1976inverse}. Full-text coefficient verification of this two-potential normalization was not obtained.}
\GSfield{Present status}{As of 10 November 2026: established in the stated scope; the notebook is exposition of this result.}
\GSfield{Assessment}{literal statement: correct with missing hypotheses; publication match: partial; verification: publisher metadata checked; development: annotated; confidence: medium}
\end{GSitem}
\begin{GSitem}{1977-443.3}{E}{The adjoint spectral identity for the recursion kernel}
\GSfield{Kind; source}{claim; \GSpage{P000778-P000779}}
\GSfield{The notes}{\(\int\widetilde\Psi\Psi\,\Lambda F=-4k^2\int\widetilde\Psi\Psi F\); equivalently the pairing with \((\Lambda+4k^2)F\) vanishes.}
\GSfield{Statement}{For test functions for which integration by parts and the nonlocal integrals have no boundary terms, the product of the two Schr\"odinger solutions is a generalized eigenfunction of the formal adjoint of \(\Lambda\), with eigenvalue \(-4k^2\). This is a pairing identity, not an assertion that \(\Lambda\) is a self-adjoint operator on a Hilbert space.}
\GSfield{Reconstruction}{Test-function domains and the adjoint distinction are supplied.}
\GSfield{Fate}{It is the mechanism behind the spectral substitution in \citet{calogerodegasperis1976inverse}; the identification is partial because the full exact scalar kernel was not read in that publication.}
\GSfield{Present status}{As of 10 November 2026: established in the stated scope; the notebook is exposition of this result.}
\GSfield{Assessment}{literal statement: correct with missing hypotheses; publication match: partial; verification: publisher metadata checked; development: annotated; confidence: high}
\end{GSitem}
\begin{GSitem}{1977-443.4}{E}{A nonlinear Backlund relation induces a rational multiplier}
\GSfield{Kind; source}{construction; \GSpage{P000778}}
\GSfield{The notes}{If \(f(\Lambda)(\widetilde q-q)+g(\Lambda)\Gamma1=0\), then \(\widetilde R=\frac{f(-4k^2)-2ikg(-4k^2)}{f(-4k^2)+2ikg(-4k^2)}R\).}
\GSfield{Statement}{Subject to the preceding Wronskian identities and a nonzero denominator, the displayed multiplier follows by adding their two left sides and solving the resulting scalar linear equation. For a full transformation of potentials, the bound-state and norming data, invertibility, and regularity must also be specified. Arbitrary functions f,g need a justified operator calculus.}
\GSfield{Reconstruction}{The spectral multiplier is a consequence under hypotheses, not a complete existence theorem for the nonlinear relation.}
\GSfield{Fate}{General B\"acklund transformations and their spectral meaning predate the seminar in \citet[abstract]{calogerodegasperis1976inverse}. The elementary algebraic implication from the recorded identities is exact, even though the paper-level normalization match is partial.}
\GSfield{Present status}{As of 10 November 2026: established in the stated scope; the notebook is exposition of this result.}
\GSfield{Assessment}{literal statement: correct with missing hypotheses; publication match: partial; verification: publisher metadata checked; development: annotated; confidence: high}
\end{GSitem}
\begin{GSitem}{1977-443.5}{E}{The infinitesimal limit and a sign-inconsistent alternative formula}
\GSfield{Kind; source}{construction; \GSpage{P000778-P000779}}
\GSfield{The notes}{Let \(\widetilde q=q(x,t+dt)\); then \(\Lambda\to L\), \(\Gamma\to 2q_x\), yielding \(q_t=\alpha(L)q_x\), \(R_t=2ik\alpha(-4k^2)R\). Gelfand writes \(L=(\partial^3+2q\partial+2\partial q)\partial^{-1}\).}
\GSfield{Statement}{At equal potentials, \(L=\partial^2-4q-2q_x\partial^{-1}\), for \(\partial^{-1}h=-Ih\). Thus \(L=(\partial^3-2q\partial-2\partial q)\partial^{-1}\). The written plus signs require replacing the potential by its negative; without that convention change they conflict with (12). Polynomial \(\alpha\) gives the isospectral KdV hierarchy.}
\GSfield{Reconstruction}{The sign correction is conditional on keeping the potential convention used throughout this lecture; no source mark is reinterpreted as a minus.}
\GSfield{Fate}{The infinitesimal spectral method is pre-existing, \citet{calogerodegasperis1976inverse}. The incompatible signs are confirmed from the manuscript and direct operator multiplication.}
\GSfield{Present status}{As of 10 November 2026: established in the stated scope; the notebook is exposition of this result.}
\GSfield{Assessment}{literal statement: false as written; publication match: partial; verification: publisher metadata checked; development: annotated; confidence: high}
\end{GSitem}
\begin{GSitem}{1977-443.6}{E}{Products of second-order solutions satisfy a third-order equation}
\GSfield{Kind; source}{claim; \GSpage{P000779}; I. M. Gelfand}
\GSfield{The notes}{Gelfand remarks that the product of two solutions of a second-order equation satisfies a third-order equation, explaining the recurrence of the KdV operator.}
\GSfield{Statement}{If \(y_1''=(q-k^2)y_1\) and \(y_2''=(q-k^2)y_2\), then \(z=y_1y_2\) satisfies \(z^{(3)}-4(q-k^2)z'-2q'z=0\). This follows by three differentiations and substitution. It is the symmetric-square equation; both factors must solve the same scalar equation.}
\GSfield{Reconstruction}{The explicit differential equation and the same-equation hypothesis are supplied.}
\GSfield{Fate}{The identity is classical and is checked directly here as the algebraic source of the equal-potential adjoint relation. The notebook's list of surnames does not determine a unique additional publication.}
\GSfield{Present status}{As of 10 November 2026: established in the stated scope; the notebook is exposition of this result.}
\GSfield{Assessment}{literal statement: correct with missing hypotheses; publication match: partial; verification: indirect; development: annotated; confidence: high}
\end{GSitem}
\begin{GSitem}{1977-443.7}{E}{Elementary soliton insertion and Riccati form}
\GSfield{Kind; source}{construction; \GSpage{P000779}}
\GSfield{The notes}{For \(g=1,f=2\rho\), \(\widetilde R=-(k+i\rho)(k-i\rho)^{-1}R\). With \(Q=Iq\), \(\widetilde Q_x+Q_x=-\tfrac12(\widetilde Q-Q)(4\rho+\widetilde Q-Q)\); \(Q=0\) gives one soliton.}
\GSfield{Statement}{The potential relation is a first-order Riccati equation. For \(\rho>0\), the regular branch from the zero seed has \(\widetilde Q=2\rho(\tanh(\rho(x-x_0))-1)\), hence \(\widetilde q=-2\rho^2\operatorname{sech}^2(\rho(x-x_0))\). Its position parameter is norming data absent from the multiplier of R. Time dependence must satisfy the chosen hierarchy equation.}
\GSfield{Reconstruction}{The regular branch and norming parameter are essential. Multiplying R=0 cannot by itself distinguish the zero seed from the resulting soliton.}
\GSfield{Fate}{This is the elementary Darboux/B\"acklund mechanism in the earlier spectral-transform framework, \citet{calogerodegasperis1976inverse}. The seed solution is directly verified in the notebook's normalization.}
\GSfield{Present status}{As of 10 November 2026: established in the stated scope; the notebook is exposition of this result.}
\GSfield{Assessment}{literal statement: correct with missing hypotheses; publication match: partial; verification: publisher metadata checked; development: annotated; confidence: high}
\end{GSitem}
\begin{GSitem}{1977-443.8}{F}{Generating all integrals by Backlund transformations}
\GSfield{Kind; source}{claim; \GSpage{P000779}}
\GSfield{The notes}{The notes say that all integrals are obtained by successive B\"acklund transformation from \(\int u^2dx\).}
\GSfield{Statement}{B\"acklund transformations can generate relations or generating functions for conserved quantities in specified integrable hierarchies. A literal finite iteration of a transformation of potentials is not, without an additional expansion or functional action, a definition that produces every conservation law from one functional.}
\GSfield{Fate}{\citet{calogerodegasperis1978conservation} is relevant to conserved quantities, but the notebook supplies no parameter expansion, action on functionals, or completeness theorem matching this sentence exactly.}
\GSfield{Present status}{As of 10 November 2026: the intended generating prescription is not identifiable precisely enough to classify as an open problem or a proved completeness result.}
\GSfield{Assessment}{literal statement: not determinable; publication match: thematic only; verification: publisher metadata checked; development: annotated; confidence: low}
\end{GSitem}
\begin{GSitem}{1977-443.9}{E}{Permutability and nonlinear superposition}
\GSfield{Kind; source}{theorem; \GSpage{P000779-P000780}}
\GSfield{The notes}{Two elementary transformations with parameters \(\rho_1,\rho_2\) commute, since the spectral multipliers commute. Formula (19) gives a rational expression for \(Q_{12}\), with an obscured sign in its denominator.}
\GSfield{Statement}{For compatible elementary transformations, including their full discrete data, and away from a zero denominator, the four Riccati relations imply \[Q_{12}=Q-\frac{(\rho_1+\rho_2)(Q_1-Q_2)}{\rho_1-\rho_2+\tfrac12(Q_1-Q_2)}.\] Subtracting the two first-level Riccati equations and the two second-level equations eliminates derivatives and gives this formula. This selects a minus between Q1 and Q2 by consistency, not by paleography.}
\GSfield{Reconstruction}{The derivative-elimination check is explicit editorial verification of an old relation, not a claimed newly discovered superposition law.}
\GSfield{Fate}{Permutability is a standard B\"acklund property in the method of \citet{calogerodegasperis1976inverse}. The algebraic reconstruction supplies the missing sign; commutativity on reflection coefficients alone would not prove equality of all reconstructed potentials.}
\GSfield{Present status}{As of 10 November 2026: established in the stated scope; the notebook is exposition of this result.}
\GSfield{Assessment}{literal statement: correct with missing hypotheses; publication match: partial; verification: publisher metadata checked; development: annotated; confidence: high}
\end{GSitem}
\begin{GSitem}{1977-443.10}{E}{Evolution as a generalized Backlund transformation}
\GSfield{Kind; source}{method; \GSpage{P000780}}
\GSfield{The notes}{Choose time-dependent f,g so that the multiplier equals \(\exp[2ik\alpha(-4k^2)(t-t_0)]\); evolution itself is then a B\"acklund transformation, analogous to a Fourier exponential.}
\GSfield{Statement}{For a time-independent polynomial \(\alpha\), the continuous scattering data evolve by the displayed exponential. Representing that multiplier as a quotient involving f,g gives a formal generalized B\"acklund relation. In general these f,g are not polynomials, so passing from a spectral identity to the operator equation needs a defined functional calculus and admissible inverse data.}
\GSfield{Reconstruction}{For time-dependent alpha the exponent is its time integral. The exact existence theorem is not inferred from a formal quotient alone.}
\GSfield{Fate}{A functional equation relating solutions at two times is explicitly advertised in \citet[abstract]{calogerodegasperis1976inverse}; the specific scalar-operator realization in the notebook was not checked from full text.}
\GSfield{Present status}{As of 10 November 2026: established in the stated scope; the notebook is exposition of this result.}
\GSfield{Assessment}{literal statement: correct with missing hypotheses; publication match: partial; verification: publisher metadata checked; development: annotated; confidence: high}
\end{GSitem}
\begin{GSitem}{1977-443.11}{A}{A polynomial PDE produces interacting-root dynamics}
\GSfield{Kind; source}{construction; \GSpage{P000780}}
\GSfield{The notes}{For \(\Psi_{tt}+\tfrac12\Psi_{xx}-x\Psi_x+n\Psi=0\), set \(\Psi=\prod_{j=1}^n(x-x_j(t))\). Then \(\ddot x_j=-x_j+\sum_{k\ne j}(1+2\dot x_j\dot x_k)/(x_j-x_k)\).}
\GSfield{Statement}{Let \(n\ge0\) and let a monic degree-n polynomial solution have distinct twice-differentiable roots on a time interval. Evaluating the PDE at each root and using logarithmic derivatives gives exactly the displayed ODE. Conversely the root ODE gives the polynomial PDE while roots remain distinct. Collisions require a polynomial or branched continuation, not the singular root coordinates.}
\GSfield{Reconstruction}{The nonnegative integer restriction, monicity, and simple-root interval are supplied. The marginal names Moser/Makkin/Ero are not used to assign exact priority for this particular ODE.}
\GSfield{Fate}{The corresponding pole/zero method was submitted 5 July 1977 and published September 1978 in \citet{calogero1978motion}. Publisher metadata and abstract identify the program; the displayed ODE is independently derived from the supplied PDE, not claimed read verbatim in the inaccessible full article.}
\GSfield{Present status}{As of 10 November 2026: the stated construction is established; the qualifications above are essential.}
\GSfield{Assessment}{literal statement: correct with missing hypotheses; publication match: partial; verification: publisher metadata checked; development: annotated; confidence: high}
\end{GSitem}
\begin{GSitem}{1977-443.12}{A}{Hermite coefficients linearize the root system}
\GSfield{Kind; source}{construction; \GSpage{P000781}}
\GSfield{The notes}{\(\Psi=2^{-n}\sum_{m=0}^n c_m(t)H_{n-m}(x)\), \(c_0=1\), gives \(\ddot c_m+mc_m=0\).}
\GSfield{Statement}{For physicists' Hermite polynomials satisfying \(\tfrac12H_r''-xH_r'+rH_r=0\), substitution diagonalizes the PDE: each coefficient satisfies the indicated oscillator equation. The root/coefficient map is locally invertible at distinct roots, so it solves the nonlinear system by linear evolution followed by finding polynomial roots. The frequencies \(\sqrt m\) are not generally commensurate.}
\GSfield{Reconstruction}{The Hermite normalization and the local limitation of root coordinates are explicit; no universal periodicity is inferred.}
\GSfield{Fate}{This belongs to the solvable many-body construction of \citet{calogero1978motion}, with the same 1977 submission/1978 publication distinction. The diagonalization follows directly from the displayed Hermite equation.}
\GSfield{Present status}{As of 10 November 2026: the stated construction is established; the qualifications above are essential.}
\GSfield{Assessment}{literal statement: correct; publication match: partial; verification: publisher metadata checked; development: annotated; confidence: high}
\end{GSitem}
\begin{GSitem}{1977-443.13}{E}{Hermite-zero equilibria and Stieltjes relations}
\GSfield{Kind; source}{theorem; \GSpage{P000781}}
\GSfield{The notes}{At equilibrium \(\bar x_j=\sum_{k\ne j}(\bar x_j-\bar x_k)^{-1}\); these are the Hermite zeros, credited to Stieltjes,1885.}
\GSfield{Statement}{The monic equilibrium polynomial solves \(\tfrac12\Psi''-x\Psi'+n\Psi=0\), hence equals \(2^{-n}H_n\). At a simple root, \(\Psi''(\bar x_j)/(2\Psi'(\bar x_j))=\sum_{k\ne j}(\bar x_j-\bar x_k)^{-1}\), proving the relations. The distinct real equilibrium is unique up to permutation.}
\GSfield{Reconstruction}{The proof and uniqueness follow from the Hermite ODE; the historical year remains a notebook attribution.}
\GSfield{Fate}{The notebook explicitly treats the root identities as an older Stieltjes result. Their use in the many-body construction is reflected in \citet{calogero1978motion}; the precise1885 bibliographic reference was not independently identified and is not fabricated.}
\GSfield{Present status}{As of 10 November 2026: established in the stated scope; the notebook is exposition of this result.}
\GSfield{Assessment}{literal statement: correct; publication match: partial; verification: indirect; development: annotated; confidence: high}
\end{GSitem}
\begin{GSitem}{1977-443.14}{A}{The integer interaction-matrix spectrum and missing identity}
\GSfield{Kind; source}{claim; \GSpage{P000781}}
\GSfield{The notes}{For \(A_{jj}=\sum_{k\ne j}(\bar x_j-\bar x_k)^{-2}\), \(A_{jk}=-(\bar x_j-\bar x_k)^{-2}\), the notes write \(\ddot\varepsilon+A\varepsilon=0\), and \(\operatorname{Spec}A=\{0,1,\ldots,n-1\}\).}
\GSfield{Statement}{Linearization of the displayed nonlinear ODE gives \(\ddot\varepsilon+(I+A)\varepsilon=0\). The Hermite-coefficient linearization has squared frequencies \(1,\ldots,n\), so \(A\) indeed has eigenvalues \(0,\ldots,n-1\), each simple. In particular, the constant vector lies in \(\ker A\), but its center-of-mass motion is an oscillator of frequency 1.}
\GSfield{Reconstruction}{The correct spectrum is retained; only its interpretation as the full dynamical matrix is repaired.}
\GSfield{Fate}{The result fits \citet{calogero1978motion}; the missing identity is established directly from the source equations. Even n=1 refutes the uncorrected linearized equation.}
\GSfield{Present status}{As of 10 November 2026: the stated construction is established; the qualifications above are essential.}
\GSfield{Assessment}{literal statement: false as written; publication match: partial; verification: publisher metadata checked; development: annotated; confidence: high}
\end{GSitem}

\GStalk{444}{1978-04-04}{D. B. Fuchs, with I. M. Gelfand}{P000782-P000785}{Lie cohomology, finite weights, and the Laplacian spectrum}
The catalogue date is retained as an identifier; the actual heading on P000782 reads 3 April 1978, while P000785 announces the next meeting on 4 April. The lecture explains Chevalley--Eilenberg homology, the finite-weight Hodge method, Goncharova's theorem, and the new integer spectrum for the positive Witt algebra. The latter appeared later in 1978 with Feigin as coauthor. The manuscript confuses eigenvalues with multiplicities and incorrectly identifies distinct eigenvalues with their partition labels; the corrected theorem permits collisions.

\begin{GSitem}{1977-444.1}{E}{The Chevalley-Eilenberg chain and cochain complexes}
\GSfield{Kind; source}{definition; \GSpage{P000782-P000783}}
\GSfield{The notes}{\(C_q(\mathfrak g)=\Lambda^q\mathfrak g\) and \(d(g_1\wedge\cdots\wedge g_q)=\sum_{s<t}(-1)^{s+t-1}[g_s,g_t]\wedge g_1\wedge\cdots\widehat g_s\cdots\widehat g_t\cdots\wedge g_q\). The dual complex defines cohomology, with continuous functionals in the topological case.}
\GSfield{Statement}{For a real Lie algebra and trivial coefficients this is the homological Chevalley--Eilenberg differential; the Jacobi identity gives \(d^2=0\). The dual differential raises degree. For formal vector fields one must use continuous cochains in the inverse-limit topology, or equivalently the indicated graded dual, rather than all algebraic functionals on an infinite product.}
\GSfield{Fate}{These standard definitions are recalled in \citet[Section~1]{gelfandfeiginfuchs1978laplace}, which explicitly specifies continuous cohomology and trivial coefficients.}
\GSfield{Present status}{As of 10 November 2026: established in the stated scope; the notebook is exposition of this result.}
\GSfield{Assessment}{literal statement: correct; publication match: exact; verification: read the result; development: annotated; confidence: high}
\end{GSitem}
\begin{GSitem}{1977-444.2}{E}{Hodge representatives on finite weight complexes}
\GSfield{Kind; source}{method; \GSpage{P000783-P000784}}
\GSfield{The notes}{An inner product identifies chains and cochains, \(\Delta=d\delta+\delta d\), and every homology class has exactly one harmonic representative; \(\Delta a=0\) iff \(da=\delta a=0\).}
\GSfield{Statement}{For a finite-dimensional chain complex with positive-definite inner products and \(\delta=d^*\), \(\langle\Delta a,a\rangle=\|da\|^2+\|\delta a\|^2\). Orthogonal decomposition gives a unique harmonic representative of each homology class. The argument applies separately to each finite weight block of \(L_1\); an arbitrary infinite-dimensional completion requires closed-range and domain hypotheses.}
\GSfield{Reconstruction}{The unrestricted opening assertion is read in the finite-weight setting justified later in the lecture, not as a general Hilbert-complex theorem.}
\GSfield{Fate}{The finite-block qualification is explicit in \citet[Section~1]{gelfandfeiginfuchs1978laplace} and \citet[Proposition~1]{fuchswilmarth2010laplacian}.}
\GSfield{Present status}{As of 10 November 2026: established in the stated scope; the notebook is exposition of this result.}
\GSfield{Assessment}{literal statement: correct with missing hypotheses; publication match: exact; verification: read the result; development: annotated; confidence: high}
\end{GSitem}
\begin{GSitem}{1977-444.3}{E}{Formal vector fields and the two vanishing-jet subalgebras}
\GSfield{Kind; source}{definition; \GSpage{P000783-P000784}}
\GSfield{The notes}{\(W_1=\mathbb R[[x]]\partial_x\), \(L_0=x\mathbb R[[x]]\partial_x\), \(L_1=x^2\mathbb R[[x]]\partial_x\). The formal diffeomorphism group has Lie algebra \(L_0\), and \(W_1/L_0\) has dimension 1.}
\GSfield{Statement}{The group of origin-preserving invertible formal changes of one variable has tangent algebra \(L_0\); its subgroup tangent to the identity has algebra \(L_1\). The constant vector field supplies the one-dimensional quotient of \(W_1\) by \(L_0\) as vector spaces. Formal topology is essential.}
\GSfield{Fate}{These are the algebras in \citet[Section~1]{goncharova1973line} and \citet[Section~1]{gelfandfeiginfuchs1978laplace}.}
\GSfield{Present status}{As of 10 November 2026: established in the stated scope; the notebook is exposition of this result.}
\GSfield{Assessment}{literal statement: correct; publication match: exact; verification: read the result; development: annotated; confidence: high}
\end{GSitem}
\begin{GSitem}{1977-444.4}{F}{Finite defect and the unqualified integration obstruction}
\GSfield{Kind; source}{claim; \GSpage{P000783-P000784}}
\GSfield{The notes}{The notes say that \(W_1\) has no local Lie group and define finite defect as the minimum codimension of a subalgebra integrating to a Lie group; \(L_0\) is the example.}
\GSfield{Statement}{In the ordinary origin-preserving formal substitution category, translation by a nonzero real constant cannot be substituted into an arbitrary formal series; hence its Lie algebra is \(L_0\), not all \(W_1\). A universal assertion that \(W_1\) integrates to no possible local group requires a specified category, topology, and regularity. The definition of defect depends on that choice.}
\GSfield{Reconstruction}{The concrete substitution obstruction is all that the supplied wording supports without a category of integration.}
\GSfield{Fate}{The distinction among the three formal vector-field algebras is established, but no exact theorem matching the notebook's unrestricted no-local-group statement was identified. Formal group functors over nilpotent bases and ordinary real local Lie groups must not be conflated.}
\GSfield{Present status}{As of 10 November 2026: the restricted substitution statement is clear; the absolute integration assertion remains insufficiently specified.}
\GSfield{Assessment}{literal statement: not determinable; publication match: none found; verification: unresolved lead; development: annotated; confidence: medium}
\end{GSitem}
\begin{GSitem}{1977-444.5}{E}{Center-free algebras and their derivation algebra}
\GSfield{Kind; source}{claim; \GSpage{P000784}}
\GSfield{The notes}{If the center is zero, the Lie algebra of the automorphism group is said to be an over-algebra or subalgebra of the original algebra; the relevant word is unclear.}
\GSfield{Statement}{The algebraic statement is an injection \(\operatorname{ad}:\mathfrak g\hookrightarrow\operatorname{Der}(\mathfrak g)\) when the center is zero. For a finite-dimensional real Lie algebra, \(\operatorname{Lie}\operatorname{Aut}(\mathfrak g)=\operatorname{Der}(\mathfrak g)\). For arbitrary infinite-dimensional algebras, existence and Lie structure of the automorphism group need further hypotheses.}
\GSfield{Reconstruction}{The algebraic injection determines the direction of containment if over-algebra was intended. The scan itself does not settle the word.}
\GSfield{Fate}{This is the standard adjoint/derivation construction. The earlier index has 1971-243.4 on automorphisms of formal vector fields, but the current fragment supplies no exact recurrence of a separately indexed problem.}
\GSfield{Present status}{As of 10 November 2026: established in the stated scope; the notebook is exposition of this result.}
\GSfield{Assessment}{literal statement: correct with missing hypotheses; publication match: partial; verification: indirect; development: annotated; confidence: high}
\end{GSitem}
\begin{GSitem}{1977-444.6}{E}{Previously computed cohomology of W1 and L0}
\GSfield{Kind; source}{remark; \GSpage{P000784}}
\GSfield{The notes}{The cohomology of \(W_1\) and \(L_0\) is said to be easily found by Gelfand and Fuchs, in contrast with \(L_1\).}
\GSfield{Statement}{For continuous real cohomology with trivial coefficients, \(H^*(L_0)\) is the exterior algebra on the one-dimensional linear-part character in degree 1. The formal one-variable algebra \(W_1\) has nonzero cohomology in degrees 0 and 3. The comparison concerns these precise coefficient and topology conventions.}
\GSfield{Reconstruction}{The note gives no dimensions here; the conventional trivial-coefficient specialization is supplied.}
\GSfield{Fate}{\citet{gelfandfuchs1970formal}, received November 1969, is the earlier formal-vector-field computation. Its metadata and abstract were checked; the specific one-dimensional specialization is a reconstructed standard consequence, not a newly claimed full-text verification.}
\GSfield{Present status}{As of 10 November 2026: established in the stated scope; the notebook is exposition of this result.}
\GSfield{Assessment}{literal statement: correct with missing hypotheses; publication match: partial; verification: publisher metadata checked; development: annotated; confidence: high}
\end{GSitem}
\begin{GSitem}{1977-444.7}{E}{The positive Witt bracket and finite weight decomposition}
\GSfield{Kind; source}{construction; \GSpage{P000784}}
\GSfield{The notes}{\([e_n,e_m]=(m-n)e_{n+m}\), \(e_n=x^{n+1}\partial_x\); the chain differential preserves \(i_1+\cdots+i_q\), so chains split into finite-dimensional complexes \(C_{q,n}\).}
\GSfield{Statement}{In \(L_1\), \(\deg e_i=i\ge1\). The finite set of strictly increasing positive tuples summing to n spans \(C_{q,n}\); the bracket preserves total weight. The formula after the bracket in notebook (11) must begin with the surviving factor \(e_{i_1}\), not a new fixed \(e_1\).}
\GSfield{Reconstruction}{The extra subscript is a mathematical repair of a faithful source typo; the weight argument confirms it.}
\GSfield{Fate}{The same finite-weight construction appears in \citet[Section~1]{goncharova1973line} and \citet[Section~1]{gelfandfeiginfuchs1978laplace}.}
\GSfield{Present status}{As of 10 November 2026: established in the stated scope; the notebook is exposition of this result.}
\GSfield{Assessment}{literal statement: false as written; publication match: exact; verification: read the result; development: annotated; confidence: high}
\end{GSitem}
\begin{GSitem}{1977-444.8}{E}{Goncharova cohomology and pentagonal weights}
\GSfield{Kind; source}{theorem; \GSpage{P000784}}
\GSfield{The notes}{For \(q\ge1\), \(\dim H_q(L_1)=2\), with generator weights \((3q^2\pm q)/2\).}
\GSfield{Statement}{Over characteristic zero, continuous trivial-coefficient \(H^q(L_1)\) is one-dimensional in each of the two weights \((3q^2-q)/2\), \((3q^2+q)/2\), and zero in all other weights; \(H^0\) is one-dimensional in weight 0. Weightwise duality gives the corresponding graded homology statement.}
\GSfield{Fate}{\citet[Introduction and main theorem]{goncharova1973line}, submitted 21 June 1972 and published 1973, proves the result; its introduction records a preceding short1972 announcement.}
\GSfield{Present status}{As of 10 November 2026: established in the stated scope; the notebook is exposition of this result.}
\GSfield{Assessment}{literal statement: correct; publication match: exact; verification: read the result; development: annotated; confidence: high}
\end{GSitem}
\begin{GSitem}{1977-444.9}{A}{The metric matters for the nonzero Laplacian spectrum}
\GSfield{Kind; source}{definition; \GSpage{P000784-P000785}}
\GSfield{The notes}{The candidates \(e_k\), \(e_k/k!\), and \(\alpha_ke_k\) are considered as orthonormal bases; the lecture chooses \(e_k\), and increasing wedge monomials are orthonormal.}
\GSfield{Statement}{Fixing \(\langle e_i,e_j\rangle=\delta_{ij}\) determines the induced inner products, adjoint differential, and nonzero spectrum of \(\Delta\). A different rescaling changes the adjoint and generally the spectrum, though harmonic dimensions still give the same cohomology on finite blocks.}
\GSfield{Fate}{This exact metric is specified by \citet[Section~1]{gelfandfeiginfuchs1978laplace} and in \citet[before Theorem~2]{fuchswilmarth2010laplacian}.}
\GSfield{Present status}{As of 10 November 2026: the stated construction is established; the qualifications above are essential.}
\GSfield{Assessment}{literal statement: correct; publication match: exact; verification: read the result; development: annotated; confidence: high}
\end{GSitem}
\begin{GSitem}{1977-444.10}{A}{Integer Laplacian eigenvalues and their corrected multiplicities}
\GSfield{Kind; source}{theorem; \GSpage{P000785}}
\GSfield{The notes}{For gap-three partitions \(I=(i_1<\cdots<i_q)\), the notes display \(E(I)=\sum_s\binom{i_s}{3}-\sum_{s<t}i_si_t\), but call it a multiplicity and claim distinct eigenvalues are counted by such partitions.}
\GSfield{Statement}{For the fixed orthonormal metric, the total weight-n chain space has eigenvalues E(I), for gap-three partitions of n. Define \(\alpha(I)\) to count the first part if \(i_1\ge3\), plus each gap strictly greater than 3. A label contributes multiplicity \(2^{\alpha(I)}\); if labels yield the same E, add these multiplicities. Thus E is the eigenvalue, not its multiplicity. At weight 25, \(I=(1,10,14)\) and \(J=(3,7,15)\) both give320, with total multiplicity \(4+8=12\).}
\GSfield{Reconstruction}{The corrected multiplicity rule comes from the later source. The explicit collision is an independent finite verification and is labelled as such in the development appendix.}
\GSfield{Fate}{\citet{gelfandfeiginfuchs1978laplace} was received 14 June 1978. \citet[Theorem~2]{fuchswilmarth2010laplacian} states the corrected version and credits an important1993 correction by Weinstein. The chapter independently checks integer CE matrices through weight 22 and the weight 25 collision; this verifies examples, not a new general proof.}
\GSfield{Present status}{As of 10 November 2026: the corrected spectral theorem is established; the notebook is not an open problem. Its distinct-eigenvalue and multiplicity wording requires the stated repairs.}
\GSfield{Assessment}{literal statement: false as written; publication match: exact; verification: read the result; development: annotated; confidence: high}
\end{GSitem}
\begin{GSitem}{1977-444.11}{A}{Nonnegativity and the two zero-eigenvalue partitions}
\GSfield{Kind; source}{problem; \GSpage{P000785}}
\GSfield{The notes}{Exercise: E(I) is nonnegative and vanishes exactly for \((1,4,7,\ldots,3q-2)\) and \((2,5,8,\ldots,3q-1)\).}
\GSfield{Statement}{For every nonempty increasing positive tuple with consecutive gaps at least3, \(E(I)\ge0\); equality holds exactly for the two displayed arithmetic progressions. Their weights are the two pentagonal numbers in Goncharova's theorem.}
\GSfield{Fate}{The assertion is explicitly given with \citet[Theorem~1]{gelfandfeiginfuchs1978laplace} and reiterated immediately before \citet[Theorem~2]{fuchswilmarth2010laplacian}. It is a resolved combinatorial exercise, not an undeveloped research problem.}
\GSfield{Present status}{As of 10 November 2026: the stated construction is established; the qualifications above are essential.}
\GSfield{Assessment}{literal statement: correct; publication match: exact; verification: read the result; development: annotated; confidence: high}
\end{GSitem}
\begin{GSitem}{1977-444.12}{A}{Vandermonde-cubed harmonic representatives}
\GSfield{Kind; source}{construction; \GSpage{P000785}}
\GSfield{The notes}{Identify \(e_i\) with \(z^i\). Multiply \(\prod_{i<j}(z_i-z_j)^3\) by \(\prod_i z_i\) or \(\prod_i z_i^2\), expand the alternating polynomial, and replace alternants by wedge monomials.}
\GSfield{Statement}{Under the alternating-polynomial realization of \(\Lambda^qL_1\), the two displayed polynomials determine nonzero harmonic chains in weights \((3q^2-q)/2\) and \((3q^2+q)/2\), respectively. They represent the two homology classes in degreeq, up to overall nonzero normalization and the sign convention for alternants.}
\GSfield{Reconstruction}{The notebook says eigenvectors; in these two cases the eigenvalue is zero. Overall alternant normalization does not affect the represented line.}
\GSfield{Fate}{Explicit harmonic representatives are supplied by \citet[Section~2]{gelfandfeiginfuchs1978laplace}. The chapter also verifies both boundary and adjoint-boundary vanish for q=1,2,3,4 with exact arithmetic, as a finite check of the historical result.}
\GSfield{Present status}{As of 10 November 2026: the stated construction is established; the qualifications above are essential.}
\GSfield{Assessment}{literal statement: correct; publication match: exact; verification: read the result; development: annotated; confidence: high}
\end{GSitem}

\GStalk{445}{}{M. D. Kruskal (adjacent lecture)}{P000786-P000787}{Catalogue heading 1978 and a lecture on surreal numbers}
Although the catalogue entry is only a year heading, both supplied pages contain substantial mathematics from Kruskal's lecture. The scan reads 30 October 1978, with November a paleographic alternative; it is adjacent later material, not a lecture of the principal 1977--1978 academic-year sequence. The pages develop sign sequences, simplicity, cuts, and addition of surreal numbers, and mention Conway's book and Knuth's narrative exposition. No research problem is posed.

\begin{GSitem}{1977-445.1}{E}{Ordinal sign sequences and their numerical order}
\GSfield{Kind; source}{definition; \GSpage{P000786}}
\GSfield{The notes}{A number is a map from a well-ordered set to two arrows. Compare arrows lexicographically; the empty path is zero and an infinite upward sequence is \(\omega\). A binary tree and periodic-arrow symbols are drawn.}
\GSfield{Statement}{A surreal sign sequence is a function \(s:\alpha\to\{-,+\}\) for a set ordinal \(\alpha\). At the first differing place, compare \(-<\text{termination}<+\). This convention orders unequal lengths as well as equal lengths. The empty sequence is zero; \(+\), \(++\), and \(+-\) represent \(1,2,1/2\), while the all-plus sequence of length \(\omega\) is \(\omega\). A periodic pattern needs its ordinal length specified.}
\GSfield{Reconstruction}{Well-ordered domains are replaced by their ordinal order types and the missing termination comparison is supplied.}
\GSfield{Fate}{This is an alternative realization of the numbers constructed by cuts in \citet[Chapter~0, pp.~4--7]{conway1976numbers}. The original book appeared in 1976; the opened text is its2001 second-edition reprint. The specific sign-sequence convention is a reconstruction of Kruskal's notes, not claimed quoted from those pages.}
\GSfield{Present status}{As of 10 November 2026: established in the stated scope; the notebook is exposition of this result.}
\GSfield{Assessment}{literal statement: correct with missing hypotheses; publication match: partial; verification: indirect; development: annotated; confidence: high}
\end{GSitem}
\begin{GSitem}{1977-445.2}{E}{Simplicity, truncations, and canonical cuts}
\GSfield{Kind; source}{definition; \GSpage{P000786-P000787}}
\GSfield{The notes}{An earlier number is on the path from 0 to the given number. Truncations split into earlier smaller and earlier greater numbers, and \(x=\{x\rfloor:x\rceil\}\).}
\GSfield{Statement}{The simplicity relation means proper initial segment of the sign sequence, not merely a smaller numerical value. Let \(L_x\) and \(R_x\) consist of proper truncations respectively less and greater than x. Then x is the unique simplest number strictly between these two sets, written \(x=\{L_x\mid R_x\}\).}
\GSfield{Reconstruction}{Canonical initial-segment simplicity is distinguished from the broader birthday ordering; different numbers can be born at the same stage.}
\GSfield{Fate}{The cut construction and distinction between a number and its many forms are explained in \citet[Chapter~0, pp.~4--6]{conway1976numbers}.}
\GSfield{Present status}{As of 10 November 2026: established in the stated scope; the notebook is exposition of this result.}
\GSfield{Assessment}{literal statement: correct; publication match: partial; verification: read the result; development: annotated; confidence: high}
\end{GSitem}
\begin{GSitem}{1977-445.3}{E}{The simplest number between two separated sets}
\GSfield{Kind; source}{theorem; \GSpage{P000787}}
\GSfield{The notes}{If A,B are sets of numbers and every a in A is less than every b in B, there exists a unique earliest x with \(a<x<b\); proof by transfinite induction.}
\GSfield{Statement}{For sets \(A,B\) of surreal numbers with \(A<B\), there is a unique surreal of minimum birthday in the interval \(\{x:A<x<B\}\); equivalently it is the simplest element of that interval. Either set may be empty. Set-sizedness is essential: no surreal exceeds every member of the proper class of ordinals.}
\GSfield{Reconstruction}{The quantifier is over sets, as the manuscript says; it is not extended to arbitrary proper classes.}
\GSfield{Fate}{The construction rule of \citet[Chapter~0, p.~4]{conway1976numbers} is exactly the separated-cut operation. The notebook's formulation uses the equivalent sign-tree simplicity theorem.}
\GSfield{Present status}{As of 10 November 2026: established in the stated scope; the notebook is exposition of this result.}
\GSfield{Assessment}{literal statement: correct; publication match: partial; verification: read the result; development: annotated; confidence: high}
\end{GSitem}
\begin{GSitem}{1977-445.4}{E}{Recursive addition and ordered additive-group laws}
\GSfield{Kind; source}{construction; \GSpage{P000787}}
\GSfield{The notes}{Define \(x+y=\{x\rfloor+y,x+y\rfloor:x\rceil+y,x+y\rceil\}\). The notes prove \(x+0=x\), assert monotonicity, commutativity and associativity, and obtain the opposite by reversing arrows.}
\GSfield{Statement}{Well-founded recursion on both arguments defines \[x+y=\{x^L+y,x+y^L\mid x^R+y,x+y^R\}.\] It makes the surreal numbers an ordered abelian group: addition is strictly increasing in either argument, zero is the identity, and the negative reverses all signs. Associativity has the usual form \((x+y)+z=x+(y+z)\). The phrase simplest monotone number abbreviates these recursive inequalities; monotonicity alone is not a uniqueness axiom for addition.}
\GSfield{Reconstruction}{The ambiguous sign in the handwritten associativity expression is interpreted contextually as plus, not declared visibly resolved.}
\GSfield{Fate}{The addition and negation formulas are explicitly given in \citet[Chapter~0, p.~5]{conway1976numbers}; the elementary zero calculation occupies pp.~6--7. No new arithmetic or unresolved associativity question is present.}
\GSfield{Present status}{As of 10 November 2026: established in the stated scope; the notebook is exposition of this result.}
\GSfield{Assessment}{literal statement: correct with missing hypotheses; publication match: exact; verification: read the result; development: annotated; confidence: high}
\end{GSitem}

\GStalk{451}{}{Unidentified preceding lecturer; I. M. Gelfand and A. A. Belavin}{P000810-P000812}{Catalogue heading 1979 with representation cohomology and instanton questions}
The catalogue year heading encloses two different pieces of adjacent 1979 material. P000810--P000811 continue an unidentified preceding lecture on extensions, continuous cohomology, smooth vectors, Harish--Chandra modules, and semidirect products. P000812 is explicitly dated 12 February 1979 and gives a physics reading programme followed by Belavin's ADHM matrix-parameter question and a short gauge/confinement fragment. The reading list is provenance evidence, not itself a collection of mathematical theorems. Its path-integral coauthor is written Gibbs, plausibly a notebook error for Hibbs.

\begin{GSitem}{1977-451.1}{E}{Infinitesimal characters and algebraic block separation}
\GSfield{Kind; source}{definition; \GSpage{P000810}}
\GSfield{The notes}{If central elements act on E by scalars, \(u e=\chi(u)e\), then \(\chi\in\operatorname{Hom}(Z(\mathfrak g),\mathbb C)\) is an infinitesimal character. The following page refers to a lemma about distinct infinitesimal characters.}
\GSfield{Statement}{Here \(Z(\mathfrak g)\) means \(Z(U(\mathfrak g))\), and \(\chi\) is a unital algebra homomorphism. If E,F are algebraic \(U(\mathfrak g)\)-modules with distinct scalar central characters, every extension of E by F splits: a central element with distinct scalar values yields a polynomial spectral projector onto F. For topological representation spaces its derived action may be unbounded, so this argument does not automatically produce a continuous splitting.}
\GSfield{Reconstruction}{The notes begin with scalar action as an assumption. An unrestricted Schur-lemma assertion for arbitrary topological irreducibles is not added.}
\GSfield{Fate}{The algebraic assertion is the standard separation of central-character blocks. The paper \citet[Section~V(C)]{delorme1977cohomologie} supplies the contemporaneous representation-cohomology context; the precise unnamed lemma is reconstructed algebraically, not identified as a numbered theorem of that paper.}
\GSfield{Present status}{As of 10 November 2026: established in the stated scope; the notebook is exposition of this result.}
\GSfield{Assessment}{literal statement: correct with missing hypotheses; publication match: partial; verification: indirect; development: annotated; confidence: high}
\end{GSitem}
\begin{GSitem}{1977-451.2}{E}{Relative first cohomology classifies split-on-h extensions}
\GSfield{Kind; source}{construction; \GSpage{P000810}}
\GSfield{The notes}{Relative cocycles satisfy \(f|_{\mathfrak h}=0\) and describe extensions trivial over \(\mathfrak h\). The relation between relative and absolute \(H^1\) is written with a membership sign. The cocycle differential is displayed.}
\GSfield{Statement}{For algebraic modules E,F, set \(M=\operatorname{Hom}_{\mathbb C}(E,F)\), with \((xT)e=x(Te)-T(xe)\). Relative1-cocycles are absolute cocycles vanishing on \(\mathfrak h\), automatically \(\mathfrak h\)-equivariant; divide by coboundaries from \(M^{\mathfrak h}\). Their classes classify extensions equipped with an \(\mathfrak h\)-splitting up to equivalence. The forgetful map \(H^1(\mathfrak g,\mathfrak h;M)\hookrightarrow H^1(\mathfrak g;M)\) is injective, with image those extensions that split over \(\mathfrak h\).}
\GSfield{Reconstruction}{The membership symbol is replaced by the canonical injection. The notes use the negative of one common differential convention; its cocycle equation is equivalent.}
\GSfield{Fate}{This is the classical degree-one cocycle/extension interpretation. Its injectivity is immediate: if a relative cocycle is \(dm\), its vanishing on \(\mathfrak h\) forces \(m\in M^{\mathfrak h}\).}
\GSfield{Present status}{As of 10 November 2026: established in the stated scope; the notebook is exposition of this result.}
\GSfield{Assessment}{literal statement: false as written; publication match: partial; verification: indirect; development: annotated; confidence: high}
\end{GSitem}
\begin{GSitem}{1977-451.3}{E}{Continuous cocycles and topologically split extensions}
\GSfield{Kind; source}{construction; \GSpage{P000810}}
\GSfield{The notes}{A continuous G-module is a topological vector space with continuous action. A cocycle \(f:G\to\operatorname{Hom}(E,F)\) is required to act jointly continuously in g and e; an alternative uses smooth vectors.}
\GSfield{Statement}{For continuous representations E,F, a cocycle with jointly continuous evaluation defines the triangular action on \(F\oplus E\). Modulo continuous changes of splitting, it classifies extensions that admit a continuous linear splitting. A general exact sequence of topological vector spaces need not split as vector spaces continuously, so these cocycles need not classify every topological extension.}
\GSfield{Reconstruction}{The required continuous splitting and a topology on the Hom space cannot be suppressed.}
\GSfield{Fate}{The continuity and cochain conventions are set out for Hilbert modules in \citet[Section~0]{delorme1977cohomologie} and for Frechet modules in \citet[Sections~2--4]{blanc1979cohomologie}. The extension category is an explicit reconstruction.}
\GSfield{Present status}{As of 10 November 2026: established in the stated scope; the notebook is exposition of this result.}
\GSfield{Assessment}{literal statement: correct with missing hypotheses; publication match: partial; verification: read the result; development: annotated; confidence: high}
\end{GSitem}
\begin{GSitem}{1977-451.4}{B}{Smooth-vector regularization and the van Est comparison}
\GSfield{Kind; source}{theorem; \GSpage{P000810}}
\GSfield{The notes}{For E equal to the trivial one-dimensional module, \(H^1(G,F)\cong H^1(G,F_\infty)\cong H^1(\mathfrak g,\mathfrak k;F_\infty)\), attributed respectively to Delorme--Blanc and van Est.}
\GSfield{Statement}{For a finite-dimensional connected Lie group G and a continuous Frechet G-module F, smoothing compares continuous group cohomology with differentiable cohomology of the Frechet smooth-vector module \(F^\infty\). Under the usual van Est hypotheses, including maximal compact K and the contractibility of \(G/K\), the latter agrees with relative Lie algebra cohomology; the K-invariance, not merely \(\mathfrak k\)-invariance, must be imposed if K is disconnected.}
\GSfield{Reconstruction}{The module class, continuous cohomology and compact-subgroup invariance are essential. The date of the preceding unidentified lecture is not fixed by P812.}
\GSfield{Fate}{The smooth-vector assertion is read in \citet[Theorem~5.2, pp.~152--153]{blanc1979cohomologie}, published 1979; an exact first-circulation date was not established. The second arrow is a classical van Est comparison; this audit verifies it only indirectly through the sources' references and framework.}
\GSfield{Present status}{As of 10 November 2026: the stated construction is established; the qualifications above are essential.}
\GSfield{Assessment}{literal statement: correct with missing hypotheses; publication match: partial; verification: read the result; development: annotated; confidence: high}
\end{GSitem}
\begin{GSitem}{1977-451.5}{E}{A real unipotent extension and the compact-period obstruction}
\GSfield{Kind; source}{construction; \GSpage{P000810}}
\GSfield{The notes}{For \(\mathbb R\), \(t\mapsto\begin{pmatrix}1&at\\0&1\end{pmatrix}\) is a nontrivial extension of the trivial module. The following sentence denies a general group/algebra correspondence and says to use cohomology trivial on the compact subgroup.}
\GSfield{Statement}{For \(a\ne0\) this is a nonsplit representation of the additive real group and of its one-dimensional Lie algebra. It does not descend to \(\mathbb R/T\mathbb Z\) for \(T\ne0\), because its value at T is not the identity. Thus it illustrates a period obstruction for a quotient group, not failure to integrate this Lie-algebra module to the simply connected group \(\mathbb R\).}
\GSfield{Reconstruction}{The ambiguous denial of correspondence is not read as a claim that the displayed real-group representation fails to exist.}
\GSfield{Fate}{The computation is elementary and consistent with the role of relative cohomology in the preceding comparison. The notebook does not name the compact quotient, so that part is explicitly an explanatory reconstruction.}
\GSfield{Present status}{As of 10 November 2026: established in the stated scope; the notebook is exposition of this result.}
\GSfield{Assessment}{literal statement: correct with missing hypotheses; publication match: partial; verification: indirect; development: annotated; confidence: high}
\end{GSitem}
\begin{GSitem}{1977-451.6}{E}{Non-Hausdorff first cohomology of the affine group}
\GSfield{Kind; source}{claim; \GSpage{P000811}}
\GSfield{The notes}{For an irreducible unitary representation of the line's affine group, \(B^1\) is dense in \(Z^1\) but \(\dim Z^1/B^1=\infty\); weak containment of the trivial representation is cited as the reason.}
\GSfield{Statement}{One must distinguish ordinary continuous \(H^1=Z^1/B^1\) from reduced \(\overline H^1=Z^1/\overline{B^1}\). For the noncharacter irreducible unitary representations of the connected affine group, the latter vanishes, while the former can be nonzero and non-Hausdorff. The manuscript does not specify its chosen irreducible realization or a proof of the stated infinite algebraic dimension.}
\GSfield{Reconstruction}{The closure is decisive. Weak containment alone is not substituted for every needed hypothesis in a general dimension theorem.}
\GSfield{Fate}{The density assertion is covered by \citet[Theorem~V.6]{delorme1977cohomologie}; its introduction p.~284 credits an earlier affine-group calculation to B. E. Johnson. The infinite-dimensional quotient claim was not independently matched to an exact full-text statement and remains a limitation of this identification.}
\GSfield{Present status}{As of 10 November 2026: established in the stated scope; the notebook is exposition of this result.}
\GSfield{Assessment}{literal statement: correct with missing hypotheses; publication match: partial; verification: read the result; development: annotated; confidence: medium}
\end{GSitem}
\begin{GSitem}{1977-451.7}{F}{Semisimple cohomology and a contradictory orthogonal-group fragment}
\GSfield{Kind; source}{claim; \GSpage{P000811}}
\GSfield{The notes}{The notes say the nonzero unitary H1 cases for semisimple groups are known, then write under nonsemisimple groups that \(SO_0(m,n)\) has no nonunitary representations with nonzero cohomology.}
\GSfield{Statement}{For a connected real simple Lie group whose Lie algebra is neither \(\mathfrak{so}(n,1)\) nor \(\mathfrak{su}(n,1)\), every irreducible unitary representation has trivial first continuous cohomology. This valid theorem does not repair the literal notebook sentence: most \(SO_0(m,n)\) are semisimple, and nonunitary is not the theorem's unitary hypothesis.}
\GSfield{Fate}{Delorme\textquotesingle s precise unitary theorem applies here \citep[Theorem~V.4]{delorme1977cohomologie}. The manuscript's two conflicting words are confirmed in the scan; several different corrections are possible, so none is silently adopted.}
\GSfield{Present status}{As of 10 November 2026: the nearby unitary theorem is established; the exact intended orthogonal-group assertion is not identifiable.}
\GSfield{Assessment}{literal statement: false as written; publication match: partial; verification: read the result; development: annotated; confidence: low}
\end{GSitem}
\begin{GSitem}{1977-451.8}{F}{An unspecified SL2 topological-extension counterexample}
\GSfield{Kind; source}{claim; \GSpage{P000811}}
\GSfield{The notes}{The infinitesimal-character lemma is false in the first continuous-Hom approach; a counterexample for \(SL(2,\mathbb R)\) is mentioned. It becomes true for smooth modules, whose cohomology is harder to calculate.}
\GSfield{Statement}{An algebraic central spectral projector need not define a continuous operator on a Hilbert globalization, while the center acts on smooth vectors. Therefore the category of extensions matters. Without the two representations, their topology, the central characters and the extension map, the alleged counterexample cannot be reconstructed uniquely.}
\GSfield{Fate}{The search through \citet{delorme1977cohomologie} and \citet{blanc1979cohomologie} established the surrounding comparison theory but did not identify this exact counterexample. No new open problem is inferred from a missing example.}
\GSfield{Present status}{As of 10 November 2026: the surviving wording does not determine a theorem with all hypotheses.}
\GSfield{Assessment}{literal statement: not determinable; publication match: thematic only; verification: unresolved lead; development: annotated; confidence: low}
\end{GSitem}
\begin{GSitem}{1977-451.9}{E}{K-finite vectors and algebraic extension computations}
\GSfield{Kind; source}{method; \GSpage{P000811}}
\GSfield{The notes}{Replace representations by their K-finite vectors. For all irreducible Harish--Chandra modules of SL2 one can compute \(H^1(\mathfrak g,\mathfrak k;\operatorname{Hom}(E_K,F_K))\), without topology.}
\GSfield{Statement}{In the admissible \((\mathfrak g,K)\)-module category, finite-dimensional K-types make K-splittings available. Relative cocycles describe first extensions, with K-equivariance and the appropriate K-finite part of Hom. For SL2 the resulting calculation depends on the two modules and their infinitesimal characters; no dimensions for individual pairs are supplied here.}
\GSfield{Reconstruction}{The category and K-invariance replace an unrestricted algebraic Hom formula.}
\GSfield{Fate}{The transition to K-finite vectors is used concretely in \citet[Section~V(C), especially p.~319]{delorme1977cohomologie}. The blanket all-pairs computability sentence is an announced method, not a supplied classification table verified by this audit.}
\GSfield{Present status}{As of 10 November 2026: established in the stated scope; the notebook is exposition of this result.}
\GSfield{Assessment}{literal statement: correct with missing hypotheses; publication match: partial; verification: read the result; development: annotated; confidence: high}
\end{GSitem}
\begin{GSitem}{1977-451.10}{E}{Semidirect products and dual orbits}
\GSfield{Kind; source}{construction; \GSpage{P000811}}
\GSfield{The notes}{For \(G=B\rtimes A\), \(B=\mathbb R^n\), the A-action determines the product; consider an orbit X in \(B^*\), with the affine group as example.}
\GSfield{Statement}{Given \(\rho:A\to GL(B)\), multiplication is \((b,a)(b',a')=(b+\rho(a)b',aa')\). The contragredient action on \(B^*\) acts on the characters \(b\mapsto e^{i\langle b,\xi\rangle}\). Orbit and stabilizer data are inputs to induction; the fragment does not state the regularity conditions or a classification theorem for all semidirect products.}
\GSfield{Reconstruction}{No unrecorded full Mackey theorem is counted as an additional notebook assertion.}
\GSfield{Fate}{Semidirect products are part of the older representation-theoretic framework discussed in \citet[Introduction and Section~V(D)]{delorme1977cohomologie}. Only the displayed setup is attributed to this fragment.}
\GSfield{Present status}{As of 10 November 2026: established in the stated scope; the notebook is exposition of this result.}
\GSfield{Assessment}{literal statement: correct with missing hypotheses; publication match: thematic only; verification: read the result; development: annotated; confidence: high}
\end{GSitem}
\begin{GSitem}{1977-451.11}{E}{The quaternionic ADHM equations and their quotient}
\GSfield{Kind; source}{construction; \GSpage{P000812}; A. A. Belavin}
\GSfield{The notes}{Belavin asks about quaternion matrices A,B of size \(N\times(N+1)\), with \(M=xA+B\) and \(MM^+\) real for every quaternion x, modulo \(M\mapsto GMV\), \(VV^+=1\).}
\GSfield{Statement}{In a consistent transposed quaternionic ADHM convention, impose the reality equations and the nondegeneracy/rank conditions needed to obtain a smooth charge-N SU2 instanton on the conformal four-sphere. The equivalence group is \(GL(N,\mathbb R)\) on the rank-N factor and \(Sp(N+1)\) on the quaternionic ambient factor. The order of quaternion multiplication must be fixed when transposing the usual \((N+1)\times N\) convention.}
\GSfield{Reconstruction}{The type of G, nondegeneracy, reality structure and transpose convention are missing from the notebook. Dropping them admits singular or spurious pairs.}
\GSfield{Fate}{The exact linear-algebra classification and these groups occur in \citet[pp.~185--186]{atiyahhitchindrinfeldmanin1978instantons}, received 16 December 1977 and published 6 March 1978, before the actual12 February 1979 notebook meeting. The printed source also proves completeness, not merely necessity of the equations.}
\GSfield{Present status}{As of 10 November 2026: established in the stated scope; the notebook is exposition of this result.}
\GSfield{Assessment}{literal statement: correct with missing hypotheses; publication match: exact; verification: read the result; development: annotated; confidence: high}
\end{GSitem}
\begin{GSitem}{1977-451.12}{F}{A beautiful parametrization of instanton matrices}
\GSfield{Kind; source}{question; \GSpage{P000812}; A. A. Belavin}
\GSfield{The notes}{There should be \(8N-3\) parameters in all; find a beautiful parametrization of such matrix pairs. The question is humorously attributed to a Chernogolovka competition.}
\GSfield{Statement}{For \(N\ge1\), the unframed smooth SU2 instanton moduli space on \(S^4\) has real dimension \(8N-3\). The ADHM quotient is an algebraic description; obtaining explicit coordinates, solving the quadratic constraints, or determining global topology are different tasks. Beautiful supplies no criterion distinguishing them, so the residual question has no single precise solved/open status.}
\GSfield{Reconstruction}{The established count is separated from the undefined requested parametrization; a research-problem ID is not manufactured.}
\GSfield{Fate}{The parameter count and full construction are on p.~186 of \citet{atiyahhitchindrinfeldmanin1978instantons}; p.~187 expressly distinguishes construction from extracting moduli geometry and topology. This does not establish an exact later solution to the notebook's unspecified aesthetic request.}
\GSfield{Present status}{As of 10 November 2026: the surviving wording does not determine a theorem with all hypotheses.}
\GSfield{Assessment}{literal statement: not determinable; publication match: partial; verification: read the result; development: annotated; confidence: medium}
\end{GSitem}
\begin{GSitem}{1977-451.13}{F}{Local gauge transformations and an unfinished confinement argument}
\GSfield{Kind; source}{remark; \GSpage{P000812}; A. A. Belavin}
\GSfield{The notes}{A heading promises exactly solvable two-dimensional field-theory questions and a naive Wilson argument about non-escape. It writes \(\delta S=0\) and \(\psi(x)\mapsto e^{i\alpha(x)}\psi(x)\).}
\GSfield{Statement}{Stationarity of an action gives its Euler--Lagrange equations. In a gauge theory a local phase transformation of matter is accompanied by a transformation of the connection so that covariant derivatives transform covariantly. These two lines alone imply neither confinement nor a Wilson-loop area law; no action, gauge dynamics or limiting regime is recorded.}
\GSfield{Fate}{The fragment is thematically related to \citet{wilson1974confinement}, but supplies no proposition matching that paper's lattice strong-coupling confinement argument. Earlier GP-1975-10 concerns a definite confinement problem; this fragment is only a thematic link and does not justify reusing its problem ID.}
\GSfield{Present status}{As of 10 November 2026: the surviving wording does not determine a theorem with all hypotheses.}
\GSfield{Assessment}{literal statement: not determinable; publication match: thematic only; verification: publisher metadata checked; development: annotated; confidence: low}
\end{GSitem}

\GStalk{457}{1978-03-05}{D. B. Fuchs; adjacent A. A. Beilinson and V. I. Ogievetsky}{P000823-P000831}{Graded Lie algebras with adjacent ambitwistor and supersymmetry notes}
The catalogue date 1978-03-05 is retained only as the envelope identifier. P000824 actually reads 5 March1979. P000823 belongs to a preceding Beilinson intervention on the flag/ambitwistor transform; P000824--P000829 are Fuchs on Cartan prolongations and Kac--Moody constructions; P000830--P000831 begin Ogievetsky's 19 March1979 lecture, headed Supergravity but recording only the super-Poincare algebra and superspace kinematics. The three speakers and dates must not be collapsed into a Fuchs lecture of 1978. The principal mathematical omissions are simplicity/irreducibility in classification, complexification in the rank-one test, transposed indices in the Serre relation, and the central quotient in loop realizations.

\begin{GSitem}{1977-457.1}{E}{The complete-flag double fibration}
\GSfield{Kind; source}{construction; \GSpage{P000823}; A. A. Beilinson}
\GSfield{The notes}{Beilinson considers full flags in four-space, projecting to \(G(4,2)\) by forgetting the line and hyperplane, and to \(Q\subset\mathbb P^3\times\mathbb P^{3*}\) by forgetting the plane.}
\GSfield{Statement}{Let \(F=\{\ell\subset P\subset H\subset\mathbb C^4:\dim(\ell,P,H)=(1,2,3)\}\), \(p:F\to\operatorname{Gr}(2,4)\), and \(q:F\to Q=\{(\ell,H):\ell\subset H\}\). A q-fiber is \(\mathbb P(H/\ell)\cong\mathbb P^1\); a p-fiber is \(\mathbb P(P)\times\mathbb P((\mathbb C^4/P)^*)\cong\mathbb P^1\times\mathbb P^1\). These are the compactified null-ray and null-direction correspondences.}
\GSfield{Reconstruction}{The subspace dimensions and both fibers are supplied explicitly; these are complex projective spaces.}
\GSfield{Fate}{The corresponding local ambitwistor construction is reviewed in \citet[Section~2]{masonskinner2006ambitwistor}, with explicit attribution to 1978 work of Witten and Isenberg--Yasskin--Green. The flag-space description itself is directly recoverable from the supplied diagram.}
\GSfield{Present status}{As of 10 November 2026: established in the stated scope; the notebook is exposition of this result.}
\GSfield{Assessment}{literal statement: correct; publication match: partial; verification: read the result; development: annotated; confidence: high}
\end{GSitem}
\begin{GSitem}{1977-457.2}{E}{Parallel sections along null curves and descent}
\GSfield{Kind; source}{construction; \GSpage{P000823}; A. A. Beilinson}
\GSfield{The notes}{A bundle with connection on the Grassmannian is canonically trivial along each CP1 because curvature vanishes in dimension 1 and the fundamental group is zero. Pull back to F, descend to Q, and obtain a holomorphic bundle trivial on CP1 times CP1.}
\GSfield{Statement}{For a holomorphic connection, curvature is a holomorphic2-form and its restriction to a complex curve vanishes. On a simply connected curve, parallel sections identify the bundle with its space of horizontal sections tensored with its structure sheaf; there is no distinguished basis of that space. On an admissible complex open region, the spaces of sections along null curves form the descended bundle, provided the fibers are connected and simply connected and the analytic descent hypotheses hold. A smooth real connection does not have zero curvature merely because the curve has complex dimension 1.}
\GSfield{Reconstruction}{Holomorphicity, the local domain and the meaning of canonical are essential additions. The notebook does not specify a nontrivial global holomorphic connection on the entire compact Grassmannian.}
\GSfield{Fate}{This is the parallel-section construction of \citet[Section~2, Theorem~1 and footnote~1]{masonskinner2006ambitwistor}. The2005 preprint/2006 publication restates the 1978 correspondence; it is not evidence of a new1979 theorem by the adjacent speaker.}
\GSfield{Present status}{As of 10 November 2026: established in the stated scope; the notebook is exposition of this result.}
\GSfield{Assessment}{literal statement: correct with missing hypotheses; publication match: partial; verification: read the result; development: annotated; confidence: high}
\end{GSitem}
\begin{GSitem}{1977-457.3}{E}{The reverse ambitwistor correspondence and missing field equations}
\GSfield{Kind; source}{claim; \GSpage{P000823}; A. A. Beilinson}
\GSfield{The notes}{The final arrow is \((E,\nabla)\longleftrightarrow\widetilde E\), with \(\widetilde E\) trivial on each CP1 times CP1; no inverse construction or field equation is written.}
\GSfield{Statement}{In a suitable local complex domain, a holomorphic ambitwistor bundle trivial on the quadrics of null directions reconstructs a space-time connection by evaluation and parallel transport. Triviality on these quadrics is not by itself the Yang--Mills equation. In the classical extension formulation the latter is a condition that the bundle extend to the third formal neighbourhood of the ambitwistor quadric.}
\GSfield{Reconstruction}{Local analytic hypotheses are supplied. No global equivalence or Yang--Mills claim is inferred from the two-headed arrow alone.}
\GSfield{Fate}{The inverse construction and the distinct field-equation condition are separated in \citet[Theorems~1--2]{masonskinner2006ambitwistor}. The notebook records only the former schematic correspondence; the latter is explanatory context, not an additional notebook theorem.}
\GSfield{Present status}{As of 10 November 2026: established in the stated scope; the notebook is exposition of this result.}
\GSfield{Assessment}{literal statement: correct with missing hypotheses; publication match: partial; verification: read the result; development: annotated; confidence: high}
\end{GSitem}
\begin{GSitem}{1977-457.4}{E}{Grading and transitivity in the adjoint action}
\GSfield{Kind; source}{definition; \GSpage{P000824}}
\GSfield{The notes}{A graded Lie algebra has \([\mathfrak g_i,\mathfrak g_j]\subset\mathfrak g_{i+j}\). The bracket gives maps to \(\operatorname{Hom}(\mathfrak g_{-1},\mathfrak g_{i-1})\) and \(\operatorname{Hom}(\mathfrak g_1,\mathfrak g_{i+1})\), assumed injective in opposite degree ranges.}
\GSfield{Statement}{The degree-zero part is a Lie algebra and each homogeneous part is its module. The stated injectivity means that a nonnegative-degree element centralizing \(\mathfrak g_{-1}\), or a nonpositive-degree element centralizing \(\mathfrak g_1\), is zero. For the classification discussed later one also imposes finite-dimensional homogeneous pieces, appropriate generation and irreducibility, and a characteristic-zero algebraically closed ground field.}
\GSfield{Reconstruction}{The additional classification assumptions are separated from the elementary definition.}
\GSfield{Fate}{The graded, irreducible and graded-simple conventions are given in \citet[Introduction and Chapter~I]{kac1968graded}, submitted 13 February1968. Its simple explicitly means absence of homogeneous ideals, which matters below.}
\GSfield{Present status}{As of 10 November 2026: established in the stated scope; the notebook is exposition of this result.}
\GSfield{Assessment}{literal statement: correct; publication match: exact; verification: read the result; development: annotated; confidence: high}
\end{GSitem}
\begin{GSitem}{1977-457.5}{E}{Four Cartan families and the missing classification hypotheses}
\GSfield{Kind; source}{claim; \GSpage{P000824}}
\GSfield{The notes}{Bounded-below grading and the stated transitivity are followed by a claim that there are four series: all polynomial vector fields, divergence-free, Hamiltonian, and contact fields.}
\GSfield{Statement}{The four simple Cartan-type families over \(\mathbb C\) are W, S, H and K, interpreted in the appropriate polynomial graded or completed formal category. The broader transitive irreducible classification also includes conformal extensions before imposing simplicity. Boundedness and transitivity alone do not imply the four-family conclusion. The contact grading has depth 2 when the contact coordinate has weight 2 and the other coordinates weight 1.}
\GSfield{Reconstruction}{Infinite dimension, irreducibility, simplicity and completion conventions are restored rather than attributed to words absent from the notebook.}
\GSfield{Fate}{Compare \citet[Chapter~III, Section~1; Theorems~4--5]{kac1968graded}. The four simple types in Theorem 4 and the six Cartan types in Theorem 5 explain why simplicity cannot be omitted.}
\GSfield{Present status}{As of 10 November 2026: established in the stated scope; the notebook is exposition of this result.}
\GSfield{Assessment}{literal statement: false as written; publication match: partial; verification: read the result; development: annotated; confidence: high}
\end{GSitem}
\begin{GSitem}{1977-457.6}{E}{The maximal Cartan prolongation and Jacobi symmetry}
\GSfield{Kind; source}{construction; \GSpage{P000824-P000825}}
\GSfield{The notes}{For depth 1, the notes say the whole algebra is recovered from \(\mathfrak g_{-1},\mathfrak g_0\), and define successive symmetric maps by the Jacobi identity.}
\GSfield{Statement}{Let V be finite-dimensional and \(\mathfrak h\subset\mathfrak{gl}(V)\). Its maximal prolongation has \(\mathfrak g_{-1}=V\), \(\mathfrak g_0=\mathfrak h\), and \[\mathfrak g_k=(S^{k+1}V^*\otimes V)\cap(S^kV^*\otimes\mathfrak h),\quad k\ge1.\] The symmetry of \(f(v_1,\ldots,v_k)v_{k+1}\) follows from the Jacobi identity and \([V,V]=0\). Every transitive algebra with these lower pieces embeds in the prolongation; equality is a maximality condition, not automatic.}
\GSfield{Reconstruction}{The notebook changes an inclusion into an equality when choosing the maximal prolongation; that choice is made explicit.}
\GSfield{Fate}{The construction is spelled out in \citet[Section~2.2]{ottazziwarhurst2011contact}, which treats it as the classical Singer--Sternberg prolongation; its earlier history includes \citet{guilleminquillensternberg1967infinite}.}
\GSfield{Present status}{As of 10 November 2026: established in the stated scope; the notebook is exposition of this result.}
\GSfield{Assessment}{literal statement: false as written; publication match: exact; verification: read the result; development: annotated; confidence: high}
\end{GSitem}
\begin{GSitem}{1977-457.7}{E}{The inductively defined bracket on a prolongation}
\GSfield{Kind; source}{construction; \GSpage{P000825}}
\GSfield{The notes}{To define \([\alpha,\beta]\), specify its bracket with every \(\delta\in\mathfrak g_{-1}\) and use Jacobi to reduce to previously defined degrees. Simplicity is an extra condition; most examples are finite-dimensional.}
\GSfield{Statement}{For homogeneous prolongation elements, define contraction of the bracket by \([[\alpha,\beta],v]=[\alpha,[\beta,v]]-[\beta,[\alpha,v]]\). Transitivity gives uniqueness, and the symmetry conditions yield existence and Jacobi. Equivalently use the commutator of their polynomial-vector-field realizations. The phrase most finite-dimensional is heuristic, not a measure-theoretic genericity theorem.}
\GSfield{Reconstruction}{The finite-dimensional faithful input and the equivalent polynomial realization justify the inductive recipe.}
\GSfield{Fate}{The recursive bracket and realization are part of the prolongation framework presented in \citet[Section~2]{ottazziwarhurst2011contact}.}
\GSfield{Present status}{As of 10 November 2026: established in the stated scope; the notebook is exposition of this result.}
\GSfield{Assessment}{literal statement: correct with missing hypotheses; publication match: partial; verification: read the result; development: annotated; confidence: high}
\end{GSitem}
\begin{GSitem}{1977-457.8}{E}{The polynomial vector-field realization}
\GSfield{Kind; source}{construction; \GSpage{P000826}}
\GSfield{The notes}{Take polynomial fields on \(\mathfrak g_{-1}\) whose linear Taylor part at every point belongs to \(\mathfrak g_0\); this is the algebra constructed by prolongation.}
\GSfield{Statement}{For \(\mathfrak h\subset\mathfrak{gl}(V)\), the polynomial fields X with \(DX(v)\in\mathfrak h\) for all v form the maximal prolongation. Comparing homogeneous coefficients gives the symmetric-tensor intersection in the preceding item. Closure under brackets follows by differentiating \([X,Y]=DY\,X-DX\,Y\) and using that derivatives of h-valued polynomials remain h-valued and h is closed under commutators.}
\GSfield{Reconstruction}{The Jacobian convention is fixed; the coefficient proof is an explanatory verification of the old construction.}
\GSfield{Fate}{This is the flat-structure interpretation explained before \citet[Theorem~2.3]{ottazziwarhurst2011contact}; the notebook gives the same coefficient criterion.}
\GSfield{Present status}{As of 10 November 2026: established in the stated scope; the notebook is exposition of this result.}
\GSfield{Assessment}{literal statement: correct; publication match: exact; verification: read the result; development: annotated; confidence: high}
\end{GSitem}
\begin{GSitem}{1977-457.9}{E}{The complex rank-one criterion for infinite type}
\GSfield{Kind; source}{theorem; \GSpage{P000826}}
\GSfield{The notes}{Wilson's lemma: the Cartan prolongation is infinite-dimensional iff some \(\xi\in\mathfrak g_0\) has rank 1; the speaker refers to Guillemin and two coauthors. Reductivity and irreducibility are separately mentioned.}
\GSfield{Statement}{For finite-dimensional complex V and \(\mathfrak h\subset\mathfrak{gl}(V)\), its prolongation has nonzero terms in arbitrarily high degree iff h contains a nonzero rank-one endomorphism. For real V apply this test to the complexification. No irreducibility or reductivity is required for this finite-type criterion; such conditions belong to classification refinements. A rank-one element \(v\otimes\xi\) directly produces \(v\otimes\xi^{k+1}\) in every prolongation degree.}
\GSfield{Reconstruction}{The complex ground field and finite-dimensional input are indispensable; the scan does not settle the historical attribution.}
\GSfield{Fate}{The theorem and proof are read in \citet[Theorem~2.3, Sections~2.3--2.4]{ottazziwarhurst2011contact}, which identifies \citet{guilleminquillensternberg1967infinite} as the first published proof. The notebook's Wilson attribution may concern a later proof; it was not independently resolved.}
\GSfield{Present status}{As of 10 November 2026: established in the stated scope; the notebook is exposition of this result.}
\GSfield{Assessment}{literal statement: correct with missing hypotheses; publication match: exact; verification: read the result; development: annotated; confidence: high}
\end{GSitem}
\begin{GSitem}{1977-457.10}{F}{An infinite-dimensional input and a Poisson-algebra suggestion}
\GSfield{Kind; source}{remark; \GSpage{P000826}}
\GSfield{The notes}{An example takes \(\mathfrak g_0=W_n\), \(\mathfrak g_{-1}=\mathbb C[x_1,\ldots,x_n]\), and says the prolongation is scarcely different from Hamiltonian fields, with a Poisson bracket.}
\GSfield{Statement}{Polynomial vector fields act on the polynomial algebra by derivations, so an algebraic prolongation recipe can be posed with this infinite-dimensional input. It lies outside the finite-dimensional rank-one criterion just stated. To identify its completion with a Poisson algebra requires the choice of algebraic versus continuous multilinear maps, grading, and the precise Poisson realization; none is supplied.}
\GSfield{Fate}{The sources checked for ordinary finite-dimensional prolongations do not fix these missing choices. The earlier Hamiltonian-vector-field cohomology problem GP-1975-21 is only thematic background, not the same question.}
\GSfield{Present status}{As of 10 November 2026: the surviving wording does not determine a theorem with all hypotheses.}
\GSfield{Assessment}{literal statement: not determinable; publication match: none found; verification: unresolved lead; development: annotated; confidence: low}
\end{GSitem}
\begin{GSitem}{1977-457.11}{E}{Generalized Cartan matrices and symmetrizability}
\GSfield{Kind; source}{definition; \GSpage{P000826-P000827}}
\GSfield{The notes}{A has diagonal2, nonpositive integer off-diagonal entries with simultaneous zeros, and can be made symmetric by \(n_{\alpha\beta}\mapsto n_{\alpha\beta}\lambda_\alpha/\lambda_\beta\). An invariant form is then mentioned.}
\GSfield{Statement}{A generalized Cartan matrix has these integrality and sign conditions. It is symmetrizable if a positive diagonal D makes DA symmetric; equivalently a positive diagonal similarity makes A symmetric, with D having entries \(\lambda_\alpha^2\) in the notebook's convention. The associated symmetrizable Kac--Moody algebra admits the standard invariant symmetric form, with the usual qualification on its center or full Cartan realization.}
\GSfield{Reconstruction}{Positive real scaling factors and the distinction between similarity and left symmetrization are stated explicitly.}
\GSfield{Fate}{These matrices, invariant forms and contragredient constructions are treated in \citet[Chapter~II, Sections~1--4]{kac1968graded}.}
\GSfield{Present status}{As of 10 November 2026: established in the stated scope; the notebook is exposition of this result.}
\GSfield{Assessment}{literal statement: correct with missing hypotheses; publication match: partial; verification: read the result; development: annotated; confidence: high}
\end{GSitem}
\begin{GSitem}{1977-457.12}{E}{Chevalley-Serre generators and an index mismatch}
\GSfield{Kind; source}{construction; \GSpage{P000827}}
\GSfield{The notes}{The notes use \([h_\beta,e_\alpha]=n_{\alpha\beta}e_\alpha\), \([e_\alpha,f_\alpha]=2h_\alpha\), and \((\operatorname{ad}e_\alpha)^{1-n_{\alpha\beta}}e_\beta=0\), with analogous f-relations.}
\GSfield{Statement}{With the recorded transposed convention for the Cartan action, the Serre exponent is \(1-n_{\beta\alpha}\), for \(\alpha\ne\beta\). The factor 2 in \([e_\alpha,f_\alpha]=2h_\alpha\) is a harmless rescaling of one root generator in characteristic zero. For singular A, the presentation by only these 3N generators describes the derived Kac--Moody algebra; the full standard algebra has an enlarged Cartan subalgebra.}
\GSfield{Reconstruction}{The scan confirms both the factor 2 and the inconsistent index order; neither is silently re-transcribed.}
\GSfield{Fate}{The contragredient presentation and its reduced quotient are part of \citet[Chapter~II]{kac1968graded}. The index repair follows by comparing the root-string exponent with the actual action of \(h_\alpha\) on \(e_\beta\), and matters for the nonsymmetric rank-two example on the next page.}
\GSfield{Present status}{As of 10 November 2026: established in the stated scope; the notebook is exposition of this result.}
\GSfield{Assessment}{literal statement: false as written; publication match: partial; verification: read the result; development: annotated; confidence: high}
\end{GSitem}
\begin{GSitem}{1977-457.13}{E}{Root-lattice and height gradings}
\GSfield{Kind; source}{definition; \GSpage{P000827-P000828}}
\GSfield{The notes}{The fine grading has degrees \(-\alpha,0,\alpha\) on \(f_\alpha,h_\alpha,e_\alpha\), and the coarse grading degrees \(-1,0,1\). Nonzero degrees with a nonzero component are roots; nonnegative simple-root combinations are called positive.}
\GSfield{Statement}{The homogeneous defining relations give a grading by the free root lattice \(Q=\bigoplus_i\mathbb Z\alpha_i\); composing with height gives the integer grading. A root is a nonzero \(\gamma\in Q\) with nonzero root space, and a positive root lies in \(Q_+\setminus\{0\}\). Not every nonnegative integral combination is a root.}
\GSfield{Reconstruction}{The zero combination is excluded and membership in the actual root support is retained.}
\GSfield{Fate}{These conventions are developed in \citet[Chapter~II]{kac1968graded}.}
\GSfield{Present status}{As of 10 November 2026: established in the stated scope; the notebook is exposition of this result.}
\GSfield{Assessment}{literal statement: correct with missing hypotheses; publication match: exact; verification: read the result; development: annotated; confidence: high}
\end{GSitem}
\begin{GSitem}{1977-457.14}{E}{Finite, affine, and indefinite growth}
\GSfield{Kind; source}{claim; \GSpage{P000827}}
\GSfield{The notes}{Positive definite A gives a finite-dimensional algebra; semidefinite singular A gives growth1 with uniformly bounded graded dimensions; a negative eigenvalue is said to give exponential growth. The type-A chain matrix gives sl(n+1).}
\GSfield{Statement}{For an indecomposable symmetrizable generalized Cartan matrix, positivity of its symmetric form characterizes finite type, and positive semidefiniteness with a null direction characterizes affine type; the latter has bounded height-graded dimensions. The indefinite case has infinite growth in Kac's finite-growth classification. The stronger quantitative exponential statement is not inferred merely from the phrase infinite growth in that source. The displayed n-by-n tridiagonal matrix is type \(A_n\), giving \(\mathfrak{sl}_{n+1}\).}
\GSfield{Reconstruction}{The symmetric form, indecomposability and meaning of growth are supplied; no false equivalence between superpolynomial and exponential growth is made.}
\GSfield{Fate}{The finite-growth alternatives are read in \citet[Chapter~II, Theorem~1 and Section~7]{kac1968graded}. A theorem-level check of the notebook's stronger exponential wording was not completed; the fate match is therefore partial.}
\GSfield{Present status}{As of 10 November 2026: established in the stated scope; the notebook is exposition of this result.}
\GSfield{Assessment}{literal statement: correct with missing hypotheses; publication match: partial; verification: read the result; development: annotated; confidence: medium}
\end{GSitem}
\begin{GSitem}{1977-457.15}{E}{The center from the kernel of the Cartan matrix}
\GSfield{Kind; source}{construction; \GSpage{P000828}}
\GSfield{The notes}{If \(\sum_\beta k_\beta n_{\alpha\beta}=0\), then \(\sum_\beta k_\beta h_\beta\) is central. Two rank-two affine matrices are displayed.}
\GSfield{Statement}{In the notebook's index convention, the displayed kernel condition makes the Cartan element commute with every e,f,h generator, hence central. The matrices \(\left(\begin{smallmatrix}2&-2\\-2&2\end{smallmatrix}\right)\) and \(\left(\begin{smallmatrix}2&-1\\-4&2\end{smallmatrix}\right)\) are the affine types \(A_1^{(1)}\) and \(A_2^{(2)}\), up to transposition. Their derived algebras have a one-dimensional center.}
\GSfield{Reconstruction}{The second matrix is not symmetric; symmetric counterpart in the translation means its transpose here.}
\GSfield{Fate}{These affine rank-two types occur in the tables and realizations of \citet[Chapter~II, Sections~5--7]{kac1968graded}. The centrality check is also immediate from the displayed relations.}
\GSfield{Present status}{As of 10 November 2026: established in the stated scope; the notebook is exposition of this result.}
\GSfield{Assessment}{literal statement: correct; publication match: exact; verification: read the result; development: annotated; confidence: high}
\end{GSitem}
\begin{GSitem}{1977-457.16}{E}{Untwisted loop algebras and the missing central quotient}
\GSfield{Kind; source}{construction; \GSpage{P000828}}
\GSfield{The notes}{The first affine rank-two matrix is said to give \(\mathfrak{sl}_2\otimes\mathbb C[t,t^{-1}]\); more generally use any finite-dimensional g, with the extra generators obtained from the longest root vectors times \(t^{\pm1}\).}
\GSfield{Statement}{For finite-dimensional simple complex g, the center quotient of the derived untwisted affine algebra is its loop algebra \(g\otimes\mathbb C[t,t^{-1}]\). The derived affine algebra is its central extension; the full affine algebra adds the degree derivation. The extra simple-root generators use the highest root and its negative, rather than an unspecified longest root. Passing to the loop quotient kills the central linear combination of h's recorded immediately above.}
\GSfield{Reconstruction}{The finite-dimensional input must be simple for the single affine type assertion; semisimple inputs give direct sums.}
\GSfield{Fate}{The distinction between contragredient and reduced algebras, and their loop realizations, is explicit in \citet[Chapter~II, Section~7]{kac1968graded}.}
\GSfield{Present status}{As of 10 November 2026: established in the stated scope; the notebook is exposition of this result.}
\GSfield{Assessment}{literal statement: false as written; publication match: partial; verification: read the result; development: annotated; confidence: high}
\end{GSitem}
\begin{GSitem}{1977-457.17}{E}{Twisted loop algebras from an involution}
\GSfield{Kind; source}{construction; \GSpage{P000829}}
\GSfield{The notes}{For \(\tau^2=1\), decompose \(\mathfrak g=\mathfrak g_+\oplus\mathfrak g_-\) and take even powers with \(\mathfrak g_+\), odd with \(\mathfrak g_-\). The example is \(\mathfrak{sl}_3\), \(\tau(X)=-X^T\).}
\GSfield{Statement}{The parity eigenspace rule defines the fixed algebra\[L(\mathfrak g,\tau)=\bigoplus_{m\in\mathbb Z}\mathfrak g_{(-1)^m}t^m.\]It is closed under brackets because eigenvalues multiply. For the nontrivial diagram involution of \(\mathfrak{sl}_3\), the corresponding central extension is twisted affine type \(A_2^{(2)}\). Again the displayed loop algebra itself is the center quotient.}
\GSfield{Reconstruction}{An arbitrary involution gives a twisted-loop presentation, but its affine-type identification depends on its outer/diagram class.}
\GSfield{Fate}{This is among the twisted realizations in \citet[Chapter~II, Section~7]{kac1968graded}.}
\GSfield{Present status}{As of 10 November 2026: established in the stated scope; the notebook is exposition of this result.}
\GSfield{Assessment}{literal statement: correct with missing hypotheses; publication match: exact; verification: read the result; development: annotated; confidence: high}
\end{GSitem}
\begin{GSitem}{1977-457.18}{E}{A one-dimensional graded bracket with period three}
\GSfield{Kind; source}{construction; \GSpage{P000829}}
\GSfield{The notes}{There are generators \(\varepsilon_i\), \(i\in\mathbb Z\), and \([\varepsilon_i,\varepsilon_j]=a_{ij}\varepsilon_{i+j}\), where \(a_{ij}\in\{-1,0,1\}\) is congruent to \(j-i\pmod3\). It is called simple.}
\GSfield{Statement}{This bracket is the principal grading of the centerless loop algebra of \(\mathfrak{sl}_2\). For \([h,e]=2e,[h,f]=-2f,[e,f]=h\), one realization is \(\varepsilon_{3r}=\tfrac12ht^r\), \(\varepsilon_{3r+1}=et^r\), \(\varepsilon_{3r-1}=\tfrac12ft^r\). It is graded-simple in the intended sense; it is not simple as an ungraded Lie algebra, since \((t-1)\mathfrak{sl}_2[t,t^{-1}]\) is a nonzero proper ideal.}
\GSfield{Reconstruction}{The normalization of e,f,h and the grading are supplied. The counterexample ideal is not homogeneous in this grading.}
\GSfield{Fate}{The distinction is already explicit in the definition of simple in \citet[Introduction]{kac1968graded}. The given basis realization verifies all nine residue cases directly; it identifies an old loop algebra, not a new classification result.}
\GSfield{Present status}{As of 10 November 2026: established in the stated scope; the notebook is exposition of this result.}
\GSfield{Assessment}{literal statement: false as written; publication match: partial; verification: read the result; development: annotated; confidence: high}
\end{GSitem}
\begin{GSitem}{1977-457.19}{F}{The incomplete three-algebra and exhaustion statements}
\GSfield{Kind; source}{claim; \GSpage{P000829}}
\GSfield{The notes}{The notes say there are three simple algebras with one generator in each integer degree, give only one bracket, and end: these two constructions exhaust everything up to grading.}
\GSfield{Statement}{No classification of all integer-graded Lie algebras follows without growth, irreducibility, generation and graded-simplicity assumptions. Kac's finite-growth theorem has a precise scope and produces reduced contragredient and Cartan families. The two missing one-dimensional brackets and the antecedent of these two constructions cannot be uniquely recovered from this page.}
\GSfield{Reconstruction}{The established finite-growth classification is separated from the incomplete last announcement.}
\GSfield{Fate}{The scoped theorem is \citet[Theorem~4]{kac1968graded}; its hypotheses are substantially more specific than the notebook's final sentence. The broader one-dimensional-component classification was not identified as an exact theorem in the supplied wording.}
\GSfield{Present status}{As of 10 November 2026: the surviving wording does not determine a theorem with all hypotheses.}
\GSfield{Assessment}{literal statement: not determinable; publication match: partial; verification: read the result; development: annotated; confidence: low}
\end{GSitem}
\begin{GSitem}{1977-457.20}{E}{Hermitian matrices, spinors and the Lorentz cover}
\GSfield{Kind; source}{construction; \GSpage{P000830}; V. I. Ogievetsky}
\GSfield{The notes}{A real four-vector is the Hermitian matrix \(X=x^0I+x^i\sigma_i\); its determinant is the Lorentz metric, and \(X\mapsto A^*XA\) gives the Lorentz action. Two Weyl spinor representations are used.}
\GSfield{Statement}{The determinant is \((x^0)^2-(x^1)^2-(x^2)^2-(x^3)^2\). The action of \(SL(2,\mathbb C)\) on Hermitian2-by-2 matrices preserves it and induces the double cover of \(SO^+(1,3)\), with kernel \(\{I,-I\}\). Dotted and undotted Weyl indices transform in conjugate representations. The inverse vector-to-matrix map uses the trace pairing and a factor 1/2, rather than an unqualified repeated Pauli-matrix contraction.}
\GSfield{Reconstruction}{The proper orthochronous component, kernel, conjugation convention and inverse normalization are made explicit.}
\GSfield{Fate}{These conventions underlie \citet[Appendix]{haaglopuszanskisohnius1975supersymmetries}. The determinant identity is also immediate from the matrix printed in the notebook.}
\GSfield{Present status}{As of 10 November 2026: established in the stated scope; the notebook is exposition of this result.}
\GSfield{Assessment}{literal statement: correct with missing hypotheses; publication match: partial; verification: read the result; development: annotated; confidence: high}
\end{GSitem}
\begin{GSitem}{1977-457.21}{E}{The real N=1 super-Poincare algebra}
\GSfield{Kind; source}{definition; \GSpage{P000830}; V. I. Ogievetsky}
\GSfield{The notes}{Ten even generators P,L and four real spinorial odd generators satisfy \(\{Q_\alpha,\bar Q_{\dot\beta}\}=2\sigma^m_{\alpha\dot\beta}P_m\), \(\{Q,Q\}=\{\bar Q,\bar Q\}=0\), and \([P,Q]=0\); dimension is written(10,4).}
\GSfield{Statement}{Add the usual Poincare brackets and the Lorentz spinor action on Q and its conjugate. These relations define the four-dimensional real N=1 super-Poincare algebra of superdimension \(10\mid4\). The odd bracket is symmetric and must satisfy super-Jacobi; it is not an ordinary Lie bracket. Its kinematic algebra by itself is not a construction of a supergravity action or field equations.}
\GSfield{Reconstruction}{Reality, conjugate indices, the even-odd brackets and the super-Jacobi condition complete the recorded relations.}
\GSfield{Fate}{This algebra and its allowed extensions are within \citet[Section~1, equations(1.1)--(1.7)]{haaglopuszanskisohnius1975supersymmetries}, received 15 November1974 and published 1975. Earlier occurrence 1975-389.1--.7 is background on supergeometry, not the same indexed open problem.}
\GSfield{Present status}{As of 10 November 2026: established in the stated scope; the notebook is exposition of this result.}
\GSfield{Assessment}{literal statement: correct with missing hypotheses; publication match: exact; verification: read the result; development: annotated; confidence: high}
\end{GSitem}
\begin{GSitem}{1977-457.22}{E}{Extended supersymmetry and the role of internal indices}
\GSfield{Kind; source}{construction; \GSpage{P000830-P000831}; V. I. Ogievetsky}
\GSfield{The notes}{With several spinors, \(\{Q_\alpha^i,\bar Q_{\dot\beta}^j\}=2\sigma^m_{\alpha\dot\beta}P_m\delta^{ij}\); i,j are internal-symmetry indices. The odd bracket is said to be unique.}
\GSfield{Statement}{With a positive Hermitian internal metric put into standard form, this is the translation-producing odd bracket for N supersymmetries. Lorentz covariance determines its form up to normalization in the N=1 case. For N greater than 1, super-Jacobi permits antisymmetric central charges in \(\{Q_\alpha^i,Q_\beta^j\}=\varepsilon_{\alpha\beta}Z^{ij}\). Thus uniqueness requires a fixed even algebra and the explicit exclusion of central extensions.}
\GSfield{Reconstruction}{The internal metric, N and the central-extension qualification are supplied; a colloquial particle/antiparticle slogan is not treated as an equation of states.}
\GSfield{Fate}{Central charges and restrictions from internal symmetry are proved in \citet[Sections~1 and~3]{haaglopuszanskisohnius1975supersymmetries}. The notebook records the zero-central-charge specialization.}
\GSfield{Present status}{As of 10 November 2026: established in the stated scope; the notebook is exposition of this result.}
\GSfield{Assessment}{literal statement: correct with missing hypotheses; publication match: exact; verification: read the result; development: annotated; confidence: high}
\end{GSitem}
\begin{GSitem}{1977-457.23}{E}{Mixing space-time and internal symmetries and the no-go scope}
\GSfield{Kind; source}{claim; \GSpage{P000831}; V. I. Ogievetsky}
\GSfield{The notes}{The notes say this is a nontrivial mixing impossible in ordinary Lie algebras, by a theorem.}
\GSfield{Statement}{The relevant no-go conclusion concerns continuous symmetries of a nontrivial relativistic S-matrix under specified assumptions, including suitable particle finiteness and regularity; it is not a theorem forbidding arbitrary semidirect products of ordinary Lie algebras. Allowing odd spinorial generators gives the supersymmetry extension. Massive and massless cases have different allowed enlargements.}
\GSfield{Reconstruction}{The statement is restricted to the physical symmetry setting of the theorem; no universal algebraic prohibition is endorsed.}
\GSfield{Fate}{The hypotheses and the extension of the Coleman--Mandula framework are explicit in \citet[Sections~1--2 and~5]{haaglopuszanskisohnius1975supersymmetries}.}
\GSfield{Present status}{As of 10 November 2026: established in the stated scope; the notebook is exposition of this result.}
\GSfield{Assessment}{literal statement: false as written; publication match: partial; verification: read the result; development: annotated; confidence: high}
\end{GSitem}
\begin{GSitem}{1977-457.24}{E}{Superspace and the finite superfield expansion}
\GSfield{Kind; source}{construction; \GSpage{P000831}; V. I. Ogievetsky}
\GSfield{The notes}{Adjoin \(\theta^\alpha,\bar\theta^{\dot\alpha}\), four anticommuting variables. A superfield expands through \(\theta^2\bar\theta^2\), and \(\varepsilon_{12}=-\varepsilon_{21}=1\) is the spinor metric.}
\GSfield{Statement}{The local superfunction algebra is the ordinary space-time function algebra tensored with an exterior algebra on four odd generators. Its degrees have dimensions \(1,4,6,4,1\), so a general scalar superfield has 16 coefficient functions, eight even and eight odd monomials. For an even superfield, spinor coefficient fields have odd parity. The mixed quadratic term is \(\theta\sigma^m\bar\theta\,v_m(x)\), with a vector index; the antisymmetric epsilon tensor raises and lowers spinor indices.}
\GSfield{Reconstruction}{The ellipsis includes the cubic spinor terms; dotted conjugate indices and the vector index repair the compressed expression.}
\GSfield{Fate}{The superspace/superfield construction predates this 1979 meeting: \citet{salamstrathdee1975superfields}, submitted 12 June1974 and published 15 March1975. Publisher abstract and metadata were checked; the precise component notation here is reconstructed from the exterior algebra, not read from that full paper.}
\GSfield{Present status}{As of 10 November 2026: established in the stated scope; the notebook is exposition of this result.}
\GSfield{Assessment}{literal statement: correct with missing hypotheses; publication match: partial; verification: publisher metadata checked; development: annotated; confidence: high}
\end{GSitem}
\begin{GSitem}{1977-457.25}{E}{Odd differential operators realizing supertranslations}
\GSfield{Kind; source}{construction; \GSpage{P000831}; V. I. Ogievetsky}
\GSfield{The notes}{With \(P_m=-i\partial_m\), the notes write \(Q_\alpha=-i\partial_{\theta^\alpha}-(\sigma^m\bar\theta)_\alpha\partial_m\) and \(\bar Q_{\dot\alpha}=i\bar\partial_{\dot\alpha}+(\theta\sigma^m)_{\dot\alpha}\partial_m\), after an obscured cancellation.}
\GSfield{Statement}{For left Grassmann derivatives these recorded signs give \(\{Q_\alpha,\bar Q_{\dot\beta}\}=-2i\sigma^m_{\alpha\dot\beta}\partial_m=2\sigma^m_{\alpha\dot\beta}P_m\), and the same-chirality anticommutators vanish. The displayed orbital \(L_{mn}=-i(x_m\partial_n-x_n\partial_m)\) must be supplemented by its action on the spinor coordinates to realize the whole super-Poincare algebra on superfields.}
\GSfield{Reconstruction}{Left derivatives, the full Lorentz action and the remaining illegible mark are distinguished. The cancellation does not prevent checking the surviving supertranslation formula.}
\GSfield{Fate}{Differential superspace realizations are the subject of \citet{salamstrathdee1975superfields}; only its abstract/metadata were available. The notebook signs are independently checked algebraically here using the Grassmann derivative rule, without guessing the cancelled mark.}
\GSfield{Present status}{As of 10 November 2026: established in the stated scope; the notebook is exposition of this result.}
\GSfield{Assessment}{literal statement: correct with missing hypotheses; publication match: partial; verification: indirect; development: annotated; confidence: high}
\end{GSitem}

\section{Developments}\label{dev:1977}
\subsection{Integral information in the formal-diffeomorphism question}
\label{dev:1977-jets}

\paragraph{The question and its scope.}
The last line of \GSpage{P000773}, independently checked against the image, asks
\ru{посчитать когомологии группы} \(\operatorname{Diff}_1(\mathbb Z)\).
The preceding page defines the inverse limit of jet groups and the displayed
series \(t+a_1t^2+a_2t^3+\cdots\). Neither the coefficient module nor the
cochain theory is specified. The strongest literal formulation is therefore a
request for a computation, with these two parameters still to be supplied.
It is not an assertion that discrete group cohomology, continuous cohomology,
and Lie-algebra cohomology coincide.

For a bounded development we explicitly choose the integral counterpart of
the finite-jet interpretation. Put
\[
 J=\{t+a_1t^2+a_2t^3+\cdots:a_i\in\mathbb Z\},\qquad
 J_N=J/\{f:f(t)\equiv t\pmod {t^{N+1}}\}\quad(N\geq2).
\]
Multiplication is composition, \(J\) has its coefficientwise topology, and
\(J_N\) is discrete. For a trivial discrete coefficient group \(A\), let
\[
 C^r_{\mathrm{jet}}(J;A)
   =\bigcup_N\operatorname{Map}(J_N^r,A),\qquad
 H^r_{\mathrm{jet}}(J;A)=H^r(C^\bullet_{\mathrm{jet}}(J;A)).
\]
Inflation makes these unions literal subgroups of the cochains on \(J\).
Every cochain in this definition depends on a common finite jet of all its
arguments. This is an editorial choice supported by the preceding inverse-limit
discussion, not a coefficient convention recoverable from the question itself.

\paragraph{What the later literature establishes.}
\citet[Theorems~1.2 and~1.5; Proposition~4.4]{babenkobogatyi2008substitutions}
identify the closed-commutator quotient with
\(\mathbb Z^2\oplus(\mathbb Z/2)^2\), and identify finite-jet real cohomology
with the continuous cohomology of \(L_1\). Their manuscript was received on
24 November 2006 and published in 2008. Thus the real version has a direct
later solution. \citet[Introduction and Theorem~1.0.2 of the preprint]{pavezcornejo2025abelianization}
extends the topological abelianization calculation to deeper substitution
subgroups. That work first circulated on 19 July 2024 and appeared in 2025;
its commutator notation explicitly denotes a closure. Neither statement
supplies a computation of the unrestricted discrete integral cohomology ring.
As of 10 November 2026, the sources checked here do not establish that stronger
computation. This is a limit of the verified match, not a proof of its absence
from all literature.

The development below proves two elementary consequences for the chosen
integral cochain theory. They are supplied here for the companion, with no
claim of research priority. They also expose precisely why replacing the
group by its Lie algebra would lose information.

\begin{lemma}[Finite-jet cochains and coefficient sequences]\ours{}
Filtered inflation gives
\(H^r_{\mathrm{jet}}(J;A)=\varinjlim_N H^r(J_N;A)\).
A short exact sequence of discrete trivial coefficient groups induces a
short exact sequence of these cochain complexes and hence its usual long
exact cohomology sequence.
\end{lemma}
\begin{proof}
Any finite list of cochains is inflated from one sufficiently large \(J_N\).
Consequently a cocycle condition or a coboundary equation in the union holds
at a finite stage. This proves the asserted equality of cohomology and the
filtered colimit. If \(B\to C\) is surjective and a \(C\)-valued cochain
factors through \(J_N^r\), choose a lift of each of its values in \(B\).
The resulting lift still factors through \(J_N^r\). Exactness at the other
terms is pointwise. The connecting homomorphisms follow from the ordinary
cochain argument.
\end{proof}

Write the first coefficients as
\[
 f(t)=t+at^2+bt^3+ct^4+dt^5+et^6+ht^7+\cdots.
\]
For clarity \(h\) here is a coefficient, not another group element. The
explicit quotient in \citet[Proposition~2.1]{babenkobogatyi2008substitutions}
has coordinates
\[
 x=a,\qquad y=b-a^2,\qquad
 \varepsilon_1=c+d+\frac{b(b+1)}2\pmod2,\qquad
 \varepsilon_2=h+db+d\pmod2.                 \tag{J1}
\]
The formula involves integer-valued expressions and depends only on the
seventh jet. Its kernel is the closed commutator subgroup by
\citet[Proposition~2.9, pp.~249--250]{babenkobogatyi2008substitutions}.
The literature input in the next proposition is
exactly this theorem; no claim about the algebraic, unclosed commutator is
needed.

\begin{proposition}[First cohomology and all torsion in degree two]\ours{}
With the definitions and trivial coefficients above,
\[
 H^1_{\mathrm{jet}}(J;\mathbb Z)=\mathbb Zx\oplus\mathbb Zy,
 \qquad
 \operatorname{Tor}H^2_{\mathrm{jet}}(J;\mathbb Z)
       \cong(\mathbb Z/2)^2.                 \tag{J2}
\]
The two torsion generators are represented by
\[
 B_i(f,g)=\frac{E_i(f)+E_i(g)-E_i(f\circ g)}2,
 \qquad i=1,2,                             \tag{J3}
\]
where \(E_i\in\{0,1\}\) is the standard lift of \(\varepsilon_i\).
\end{proposition}
\begin{proof}
Degree-one cocycles with trivial coefficients are homomorphisms; the
finite-jet condition is equivalent to continuity to a discrete target.
Their kernels are closed and contain every commutator, hence contain its
closure. Conversely every homomorphism from
\(A_0=\mathbb Z^2\oplus(\mathbb Z/2)^2\) is continuous after composition
with (J1), since that quotient map already factors through a finite jet.
Therefore
\[
 H^1_{\mathrm{jet}}(J;D)=\operatorname{Hom}(A_0,D)
\]
for each discrete trivial group \(D\). Taking \(D=\mathbb Z\) gives the
first assertion.

Apply the lemma to \(0\to\mathbb Z\to\mathbb Q\to\mathbb Q/\mathbb Z\to0\).
The relevant part of the exact sequence is
\[
 \operatorname{Hom}(A_0,\mathbb Q)\longrightarrow
 \operatorname{Hom}(A_0,\mathbb Q/\mathbb Z)\longrightarrow
 H^2_{\mathrm{jet}}(J;\mathbb Z)\longrightarrow
 H^2_{\mathrm{jet}}(J;\mathbb Q).
\]
The first map is surjective on the two free coordinates, while its image
on the two order-two coordinates is zero. Its cokernel is exactly
\((\mathbb Z/2)^2\), which therefore injects into integral degree two as
the kernel of the coefficient map to \(\mathbb Q\). Rational cochains
form a complex of rational vector spaces, so their cohomology is
torsion-free. All integral torsion lies in that kernel; conversely the
kernel just found is finite, so it consists entirely of torsion. This
proves the second assertion, including the absence of odd torsion and of
elements of order four or higher in this degree and cochain theory.

Under \(\operatorname{Hom}(A_0,\mathbb Q/\mathbb Z)\), the torsion
coordinates are the characters \(\varepsilon_i/2\). Their connecting
cocycles are exactly (J3). Additivity modulo two makes (J3) integral, and
\(\delta^2=0\) gives \(\delta B_i=0\). Also \(2B_i=\delta E_i\).
Neither character lifts to a rational homomorphism on the torsion summand,
and they are independent there. Their two classes thus have exact order
two and generate the kernel calculated above.
\end{proof}

\begin{proposition}[An integral primitive for the degree-one product]\ours{}
For \(w(f)=c-3ab+2a^3\), the inhomogeneous coboundary is
\[
 (\delta w)(f,g)=x(f)y(g).                  \tag{J4}
\]
In particular \([x]\smile[y]=0\) over \(\mathbb Z\), both for finite-jet
cochains and in the ordinary cohomology of the underlying discrete group.
The latter assertion concerns these particular classes only.
\end{proposition}
\begin{proof}
For \(g(t)=t+At^2+Bt^3+Ct^4+\cdots\), direct substitution gives
\[
 a(f\circ g)=a+A,\quad
 b(f\circ g)=b+B+2aA,\quad
 c(f\circ g)=c+C+aA^2+2aB+3bA.
\]
Substitute these three identities in
\(\delta w(f,g)=w(g)-w(f\circ g)+w(f)\).
The cubic and mixed terms cancel, leaving \(a(B-A^2)\), as required.
All functions used depend on the fourth jet, so the equation holds in
both cochain complexes claimed.
\end{proof}

There is a second algebraic check. Put
\(z=c-a^3-2a(b-a^2)\). The fourth-jet composition law in the integer
coordinates \((x,y,z)\) is
\[
 (x,y,z)(X,Y,Z)=(x+X,y+Y,z+Z+yX).
\]
The coordinate change has an integer inverse, so this quotient is the
integer Heisenberg group. Since \(w=z-xy\), its displayed law reproduces
(J4), independently of the direct expansion. The degree-one squares also
vanish: for an integer homomorphism \(u\),
\(\delta\binom{u}{2}=-u\smile u\). Together with graded commutativity,
this checks that every product of two classes in the first group of (J2)
is zero.

\paragraph{Checks and their exact force.}
\texttt{code/check\_jets.py} uses symbolic substitution for (J4), a separate
integer-convolution implementation for composition, 1,500 deterministically
seeded triples with coefficients between \(-50\) and \(50\), and exhaustive
residue checks of the two parity characters (1,024 and 4,096 pairs).
All identities, including the two cocycle equations, passed. The proof of
the torsion assertion was checked a second way with
\(0\to\mathbb Z\xrightarrow{\,2\,}\mathbb Z\to\mathbb F_2\to0\):
\(H^1_{\mathrm{jet}}(J;\mathbb F_2)\) has dimension four, the reduction of
integral degree one has dimension two, and the resulting Bockstein image
has dimension two. The rational-coefficient argument is still needed to
exclude higher and odd torsion; the modulo-two check alone cannot do that.
The computational output is included. A separate review re-derived the
cup-product sign and checked the exact-sequence proof against the published
kernel statement; it found no mathematical defect. It did sharpen the
locator from the displayed quotient map to Proposition~2.9, which proves
that its kernel is the closed commutator. These checks concern formulas and
logical implications, not an independent proof of that published classification.

\paragraph{What remains, and a bounded next step.}
These results determine degree one, the torsion in degree two, and its two
explicit cocycles. They do not determine the torsion-free integral part of
degree two, higher integral groups or their products. They also do not
identify finite-jet cohomology with arbitrary continuous or discrete
cohomology. The remaining obstacle is control of integral inflation maps
along the central extensions \(J_{N+1}\to J_N\); rational Lie-algebra
comparison discards precisely that arithmetic information. A further
bounded run could construct integral resolutions for several successive
jet quotients, compute their degree-two inflation maps, and decide which
integral lattices and denominators survive. A finite range would be
evidence for a stabilization assertion, not a proof of one. Nothing in
the notebook or the sources checked establishes why the original speaker
did not pursue the full integral question; the documented topology and
coefficient distinctions explain a mathematical obstacle without inventing
a historical motive.

\subsection{Verification annotation: the Laplace spectrum on the positive Witt algebra}
\label{dev:1977-laplacian-check}
This is a check of a published class-A thread, not a newly selected open
problem. On \GSpage{P000784--P000785}, take
\([e_i,e_j]=(j-i)e_{i+j}\) for \(i,j\geq1\), declare the increasing wedge
monomials orthonormal, and write
\(\Delta_q=d_q^{\mathsf T}d_q+d_{q+1}d_{q+1}^{\mathsf T}\).
Each weight space is finite-dimensional. The corrected spectral statement
is given by \citet[Theorem~2, p.~324]{fuchswilmarth2010laplacian}; the following checks locate a
specific error in the notebook's assertion about distinct eigenvalues.

\begin{example}[A collision of the displayed spectral labels]\oursComputed{}
The two gap-three partitions of weight 25
\[
 (1,10,14),\qquad(3,7,15)
\]
both have
\[
 E(I)=\sum_s\binom{i_s}{3}-\sum_{s<t}i_si_t=320.
\]
For the first, \(120+364-(10+14+140)=320\); for the second,
\(1+35+455-(21+45+105)=320\). Thus partitions cannot count
distinct spectral values. Under the corrected multiplicity rule their
contributions are four and eight, and must be added.
\end{example}

\begin{remark}[Exact finite checks]\oursComputed{}
\texttt{code/check\_laplacian.py} constructs the Chevalley--Eilenberg
boundary with its alternating signs, checks \(d^2=0\), and factors exact
integer characteristic polynomials for every degree of weights 1--22 and
25. A separately implemented gap-partition enumeration gives exactly the
same eigenvalues and total multiplicities, including multiplicity 12 at
eigenvalue 320 in weight 25. It also expands
\[
 \prod_{i<j}(z_i-z_j)^3\prod_i z_i^a,\qquad a=1,2,
\]
as alternating polynomials and checks both boundary and adjoint boundary
vanish for exterior degrees 1--4. No numerical eigenvalue rounding is
used. The second check, partition enumeration, shares no boundary-matrix
code with the first. These computations confirm a finite range and the
explicit correction; they do not replace the all-weight published theorem.
An independent reviewer also built the weight-25 boundaries with a different
sign algorithm and computed exact ranks of \(\Delta-320I\). The nullities
in exterior degrees 3, 4, 5 and 6 are respectively 2, 5, 4 and 1, totaling
12. Its separate enumeration checks every weight through 25, including
23 and 24, before identifying the first collision. The review code and
output are included.
\end{remark}

\subsection{Verification annotation: the literal KP normalization}
\label{dev:1977-kp-check}
Items 1977-434.3--1977-434.5 belong to the published rational-KP thread
\citep[Theorem~3 and its corollaries]{krichever1978rational}. The following
calculation makes the editorial choice of signs and momenta explicit. It
does not claim an error in a published theorem stated with other conventions.

\begin{proposition}[Pole equations in the notebook's convention]\ours{}
Let the functions \(q_i(y,t)\), \(1\leq i\leq N\), be smooth and pairwise
distinct on an open parameter domain, and put
\[
 u(x,y,t)=-2\sum_{i=1}^N(x-q_i(y,t))^{-2}.
\]
As an identity of meromorphic functions of \(x\),
\[
 3u_{yy}+(4u_t+6uu_x+u_{xxx})_x=0
\]
holds if and only if, for every \(i\),
\[
 q_{i,yy}=-8\sum_{j\ne i}(q_i-q_j)^{-3},\qquad
 q_{i,t}=\frac34q_{i,y}^{\,2}-3\sum_{j\ne i}(q_i-q_j)^{-2}.
 \tag{K1}
\]
\end{proposition}
\begin{proof}
At a fixed pole write \(s=x-q_i\) and
\(u=-2s^{-2}+R_0+R_1s+O(s^2)\), where
\[
 R_0=-2\sum_{j\ne i}(q_i-q_j)^{-2},\qquad
 R_1=4\sum_{j\ne i}(q_i-q_j)^{-3}.
\]
The leading sixth-order terms cancel. The only remaining negative powers
in the differential equation have coefficients
\[
 [s^{-4}]=48q_{i,t}-72R_0-36q_{i,y}^{\,2},\qquad
 [s^{-3}]=-24R_1-12q_{i,yy}.
\]
Their vanishing is exactly (K1). Conversely, if (K1) holds, the left-hand
side is a rational function of \(x\) with no finite poles. It tends to
zero at infinity, so it is identically zero. The distinct-pole hypothesis
is used in the expansion and is not omitted at collisions.
\end{proof}

\begin{proposition}[Compatible canonical Hamiltonians]\ours{}
With canonical coordinates \((q_i,p_i)\) on the distinct-position locus,
define
\[
 H=\frac12\sum_i p_i^2-4\sum_{i<j}(q_i-q_j)^{-2},\qquad
 K=\frac14\sum_i p_i^3-3\sum_{i<j}(p_i+p_j)(q_i-q_j)^{-2}.
 \tag{K2}
\]
Then \(\{H,K\}=0\). Taking the \(H\)-flow as \(y\) and the \(K\)-flow
as \(t\) gives \(p_i=q_{i,y}\) and (K1), hence the displayed rational KP
solution locally away from collisions.
\end{proposition}
\begin{proof}
Write \(d_{ij}=q_i-q_j\). Direct differentiation in
\(\{H,K\}=\sum_i(H_{q_i}K_{p_i}-H_{p_i}K_{q_i})\) cancels the
\(p_i^2\) terms. The remaining momentum terms are
\(-6\sum_{i\ne j}p_ip_jd_{ij}^{-3}=0\), by exchanging \(i,j\).
What remains is
\[
 -24\sum_i\sum_{j\ne i}\sum_{k\ne i}d_{ij}^{-3}d_{ik}^{-2}.
\]
Terms with \(j=k\) cancel in pairs. Group the others by unordered triples
and set \(a=q_i-q_j\), \(b=q_j-q_k\), \(c=q_k-q_i\), so
\(a+b+c=0\). Multiplication of a triple's sum by \(a^3b^3c^3\) gives
\[
 a^3(b-c)+b^3(c-a)+c^3(a-b)=0.
\]
This last polynomial vanishes upon substituting \(c=-a-b\). Thus the
bracket is zero. Hamilton's equations give
\(q_{i,y}=p_i\), \(p_{i,y}=-8\sum_{j\ne i}d_{ij}^{-3}\), and the
second formula in (K1) for \(q_{i,t}\). Commuting smooth vector fields have
local compatible flows on this open locus; no global continuation through
a collision is claimed.
\end{proof}

\begin{remark}[The lower coefficient in the third-order operator]\oursComputed{}
Let \(D=\partial_x\), \(L=D^2+u\) and
\(M=D^3+\frac32uD+w\). The coefficient of \(D\) in
\([\partial_y-L,\partial_t+M]\) vanishes precisely when
\(w_x=\frac34(u_{xx}+u_y)\). The remaining scalar equation is
\[
 u_t+w_y-w_{xx}+u_{xxx}+\frac32uu_x=0.
\]
Differentiating it in \(x\) and eliminating derivatives of \(w\) gives
the notebook's KP equation. Locally one may choose the function of
\((y,t)\) arising on integrating \(w_x\) to satisfy this scalar equation;
global boundary or periodicity conditions must be supplied separately.
Using \(\partial_t-M\) instead reverses the KP time convention. Thus the
final coefficient in \(M\) cannot in general be the same \(u\) as in \(L\).
\end{remark}

The Laurent expansion was derived independently by the record editor and
the chapter editor. \texttt{code/check\_kp.py} reproduces it symbolically,
constructs differential-operator composition by the Leibniz rule, and
checks the coefficient elimination. The triple cancellation in the
Hamiltonian proof is checked as a polynomial identity, with independent
exact evaluations of the full Poisson bracket. These checks found the
same normalization. They do not assert a classification beyond the
explicit pole ansatz and local distinct-position domain just proved.


% No new C/D theorem is claimed for records419--427.
\subsection{An annotation on the thirteen-coin convention}

This is a check of an expository puzzle (1977-424.2), not a proposed new
solution to a class C or D research question.  The classical subject appears
in \citet{dyson1946pennies}; the publisher record was checked, but the proof
below is a self-contained reconstruction.  The difference between identifying
a counterfeit and determining its sign is essential to the notebook's number.

\begin{proposition}[Identity only, with equal pan counts]\ours{}
Suppose exactly one of \(m\) coins has a different mass, either heavier or
lighter, and all other coins have a common mass.  No genuine coin is supplied
externally.  Every weighing compares equal numbers of coins on an ideal
three-outcome balance.  For \(N\geq1\), the largest \(m\) whose counterfeit
identity can always be determined in at most \(N\) weighings is
\((3^N-1)/2\).  Its sign need not always be determined.
\end{proposition}

\begin{proof}
For the upper bound let the first weighing compare \(a\) coins against
\(a\), leaving \(b=m-2a\) aside.  Either unbalanced branch contains
\(2a\) possible signed counterfeits.  They must be distinguished by
\(N-1\) subsequent ternary observations, so
\(2a\leq3^{N-1}\).  Since the left side is even,
\(2a\leq3^{N-1}-1\).  On the balanced branch the \(b\) unweighed coins
have unknown signs.  For any fixed coin, its two signs can follow the same
history only if that coin is never weighed on that history; until its first
weighing both signs produce balance, whereas its first appearance produces
opposite imbalances.  Consequently only the all-balance history can fail to
determine the sign, and it can represent at most one identity.  At least
\(2b-1\) terminal histories are therefore necessary when \(b>0\).
Hence \(b\leq(3^{N-1}+1)/2\); the same bound is trivial if \(b=0\).
Adding gives \(m\leq(3^N-1)/2\).  The argument permits adaptive choices
of the subsequent weighings.

For attainment, use the alphabet \(\{-1,0,1\}\), in cyclic order
\(-1\to0\to1\to-1\).  For every nonconstant word of length \(N\),
find its first unequal adjacent pair.  Retain the word precisely when that
pair follows the cyclic order.  Negation reverses the order, so exactly one
of each opposite pair is retained.  There are
\(M=(3^N-3)/2\) retained words.  Cyclically permuting all three alphabet
symbols preserves the selection rule, so each coordinate takes each symbol
equally often among the retained words.  Their coordinate sums are therefore
zero.  Add the all-zero word, giving \(M+1=(3^N-1)/2\) columns.

Assign one column to each coin.  In weighing \(j\), put entries \(-1\)
on the left, \(1\) on the right, and leave entries \(0\) aside.  Every
weighing has equal pan counts.  A heavy counterfeit produces its column as
the outcome word, with right-heavy denoted \(1\); a light one produces the
negative column.  Distinct nonzero columns are never opposites, so these
histories identify their coins and signs.  The all-zero history identifies
the last coin, with either sign.  For \(N=1\) the retained set is empty and
the single coin is identified without weighing, as required.
\end{proof}

\begin{remark}[If the sign is required]\ours{}
When both identity and sign must be recovered, the balanced first branch
must distinguish all \(2b\) states, hence
\(b\leq(3^{N-1}-1)/2\).  Together with the same unbalanced-branch bound
this gives \(m\leq(3^N-3)/2\).  The nonzero retained columns above attain
that bound for \(N\geq2\).  Thus three weighings permit 12 coins with
both outputs, or 13 with identity alone.  This distinction, not a numerical
misprint, accounts for the notebook's 13.
\end{remark}

\begin{example}[Exhaustive finite checks]\oursComputed{}
The accompanying \texttt{coin\_check.py} generates the columns and computes
actual left and right pan masses, independently of the signed-column
formula.  For \(N=1,\ldots,6\), it checks respectively
\(1,4,13,40,121,364\) coins.  Every coordinate balances, and every outcome
determines a unique coin.  The only two-state collision is the sign of the
zero-column coin.  In the three-weighing case all 26 signed hypotheses were
enumerated, yielding 25 histories and no collision between distinct
identities.  The proof establishes the general statement; finite enumeration
only checks the implementation and the notebook's small case.
\end{example}


% No independent C/D result is claimed by this record editor.
% Items 434.3--434.5 contain normalized KP verification annotations; the assembler independently checks them.
% Item 436.5 remains an annotated partial result: breathers are already explicitly present in the notebook.
% The chapter assembler supplies any selected deep development separately.


\subsection{Independent checks of the spectral and algebraic formulas}

The following computations explain emendations in occurrences
1977-442.4, 1977-442.6--.7, 1977-443.9 and 1977-443.12--.14, and two later
algebraic formulas. They are checks made for this companion, with no claim
of historical novelty. The relevant published subjects are described by
\citet{calogerodegasperis1976inverse,calogerodegasperis1978extension,calogero1978motion}; only their publisher metadata and abstracts or
editorial footnotes were available for the present comparison. Thus the
proofs below verify the displayed notebook equations independently of a
claim to have read their full published normalization conventions.

\begin{proposition}[Transmission phase in the notebook convention]\oursComputed{}
Let \(u\in C(\mathbb R)\cap L^1(\mathbb R)\) be real. Use the
notebook's physical scattering normalization for
\(\psi''=(u-k^2)\psi\): at negative infinity,
\(\psi\sim T(k)e^{-ikx}\), and at positive infinity,
\(\psi\sim e^{-ikx}+R(k)e^{ikx}\). With
\(T=|T|e^{2i\varphi}\) and the branch \(\varphi\to 0\) as
\(k\to+\infty\), its leading high-energy phase is
\[
 \varphi(k)=-\frac{1}{4k}\int_{\mathbb R}u(x)\,dx+o(k^{-1}).
\]
Consequently, the source's coefficient \(-c_0/2\), where
\(c_0=\int u\), is incompatible with its own factor 2 in the phase of T.
\end{proposition}
\begin{proof}
Normalize the left Jost solution f by \(f\sim e^{-ikx}\) at negative
infinity. Its Volterra equation is
\[
 f(x)=e^{-ikx}+\frac1k\int_{-\infty}^x
          \sin(k(x-y))u(y)f(y)\,dy.
\]
For real \(k>0\), Gronwall's inequality gives
\(\sup_x|f(x)|\le\exp(\|u\|_1/k)\) and
\(\sup_x|f(x)-e^{-ikx}|\le
\|u\|_1\exp(\|u\|_1/k)/k\).
Splitting the sine into exponentials shows that the coefficient of
\(e^{-ikx}\) at positive infinity is
\[
 A(k)=1-\frac1{2ik}\int_{\mathbb R}e^{iky}u(y)f(y)\,dy
     =1+\frac{i c_0}{2k}+O(k^{-2}).
\]
The error bound follows from the uniform estimate and integrability
of u. For all sufficiently large k, \(A(k)\ne0\).
The physical solution is \(\psi=f/A\), so \(T=A^{-1}\), and
\(\log T=-i c_0/(2k)+O(k^{-2})\). Taking half the imaginary part
on the branch near1 proves the formula. No convergence of a full
asymptotic series is asserted.

An independent exact check uses
\(u=-2r^2\operatorname{sech}^2(rx)\), \(r>0\), and
\[
 \psi=e^{-ikx}\frac{ik+r\tanh(rx)}{ik+r},\qquad
 T(k)=\frac{k+ir}{k-ir}.
\]
Direct differentiation gives \(\psi''=(u-k^2)\psi\). Its multiplying
factor tends to 1 at positive infinity and to
\((ik-r)/(ik+r)=(k+ir)/(k-ir)=T(k)\) at negative infinity,
so it has exactly the stated physical normalization, with \(R=0\).
Here \(c_0=-4r\) and \(\varphi=\arctan(r/k)=r/k+O(k^{-3})\),
again giving \(-c_0/4\). The symbolic residual and expansion are
recorded in the accompanying verification output. An independent
review caught an earlier reciprocal error in the stated boundary
normalization; these explicit endpoint checks correct it without
altering the final coefficient \(-c_0/4\).
\end{proof}

\begin{proposition}[Weights of the integrated recursion moments]\oursComputed{}
Let \(u\in\mathcal S(\mathbb R)\), and for
\(f\in\mathcal S(\mathbb R)\) define
\[
 L_u f=f''-4uf+2u'\int_x^\infty f(s)\,ds,\qquad
 S_u=xu'+2u,\qquad c_n[u]=\int_{\mathbb R}L_u^n S_u\,dx,
\]
The upper-tail primitive of a Schwartz function is bounded; its
product with \(u'\) is Schwartz. Thus L preserves this space,
all the displayed iterates and integrals exist, and the boundary
terms used below vanish.
For \(u_\lambda(x)=\lambda^2u(\lambda x)\), \(\lambda>0\),
\[
 c_n[u_\lambda]=\lambda^{2n+1}c_n[u].
\]
Thus a pure scaling flow \(u_t=b(xu_x+2u)\), for constant \(b\in\mathbb R\),
has \(\dot c_n=(2n+1)bc_n\), rather than \(nbc_n\).
\end{proposition}
\begin{proof}
If \(h_\lambda(x)=\lambda^a h(\lambda x)\), then its upper-tail
primitive has weight \(a-1\). Each of the three terms in
\(L_{u_\lambda}h_\lambda\) therefore has weight \(a+2\). Since
\(S_{u_\lambda}(x)=\lambda^2S_u(\lambda x)\), induction gives
weight \(2n+2\) before integration and \(2n+1\) afterwards.
The scaling-flow solution is obtained by \(\lambda=e^{bt}\).

Independently, integration by parts gives \(c_0=\int u\).
Put \(U(x)=\int_x^\infty u(s)\,ds\). Then
\(\int_x^\infty S_u=U-xu\), and the formula for L gives
\(c_1=-3\int u^2\). For example, on
\(u_\lambda=\lambda^2e^{-\lambda^2x^2}\), the two quantities are
\(\lambda\sqrt\pi\) and \(-3\lambda^3\sqrt{\pi/2}\).
These exact integrals independently test weights 1 and 3, already
refuting the source's proposed weights 0 and 1.
\end{proof}

\begin{proposition}[The sign in the Bianchi quadrilateral]\oursComputed{}
Let \(p,q\in\mathbb C\) be constants, and let U,A,B be
complex-valued \(C^1\) functions on a real interval I. Suppose the
two lower B\"acklund edges satisfy
\[
 U_x+A_x=-\tfrac12(A-U)(4p+A-U),
\]
and the same equation with \((A,p)\) replaced by \((B,q)\).
On any subinterval where the displayed denominator is nonzero,
the compatible upper value is the \(C^1\) function
\[
 C=U-\frac{(p+q)(A-B)}{p-q+\tfrac12(A-B)}.
\]
In particular, the denominator uses \(A-B\), with the same order
as the numerator in these conventions.
\end{proposition}
\begin{proof}
The upper edges require
\[
 A_x+C_x=-\tfrac12(C-A)(4q+C-A),\qquad
 B_x+C_x=-\tfrac12(C-B)(4p+C-B).
\]
Subtracting them and eliminating \(A_x-B_x\) by the two lower
edges gives a linear equation for C. Solving it yields the displayed
formula. This elimination is one check.

For a second check, differentiate that rational formula using the two
lower equations, leaving \(U_x\) arbitrary. Substitution into either
upper equation gives zero identically. The symbolic verification file
performs these as two separate residual calculations. The computation
only concerns this local differential quadrilateral; it neither fills
in missing scattering norming data nor proves a global permutability
statement at zeros of the denominator.
\end{proof}

\begin{proposition}[The Hermite interaction matrix]\oursComputed{}
Let \(\bar x_1,\ldots,\bar x_n\) be the distinct zeros of the
physicists' Hermite polynomial \(H_n\), and put
\[
 A_{jk}=-\frac1{(\bar x_j-\bar x_k)^2}\quad(j\ne k),\qquad
 A_{jj}=\sum_{k\ne j}\frac1{(\bar x_j-\bar x_k)^2}.
\]
For the source's root dynamics
\[
 \ddot x_j=-x_j+\sum_{k\ne j}
       \frac{1+2\dot x_j\dot x_k}{x_j-x_k},
\]
the equilibrium linearization is
\(\ddot\varepsilon+(I+A)\varepsilon=0\).
Moreover, \(\operatorname{Spec}(A)=\{0,1,\ldots,n-1\}\), with
each eigenvalue simple.
\end{proposition}
\begin{proof}
The derivative of each inverse difference contributes
\(-(\varepsilon_j-\varepsilon_k)/(\bar x_j-\bar x_k)^2\).
The velocity products have zero linear part at a stationary point.
The independent term \(-x_j\) contributes \(-\varepsilon_j\),
giving the stated matrix \(I+A\).

For the spectrum, use the notebook polynomial equation
\[
 \Psi_{tt}+\tfrac12\Psi_{xx}-x\Psi_x+n\Psi=0,
 \qquad
 \Psi=2^{-n}\sum_{m=0}^n c_m(t)H_{n-m}(x),\quad c_0=1.
\]
The Hermite differential equation gives
\(\ddot c_m+mc_m=0\) for \(1\le m\le n\).
At a polynomial with distinct roots, the differential of the
root-to-coefficient map is invertible: a variation has
\(\delta\Psi(\bar x_j)=-\Psi'(\bar x_j)\delta x_j\), and
evaluation at n distinct points is invertible on polynomials of
degree at most \(n-1\). Therefore \(I+A\) is conjugate to
\(\operatorname{diag}(1,\ldots,n)\).

As a second explicit check, for n=3 the roots are
\(-\sqrt{3/2},0,\sqrt{3/2}\) and
\[
 A=\begin{pmatrix}
 5/6&-2/3&-1/6\\
 -2/3&4/3&-2/3\\
 -1/6&-2/3&5/6
 \end{pmatrix},\qquad
 \det(\lambda I-A)=\lambda(\lambda-1)(\lambda-2).
\]
The verification code independently differentiates the full nonlinear
force and obtains its Jacobian \(-I-A\). It also checks the Hermite
mode identity exactly for all 36 pairs \(1\le m\le n\le8\).
Already n=1 shows why the identity cannot be omitted: A is zero but
the single root satisfies \(\ddot x=-x\).
\end{proof}

\begin{remark}[Two later algebraic checks]\oursComputed{}
For occurrence 1977-457.18, take the ordinary \(\mathfrak{sl}_2\)
relations and define
\[
 \varepsilon_{3r}=\tfrac12ht^r,\quad
 \varepsilon_{3r+1}=et^r,\quad
 \varepsilon_{3r-1}=\tfrac12ft^r.
\]
Matrix commutators verify all nine residue cases of
\([\varepsilon_i,\varepsilon_j]=a_{ij}\varepsilon_{i+j}\),
where \(a_{ij}\in\{-1,0,1\}\) is congruent to \(j-i\) modulo 3.
Separately, the 27 residue triples satisfy Jacobi directly. The
nonhomogeneous ideal \((t-1)\mathfrak{sl}_2[t,t^{-1}]\) explains why
the classification convention is graded simplicity, as already stated
in \citet[Introduction]{kac1968graded}.

For occurrence 1977-457.25, left Grassmann differentiation satisfies
\(\{\partial_{\theta^a},\theta^b\}=\delta_a^b\), while distinct
odd multiplications or derivatives anticommute. Hence the two cross
terms in the notebook's Q and \(\bar Q\) each contribute
\(-i\sigma^m\partial_m\), giving
\(\{Q_\alpha,\bar Q_{\dot\beta}\}=-2i\sigma^m_{\alpha\dot\beta}
\partial_m\); the same-chirality brackets vanish. Independently,
the code realizes multiplication and differentiation as exact 16-by-16
matrices on the exterior algebra of four generators and checks all 12
anticommutator pairs with four independent commuting derivative
symbols. These calculations validate the surviving manuscript signs;
they do not resolve the illegible cancelled mark or supply the missing
Lorentz action on the spinor coordinates.
\end{remark}

\section{Problem index}\label{problems:1977}
Earlier-year occurrence identifiers are retained as text; current-year identifiers link to their entries.

\subsection{GP-1977-01: Nonarithmetic lattices in complex hyperbolic dimension at least four}\label{prob:GP-1977-01}
For which \(n\geq4\), if any, is there a nonarithmetic lattice in \(PU(n,1)\)? This is the modern unresolved part of the notebook\textquoteright s broader rank-one arithmeticity question.
\par\noindent\textit{Occurrences:} \hyperref[item:1977-421.14]{1977-421.14}.
\par\noindent\textit{Status (2026-11):} As of 10 November 2026, open in the literature located. The 2023 paper explicitly states the question for \(n\geq4\), and the introduction of \citet{yu2026mirror}, first submitted 30 October 2026, still reports known examples only in dimensions two and three. No higher-dimensional solution was found.

\subsection{GP-1977-02: A conceptual proof of the first Gelfand--Dickey Jacobi identity}\label{prob:GP-1977-02}
Explain and prove the Jacobi identity for the coefficient-linear Gelfand--Dickey Hamiltonian operator through a structural construction, rather than a direct binomial calculation.
\par\noindent\textit{Occurrences:} \hyperref[item:1977-426.10]{1977-426.10}.
\par\noindent\textit{Status (2026-11):} Resolved by the formal Volterra-group/coadjoint construction: \citet[\S\S5--6]{lebedevmanin1979volterra}, received November 1978 and published 1979. The full proof was read; the object is the first Gelfand--Dickey bracket.

\subsection{GP-1977-03: Local algebraization of Harish–Chandra categories}\label{prob:GP-1977-03}
For a real linear connected semisimple group and bounded central-support domain \(D\), construct an algebra finite over the center realizing \(\mathcal C_D\); near a regular infinitesimal character, realize it by an explicit nice analytic matrix order.
\par\noindent\textit{Occurrences:} \hyperref[item:1977-428.2]{1977-428.2}, \hyperref[item:1977-428.6]{1977-428.6}, \hyperref[item:1977-429.7]{1977-429.7}, \hyperref[item:1977-430.7]{1977-430.7}.
\par\noindent\textit{Status (2026-11):} Solved in the stated scope by Bernstein, I. M. Gelfand and S. I. Gelfand, presented 9 January 1978 and published March 1978; the nice-algebra theorem requires a regular infinitesimal character and a sufficiently small neighborhood. As of 10 November 2026, no extension to arbitrary singular blocks is asserted here.

\subsection{GP-1977-04: Finite-action Yang–Mills critical connections beyond instantons}\label{prob:GP-1977-04}
Classify smooth finite-action SU(2) Yang--Mills connections on Euclidean four-space up to gauge, equivalently critical connections on bundles over the round four-sphere after removable-singularity compactification.
\par\noindent\textit{Occurrences:} \hyperref[item:1977-432.4]{1977-432.4}.
\par\noindent\textit{Status (2026-11):} As of 10 November 2026, ADHM-type data classify the self-dual or anti-self-dual minima, while nonminimal critical points also exist. The checked 2026 gap theorem controls an energy range and does not give a complete classification of all critical points.

\subsection{GP-1977-05: Fixed-charge instanton classification and connectedness}\label{prob:GP-1977-05}
For each charge \(k\), describe all \(SU(2)\) instantons on \(S^4\) up to gauge and determine whether their moduli space is path-connected.
\par\noindent\textit{Occurrences:} \hyperref[item:1977-432.8]{1977-432.8}, \hyperref[item:1977-432.12]{1977-432.12}, \hyperref[item:1977-435.13]{1977-435.13}.
\par\noindent\textit{Status (2026-11):} As of 10 November 2026, as of 10 November 2026, as of 10 November 2026, the finite-dimensional construction was announced in 1977 and published in 1978; path-connectedness for SU(2) was proved by Taubes in 1984. Dayaprema--Waldron (2026), Theorem 1.2 and section 3.3, provide a new proof and explain the extension to SU(r).

\subsection{GP-1977-06: Selecting combinations with prescribed inner divisors}\label{prob:GP-1977-06}
For fixed inner \(\omega\) and a finite family in \(H^\infty\), find a constant-coefficient linear combination whose common inner divisor with \(\omega\) equals that of the whole family; determine the exceptional parameters for a one-parameter family.
\par\noindent\textit{Occurrences:} \hyperref[item:1977-433.6]{1977-433.6}, \hyperref[item:1977-433.7]{1977-433.7}.
\par\noindent\textit{Status (2026-11):} Known before the talk: Sherman\textquotesingle s two-function lemma is printed with proof in Sz.-Nagy (1976), and the same paper proves a simultaneous countable-row extension with summable coefficients. As of 10 November 2026, no new open problem is inferred from the expository statement.

\subsection{GP-1977-07: Rational decaying KP solutions and commuting pole dynamics}\label{prob:GP-1977-07}
Classify solutions of \(3u_{yy}+(4u_t+6uu_x+u_{xxx})_x=0\) rational in \(x\) and vanishing at infinity, through finite-dimensional commuting equations for their poles.
\par\noindent\textit{Occurrences:} \hyperref[item:1977-434.2]{1977-434.2}, \hyperref[item:1977-434.3]{1977-434.3}, \hyperref[item:1977-434.4]{1977-434.4}.
\par\noindent\textit{Status (2026-11):} As of 10 November 2026, as of 10 November 2026, as of 10 November 2026, krichever\textquotesingle s 1978 Theorem 3 and corollaries solve the rational classification in its stated sense. Distinct-pole coordinates require local treatment across collisions, and the source\textquotesingle s particle coefficients must be normalized consistently. Earlier particle-integrability occurrences 1975-372.7 and 1975-372.12 are related background, not the same rational-KP question.

\subsection{GP-1977-08: Inducing a unitary connection by orthogonal projection}\label{prob:GP-1977-08}
Given a smooth Hermitian vector bundle with a compatible connection over a finite-dimensional smooth manifold, realize it isometrically inside a finite-dimensional trivial Hermitian bundle so that its connection is the orthogonal projection of the trivial connection.
\par\noindent\textit{Occurrences:} \hyperref[item:1977-435.6]{1977-435.6}.
\par\noindent\textit{Status (2026-11):} Affirmative before the seminar: the Narasimhan--Ramanan universal-connection theorem was announced in 1961 and is stated in the 1962 Prague proceedings. As of 10 November 2026, this is an established theorem, not an open problem.

\subsection{GP-1977-09: Soliton resolution with the correct coherent structures and radiation}\label{prob:GP-1977-09}
In a specified class of global finite-energy sine-Gordon solutions, establish a large-time decomposition into finitely many localized coherent structures and radiation, with an explicit convergence topology; determine whether breathers and modified rather than free scattering are necessary.
\par\noindent\textit{Occurrences:} \hyperref[item:1977-436.5]{1977-436.5}.
\par\noindent\textit{Status (2026-11):} As of 10 November 2026, Chen--Liu--Lu prove a precise generic weighted-Sobolev version with kinks, antikinks, breathers and modified radiation (preprint 2020, version 2022). The source\textquotesingle s unrestricted class and convergence topology are not fixed. Breathers are explicitly recognized in the notebook itself; no counterexample to that broader wording is claimed here.

\subsection{GP-1977-10: Constructing a quantum inverse-scattering formalism}\label{prob:GP-1977-10}
For a specified classically integrable field equation, construct quantum scattering-data operators, their exchange algebra, commuting quantum integrals and eigenstates.
\par\noindent\textit{Occurrences:} \hyperref[item:1977-436.16]{1977-436.16}.
\par\noindent\textit{Status (2026-11):} An early concrete solution is the Sklyanin--Faddeev announcement of 1978, especially its explicit nonlinear-Schrödinger construction. Its sine-Gordon comparison is expressly less complete. As of 10 November 2026, quantum inverse scattering is an established model-specific method; the notebook does not state a universal quantization theorem.

\subsection{GP-1977-11: Describe projective bundles by finite exterior modules}\label{prob:GP-1977-11}
Can algebraic vector bundles on \(\mathbb P^n\) be represented and compared by graded modules of finite dimension over an exterior algebra? The notebook's broad classification question is indexed in the specific form answered by its theorem.
\par\noindent\textit{Occurrences:} \hyperref[item:1977-438.5]{1977-438.5}.
\par\noindent\textit{Status (2026-11):} As of November 2026, the representation and stable-isomorphism criterion are established by BGG. A full explicit moduli classification is a different, substantially broader question and is not claimed solved here.

\subsection{GP-1977-12: Why homogeneous functions define a line bundle}\label{prob:GP-1977-12}
Why do homogeneous rational functions of integer degree \(j\), regular above a projective open set, form a locally free sheaf of rank one?
\par\noindent\textit{Occurrences:} \hyperref[item:1977-438.6]{1977-438.6}.
\par\noindent\textit{Status (2026-11):} Resolved classically: division by \(x_i^j\) gives an explicit local trivialization on \(x_i\ne0\). This is an exercise, not an open research question.

\subsection{GP-1977-13: Cohomology of the integral formal substitution group}\label{prob:GP-1977-13}
Compute the cohomology of \(J=\{t+\sum_{i\ge1}a_it^{i+1}:a_i\in\mathbb Z\}\). The notebook leaves coefficients and topology unstated. This chapter distinguishes finite-jet cohomology \(\varinjlim_NH^*(J/(t^{N+1});A)\), continuous cochains, and unrestricted cohomology of the abstract group.
\par\noindent\textit{Occurrences:} \hyperref[item:1977-441.7]{1977-441.7}.
\par\noindent\textit{Status (2026-11):} As of November 2026, Babenko--Bogatyi's real uniformly locally constant-support theory is computed and its continuous abelianization is known. For the explicit integral finite-jet refinement, the chapter derives torsion \((\mathbb Z/2)^2\) in degree two and a vanishing product of the two integral degree-one generators. It does not compute the full integral ring or the abstract discrete theory. Pavez-Cornejo's 2024 preprint/2025 publication uses closed commutators, not ordinary abstract commutators.

\section{Corrections to the transcription and translation}\label{corr:1977}
The evidence distinguishes a transcription correction from a mathematical emendation of wording faithfully preserved in the notebook.
\subsection*{P000705 (both; provenance): Opening date}
\noindent\textit{Transcription reads.} Handwritten 5/IV--77; catalogue date 1977-09-05\par
\noindent\textit{Critical reading.} Retain the catalogue date with a provenance note; do not silently change the visible IV into IX.\par
\noindent\textit{Evidence.} P705 scan visually inspected; the supplied catalogue and handwriting conflict. No independent chronological document resolves this conflict.\par
\subsection*{P000706 (both; faithful error): Instanton/bundle correspondence}
\noindent\textit{Transcription reads.} Self-dual connections correspond to holomorphic bundles on $\mathbb{CP}^{3}$.\par
\noindent\textit{Critical reading.} Specify rank, determinant, unitary reality data, triviality on real twistor lines, and the relevant equivalence classes.\par
\noindent\textit{Evidence.} The notebook statement is abbreviated. Atiyah--Ward 1977 publisher abstract supports the correspondence, but its full detailed theorem was not obtained; supplied conditions are a modern reconstruction.\par
\subsection*{P000706 (both; faithful error): Projective quotient notation}
\noindent\textit{Transcription reads.} $\mathbb{C}^{4}/\mathbb{C}^{*}$ and $\mathbb{H}^{2}/\mathbb{H}^{*}$\par
\noindent\textit{Critical reading.} Remove the origin before projectivizing.\par
\noindent\textit{Evidence.} P706 scan and both transcriptions preserve the shorthand; the zero vector is not a projective point.\par
\subsection*{P000706-P000707 (both; ambiguity): Quaternionic real structure}
\noindent\textit{Transcription reads.} The quaternionic antilinear map is called an involution.\par
\noindent\textit{Critical reading.} On $\mathbb{C}^{4}$ the lift $J$ has $J^{2}=-1$; the induced projective map has square $1$.\par
\noindent\textit{Evidence.} P706--P707 scans and displayed coordinate map checked; this is a distinction between lift and projectivization, not evidence of a transcription mistake.\par
\subsection*{P000707 (both; ambiguity): Genus of the spectral curve}
\noindent\textit{Transcription reads.} Rational solutions correspond to genus zero.\par
\noindent\textit{Critical reading.} Specify geometric genus of the normalization when a singular rational degeneration is intended.\par
\noindent\textit{Evidence.} Krichever 1977, §§1--2 and Appendix 2; singular rational curves can retain positive arithmetic genus.\par
\subsection*{P000707 (both; ambiguity): Manin algebraic-type criterion}
\noindent\textit{Transcription reads.} $L_{3}\notin\mathbb{C}[L_{1},L_{2}]$\par
\noindent\textit{Critical reading.} Retain the expression with an ambiguity note: $L_{2}$ is a time-evolution operator in a Lax pair and need not commute with $L_{1}$.\par
\noindent\textit{Evidence.} P707 scan visually inspected; the commuting spectral ring must be distinguished from the displayed evolution pair.\par
\subsection*{P000707 (both; ambiguity): Polynomial relation attribution}
\noindent\textit{Transcription reads.} Jacobson, 1930s, is named for a polynomial relation of operators.\par
\noindent\textit{Critical reading.} Preserve the historical attribution as unresolved; distinguish the Burchnall--Chaundy theorem for commuting scalar ordinary differential operators.\par
\noindent\textit{Evidence.} P707 scan and Krichever 1978 introduction checked. No exact Jacobson source was identified; no replacement attribution is claimed to correct handwriting.\par
\subsection*{P000710 (both; faithful error): Algebraic matrix-entry rigidity}
\noindent\textit{Transcription reads.} Every lattice in $\mathrm{SL}(n,\mathbb{R})$ is conjugate to matrices with algebraic entries.\par
\noindent\textit{Critical reading.} Impose the higher-rank hypotheses, in particular $n\geq3$ for this formulation.\par
\noindent\textit{Evidence.} P710 scan; Tits 1977 superrigidity/arithmeticity statement. The $n=2$ formulation admits deformations and is false without qualification.\par
\subsection*{P000711 (both; faithful error): Nonarchimedean places}
\noindent\textit{Transcription reads.} Use $\mathbb{Q}_{p}$ for the $p$-adic realization of the number-field coefficients.\par
\noindent\textit{Critical reading.} For a coefficient field $K$ use all nonarchimedean completions $K_{v}$, or the equivalent restriction-of-scalars representation over $\mathbb{Q}_{p}$.\par
\noindent\textit{Evidence.} P711 scan and Tits 1977 §5 deduction; primes and embeddings of a general number field cannot be suppressed in the exact argument.\par
\subsection*{P000712 (both; provenance): Catalogue record 422 overlap}
\noindent\textit{Transcription reads.} Kirzhnits record spans pages also assigned to Drinfeld and the subsequent physics report.\par
\noindent\textit{Critical reading.} Treat the bracketed notice of Kirzhnits as a heading; assign P712 formulas to record 423 and the distinct later contributions to records 424--425.\par
\noindent\textit{Evidence.} P712--P714 scans and both transcriptions read; no second independent Kirzhnits exposition occurs at the heading.\par
\subsection*{P000713 (both; ambiguity): Unknown-sign counterfeit puzzle}
\noindent\textit{Transcription reads.} 13 coins in three weighings\par
\noindent\textit{Critical reading.} Keep 13 when identifying the counterfeit coin alone; 12 is the bound if its heavier/lighter sign must also be determined.\par
\noindent\textit{Evidence.} P713 scan confirms 13. Complete independent balanced-ternary construction, upper bound and physical-pan check are in \texttt{coin\_annotation.tex} and \texttt{coin\_check\_output.json}; this is an ambiguity in the success criterion, not a digit error.\par
\subsection*{P000714 (both; faithful error): Acquiring information}
\noindent\textit{Transcription reads.} The wording says that acquiring information decreases information.\par
\noindent\textit{Critical reading.} It decreases the remaining uncertainty or entropy, under the stated conditioning model.\par
\noindent\textit{Evidence.} P714 scan and adjoining conditional-information discussion; Shannon 1948 §6 properties 5--6. This is a conceptual slip preserved by the transcription.\par
\subsection*{P000714 (both; faithful error): Binary entropy}
\noindent\textit{Transcription reads.} $H=p\log p+(1-p)\log(1-p)$\par
\noindent\textit{Critical reading.} $H=-p\log p-(1-p)\log(1-p)$\par
\noindent\textit{Evidence.} P714 scan confirms the missing leading negative sign in the original notebook. Shannon 1948 §6 and positivity at $p=1/2$ verify the emendation.\par
\subsection*{P000714 (both; ambiguity): Ideal-gas entropy schematic formula}
\noindent\textit{Transcription reads.} $S=\ln(V/N)$\par
\noindent\textit{Critical reading.} Treat this as the volume-dependent entropy per particle in units $k_{B}$, up to temperature-dependent and additive terms; use $\Delta S/N=k_{B}\ln(V_{f}/V_{i})$ at fixed temperature.\par
\noindent\textit{Evidence.} P714 scan; source omits normalization and fixed-temperature assumptions. The item supplies the elementary thermodynamic interpretation rather than replacing the handwritten formula wholesale.\par
\subsection*{P000715 (both; faithful error): Area increase}
\noindent\textit{Transcription reads.} Black-hole area always increases.\par
\noindent\textit{Critical reading.} Apply the classical area theorem under its energy and global-causality hypotheses; quantum evaporation is outside that claim.\par
\noindent\textit{Evidence.} P715 scan; Hawking 1971 abstract and Hawking 1975 §1 and §3 distinguish classical area increase from evaporation.\par
\subsection*{P000715 (both; faithful error): Bekenstein supervisor}
\noindent\textit{Transcription reads.} Bekenstein is described as Hawking's student.\par
\noindent\textit{Critical reading.} Bekenstein was John Archibald Wheeler's doctoral student.\par
\noindent\textit{Evidence.} P715 scan confirms the attribution. Wald 2015 personal recollection explicitly describes their graduate study under Wheeler.\par
\subsection*{P000715 (both; faithful error): Black-hole parameters}
\noindent\textit{Transcription reads.} Mass, radius and charge are the only remembered data.\par
\noindent\textit{Critical reading.} For the appropriate stationary four-dimensional electrovac uniqueness theorem, use mass, angular momentum and electric charge; radius is not an independent fourth parameter.\par
\noindent\textit{Evidence.} P715 scan; Chrusciel--Costa--Heusler 2012 Theorem 3.3 read with hypotheses. The broad notebook slogan is not a theorem for all black holes.\par
\subsection*{P000715 (both; ambiguity): Ground-state ket}
\noindent\textit{Transcription reads.} $H\lvert\Phi\rangle=E_{0}\lvert\phi\rangle$\par
\noindent\textit{Critical reading.} Use the same state symbol on both sides of the eigenvalue equation.\par
\noindent\textit{Evidence.} P715 scan and transcriptions inspected; capitalization changes in the displayed equation should not be interpreted as two eigenstates.\par
\subsection*{P000715 (both; faithful error): Solar-mass Schwarzschild radius}
\noindent\textit{Transcription reads.} Approximately 7 km\par
\noindent\textit{Critical reading.} Approximately 2.95 km for one solar mass.\par
\noindent\textit{Evidence.} P715 scan confirms 7. Evaluation $2(GM_{\odot})/c^{2}$ with nominal $GM_{\odot}$ from Prsa et al. 2016 Table 1 gives 2.95325 km; the notebook error is faithfully transcribed.\par
\subsection*{P000716 (both; faithful error): Acceleration temperature}
\noindent\textit{Transcription reads.} $kT=ha/c$\par
\noindent\textit{Critical reading.} $k_{B}T=\hbar a/(2\pi c)$, with $\hbar$ explicitly distinguished from Planck $h$.\par
\noindent\textit{Evidence.} P716 scan and transcribed TikZ formula inspected. Crispino--Higuchi--Matsas 2008 thermal formulas verify the coefficient; if $h$ means ordinary Planck constant the denominator is $4\pi^{2}c$.\par
\subsection*{P000716 (both; faithful error): Cosmological horizon}
\noindent\textit{Transcription reads.} Light has not yet arrived from a region, so it is called an event horizon and assigned a temperature.\par
\noindent\textit{Critical reading.} Distinguish a particle horizon from a future event horizon; a temperature requires a specified spacetime, state and observer.\par
\noindent\textit{Evidence.} P716 scan. The definitions have different time limits; no definite cosmological model in the passage resolves the proposed temperature. This is a conceptual correction, not an invented de Sitter assertion.\par
\subsection*{P000717 (both; faithful error): Schwartz local-functional exercise}
\noindent\textit{Transcription reads.} Integral over $\mathbb{R}$ of a differential polynomial is used without restricting its constant term.\par
\noindent\textit{Critical reading.} Require \(f(0)=0\) for this integral on all Schwartz tuples. The algebraic kernel is\[\ker(\text{Euler derivative})=\mathbb C+D(A),\]so constants must be separated.\par
\noindent\textit{Evidence.} P717 scan and Gelfand--Dikii 1977 §1 read. A nonzero constant density is not integrable on $\mathbb{R}$; the formal kernel also includes constants.\par
\subsection*{P000718 (both; provenance): End of record 426}
\noindent\textit{Transcription reads.} Catalogue page range ends at P717.\par
\noindent\textit{Critical reading.} Include P718 in record 426; the next speaker starts at P719.\par
\noindent\textit{Evidence.} Both transcriptions and P717--P718 scans read; the definitions, theorem, simpler-proof request and Novikov intervention form one continuation.\par
\subsection*{P000718 (both; ambiguity): Hamiltonian-matrix line continuation}
\noindent\textit{Transcription reads.} A subtraction sign appears at the end of one line and again at the start of the continuation.\par
\noindent\textit{Critical reading.} Read a single subtraction between the two binomial differential-operator terms.\par
\noindent\textit{Evidence.} P718 scan compared with the displayed formula and Gelfand--Dikii 1977 §2 Eq. 8; treating the line-wrap marks as two negatives would violate formal skew-adjointness.\par
\subsection*{P000718-P000719 (both; provenance): Record 428 boundary}
\noindent\textit{Transcription reads.} The catalogue routing window includes P000718\par
\noindent\textit{Critical reading.} S. I. Gelfand begins on P000719; P000718 belongs to the preceding Dickey--Novikov material\par
\noindent\textit{Evidence.} Heading and continuity of the formulas; independently checked with the early-record editor.\par
\subsection*{P000722 (both; provenance): Uploaded scan integrity}
\noindent\textit{Transcription reads.} The supplied PNG ends after 10,722 bytes\par
\noindent\textit{Critical reading.} Use the complete 169,042-byte original recovered from the archive\par
\noindent\textit{Evidence.} Recovered \url{https://old.mccme.ru/gelfand/notes/chapter_011/11_part_000018.png}; the supplied file is its exact prefix, and downloaded neighboring P000721/P000723 match the supplied files byte for byte. The full recovered page was visually read; see \texttt{work/code/output/scan\_722\_recovery.json}.\par
\subsection*{P000722; P000726 (both; faithful error): Quadratic central element}
\noindent\textit{Transcription reads.} \(4E_-E_+-H^2+2H\)\par
\noindent\textit{Critical reading.} \(4E_-E_++H^2+2H\)\par
\noindent\textit{Evidence.} The complete recovered P000722 and P000726 scans retain the minus sign. Standard positive matrix generators on P000722 give \([E_+,E_-]=H\); a direct commutator check requires \(+H^2\). This is a notebook error, not a transcription correction.\par
\subsection*{P000723 (both; faithful error): Right pair of arrows in the three-vertex diagram}
\noindent\textit{Transcription reads.} \(e_{13},e_{31}\) arrow labels\par
\noindent\textit{Critical reading.} For left modules, \(e_{ij}:L_j\to L_i\); reverse the inconsistent right-hand labels\par
\noindent\textit{Evidence.} The right pair conflicts with the displayed matrix units while the left pair uses the left-module convention. A global transpose does not repair both pairs. Scan inspected.\par
\subsection*{P000724 (both; ambiguity): Upper-left matrix order}
\noindent\textit{Transcription reads.} Vanishing positions transposed relative to P000723\par
\noindent\textit{Critical reading.} State the left/right-module convention and transpose the entire order if changing it\par
\noindent\textit{Evidence.} P000723 and P000724 scans give different off-diagonal vanishing patterns. They can describe opposite/transposed orders but are not literally the same presentation without a convention change.\par
\subsection*{P000726 (both; faithful error): Compact generator spectrum}
\noindent\textit{Transcription reads.} \(\operatorname{Spec}(X_0)\subset2\mathbb{Z}\)\par
\noindent\textit{Critical reading.} \(\operatorname{Spec}(X_0)\subset2i\mathbb{Z}\)\par
\noindent\textit{Evidence.} For \(X_0=E_+-E_-\), the compact generator acts skew-Hermitian in the usual unitary convention. P000729 explicitly gives \(2ik\) on the Fourier vectors, confirming the missing \(i\).\par
\subsection*{P000727 (both; faithful error): Normalizer of the left ideal}
\noindent\textit{Transcription reads.} Smallest subalgebra in which \(J\) is a two-sided ideal\par
\noindent\textit{Critical reading.} Largest such subalgebra: the normalizer \(N(J)\)\par
\noindent\textit{Evidence.} The displayed normalizer definition and quotient require the largest subalgebra preserving the ideal. Scan confirms the adjective.\par
\subsection*{P000728 (both; faithful error): Casimir on functions homogeneous of degree \(s\)}
\noindent\textit{Transcription reads.} \(-s^2+2s\)\par
\noindent\textit{Critical reading.} \(s(s+2)\) in the displayed homogeneous-degree convention\par
\noindent\textit{Evidence.} Substituting \(E_+=x_1\partial_{x_2}\), \(E_-=x_2\partial_{x_1}\), and \(H=x_1\partial_{x_1}-x_2\partial_{x_2}\) into the corrected central element gives \(s(s+2)\). Scan checked.\par
\subsection*{P000728-P000729 (both; faithful error): Second parameter at the trivial central character}
\noindent\textit{Transcription reads.} \(z=0\) at \(s=0\) and \(s=2\); the dual reducibility point is \(s=2\)\par
\noindent\textit{Critical reading.} The second point is \(s=-2\) in the stated homogeneous-degree convention\par
\noindent\textit{Evidence.} The corrected Casimir eigenvalue \(s(s+2)\) vanishes at \(0\) and \(-2\). The dual principal-series extension occurs at degree \(-2\); the notation \(s=2\) uses the opposite parameter convention.\par
\subsection*{P000729 (both; ambiguity): Weyl reflection on the principal-series parameter}
\noindent\textit{Transcription reads.} \(s\mapsto2-s\)\par
\noindent\textit{Critical reading.} \(s\mapsto-s-2\) for functions homogeneous of degree \(s\)\par
\noindent\textit{Evidence.} The notebook mixes the homogeneous-degree parameter with its negative; \(2-s\) is the reflection for the latter convention. It cannot be combined unchanged with the preceding formulas.\par
\subsection*{P000730 (both; faithful error): First gauge-group example}
\noindent\textit{Transcription reads.} \(\operatorname{SU}(1)\)\par
\noindent\textit{Critical reading.} \(\operatorname{U}(1)\)\par
\noindent\textit{Evidence.} \(\operatorname{SU}(1)\) is trivial; the electromagnetic abelian example is \(\operatorname{U}(1)\). The scan confirms \(\operatorname{SU}(1)\), so the notebook itself is corrected.\par
\subsection*{P000730 (both; faithful error): Topological charge formula}
\noindent\textit{Transcription reads.} \(\int F\wedge F\)\par
\noindent\textit{Critical reading.} Include an invariant trace and normalization; for example, \(-\frac{1}{8\pi^2}\int\operatorname{tr}(F\wedge F)\) in a fixed anti-Hermitian convention\par
\noindent\textit{Evidence.} Matrix-valued curvature must be traced to obtain a scalar characteristic number. Orientation and self-duality conventions determine its sign.\par
\subsection*{P000731 (both; faithful error): Action minimizers}
\noindent\textit{Transcription reads.} Strict local minima at self-dual fields\par
\noindent\textit{Critical reading.} Absolute minima within a fixed topological sector, with gauge and moduli zero directions\par
\noindent\textit{Evidence.} Gauge-equivalent connections and distinct instantons of the same charge have the same action. The scan retains strict; the Bogomolny decomposition gives the qualified statement.\par
\subsection*{P000731 (both; faithful error): Multi-instanton parameter attribution}
\noindent\textit{Transcription reads.} \(5n+4\) parameters attributed to 't Hooft\par
\noindent\textit{Critical reading.} Jackiw--Nohl--Rebbi family; distinguish raw parameters, redundancies and gauge quotient\par
\noindent\textit{Evidence.} Atiyah--Hitchin--Singer, PNAS 74 (1977), p.2662 explicitly distinguish the \(5n\) family of 't Hooft from the \(5n+4\) family of Jackiw et al.; both primary pages read.\par
\subsection*{P000731 (both; faithful error): Pullback of an instanton}
\noindent\textit{Transcription reads.} An arbitrary continuous map of the sphere pulls an instanton back to an instanton\par
\noindent\textit{Critical reading.} Conformal diffeomorphisms preserve the relevant self-duality equation; arbitrary continuous maps do not\par
\noindent\textit{Evidence.} Continuity alone does not define a pulled-back smooth connection, and a general smooth pullback need not preserve the Hodge star. Scan checked.\par
\subsection*{P000732 (both; faithful error): Self-dual Lax pair}
\noindent\textit{Transcription reads.} The spectral parameter occurs with the same relative signs in both factors\par
\noindent\textit{Critical reading.} Choose compatible opposite relative signs for the stated self-duality convention\par
\noindent\textit{Evidence.} Direct expansion of the commutator and comparison with the corrected pair on P000741 reveal the relative-sign inconsistency; absolute signs depend on the chosen complex coordinates.\par
\subsection*{P000732 (both; faithful error): Twistor fibre}
\noindent\textit{Transcription reads.} All linear complex structures on \(\mathbb{R}^4\) form \(\mathbb{CP}^1\)\par
\noindent\textit{Critical reading.} Orthogonal complex structures compatible with a fixed orientation form \(\mathbb{CP}^1\)\par
\noindent\textit{Evidence.} Without orthogonality the space is \(\operatorname{GL}^{+}(4,\mathbb{R})/\operatorname{GL}(2,\mathbb{C})\), not a two-sphere. The metric/orientation hypotheses are missing from the notebook.\par
\subsection*{P000733 (both; faithful error): Fibre label in the bottom diagram}
\noindent\textit{Transcription reads.} \(\mathbb{CP}^2\)\par
\noindent\textit{Critical reading.} \(\mathbb{CP}^1\)\par
\noindent\textit{Evidence.} The fibre is the projective line of complex lines in one quaternionic line; scan inspected.\par
\subsection*{P000733 (both; faithful error): Top projection arrow}
\noindent\textit{Transcription reads.} \(\mathbb{CP}^3\leftarrow S^4\)\par
\noindent\textit{Critical reading.} \(\mathbb{CP}^3\longrightarrow S^4\)\par
\noindent\textit{Evidence.} The twistor fibration projects complex lines to quaternionic lines. The scan confirms the reversed arrow.\par
\subsection*{P000734 (both; faithful error): Integrability criterion}
\noindent\textit{Transcription reads.} Two arbitrary spanning fields commute\par
\noindent\textit{Critical reading.} The distribution is involutive; a suitable local coordinate frame can be chosen commuting\par
\noindent\textit{Evidence.} Vanishing of the commutator of every chosen spanning pair is stronger than Frobenius involutivity. The source is abbreviated.\par
\subsection*{P000735 (both; ambiguity): Definition of an inner function}
\noindent\textit{Transcription reads.} \(\lvert u\cdot e^{it}\rvert\)\par
\noindent\textit{Critical reading.} \(\lvert u(e^{it})\rvert\)\par
\noindent\textit{Evidence.} The definition of an inner function uses boundary values of \(u\) on the unit circle. The handwritten juxtaposition denotes evaluation; literal multiplication would make the formula incorrect.\par
\subsection*{P000736 (both; faithful error): Polynomial linear combination}
\noindent\textit{Transcription reads.} The final coefficient is \(u_{r+1}\)\par
\noindent\textit{Critical reading.} The final coefficient is \(u_r\) for the displayed list \(u_1,\ldots,u_r\)\par
\noindent\textit{Evidence.} The number of listed functions and degree \(r-1\) determine the final index; scan inspected.\par
\subsection*{P000736 (both; faithful error): Singular inner factor}
\noindent\textit{Transcription reads.} Positive sign in the singular-measure exponential\par
\noindent\textit{Critical reading.} Negative sign for a positive singular measure\par
\noindent\textit{Evidence.} The Poisson kernel has positive real part, so the singular inner factor is \(\exp\bigl(-\int\frac{\zeta+z}{\zeta-z}\,d\mu(\zeta)\bigr)\); the scan confirms the positive sign.\par
\subsection*{P000737 (both; faithful error): Hamiltonian of pole motion}
\noindent\textit{Transcription reads.} \(\frac12\sum_i p_i^2+\sum_{i<j}(q_i-q_j)^{-2}\)\par
\noindent\textit{Critical reading.} \(\frac12\sum_i p_i^2-4\sum_{i<j}(q_i-q_j)^{-2}\) when \(p_i=q_{i,y}\)\par
\noindent\textit{Evidence.} Laurent expansion for precisely the notebook KP equation and \(u=-2\sum_j(x-q_j)^{-2}\) gives \(q_{i,yy}=-8\sum_{j\ne i}(q_i-q_j)^{-3}\). Independently verified by the assembler in \texttt{work/code/check\_kp.py}. Other time/Poisson conventions must be explicit.\par
\subsection*{P000737 (both; faithful error): Lax operator and time sign}
\noindent\textit{Transcription reads.} \(L_2=D^3+\frac32uD+u\); \([\partial_y-L_1,\partial_t-L_2]=0\)\par
\noindent\textit{Critical reading.} \(M=D^3+\frac32uD+w\), \(w_x=\frac34(u_{xx}+u_y)\); \([\partial_y-L_1,\partial_t+M]=0\) for the notebook equation\par
\noindent\textit{Evidence.} The scan unmistakably ends with \(+u\). Direct operator expansion yields the displayed correction for the notebook KP sign; changing the time sign changes the PDE. Independent symbolic check saved by the assembler.\par
\subsection*{P000738 (both; ambiguity): Equation (2), duality convention}
\noindent\textit{Transcription reads.} \(F_{\mu\nu}=\varepsilon^{\mu\nu\rho\sigma}F_{\rho\sigma}\)\par
\noindent\textit{Critical reading.} Include \(1/2\) if summing over both ordered complementary indices; no factor if a single ordered complementary pair is intended\par
\noindent\textit{Evidence.} The preceding sentence describes a complementary pair, while the formula resembles Einstein summation. Record the convention instead of silently adding a factor.\par
\subsection*{P000738 (both; faithful error): Equation (3) summation index}
\noindent\textit{Transcription reads.} \(\sum_{i=1}^r A_\mu^{ij}s_j\)\par
\noindent\textit{Critical reading.} \(\sum_{j=1}^r A_\mu^{ij}s_j\)\par
\noindent\textit{Evidence.} The free index is \(i\) and the repeated vector index is \(j\); the scan confirms lower summation index \(i\).\par
\subsection*{P000738 (both; provenance): Record 435 date}
\noindent\textit{Transcription reads.} Catalogue: 21 September 1977\par
\noindent\textit{Critical reading.} Source heading: 21 November 1977 (21/XI/77)\par
\noindent\textit{Evidence.} The manuscript XI is clear and underlined; transcriptions retain it. Catalogue date is retained only as catalogue metadata.\par
\subsection*{P000740 (both; faithful error): Pulled-back curvature criterion}
\noindent\textit{Transcription reads.} \(D^*F=0\)\par
\noindent\textit{Critical reading.} \((D^*F)^{2,0}=0\), together with the conjugate condition for a unitary connection\par
\noindent\textit{Evidence.} The following coordinate equations and the twistor construction require vanishing of a type component, not the entire nonzero instanton curvature.\par
\subsection*{P000740 (both; faithful error): Quaternionic conjugation notation}
\noindent\textit{Transcription reads.} The antilinear map on \(\mathbb{C}^4\) is called an involution\par
\noindent\textit{Critical reading.} It squares to \(-1\) on \(\mathbb{C}^4\); its projectivization is a fixed-point-free involution\par
\noindent\textit{Evidence.} Substitution in the written coordinate formula gives \(\sigma^2=-\operatorname{id}\), while multiplication by \(-1\) acts trivially projectively.\par
\subsection*{P000743 (both; faithful error): Basis of \(I\)}
\noindent\textit{Transcription reads.} \(e_1,\ldots,e_r\)\par
\noindent\textit{Critical reading.} \(e_1,\ldots,e_n\)\par
\noindent\textit{Evidence.} The declared dimension of \(I\) is \(n\), and the vectors \(e_{ij}\) run through \(n\) rows. Scan checked.\par
\subsection*{P000743 (both; faithful error): Composition with quaternionic antilinear map}
\noindent\textit{Transcription reads.} \(\ell\circ\sigma\)\par
\noindent\textit{Critical reading.} \(\overline{\ell}\circ\sigma\) with the appropriate conjugate target\par
\noindent\textit{Evidence.} The source composition of a complex-linear map with \(\sigma\) is antilinear; target conjugation is required for the claimed complex-linear monad map. Scan checked.\par
\subsection*{P000743 (both; faithful error): Upper tensor-product arrow}
\noindent\textit{Transcription reads.} \(H\otimes\mathbb{C}\)\par
\noindent\textit{Critical reading.} \(H\otimes\mathbb{C}^4\)\par
\noindent\textit{Evidence.} The next construction uses four columns for every basis vector and a \(4n\) by \(4n\) Gram matrix; the exponent \(4\) is absent in the scan.\par
\subsection*{P000744 (both; faithful error): Frame of the orthogonal complement}
\noindent\textit{Transcription reads.} \(L^\perp\) is generated globally by \(2n\) sections over \(S^4\)\par
\noindent\textit{Critical reading.} The displayed affine formulas give a frame on the \(\mathbb{R}^4\) chart; no such global frame on \(S^4\) is implied\par
\noindent\textit{Evidence.} For nonzero instanton charge the complement has nontrivial Chern class and cannot be globally framed by exactly its rank many independent sections. The formulas use affine \(x\) coordinates.\par
\subsection*{P000746 (both; faithful error): Equation (7), occupation density}
\noindent\textit{Transcription reads.} \(a^*(k^{\prime})a(k)\)\par
\noindent\textit{Critical reading.} \(a^*(k)a(k)\)\par
\noindent\textit{Evidence.} The extra prime is visible in the scan but no independent \(k^{\prime}\) is integrated or defined; the free number density needs equal arguments.\par
\subsection*{P000748 (both; faithful error): Breather endpoint}
\noindent\textit{Transcription reads.} \(0<\alpha\leq\pi/2\)\par
\noindent\textit{Critical reading.} \(0<\alpha<\pi/2\) for a finite-period breather\par
\noindent\textit{Evidence.} At \(\alpha=\pi/2\), the frequency \(m\cos\alpha\) vanishes. It is a separatrix limit; the mass formula alone does not make the endpoint periodic.\par
\subsection*{P000748 (both; faithful error): Equation (18), relativistic energy}
\noindent\textit{Transcription reads.} \(\sqrt{p_k^2+M(\alpha_k)}\)\par
\noindent\textit{Critical reading.} \(\sqrt{p_k^2+M(\beta_k)^2}\) in the preceding parameter convention\par
\noindent\textit{Evidence.} The preceding text makes \(\beta\) the internal period parameter and \(\alpha\) the initial phase. The mass must be squared and cannot depend on the time phase. Scan inspected.\par
\subsection*{P000748 (both; faithful error): Equation (21), Poisson bracket order}
\noindent\textit{Transcription reads.} \(\{u_t(x),u(y)\}=\gamma\delta(x-y)\)\par
\noindent\textit{Critical reading.} \(\{u(x),u_t(y)\}=\gamma\delta(x-y)\) in the convention of equation (3)\par
\noindent\textit{Evidence.} The notebook reverses the order relative to its initial canonical bracket without changing the sign.\par
\subsection*{P000749 (both; faithful error): Equation (26), commutator convention}
\noindent\textit{Transcription reads.} \([\pi(x),u(y)]=i\delta(x-y)\)\par
\noindent\textit{Critical reading.} \([u(x),\pi(y)]=i\delta(x-y)\) in the usual convention matching equation (3)\par
\noindent\textit{Evidence.} The printed order is reversed relative to the earlier Poisson bracket. A different commutator convention would have to be stated.\par
\subsection*{P000749 (both; ambiguity): Unnumbered interaction after equation (26)}
\noindent\textit{Transcription reads.} \(m^2\gamma^{-1}(1-\cos\gamma u)\)\par
\noindent\textit{Critical reading.} Specify the field and coupling rescaling; canonical normalization is \(m^2\beta^{-2}(1-\cos\beta\varphi)\)\par
\noindent\textit{Evidence.} The displayed potential has quadratic term \(m^2\gamma u^2/2\) and hence does not preserve mass \(m\) under the usual canonical kinetic term. A convention ambiguity is retained.\par
\subsection*{P000750 (both; faithful error): Equation (32), stationary phase}
\noindent\textit{Transcription reads.} \(\widetilde{\Psi}(mx/t)e^{ix^2/(2mt)}\)\par
\noindent\textit{Critical reading.} \(\sqrt{m/(2\pi t)}e^{-i\pi/4}\widehat{\Psi}(mx/t)e^{imx^2/(2t)}\) for the stated Fourier normalization\par
\noindent\textit{Evidence.} Critical point \(p=mx/t\) and phase value \(mx^2/(2t)\) follow directly. The missing amplitude depends on the Fourier convention but includes \(t^{-1/2}\). Scan checked.\par
\subsection*{P000751 (both; ambiguity): No-higher-corrections conjecture}
\noindent\textit{Transcription reads.} \(M=8m/\gamma-m/\pi\) exactly\par
\noindent\textit{Critical reading.} Specify the renormalized mass convention before deciding the claim\par
\noindent\textit{Evidence.} The physical lightest-breather mass gives \(M/m_1=[2\sin(\pi\xi/2)]^{-1}\), \(\xi=\gamma/(8\pi-\gamma)\), whose expansion has a nonzero \(\gamma/192\) term. This does not determine a differently renormalized parameter \(m\).\par
\subsection*{P000751 (both; faithful error): Stationary sine-Gordon fluctuation potential}
\noindent\textit{Transcription reads.} \(q(x)=1/\cosh^2x\)\par
\noindent\textit{Critical reading.} \(K=-\partial_x^2+m^2-2m^2\operatorname{sech}^2(mx)\)\par
\noindent\textit{Evidence.} Korepin--Faddeev 1975, equation (16), and the second variation of the stated kink give the negative coefficient, factor two and asymptotic mass. The translation zero mode must be excluded from the determinant.\par
\subsection*{P000752 (both; transcription): Final handwritten sentence}
\noindent\textit{Transcription reads.} \ru{Массивные частицы появляются как солитоны} / Massive particles appear as solitons\par
\noindent\textit{Critical reading.} \ru{Массивные частицы рассеиваются как солитоны} / Massive particles scatter like solitons\par
\noindent\textit{Evidence.} The verb \ru{рассеиваются} is legible in the scan, independently read by both the record editor and chapter assembler. This changes the meaning to a scattering comparison.\par
\subsection*{P000753-P000754 (both; provenance): Record 436 envelope endpoint}
\noindent\textit{Transcription reads.} Catalogue routing window includes these pages in Faddeev material\par
\noindent\textit{Critical reading.} These pages begin S. I. Gelfand, 13 February 1978; combine with the continuation under record 438\par
\noindent\textit{Evidence.} Date and speaker heading on P000753; mathematical continuation on P000756. The late-record editor handles this adjacent precursor explicitly.\par
\subsection*{P000754 (both; ambiguity): regular-module example}
\noindent\textit{Transcription reads.} free modules lead to trivial bundles\par
\noindent\textit{Critical reading.} free modules lead to the zero bundle\par
\noindent\textit{Evidence.} The same paragraph computes all evaluation homology as zero; a positive-rank trivial bundle is different.\par
\subsection*{P000756 (both; faithful error): Serre bibliography}
\noindent\textit{Transcription reads.} 1953\par
\noindent\textit{Critical reading.} 1955 for Faisceaux algebriques coherents\par
\noindent\textit{Evidence.} Opened original Serre publication metadata and text; the quoted projective sheaf theorem belongs to FAC 1955.\par
\subsection*{P000759 (both; faithful error): section-module inverse}
\noindent\textit{Transcription reads.} direct sum over all integers is finitely generated\par
\noindent\textit{Critical reading.} use sufficiently high nonnegative truncation for an arbitrary coherent sheaf\par
\noindent\textit{Evidence.} A skyscraper sheaf has nonzero global sections in every negative twist; this cannot be generated in finitely many graded degrees.\par
\subsection*{P000760 (both; faithful error): total differential}
\noindent\textit{Transcription reads.} \(d=\partial+\sum_i\xi_i x_i\), with no totalization sign\par
\noindent\textit{Critical reading.} insert the usual degree-dependent totalization sign or state anticommuting conventions\par
\noindent\textit{Evidence.} The scan shows only a plus sign; commuting differentials give the cross term \(2\partial d\) in characteristic zero.\par
\subsection*{P000761 (both; faithful error): comparison evaluated at 1}
\noindent\textit{Transcription reads.} \(\varphi(v)(1)=\sum_i x_i\otimes\xi_i v\)\par
\noindent\textit{Critical reading.} A unit comparison must include \(1\otimes v\).\par
\noindent\textit{Evidence.} Scan confirms sum; on a trivial module the displayed value vanishes, contradicting the asserted quasi-isomorphism.\par
\subsection*{P000764 (both; faithful error): decompositions (7)--(10)}
\noindent\textit{Transcription reads.} orthogonal support subspaces used as dual isotropic CAR subspaces\par
\noindent\textit{Critical reading.} distinguish the support decomposition from a complex isotropic polarization\par
\noindent\textit{Evidence.} Conserved pairing (4) makes opposite-support data orthogonal; (9) instead pairs the two spaces nontrivially. Missing identification is mathematical, not a legibility repair.\par
\subsection*{P000765 (both; faithful error): definition of \(\operatorname{Spin}\)}
\noindent\textit{Transcription reads.} all Clifford elements normalizing \(W\)\par
\noindent\textit{Critical reading.} invertible Clifford group; normalized even elements for \(\operatorname{Spin}\)\par
\noindent\textit{Evidence.} Scan has no parity or norm-one condition and explicitly observes a noninjective conjugation map.\par
\subsection*{P000766 (both; faithful error): Painleve classification}
\noindent\textit{Transcription reads.} only six types have fixed singular points\par
\noindent\textit{Critical reading.} six irreducible Painleve types within the classified class, with no movable critical points\par
\noindent\textit{Evidence.} Movable poles are allowed; Garnier 1912 and the standard isomonodromic formulation concern critical branching.\par
\subsection*{P000766 (both; faithful error): Physical Review B citation}
\noindent\textit{Transcription reads.} Wu, McCoy, Tracy\par
\noindent\textit{Critical reading.} Wu, McCoy, Tracy, and Barouch\par
\noindent\textit{Evidence.} Publisher record DOI 10.1103/PhysRevB.13.316 identifies four authors.\par
\subsection*{P000766 (both; faithful error): equation (23)}
\noindent\textit{Transcription reads.} \((m-\Gamma)w=0\), with \(\Gamma\) the constant exchange matrix\par
\noindent\textit{Critical reading.} intended differential operator unresolved; do not silently supply it\par
\noindent\textit{Evidence.} The scan was inspected: \(m\) is scalar and no differential sign is present. The literal matrix has no kernel for \(m\ne\pm1\).\par
\subsection*{P000768 (both; faithful error): Appell \(F_4\) series (2)}
\noindent\textit{Transcription reads.} coefficient sum without \(x^m y^n\)\par
\noindent\textit{Critical reading.} insert \(x^m y^n\) after the coefficient\par
\noindent\textit{Evidence.} The scan omits the factor and the printed expression is otherwise independent of both variables; this is a faithful notebook omission.\par
\subsection*{P000769 (both; faithful error): determinant and discriminant}
\noindent\textit{Transcription reads.} last row \(i_n\); \(Y\) defined by vanishing minors without a union quantifier\par
\noindent\textit{Critical reading.} last row \(i_{n+1}\); at least one relevant minor vanishes\par
\noindent\textit{Evidence.} There are \(n+1\) columns; general-position failure is the union of minor-zero loci, whereas vanishing of all maximal minors is a stronger rank condition.\par
\subsection*{P000770 (both; faithful error): nonresonance and vanishing}
\noindent\textit{Transcription reads.} subset sums outside the strictly positive integers imply \(H^p=0\) for \(p\ne n\)\par
\noindent\textit{Critical reading.} exclude the trivial local system and impose sufficient nonresonance for concentration\par
\noindent\textit{Evidence.} The scan explicitly defines \(\mathbb Z_+\) as integers greater than zero. All residues zero satisfy the written hypothesis and give nonzero \(H^0\).\par
\subsection*{P000770-P000771 (both; faithful error): Euler characteristic}
\noindent\textit{Transcription reads.} \(H^n=\chi\) and \(\chi=\binom{m-1}{n}\)\par
\noindent\textit{Critical reading.} \(\dim H^n=(-1)^n\chi=\binom{m-1}{n}\), under concentration\par
\noindent\textit{Evidence.} For \(n=1\), the complement of \(m\) points has \(\chi=1-m\). The scan has neither \(\dim\) nor the parity factor.\par
\subsection*{P000773 (both; faithful error): all natural operations}
\noindent\textit{Transcription reads.} Landweber-Novikov operations are all natural cohomological operations\par
\noindent\textit{Critical reading.} specify the Landweber--Novikov subalgebra \(S\) and the cohomology theory\par
\noindent\textit{Evidence.} The full complex-bordism stable operation algebra also contains coefficient operations; the opened 1978 paper states the precise \(S\) identification.\par
\subsection*{P000775 (both; faithful error): Novikov operator factor order}
\noindent\textit{Transcription reads.} \(L=\partial^{-1}L_3\) identified with the earlier vector-field recursion\par
\noindent\textit{Critical reading.} distinguish gradient recursion \(\partial^{-1}J_2\) from vector-field recursion \(J_2\partial^{-1}\)\par
\noindent\textit{Evidence.} The scanned inverse derivative is on the left; applying the earlier explicit \(L\) to a test function fixes the opposite order for vector fields.\par
\subsection*{P000775 (both; faithful error): transmission-phase formula (8)}
\noindent\textit{Transcription reads.} \(\varphi^{(0)}=-c^{(0)}/2\) when \(T=|T|\exp(2i\varphi)\)\par
\noindent\textit{Critical reading.} \(\varphi^{(0)}=-c^{(0)}/4\) for the displayed Schr\"odinger normalization\par
\noindent\textit{Evidence.} The leading term is \(\log T=-i\int_{\mathbb R}u/(2k)\), with \(c^{(0)}=\int_{\mathbb R}u\); the left-Jost Volterra estimate independently verifies this normalization. The formula requires a normalization repair, not a transcription change.\par
\subsection*{P000776 (both; provenance): lower-page lecture boundary}
\noindent\textit{Transcription reads.} included within Calogero envelope\par
\noindent\textit{Critical reading.} V. P. Palamodov at the Mathematical Society, 21 March 1978\par
\noindent\textit{Evidence.} Fresh dated Russian heading and speaker visibly occur below the end of the Calogero notes.\par
\subsection*{P000776 (both; faithful error): moment equation (11)}
\noindent\textit{Transcription reads.} \(dc^{(n)}/dt=n\beta_0c^{(n)}\) in the pure scaling case\par
\noindent\textit{Critical reading.} \(dc^{(n)}/dt=(2n+1)\beta_0c^{(n)}\)\par
\noindent\textit{Evidence.} The scan clearly has \(n\). Direct scaling \(u(x,t)=e^{2\beta_0t}u_0(e^{\beta_0t}x)\) gives the stated weight; \(n=0\) already contradicts the notebook.\par
\subsection*{P000777 (both; faithful error): equation (6) second scattering factor}
\noindent\textit{Transcription reads.} \(\Psi(x,t)\)\par
\noindent\textit{Critical reading.} \(\Psi(x,k)\)\par
\noindent\textit{Evidence.} The scanned expression uses \(t\), but this is the stationary two-potential identity with common spectral parameter \(k\); mathematical emendation.\par
\subsection*{P000777 (both; faithful error): reflection uniqueness}
\noindent\textit{Transcription reads.} potential determined by \(R\) alone and discrete spectrum by its poles\par
\noindent\textit{Critical reading.} retain bound-state eigenvalues and norming constants, or specify a restrictive analytic class\par
\noindent\textit{Evidence.} Zero and nonzero reflectionless potentials have \(R\equiv0\); the notebook first lists discrete information separately.\par
\subsection*{P000778 (both; faithful error): alternative recursion operator}
\noindent\textit{Transcription reads.} \(\partial^3+2q\partial+2\partial q\)\par
\noindent\textit{Critical reading.} minus signs for both \(q\) terms when retaining the earlier potential convention\par
\noindent\textit{Evidence.} The scan has plus signs. Multiplication by \(\partial^{-1}\) otherwise gives \(+4q\), contrary to the explicitly written \(-4q\) in (12).\par
\subsection*{P000780 (both; ambiguity): superposition denominator}
\noindent\textit{Transcription reads.} \(Q_1\) [unclear sign] \(Q_2\)\par
\noindent\textit{Critical reading.} \(Q_1-Q_2\), conditional on compatibility with the four displayed Riccati relations\par
\noindent\textit{Evidence.} The overwritten scan mark remains unclear; eliminating derivatives algebraically forces the minus sign in this normalization.\par
\subsection*{P000781 (both; faithful error): small-oscillation equation}
\noindent\textit{Transcription reads.} \(\ddot\varepsilon+A\varepsilon=0\)\par
\noindent\textit{Critical reading.} \(\ddot\varepsilon+(I+A)\varepsilon=0\)\par
\noindent\textit{Evidence.} The scan confirms the omitted identity; linearization of the \(-x_j\) force contributes \(-\varepsilon_j\), also visible for \(n=1\).\par
\subsection*{P000782-P000785 (both; provenance): catalogue session date}
\noindent\textit{Transcription reads.} 4 April 1978\par
\noindent\textit{Critical reading.} heading 3 April 1978; next meeting 4 April 1978\par
\noindent\textit{Evidence.} The P782 heading reads 3/IV-78 and P785 explicitly labels 4/IV as the next meeting. The catalogue date is retained in the date field solely for envelope identity.\par
\subsection*{P000784 (both; faithful error): chain differential (11)}
\noindent\textit{Transcription reads.} \(e_1\) after the bracket term\par
\noindent\textit{Critical reading.} \(e_{i_1}\), with the indicated omitted factors\par
\noindent\textit{Evidence.} The general CE differential on P782 and preservation of weight require the indexed original first factor.\par
\subsection*{P000785 (both; faithful error): spectral formula description}
\noindent\textit{Transcription reads.} \(E(I)\) is the multiplicity; labels count distinct eigenvalues\par
\noindent\textit{Critical reading.} \(E(I)\) is an eigenvalue; multiplicities \(2^{\alpha}\) add at collisions\par
\noindent\textit{Evidence.} The full corrected theorem is Fuchs--Wilmarth (2010), Theorem 2. The exact weight-25 example \(E(1,10,14)=E(3,7,15)=320\) gives multiplicity 12.\par
\subsection*{P000786-P000787 (both; provenance): catalogue heading versus embedded lecture}
\noindent\textit{Transcription reads.} 1978 heading without mathematics\par
\noindent\textit{Critical reading.} adjacent Kruskal lecture, probably 30 October 1978, possibly 30 November\par
\noindent\textit{Evidence.} The scan includes named speaker, date and two pages of surreal-number mathematics. Month uncertainty retained.\par
\subsection*{P000787 (both; ambiguity): associativity sign}
\noindent\textit{Transcription reads.} \(x+(y\mathbin{?}z)\)\par
\noindent\textit{Critical reading.} \(x+(y+z)\) as the intended additive-group law\par
\noindent\textit{Evidence.} The sign is not paleographically certain; surrounding addition properties determine the mathematical interpretation.\par
\subsection*{P000810 (both; faithful error): relative/absolute cohomology sign}
\noindent\textit{Transcription reads.} \(H^1\) of the relative complex is a member of \(H^1\) of the absolute complex\par
\noindent\textit{Critical reading.} canonical injection in degree one, with image the extensions splitting over \(\mathfrak h\)\par
\noindent\textit{Evidence.} The scan confirms the membership sign. A relative coboundary is an absolute coboundary from an \(\mathfrak h\)-invariant element.\par
\subsection*{P000811 (both; ambiguity): orthogonal-group cohomology sentence}
\noindent\textit{Transcription reads.} nonsemisimple \(SO_0(m,n)\), no nonunitary representations with nonzero cohomology\par
\noindent\textit{Critical reading.} intended statement unresolved; do not silently replace both negations\par
\noindent\textit{Evidence.} The scan confirms both words; Delorme (1977), Theorem V.4, concerns unitary representations of simple groups excluding \(\mathfrak{so}(n,1)\) and \(\mathfrak{su}(n,1)\).\par
\subsection*{P000812 (both; provenance): catalogue heading and embedded date}
\noindent\textit{Transcription reads.} YEAR 1979 without a dated lecture\par
\noindent\textit{Critical reading.} 12 February 1979 for the Gelfand/Belavin page; the earlier two pages are undated continuations\par
\noindent\textit{Evidence.} The heading 12/II-79 is plainly visible; the year heading is retained as the catalogue envelope.\par
\subsection*{P000812 (both; faithful error): path-integral coauthor}
\noindent\textit{Transcription reads.} Gibbs\par
\noindent\textit{Critical reading.} probably Hibbs\par
\noindent\textit{Evidence.} The manuscript explicitly has Gibbs; the Feynman path-integral coauthor is A.R.Hibbs. Bibliographic interpretation, not a claim that the scan says Hibbs.\par
\subsection*{P000823 (both; faithful error): curvature on \(\mathbb{CP}^1\)}
\noindent\textit{Transcription reads.} curvature zero because the dimension is one\par
\noindent\textit{Critical reading.} true for a holomorphic connection on a complex curve; false for a general smooth connection\par
\noindent\textit{Evidence.} The manuscript does not state holomorphic at the input but asks for a holomorphic output; complex dimension one does not annihilate smooth real two-forms.\par
\subsection*{P000824 (both; faithful error): Cartan classification hypotheses}
\noindent\textit{Transcription reads.} bounded below and transitive imply four series\par
\noindent\textit{Critical reading.} state characteristic zero, finite pieces, infinite dimension, irreducibility and simplicity/category\par
\noindent\textit{Evidence.} Kac (1968) distinguishes four simple types in Theorem 4 from six transitive irreducible Cartan types in Theorem 5.\par
\subsection*{P000824-P000825 (both; faithful error): recovery from degrees \(-1\) and \(0\)}
\noindent\textit{Transcription reads.} everything recovered\par
\noindent\textit{Critical reading.} maximal prolongation is recovered; another transitive algebra embeds in it\par
\noindent\textit{Evidence.} The definitions explicitly change containment to equality; maximality is the missing condition.\par
\subsection*{P000824-P000831 (both; provenance): catalogue year and speaker boundaries}
\noindent\textit{Transcription reads.} 5 March 1978, Fuchs over P823--P831\par
\noindent\textit{Critical reading.} P823: preceding Beilinson; P824--P829: Fuchs, 5 March 1979; P830--P831: Ogievetsky, 19 March 1979\par
\noindent\textit{Evidence.} All three speaker headings and both 1979 dates are visible in the scans. The catalogue date is retained as the required envelope identifier.\par
\subsection*{P000826 (both; faithful error): rank-one criterion}
\noindent\textit{Transcription reads.} rank-one element over an unspecified field\par
\noindent\textit{Critical reading.} finite-dimensional complex input; for real input test its complexification\par
\noindent\textit{Evidence.} Ottazzi--Warhurst (2011), Theorem 2.3 and Section 2.4, give the exact condition and explain the field dependence.\par
\subsection*{P000827 (both; faithful error): Cartan/Serre index convention}
\noindent\textit{Transcription reads.} \([h_\beta,e_\alpha]=n_{\alpha\beta}e_\alpha\), but the exponent is \(1-n_{\alpha\beta}\)\par
\noindent\textit{Critical reading.} with that bracket convention, exponent \(1-n_{\beta\alpha}\)\par
\noindent\textit{Evidence.} The scan confirms the inconsistent ordering; the root-string length uses the action of \(h_\alpha\) on \(e_\beta\). The factor 2 in \([e,f]\) is only a rescaling.\par
\subsection*{P000828 (both; faithful error): loop realization}
\noindent\textit{Transcription reads.} affine algebra equals \(\mathfrak g\otimes\mathbb C[t,t^{-1}]\)\par
\noindent\textit{Critical reading.} center quotient of derived affine algebra; restore central extension/full Cartan when needed\par
\noindent\textit{Evidence.} The preceding kernel-of-\(A\) line gives a nonzero center, whereas the loop algebra for simple \(\mathfrak g\) has zero center.\par
\subsection*{P000829 (both; faithful error): simplicity of periodic bracket}
\noindent\textit{Transcription reads.} simple Lie algebra\par
\noindent\textit{Critical reading.} graded-simple in the finite-growth classification sense\par
\noindent\textit{Evidence.} The explicit loop realization has the proper ungraded ideal \((t-1)\mathfrak g[t,t^{-1}]\); Kac (1968) defines simple as having no homogeneous ideals.\par
\subsection*{P000831 (both; ambiguity): superfield mixed quadratic term}
\noindent\textit{Transcription reads.} \(\theta\sigma^\alpha\bar\theta h_\alpha\)\par
\noindent\textit{Critical reading.} \(\theta\sigma^m\bar\theta v_m\), if the index is genuinely Greek\par
\noindent\textit{Evidence.} The two-spinor product transforms as a vector. Latin \(a\) remains a paleographic possibility, so this is a conditional notation repair.\par
